\chaptertitle{Lisp: Lists and Recursion}


                      Lisp: Lists and Recursion

                                     March, 1983




SINCE I ended the previous column with a timely newsbreak about the homely
Glazunkian porpuquine, I felt it only fitting to start off the present column with more
about that little-known but remarkable beast. As you may recall, the quills on any
porpuquine (except for the tiniest ones) are smaller porpuquines. The tiniest
porpuquines have no quills but do have a nose, and a very important nose at that,
since the Glazunkians base their entire monetary system on that little nose. Consider,
for instance, the value of 3-inch porpuquines in Outer Glazunkia. Each one always
has nine quills (contrasting with their cousins in Inner Glazunkia, which always have
seven); thus each one has nine 2-inch porpuquines sticking out of its body. Each of
those in turn sports nine 1-inch porpuquines, out of each of which sprout nine zero-
inch porpuquines, each of which has one nose. All told, this comes to 9 X 9 X 9 X 1
noses, which means that a 3-inch porpuquine in Outer Glazunkia has a buying power
of 729 noses. If, by contrast, we had been in Inner Glazunkia and had started with a 4-
incher, that porpuquine would have a buying power of 7 X 7 X 7 X 7 X 1=2,401
noses.
        Let's see if we can't come up with a general recipe for calculating the buying
power (measured in noses) of any old porpuquine. It seems to me that it would go
something like this:

The buying power of a porpuquine with a given quill count and size is:
      if its size = 0, then 1;
      otherwise, figure out the buying power of a porpuquine with
                       the same quill count but of the next smaller size,
               and multiply that by the quill count.

We can shorten this recipe by adopting some symbolic notation. First, let q stand for
the quill count and s for the size. Then let cond stand for "if" and t for "otherwise".
Finally, use a sort of condensed algebraic notation in which the English names of
operations are placed to the left of their operands, inside parentheses. We get
something like this:




Lisp: Lists and Recursion                                                          410

(buying-power q s) Is:
      cond (eq s 0) 1;
             t (times q (buying-power q (next-smaller s)))

This is an exact translation of the earlier English recipe into a slightly more symbolic
form. We can make it a little more compact and symbolic by adopting a couple of
new conventions. Let each of the two cases (the case where s equals zero and the
"otherwise" case) be enclosed in parentheses; in general, use parentheses to enclose
each logical unit completely. Finally, indicate by the words def and lambda that this is
a definition of a general notion called buying-power with two variables (quill count q
and size s). Now we get:

(def buying-power (lambda (q s)
       (cond ((eq s 0) 1)
              (t (times q (buying-power q (next-smaller s)))))))

I mentioned above that the buying power of a 9-quill, 3-inch porpuquine is 729 noses.
This could be expressed by saying that (buying-power 9 3) equals 729. Similarly,
(buying-power 7 4) equals 2,401.

                                        * * *

       Well, so much for porpuquines. Now let's limp back to Lisp after this rather
long digression. I had posed a puzzle, toward the end of last month's column, in which
the object was to write a Lisp function that subsumed a whole family of related
functions called square, cube, 4th-power, 5th-power, and so on. I asked you to
come up with one general function called power, having two variables, such that
(power 9 3) gives 729, (power 7 4) gives 2,401, and so on. I had presented a "tower
of power"-that is; an infinitely tall tower of separate Lisp definitions, one for each
power, connecting it to the preceding power. Thus a typical floor in this tower would
be:

(def 102nd-power (lambda (q) (times q (101st-power q))))

Of course, 101st-power would refer to 100th-power in its definition, and so on, thus
creating a rather long regress back to the simplest, or "embryonic", case. Incidentally,
that very simplest case, rather than square or even 1st-power, is this:

(def 0th-power (lambda (q) 1))

I told you that you had all the information necessary to assemble the proper definition.
All you needed to observe is, of course, that each floor of the




Lisp: Lists and Recursion                                                           411

tower rests on the "next-smaller" floor (except for the bottom floor, which is a "stand-
alone" floor). By "next-smaller", I mean the following:

(def next-smaller (lambda (s) (difference s 1)))

Thus (next-smaller 102) yields 101. Actually, Lisp has a standard name for this
operation (namely, subl) as well as for its inverse operation (namely, addl ). If we put
all our observations together, we come up with the following universal definition:

(def power (lambda (q s)
       (cond ((eq s 0) 1)
              (t (times q (power q (next-smaller s)))))))

This is the answer to the puzzle I posed. Hmmm, that's funny ... I have the strangest
sense of deja vu. I wonder why!

                                        *   *   *

    The definition presented here is known as a recursive definition, for the reason
that inside the definiens, the definiendum is used. This is a fancy way of saying that I
appear to be defining something in terms of itself, which ought to be considered
gauche if not downright circular in anyone's book! To see whether the Lisp genie
looks askance upon such trickery, let's ask it to figure out (power 9 3):

    (power 9 3)
    729
   


Well, fancy that! No complaints? No choking? How can the Lisp genie swallow such
nonsense?
         The best explanation I can give is to point out that no circularity is actually
involved. While it is true that the definition of power uses the word power inside
itself, the two occurrences are referring to different circumstances. In a nutshell,
(power q s) is being defined in terms of a simpler case, namely, (power q (next-
smaller s)). Thus I am defining the 44th power in terms of the 43rd power, and that in
terms of the next-smaller power, and so on down the line until we come to the
"bottom line", as I call it-the 0th power, which needs no recursion at all. It suffices to
tell the genie that its value is 1. So when you look carefully, you see that this
recursive definition is no more circular than the "tower of power" was-and you can't
get any straighter than an infinite straight line! In fact, this one compact definition
really is just a way of getting the whole tower of power into one finite expression. Far
from being circular, it is just a handy summary of infinitely many different
definitions, all belonging to one family.




Lisp: Lists and Recursion                                                             412

       In case you still have a trace of skepticism about this sleight of hand, perhaps I
should let you watch what the Lisp genie will do if you ask for a "trace" of the
function, and then ask it once again to evaluate (power 9 3).

-> (power 9 3).
   ENTERING power (q=9, s=3)
    ENTERING power (q=9, s=2)
      ENTERING power (q=9, s=1)
       ENTERING power (q=9, s=0)
       EXITING power (value: 1)
     EXITING power (value: 9)
   EXITING power (value: 81)
EXITING power (value: 729)
->

On the lines marked ENTERING, the genie prints the values of the two arguments,
and on the lines marked EXITING, it prints the value it has computed and is
returning. For each ENTERING line there is of course an EXITING line, and the
two are aligned vertically-that is, they have the same amount of indentation.
        You can see that in order to figure out what (power 9 3) is, the genie must first
calculate (power 9 2). But this is not a given; instead it requires knowing the value of
(power 9 1), and this in turn requires (power 9 0). Ah! But we were given this one-it
is just 1. And now we can bounce back "up", remembering that in order to get one
answer from the "deeper" answer, we must multiply by 9. Hence we get 9, then 81,
then 729, and we are done.
        I say "we", but of course it is not we but the Lisp genie who must keep track of
these things. The Lisp genie has to be able to suspend one computation to work on
another one whose answer was requested by the first one. And the second
computation, too, may request the answer to a third one, thus putting itself on hold-as
may the third, and so on, recursively. But eventually, there will come a case where the
buck stops that is, where a process runs to completion and returns a value-and that
will enable other stacked-up processes to finally return values, like stacked-up
airplanes that have circled for hours finally getting to land, each landing opening up
the way for another landing.
        Ordinarily, the Lisp genie will. not print out a trace of what it is thinking
unless you ask for it. However, whether you ask to see it or not, this kind of thing is
going on behind the scenes whenever a function call is evaluated. One of the
enjoyable things about Lisp is that it can deal with such recursive definitions without
getting flustered.

                                       *   *   *




Lisp: Lists and Recursion                                                            413

        I am not so naive as to expect that you've now totally got the hang of recursion
and could go out and write huge recursive programs with the greatest of ease. Indeed,
recursion can be a remarkably subtle means of defining functions, and sometimes
even an expert can have trouble figuring out the meaning of a complicated recursive
definition. So I thought I'd give you some practice in working with recursion.
Let me give a simple example based on this silly riddle: "How do you make a pile of
13 stones?" Answer: "Put one stone on top of a pile of 12 stones." (Ask a silly
question and get an answer 12/13 as silly.) Suppose we want to make a Lisp function
that will give us not a pile of 13 stones, but a list consisting of 13 copies of the atom
stone-or in general, n copies of that atom. We can base our answer on the riddle's
silly-seeming yet correct recursive answer. The general notion is to build the answer
for n out of the answer for n's predecessor. Build how? Using the list-building
function cons, that's how. What's the embryonic case? That is, for which value of
• does this riddle present absolutely no problem at all? That's easy: when
• equals 0, our list should be empty, which means the answer is nil. We can now put
our observations together as follows:

(def bunch-of-stones (lambda (n
       (cond ((eq n 0) nil)
              (t (cons 'stone (bunch-of-stones (next-smaller n)))))))

Now let's watch the genie put together a very small bunch of stones (with trace on,
just for fun):

-> (bunch-of-stones 2)
       ENTERING bunch-of-stones (n=2)
              ENTERING bunch-of-stones (n=1)
                     ENTERING bunch-of-stones (n=0)
                     EXITING bunch-of-stones (value: nil)
              EXITING bunch-of-stones (value: (stone))
       EXITING bunch-of-stones (value: (stone stone))
(stone stone)
->

        This is what is called "consing up a list". Now let's try another one. This one is
an old chestnut of Lisp and indeed of recursion in general. Look at the definition and
see if you can figure out what it's supposed to do; then read on to see if you were
right.

-> (def wow (lambda (n)
        (cond ((eq n 0) 1)
               (t (times n (wow (subs n)))))))

Remember, subl means the same as next-smaller. For a lark, why don't




Lisp: Lists and Recursion                                                             414

you calculate the value of (wow 100)? (If you ate your mental Wheaties this morning,
try it in your head.)
         It happens that Lisp genies often mumble out loud while they are executing
wishes, and I just happen to have overheard this one as it was executing the wish
(wow 100). Its soliloquy ran something like this:

            Hmm ... (wow 100), eh? Well, 100 surely isn't equal to 0, so I guess the
   answer has to be 100 times what it would have been, had the problem been
   (wow 99). All rightie-now all I need to do is figure out what (wow 99) is. Oh,
   this is going to be a piece of cake! Let's see, is 99 equal to 0? No, seems not to
   be, so I guess the answer to this problem must be 99 times what the answer
   would have been, had the problem been (wow 98). Oh, this is going to be child's
   play! Let's see ...

At this point, the author, having some pressing business at the bank, had to leave the
happy genie, and did not again pass the spot until some milliseconds afterwards.
When he did so, the genie was just finishing up, saying:

       ... And now I just need to multiply that by 100, and I've got my final answer.
       Easy as pie! I believe it comes out to be 93326215443944152681699238
       00000000-if I'm not mistaken.

Is that the answer you got, dear reader? No? Ohhh, I see where you went wrong. It
was in your multiplication by 52. Go back and try it again from that point on, and be a
little more careful in adding those long columns up. I'm quite sure you'll get it right
this time.

                                        *   *   *

        This wow function is ordinarily called factorial,- n factorial is usually defined
to be the product of all the numbers from 1 through n. But a recursive definition looks
at things slightly differently: speaking recursively, n factorial is simply the product of
n and the previous factorial. It reduces the given problem to a simpler sort of the same
type. That simpler one will in turn be reduced, and so on down the line, until you
come to the simplest problem of that type, which I call the "embryonic case" or the
"bottom line". People often speak, in fact, of a recursion "bottoming out".
        A New Yorker cartoon from a few years back illustrates the concept perfectly.
It shows a fifty-ish man holding a photograph of himself roughly ten years earlier. In
that photograph, he is likewise holding a photograph of himself, ten years earlier than
that. And on it goes, until eventually it "bottoms out"-quite literally-in a photograph
of a bouncy baby boy in his birthday suit (bottom in the air). This idea of recursive
photos catching you as you grow up is quite appealing. I wish my parents had thought
of it!




Lisp: Lists and Recursion                                                             415

Contrast it with the more famous Morton Salt infinite regress, in which the Morton
Salt girl holds a box of Morton Salt with her picture on it-but as the girl in the picture
is no younger, there is no bottom line and the regress is endless, at least theoretically.
Incidentally, the Dutch cocoa called "Droste's" has a similar illustration on its boxes,
and very likely so do some ,other products.
        The recursive approach works when you have a family of related problems, at
least one of which is so simple that it can be answered immediately. This I call the
embryonic case. (In the factorial example, that's the (eq n 0) case, whose answer is 1.)
Each problem ("What is 100 factorial?", for instance) can be viewed as a particular
case of one general problem ("How do you calculate factorials?"). Recursion takes
advantage of the fact that the answers to various cases are related in some logical way
to each other. (For example, I could very easily tell you the value of 100 factorial if
only somebody would hand me the value of 99 factorial-all I need to do is multiply by
100.) You could say that the "Recursioneer's Motto" is: "Gee, I could solve this case if
only someone would magically hand me the answer to the case that's one step closer
to the embryonic case." Of course, this motto presumes that certain cases are, in some
sense, "nearer" to the embryonic case than others are-in fact, it presumes that there is
a natural pathway leading from any case through simpler cases all the way down to
the embryonic case, a pathway whose steps are clearly marked all along the way.
        As it turns out, this is a very reasonable assumption to make in all sorts of
circumstances. To spell out the exact nature of this recursion-guiding pathway, you
have to answer two Big Questions:

       (1) What is the embryonic case?
       (2) What is the relationship of a typical case to the next simpler case?

Now actually, both of these Big Questions break up into two subquestions (as befits
any self-respecting recursive question!), one concerning how you recognize where
you are or how to move, the other concerning what the answer is at any given stage.
Thus, spelled out more explicitly, our Big Questions are:

       (la) How can you know when you've reached the embryonic case? (lb) What is
the embryonic answer?

        (2a) From a typical case, how do you take exactly one step toward the
embryonic case?
        (2b) How do you build this case's answer out of the "magically given" answer
to the simpler case?
        ' Question (2a) concerns the nature of the descent toward the embryonic case,




Lisp: Lists and Recursion                                                             416

or bottom line. Question (2b) concerns the inverse aspect, namely, the ascent that
carries you back up from the bottom to the top level. In the case of the factorial, the
answers to the Big Questions are:

       (1a) The embryonic case occurs when the argument is 0.
       (lb) The embryonic answer is 1.

        (2a) Subtract 1 from the present argument.
        (2b) Multiply the "magic" answer by the present argument.
        Notice how the answers to these four questions are all neatly incorporated in
the recursive definition of wow.

                                          *   *   *

        Recursion relies on the assumption that sooner or later you will bottom out.
One way to be sure you'll bottom out is to have all the simplifying or "descending"
steps move in the same direction at the same rate, so that your pathway is quite
obviously linear. For instance, it's obvious that by subtracting 1 over and over again,
you will eventually reach 0, provided you started with a positive integer. Likewise, it's
obvious that by performing the list-shortening operation of cdr, you will eventually
reach nil, provided you started with a finite list. For this reason, recursions using
SUM or cdr to define their pathway of descent toward the bottom are commonplace.
I'll show a cdr-based recursion shortly, but first I want to show a funny numerical
recursion in which the pathway toward the embryonic case is anything but linear and
smooth. In fact, it is so much like a twisty mountain road that to describe it as moving
"towards the embryonic case" seems hardly accurate. And yet, just as mountain roads,
no matter how many hairpin turns they make, eventually do hit their destinations, so
does this path.
        Consider the famous "3n -f 1" problem, in which you start with any positive
integer, and if it is even, you halve it; otherwise, you multiply it by 3 and add 1. Let's
call the result of this operation on n (hotpo n) (standing for "half or triple plus one").
Here is a Lisp definition of hotpo:

(def hotpo (lambda (n)
       (cond ((even n) (half n))
              (t (addl (times 3 n))))))

This definition presumes that two other functions either have been or will be defined
elsewhere for the Lisp genie, namely even and half (addl and times being, as
mentioned earlier, intrinsic parts of Lisp). Here are the lacking definitions:

(def even (lambda (n) (eq (remainder n 2) 0)))
(def half (lambda (n) (quotient n 2)))




Lisp: Lists and Recursion                                                             417

        What do you think happens if you begin with some integer and perform hotpo
over and over again? Take 7, for instance, as your starting point. Before you do the
arithmetic, take a guess as to what sort of behavior might occur.
        As it turns out, the pathway followed is often surprisingly chaotic and bumpy.
For instance, if we begin with 7, the process leads us to 22, then 11, then 34, 17, 52,
26, 13, 40, 20, 10, 5, 16, 8, 4, 2, 1, 4, 2, 1, 4, 2, 1.... Note that we have wound up in a
short loop (a 3-cycle, in the terminology of Chapter 16). Suppose we therefore agree
that if we ever reach 1, we have "hit bottom" and may stop. You might well ask,
"Who says we will hit 1? Is there a guarantee?" (Again in the terminology of Chapter
16, we could ask, "Is the 1-4-2-1 cycle an attractor?") Indeed, before you try it out in a
number of cases, you have no particular reason to suspect that you will ever hit 1, let
alone always. (It would be very surprising if someone correctly anticipated what
would happen in the case of, say, 7 before trying it out.) However, numerical
experimentation reveals a remarkable reliability to the process; it seems that no matter
where you start, you always do enter the 1-4-2-1 cycle sooner or later. (Try starting
with 27 as seed if you want a real roller-coaster ride!)
        Can you write a recursive function to reveal the pathway followed from an
arbitrary starting point "down" to 1? Note that I say "down" advisedly, since many of
the steps are in fact up! Thus the pathway starting at 3 would be the list (3 10 5 16 8 4
2 1). In order to solve this puzzle, you need to go back and answer for yourself the
two Big Questions of Recursion, as they apply here. Note:

(cond ((not (want help)) (not (read further)))
       (t (read further)))

                                        *   *   *

        First-about the embryonic case. This is easy. It has already been defined as the
arrival at 1; and the embryonic, or simplest possible, answer is the list (1), a tiny but
valid pathway from 1 to 1.
        Second-about the more typical cases. What operation will carry us from
typical 7 one step closer to embryonic 1? Certainly not the subl operation. No-by
definition it's the function hotpo itself that brings you ever "nearer" to 1-even when it
carries you up! This teasing quality is of course the whole point of the example. What
about (2b)-how to recursively build a list documenting our wildly oscillating
pathway? Well, the pathway belonging to 7 is gotten by tacking (i.e., consing) 7 onto
the shorter pathway belonging to (hotpo 7), or 22. After all, 22 is one step closer to
being embryonic than 7 is!
        These answers enable us to write down the desired function definition, using
tato as our dummy variable (tato being a well-known acronym for




Lisp: Lists and Recursion                                                              418

tato (and tato only), which recursively expands to tato (and tato only) (and tato
(and tato only) only)-and so forth).

(def pathway-to-1 (lambda (tato)
       (cond ((eq tato 1) '(1))
              (t (cons tato (pathway-to-1 (hotpo tato)))))))

Look at the way the Lisp genie "thinks" (as revealed when the trace feature is on):

-> (pathway-to-i 3)
   ENTERING pathway-to-i (tato=3)
      ENTERING pathway-to-1 (tato=10)
           ENTERING pathway-to-1 (tato=5)
               ENTERING pathway-to-1 (tato=16)
                  ENTERING pathway-to-1 (tato=8)
                    ENTERING pathway-to-i (tato=4)
                       ENTERING pathway-tai (tato=2)
                         ENTERING pathway-to-1 (tato=1)
                         EXITING pathway-to-1 (value: (1))
                       EXITING pathway-to-1 (value: (2 1))
                    EXITING pathway-to-i (value: (4 2 1))
                  EXITING pathway-to-1 (value: (8 4 2 1))
              EXITING pathway-to-1 (value: (16 8 4 2 1))
          EXITING pathway-to-1 (value: (5 16 8 4 2 1))
       EXITING pathway-to-1 (value: (10 5 16 8 4 2 1))
   EXITING pathway-to-1 (value: (3 10 5 16 8 4 2 1))
(3 10 5 16 8 4 2 1)
->

        Notice the total regularity (the sideways `V' shape) of the left margin of the
trace diagram, despite the chaos of the numbers involved. Not all recursions are so
geometrically pretty, when traced. This is because some problems request more than
one subproblem to be solved. As a practical real-life example of such a problem,
consider how you might go about counting up all the unicorns in Europe. This is
certainly a nontrivial undertaking, yet there is an elegant recursive answer: Count up
all the unicorns in Portugal, also count up all the unicorns in the other 30-odd
countries of Europe, and finally add those two results together.
        Notice how this spawns two smaller unicorn-counting subproblems, which in
turn will spawn two subproblems each, and so on. Thus, how can one count all the
unicorns in Portugal? Easy: Add the number of unicorns in the Estremadura region to
the number of unicorns in the rest of Portugal! And how do you count up the unicorns
in Estremadura (not to mention those in the remaining regions of Portugal)? By
further breakup, of course.




Lisp: Lists and Recursion                                                             419

But what is the bottom line? Well, regions can be broken up into districts,

        districts into square kilometers, square kilometers into hectares, hectares into
square meters-and presumably we can handle each square meter without further
breakup.
        Although this may sound rather arduous, there really is no other way to
conduct a thorough census than to traverse every single part on every level of the full
structure that you have, no matter how giant it may be. There is a perfect Lisp
counterpart to this unicorn census: it is the problem of determining how many atoms
there are inside an arbitrary list. How can we write a Lisp function called atomcount
that will give us the answer 15 when it is shown the following strange-looking list
(which we'll call brahma)?

       (((ac ab cb) ac (ba be ac)) ab ((cb ca ba) cb (ac ab cb)))

One method, expressed recursively, is exactly parallel to that for ascertaining the
unicorn population of Europe. See if you can come up with it on your own.

                                       *   *   *

       The idea is this. We want to construct the answer-namely, 15-out of the
answers to simpler atom-counting problems. Well, it is obvious that one simpler
atom-counting problem than (atomcount brahma) is (atomcount (car brahma)).
Another one is (atomcount (cdr brahma)). The answers to these two problems are,
respectively, 7 and 8. Now clearly, 15 is made out of 7 and 8 by addition-which
makes sense, after all, since the total number of atoms must be the number in the car
plus the number in the cdr. There's nowhere else for any atoms to hide! Well, this
analysis gives us the following recursive definition, with s as the dummy variable:

(def atomcount (lambda (s)
       (plus (atomcount (car s)) (atomcount (cdr s)))))

        It looks very simple, but it has a couple of flaws. First, we have written the
recursive part of the definition, but we have utterly forgotten the other equally vital
half-the "bottom line". It reminds me of the Maryland judge I once read about in the
paper, who ruled: "A horse is a four-legged animal that is produced by two other
horses." This is a lovely definition, but where does it bottom out? Similarly for
atomcount. What is the simplest case, the embryonic case, of atomcount? Simple: It
is when we are asked to count the atoms in a single atom. The answer, in such a case,
is of course 1. But how can we know when we are looking at an atom? Fortunately,
Lisp has a built-in function called atom that returns t (meaning "true") whenever we
are looking at an atom, and nil otherwise. Thus (atom 'plop) returns t, while (atom
'(a b c)) returns nil. Using that, we can patch up our definition:




Lisp: Lists and Recursion                                                           420

(def atomcount (lambda (s)
       (cond ((atom s) 1)
              (t (plus (atomcount (car s)) (atomcount (cdr s)))))))

       Still, though, it is not quite right. If we ask the genie for atomcount of (a b c),
instead of getting 3 for an answer, we will get 4. Shocking! How come this happens?
Well, we can pin the problem down by trying an even simpler example: if we ask for
(atomcount '(a)), we find we get 2 instead of 1. Now the error should be clearer: 2 =1
+ 1, with 1 each coming from the car and cdr of (a). The car is the atom a which
indeed should be counted as 1, but the cdr is nil, which should not. So why does nil
give an atomcount of 1? Because nil is not only an empty list, it is also an atom! To
suppress this bad effect, we simply insert another cond clause at the very top:

(def atomcount (lambda (s)
       (cond ((null s) 0)
              ((atom s) 1)
              (t (plus (atomcount (car s)) (atomcount (cdr s)))))))

I wrote (null s), which is just another way of saying (eq s nil). In general, if you want
to determine whether the value of some expression is nil or not, you can use the in-
built function null, which returns t if yes, nil if no. Thus, for example, (null (null
nil)) evaluates to nil, since the inner function call evaluates to t, and t is not nil!
        Notice in this recursion that we have more than one type of embryonic case
(the null case and the atom case), and more than one way of descending toward the
embryonic case (via both car and cdr). Thus, our Big Questions can be revised a bit
further:

       (la) Is there just one embryonic case, or are there several, or even an
                infinite class of them?
       (1 b) How can you know when you've reached an embryonic case?
       (1c) What are the answers to the various embryonic cases?

       (2a) From a typical case, is there exactly one way to step toward an
              embryonic case, or are there various possibilities?
       (2b) From a typical case, how do you determine which of the various
              routes toward an embryonic case to take?
       (2c) How do you build this case's answer out of the "magically given"
              answers to one or more simpler cases?

      Now what happens when we trace our function as it counts the atoms in
brahma, our original target? The result is shown in Figure 18-1. Notice the more
complicated topography of this recursion, with its many ins and outs.




Lisp: Lists and Recursion                                                             421

Lisp: Lists and Recursion   422

         Whereas the previous 'V'-shaped recursion looked like a simple descent into a
 smooth-walled canyon and then a simple climb back up the other side, this recursion
  looks like a descent into a much craggier canyon, where on your way up and down
  each wall you encounter various "subcanyons" that you must treat in the same way-
 and who knows how many levels of such structure you will be called on to deal with
                                   in your exploration?
        Shapes with substructure that goes on indefinitely like that, never bottoming
out in ordinary curves, are called fractals. Their nature forms an important area of
inquiry in mathematics today. An excellent introduction can be found in Martin
Gardner's "Mathematical Games" for April, 1978, and a much fuller treatment in
Benoit Mandelbrot's splendid book The Fractal Geometry of Nature. For a dynamic
view of a few historically revolutionary fractals, there is Nelson Max's marvelous film
Space-Filling Curves, where so-called "pathological" shapes are constructed step by
step before your eyes, and their mathematical significance is geometrically presented.
Then, as eerie electronic music echoes all about, you start shrinking like Alice in
Wonderland-but unlike her, you can't stop, and as you shrink towards oblivion, you
get to see ever more microscopic views of the infinitely detailed fractal structures. It's
a great visual adventure, if you're willing to experience infinity-vertigo!

                                        *   *   *

        One of the most elegant recursions I know of originates with the famous disk-
moving puzzle known variously as "Lucas' Tower", the "Tower of Hanoi", and the
"Tower of Brahma". Apparently it was originated by the French mathematician
Edouard Lucas in the nineteenth century. The legend that is popularly attached to the
puzzle goes like this:
        In the great Temple of Brahma in Benares, on a brass plate beneath the dome
that marks the Center of the World, there are 64 disks of pure gold which the priests
carry one at a time between three diamond needles according to Brahma's immutable
law: No disk may be placed on a smaller disk. In the Beginning of the World, all 64
disks formed the Tower of Brahma on one needle. Now, however, the process of
transfer of the tower from one needle to another is in midcourse. When the last disk is
finally in place, once again forming the Tower of Brahma but on a different needle,
then will come the End
        of the World, and All will turn to dust.
        A picture of the puzzle is shown in Figure 18-2. In it, the three needles are
labeled a, b, and c.
        If you work at it, you certainly can discover the systematic method that the
priests must follow in order to get the disks from needle a to needle b. For only three
disks, for instance, it is very easy to write down the order in which the moves go:
                ab ac be ab ca cb ab




Lisp: Lists and Recursion                                                             423

FIGURE 18-2. The Tower of Brahma puzzle, with 64 disks to be transferred.
[Drawing by David Moser. ]

Here, the Lisp atom ab represents a jump from needle a to needle b. There is a
structure to what is going on, however, that is not revealed by a mere listing of such
atoms. It is better revealed if one groups the atoms as follows:

               ab ac be ab ca cb ab

         The first group accomplishes a transfer of a 2-tower from needle a to needle c,
thus freeing up the largest disk. Then the middle move, ab, picks up that big heavy
disk and carries it over from needle a to needle b. The final group is very much like
the initial group, in that it transfers the 2-tower back from needle c to needle b. Thus
the solution to moving three depends on being able to move two. Similarly, the
solution to moving 64 depends on being able to move 63. Enough said? Now try to
write a Lisp function that will give you a solution to the Tower of Brahma for n disks.
(You may prefer to label the three needles with digits rather than letters, so that moves
are represented by two-digit numbers such as 12.) I will present the solution in the
next column-unless, of course, the dedicated priests, working by day and by night to
bring about the end of the world, should chance to reach their cherished goal before
then ...




Lisp: Lists and Recursion                                                            424

                  Lisp: Recursion and Generality
                                      April, 1983




IN the preceding column, I described Edouard Lucas' Tower of Brahma puzzle, in
which the object is to transfer a tower of 64 gold disks from one diamond needle to
another, making use of a third needle on which disks can be placed temporarily. Disks
must be picked up and moved one at a time, the only other constraint being that no
disk may ever sit on a smaller one. The problem I posed for readers was to come up
with a recursive description, expressed as a Lisp function, of how to accomplish this
goal (and thereby end the world).
        I pointed out that the recursion is evident enough: to transfer 64 disks from
one needle to another (using a third), it suffices to know how to transfer 63 disks from
one needle to another (using a third). To recap it, the idea is this. Suppose that the 64-
disk tower of Brahma starts out on needle a. Figure 19-1 shows a schematic picture,
representing all 64 disks by a mere 4. First of all, using your presumed 63-disk-
moving ability, transfer 63 disks from needle a to needle c, using needle b as your
"helping needle". Figure 19-1b shows how the set-up now looks. (Note: In the figure,
4 plays the role of 64, so 3 plays the role of 63, but for some reason, 1 doesn't play the
role of 61. Isn't that peculiar?) All right. Now simply pick up the one remaining a-
disk-the biggest disk of all-and plunk it down on needle b, as is shown in Figure 19-
1c. Now you can see how easy it will be to finish up-simply re-exploit your 63-disk
ability so as to transfer that pile on c back to b, this time using a as the "helping
needle". Notice how in this maneuver, needle a plays the helping role that needle c
played in the previous 63-disk maneuver. Figure 19-1d shows the situation just a split
second before the last disk is put in place. Why not after it's in place? Simple:
Because the entire world then turns to dust, and it's too hard to draw dust.

                                        *   * *

Now someone might complain that I left out all the hard parts: "You just magically
assumed an ability to move 63 disks!" So it might seem, but there's nothing magical
about such an assumption. After all, to move 63, you




Lisp: Recursion and Generality                                                        425

FIGURE 19-1. A smaller Tower of Brahma puzzle. At the top, the starting position.
Below it is shown in an intermediate stage, in which a three-high pile has been
transferred from needle a to needle c. At this point, the biggest disk has become free,
and can be jumped to needle b. Then all that is left is to re-transfer the three-high pile
from c to b. When this is done, the world will end. Thus, the final picture shows an
artist's conception of the world a mere split-second before it all turns to dust.

merely need to know how to move 62. And to move 62, you merely need to know
how to move 61. On it goes down the line, until you "bottom out" at the "embryonic
case" of the Tower of Brahma puzzle, the 1-disk puzzle. Now, I'll admit that you have
to keep track of where you are in this process, and that may be a bit tedious-but that's
merely bookkeeping. In principle, you now could actually carry the whole process
out-if you were bent on seeing the world end!
        As our first approximation to a Lisp function, let's write an English description
of the method. Let's call the three needles sn, dn, and hn, standing for "source-
needle", "destination-needle", and "helping-needle". Here goes:

To move a tower of height n from sn to do making use of hn:
     If n = 1, then just carry that one disk directly from an to dn;
     otherwise, do the following three steps:




Lisp: Recursion and Generality                                                        426

       (1) move a tower of height n-1 from sin to hn making use of dn;
       (2) carry 1 disk from sn to dn;
       (3) move a tower of height n-1 from hn to do making use of sn.

Here, lines (1) and (3) are the two recursive calls; skirting paradox, they seem to call
upon the very ability they are helping to define. The saving feature is that they involve
n-1 disks instead of n. Note that in line (1), hn plays the "destination" role while dn
plays the "helper" role. And in (3), hn plays the "source" role while sn plays the
"helper" role. Since the whole thing is recursive, every needle will be switching hats
many times over during the course of the transfer. That's the beauty of this puzzle and
in a way it's the beauty of recursion.
        Now how do we make the transition from English to Lisp? It's quite simple:

(def move-tower (lambda (n sn do hn)
      (cond ((eq n 1) (carry-one-disk sn dn))
             (t (move-tower (sub1 n) sn hn dn)
                    (carry-one-disk sn dn)
                    (move-tower (subl n) hn do sn)))))

        Where are the Lisp equivalents of the English words "from", "to", and
"making use of"? They seem to have disappeared! So how can the genie know which
needle is to play which role at each stage? The answer is, this information is conveyed
positionally. There are, in this function definition, four parameters: one integer and
three "dummy needles". The first of these three is the source, the second the
destination, the third the helper. Thus in the initial list of parameters (following the
lambda) they are in the order sn dn hn. In the first recursive call, the Lisp translation
of line (1), they are in the order sn hn dn, showing how hn and dn have switched
hats. In the second recursive call, the Lisp translation of line (3), you can see that hn
and sn have switched hats.
        The point is that the atom names sn, dn, and hn carry no intrinsic meaning to
the genie. They could as well have been apple, banana, and cherry. Their meanings
are defined operationally, by the places where they appear in the various parts of the
function definition. Thus it would have been a gross blunder to have written, for
instance, (move-tower (subl n) sn do hn) as Lisp for line (1), because this contains
no indication that hn and dn must switch roles in that line.

                                       *   *   *

An important question remains. What happens when that friendly Lisp genie comes to
a line that says carry-one-disk? Does it suddenly zoom off




Lisp: Recursion and Generality                                                       427

to the Temple at Benares and literally heft a solid gold disk? Or, more prosaically,
does it pick up a plastic disk on a table and transfer it from one plastic needle to
another? In other words, does some physical action, rather than a mere computation,
take place?
        Well, in theory that is quite possible. In fact, even in practice the -execution of
a Lisp command could actually cause a mechanical arm to move to a specific
location, to pick up whatever its mechanical hand grasps there, to carry that object to
another specific location, and then to release it there. In these days of industrial
robots, there is nothing science-fictional about that. However, in the absence of a
mechanical arm and hand to move physical disks, what could it mean?
        One obvious and closely related possibility is to have there be a display of the
puzzle on a screen, and for carry-one-disk to cause a picture of a hand to move, to
appear to grasp a disk, pick it up, and replace it somewhere else. This would amount
to simulating the puzzle graphically, which can be done with varying degrees of
realism, as anyone who has seen science-fiction films using state-of-the-art computer
graphics techniques knows.
        However, suppose that we don't have fancy graphics hardware or software at
our disposition. Suppose that all we want is to create a printed recipe telling us how to
move our own soft, organic, human hands so as to solve the puzzle. Thus in the case
of a three-disk puzzle, the desired recipe might read ab ac be ab ca cb ab or equally
well 1213 23 12 3132 12. What could help us reach this humbler goal?
        If we had a program that moved an arm, we would be concerned not with the
value it returned, but with the patterned sequence of "side effects" it carried out. Here,
by contrast, we are most concerned with the value that our program is going to return-
a patterned list of atoms. The list for n = 3 has got to be built up from two lists for
n=2. This idea was shown at the end of last month's column:

                       ab ac be ab ca cb ab

        In Lisp, to set groups apart, rather than using wider spaces, we use
parentheses. Thus our goal for n = 3 might be to produce the sandwich-like list ((ab
ac bc) ab (ca cb ab)). One way to produce a list out of its components is to use cons
repeatedly. Thus if the values of the atoms apple, banana, and cherry are 1, 2, and 3,
respectively, then the value returned by (cons apple (cons banana (cons cherry
nil))) will be the list (1 2 3). However, there is a shorter way to get the same result,
namely, to write (list apple banana cherry). It returns the same value. Similarly, if
the atoms sn and dn are bound to a and b respectively, then execution of the
command (list sn dn) will return (a b). The function list is an unusual function in that
it takes any number of arguments at all-even none, so that (list) returns the value of
nil !
        Now let us tackle the problem of what value we want the function called `
carry-one-disk to return. It has two parameters that represent needles, and




Lisp: Recursion and Generality                                                         428

ideally we'd like it to return a single atom made out of those needles' names, such as
ab or 12. For the moment, it'll be easier if we assume that the needle names are the
numbers 1, 2, and 3. In this case, to make the number 12 out of 1 and 2, it suffices to
do a little bit of arithmetic: multiply the first by 10 and add on the second. Here is the
Lisp for that:

(def carry-one-disk (lambda (sn dn) (plus (times 10 sn) dn)))

On the other hand, if the needle names are nonnumeric atoms, then we can use a
standard Lisp function called concat, which takes the values of its arguments (any
number, as with list) and concatenates them all so as to make one big atom. Thus
(concat 'con 'cat 'e 'nate) returns the atom concatenate. In such a case, we could
write:

(def carry-one-disk (lambda (sn dn) (concat sn dn)))

Either way, we have solved the "bottom line" half of the move-tower problem.
        The other half of the problem is what the recursive part of move-tower will
return. Well, that is pretty simple. We simply would like it to return a sandwich-like
list in which the values of the two recursive calls flank the value of the single call to
carry-one-disk. So we can modify our previous recursive definition very slightly, by
adding the word list:

(def move-tower (lambda (n sn do hn)
      (cond ((eq n 1) (carry-one-disk sn dn))
             (t (list (move-tower (subl n) sn hn dn)
                       (carry-one-disk sn dn)
                       (move-tower (subl n) hn do sn))))))

Now let's conduct a little exchange with the Lisp genie:

-> (move-tower 4 'a 'b 'c)
(((ac ab cb) ac (ba be ac)) ab ((cb ca ba) cb (ac ab cb)))

Smashing! It actually works! Isn't that pretty? In last month's column, this list was
called brahma.

                                        *   *   *

Suppose we wished to suppress all the inner parentheses, so that just a long
uninterrupted sequence of atoms would be printed out. For instance, we would get (ac
ab cb ac ba be ac ab cb ca ba cb ac ab cb) instead of the intricacy of brahma. This
would be slightly less informative, but it would be more impressive in its opaqueness.
In this case, we would not want to




Lisp: Recursion and Generality                                                        429

use the function list to make our sandwich of three values, but would have to use
some other function that removed the parentheses from the two flanking recursive
values.
        This is a case where the Lisp function append comes in handy. It splices any
number of lists together, dropping their outermost parentheses as it does so. Thus
(append '(a (b)), ‘(c) nil '(d e)) yields the five-element list (a (b) c d e) rather than
the four-element list ((a (b)) (c) nil (d e)), which would be yielded if list rather than
append appeared in the function position. Using append and a slightly modified
version of carry-one-disk to work with it, we can formulate a final definition of
move-tower that does what we want:

(def move-tower (lambda (n sn do hn)
      (cond ((eq n 1) (carry-one-disk sn dn))
             (t (append (move-tower (subl n) sn hn dn)
                    (carry-one-disk sn dn)
                    (move-tower (subl n) hn do sn))))))

(def carry-one-disk (lambda (sn dn) (list (concat sn dn))))

To test this out, I asked the Lisp genie to solve a 9-high Tower of Brahma puzzle.
Here is what it shot back at me, virtually instantaneously:

 (ab c be ab ca cb ab ac be be ca be ab ac be ab ca cb ab ca be be ca cb ab ac be
ab ca cb ab ac be ba ca be ab ac be ba ca cb ab ca be ba ca be ab ac bc ab ca cb
ab ac bc ba ca bc ab ac bc ab ca cb ab ca bc ba ca cb ab ac be ab ca cb ab ca be
ba ca be ab ac be ba ca cb ab ca be ba ca cb ab ac bc ab ca cb ab ac bc ba ca bc
ab ac bc ab ca cb ab ca bc be ca cb ab ac be ab ca cb ab ac be be ca be ab ac be
ba ca cb ab ca be be ca be ab ac be ab ca cb ab ac be ba ca be ab ac be ba ca cb
ab ca be be ca cb ab ac bc ab ca cb ab ca bc be ca bc ab ac bc be ca cb ab ca bc be
ca bc ab ac be ab ca cb ab ac be be ca be ab ac be ab ca cb ab ca be be ca cb ab ac
be ab ca cb ab ac be ba ca be ab ac be ba ca cb ab ca be ba ca be ab ac be Ab ca
cb ab ac be ba ca be ab ac be ab ca cb ab ca be ba ca cb ab ac be ab ca cb ab ca
be be ca be ab ac be be ca cb ab ca be be ca cb ab ac be ab ca cb ab ac be ba ca be
ab ac be ab ca cb ab ca be be ca cb ab ac be ab ca cb' ab ca be ba ca be ab ac be
ba ca cb ab ca be ba ca be ab ac be ab ca cb ab ac be be ca be ab ac be be ca cb
ab ca be be ca cb ab ac be ab ca cb ab ca be ba ca be ab ac be ba ca cb ab ca be
be ca cb ab ac be ab ca cb ab ac be be ca be ab ac be ab ca cb ab ca be ba ca cb
ab ac be ab ca cb ab ac be be ca be ab ac be ba ca cb ab ca be ba ca be ab ac be
ab ca cb ab ac be ba ca be ab ac be ab ca cb ab ca be ba ca cb ab ac be ab ca cb
ab ca be ba ca be ab ac be ba ca cb ab ca be ba ca cb ab ac be ab ca cb ab ac be
be ca be ab ac be ab ca cb ab ca be be ca cb ab ac be ab ca cb ab)

Now that's the, kind of genie I like!




Lisp: Recursion and Generality                                                       430

                                         *   *   *

        Congratulations! You have just been through a rather sophisticated and brain-
taxing example of recursion. Now let us take a look at a recursion that offers us a
different kind of challenge. This recursion comes from an offhand remark I made last
column. I used the odd variable name tato, mentioning that it is a recursive acronym
standing for tato (and tato only). Using this fact you can expand tato any number of
times. The sole principle is that each occurrence of tato on a given level is replaced by
the two-part phrase tato (and tato only) to make the next level. Here is a short table:

n=0: tato
n=1: tato
       (and tato only)
n=2: tato
       (and tato only)
          (and tato (and tato only) only)
n=3: tato
       (and tato only)
         (and tato (and tato only) only)
           (and tato (and tato only) (and tato (and tato only) only) only)

        For us the challenge is to write a Lisp function that returns what tato becomes
after n recursive expansions, for any n. Irrelevant that for n much bigger than 3 the
whole thing gets ridiculously large. We're theoreticians!
        There is only one problem. Any Lisp function must return a single Lisp
structure (an atom or a list) as its value; however, the entries in our table do not satisfy
this criterion. For instance, the one for n =2 consists of one atom followed by two
lists. To fix this, we can turn each of the entries in the table into a list by enclosing it
in one outermost pair of parentheses. Now our goal is consistent with Lisp. How do
we attain it?
        Recursive thinking tells us that the bottom line, W embryonic case, occurs
when n=0, and that otherwise, the nth' line is made from the line before it by
replacing the atom tato, wherever it occurs, by the list (tato (and tato only)), only
without its outermost parentheses. We can write this up right away.

(def tato-expansion (lambda (n)
  (cond ((eq n 0) '(tato))      ,
        (t (replace 'tato '(tato (and tato only)) (tato-expansion (sub 1 n)))))))

The only thing is, we have not specified what we mean by replace. We must be very
careful in defining how to carry out our replace operation. Look at any. of the lines of
the tato table, and you will see that it contains




Lisp: Recursion and Generality                                                         431

one element more than the preceding line. Why is this? Because the atom tato gets
replaced each time by a two-element list whose parentheses, as I pointed out earlier,
are dropped during the act of replacement. It's this parenthesis-dropping that is the
sticky point. A less tricky example of such parenthesis-dropping replacement than the
recursive one involving tato would be this: (replace 'a '(1 2 3) '(a b a)), whose value
should be (12 3 b 1 2 3) rather than ((1 2 3) b (12 3)). Rather than exact substitution
of a list for an atom, this kind of replacement involves splicing or appending a list
inside a longer list.
        Let's try to specify in Lisp-using recursion, as usual just what we mean by
replace-ing all occurrences of the atom atm by a list called lyst, inside a long list
called longlist. This is a good puzzle for you to try. A hint: See how the answer for
argument (a b a) is built out of the answer for argument (b a). Also look at other
simple cases like that, moving back down toward the embryonic case.

                                        *   *   *

       The embryonic case occurs when longlist is nil. Then, of course, nothing
happens so our answer should be nil.
       The recursive case involves building a more complex answer from a simpler
one assumed given. We can fall back on our (a b a) example for this. We can build
the complex answer (12 3 b 1 2 3) out of the simpler answer (b 1 2 3) by appending
(12 3) onto it. On the other hand, we could consider (b 1 2 3) itself to be a complex
answer built from the simpler answer (1 2 3) by consing b onto it. Why does one
involve appending and the other involve consing? Simple: Because the first case
involves the atom a, which does get replaced, while the second involves the atom b,
which does not get replaced. This observation allows us to attempt to write down an
attempt at a recursive definition of replace, as follows:

(def replace (lambda (atm lyst longlist)
       (cond ((null longlist) nil)
               ((eq (car longlist) atm)
               (append lyst (replace atm lyst (cdr longlist))))
               (t (cons (car longlist) (replace atm lyst (cdr longlist)))))))

As you can see, there is an embryonic case (where longlist equals nil), and then one
recursive case featuring append and one recursive case featuring cons. Now let's try
out this definition on a new example.

-> (replace 'a '(1 2 3) '(a (a) b a))
(1 2 3 (a) b 1 2 3)
->

Whoops! It almost worked, except that one of the occurrences of a was completely
missed. This means that in our definition of replace, we must




Lisp: Recursion and Generality                                                     432

have overlooked some eventuality. Indeed, if you go back, you will see that an
unwarranted assumption slipped in right under our noses-namely, that the elements of
longlist are always atoms. We ignored the possibility that longlist might contain
sublists. And what to do in such a case? Answer: Do the replacement inside those
sublists as well. And inside sublists of sublists, too-and so on. Can you figure out a
way to fix up the ailing definition?

                                       *   *   *

We've seen a recursion before in which all parts on all levels of a structure needed to
be explored; it was the function atomcount last month, in which we did a
simultaneous recursion on both car and cdr. The recursive line ran (plus (atomcount
(car s)) (atomcount (cdr s))). Here it will be quite analogous. We'll have a recursive
line featuring two calls on replace, one involving the car of longlist and one involving
the cdr, instead of just one involving the cdr. And this makes perfect sense, once you
think about it. Suppose you wanted to replace all the unicorns in Europe by
porpuquines. One way to achieve this nefarious goal would be to split Europe into
two pieces: Portugal (Europe's car), and all the rest (its cdr). After replacing all the
unicorns in Portugal by porpuquines, and also all the unicorns in the rest of Europe by
porpuquines, finally you would recombine the two new pieces into a reunified Europe
(this is supposed to suggest a cons operation). Of course, to carry out this dastardly
operation on Portugal, an analogous splitting and rejoining would have to take place-
and so on. This suggests that our recursive line will look like this:

(cons (replace 'unicorn '(porpuquine) (car geographical-unit))
       (replace 'unicorn '(porpuquine) (cdr geographical-unit)))

or, more elegantly and more generally,

(cons (replace atm lyst (car longlist)) (replace atm lyst (cdr longlist)))

This cons line will cover the case where longlist's car is nonatomic, as well as the
case where it is atomic but not equal to atm. In order to make this work, we need to
augment the embryonic case slightly: we'll say that when longlist is not a list but an
atom, then replace has no effect on longlist at all. Conveniently, this subsumes the
earlier null line, so we can drop that one. If we put all this together, we come up with
a new, improved definition:

(def replace (lambda (atm lyst longlist)
       (cond ((atom longlist) longlist)
               ((eq (car longlist) atm)
               (append lyst (replace atm lyst (cdr longlist))))
               (t (cons (replace atm lyst (car longlist))
                       (replace atm lyst (cdr longlist)))))))




Lisp: Recursion and Generality                                                      433

Now when we say (tato-expansion 2) to the Lisp genie, it will print out for us the list
(tato (and tato only) (and tato (and tato only) only)).

                                       *   *   *

       Well, well. Isn't this a magnificent accomplishment? If it seems less than
magnificent, perhaps we can carry it a step further. A recursive acronymone
containing a letter standing for the acronym itself-can be amusing, but what of
mutually recursive acronyms? This could mean, for instance, two acronyms, each of
which contains a letter standing for the other acronym. An example would be the pair
of acronyms NOODLES and LINGUINI, standing for:

       NOODLES (oodles of delicious LINGUINI), elegantly served

                                     and

       luscious Itty-bitty NOODLES gotten usually in naples, Italy

respectively. Notice, incidentally, that NOODLES is not only indirectly but also
directly recursive. There's nothing wrong with that.
        In general, the notion of mutual recursion means a system of arbitrarily many
interwoven structures, each of which is defined in terms of one or more members of
the system (possibly including itself). If we are speaking of a family of mutually
recursive acronyms, then this means a collection of words, letters in any one of which
can stand for any word in the family.
I have to admit that this specific notion of mutually recursive acronyms is not
particularly useful in any practical sense. However, it is quite useful as a droll
example of a very common abstract phenomenon. Who has not at some time mused
about the inevitable circularity of dictionary definitions? Anyone can see that all
words eventually are defined in terms of some fundamental set that is not further
reducible, but simply goes round and round endlessly. You can amuse yourself by
looking up the definition of a common word in a dictionary and replacing the main
words in it by their definitions. I once carried this process out for "love" (defined as
"A strong affection for or attachment or devotion to a person or persons"), substituting
for "strong", "affection", "attachment", "devotion", and "person", and coming up with
this concoction:

  A morally powerful mental state or tendency, having strength of character or
  will for, or affectionate regard, or loyalty, faithfulness, or deep affection to, a
  human being -or beings, especially as distinguished from a thing or lower
  animal.

But not being satisfied with that, I carried the whole process one step further. This
was my result:




Lisp: Recursion and Generality                                                      434

  A set of circumstances or attributes characterizing a person or thing at a given
  time in, with, or by the conscious or unconscious together as a unit full of or
  having a specific ability or capacity in a manner relating to, dealing with, or
  capable of making the distinction between right and wrong in conduct, or an
  inclination to move or act in a particular direction or way, having the state or
  quality of being strong in moral strength, self-discipline, or fortitude, or the act
  or process of volition for, or consideration, attention, or concern full of fond or
  tender feeling for, or the quality, state, or instance of being faithful to, those
  persons or ideals that one is under obligation to defend or support, or the
  condition, quality or state of being worthy of trust, or a strongly felt fond or
  tender feeling to a creature or creatures of or characteristic of a person or
  persons, that lives or exists, or is assumed to do so, particularly as separated or
  marked off by differences from that which is conceived, spoken of, or referred
  to as existing as an individual entity, or from any living organism inferior in
  rank, dignity, or authority, typically capable of moving about but not of making
  its own food by photosynthesis.

Isn't it romantic? It certainly makes "love" ever more mysterious. Stuart Chase, in his
lucid classic on semantics, The Tyranny of Words, does a similar exercise for "mind"
and shows its opacity equally well. But of course concrete words as well as abstract
ones get caught in this vortex of confusion. My favorite example is one I discovered
while looking through a French dictionary many years ago. It defined the verb clocher
("to limp") as marcher en boitant ("to walk while hobbling", roughly), and boiler ("to
hobble") as clocher en marchant ("to limp while walking"). This eager learner of
French was helped precious little by that particular pair of definitions.

                                       *   *   *

But let us return to mutually recursive acronyms. I put quite a bit of effort into
working out a family of them, and to my surprise, they wound up dealing mostly
(though by no means exclusively!) with Italian food. It all began when, inspired by
tato, I chose the similar word tomato- and then decided to use its plural, coming up
with this meaning for tomatoes:

TOMATOES on MACARONI (and TOMATOES only), exquisitely SPICED.

The capitalized words here are those that are also acronyms. Here is the rest of my
mutually recursive family:

MACARONI:
    MACARONI and CHEESE (a REPAST of Naples, Italy)
REPAST:
    rather extraordinary PASTA and SAUCE, typical




Lisp: Recursion and Generality                                                       435

CHEESE:
      cheddar, havarti, emmenthaler (especially SHARP
      emmenthaler)
SHARP:
      strong, hearty, and rather pungent
SPICED:
      sweetly pickled in CHEESE ENDIVE dressing
ENDIVE:
      egg NOODLES, dipped in vinegar eggnog
NOODLES:
      NOODLES (oodles of delicious LINGUINI), elegantly served
LINGUINI:
      LAMBCHOPS (including NOODLES), gotten usually in Northern Italy
PASTA:
      PASTA and SAUCE (that's ALL!)
ALL!:
      a luscious lunch
SAUCE:
      SHAD and unusual COFFEE (eccellente!)
SHAD:
      SPAGHETTI, heated al dente
SPAGHETTI:
      standard PASTA, always good, hot especially (twist, then ingest)
COFFEE:
      choice of fine flavors, especially ESPRESSO
ESPRESSO:
      excellent, strong, powerful, rich, ESPRESSO, suppressing sleep
outrageously
BASTAI:
      belly all stuffed (tummy achel)
LAMBCHOPS:
      LASAGNE and meat balls, casually heaped onto PASTA SAUCE
LASAGNE:
      LINGUINI and SAUCE and GARLIC (NOODLES everywhere!)
RHUBARB:
      RAVIOLI, heated under butter and RHUBARB (BASTA!)




Lisp: Recursion and Generality                                    436

RAVIOLI:
     RIGATONI and vongole In oil, lavishly introduced
RIGATONI:
     rich Italian GNOCCHI and TOMATOES (or NOODLES
     instead)
GNOCCHI:
     GARLIC NOODLES over crisp CHEESE, heated Immediately
GARLIC:
     green and red LASAGNE in CHEESE

Any gourmet can see that little attempt has been made to have each term defined by
its corresponding phrase; it is simply associated more or less arbitrarily with the
phrase.
         Now what happens if we begin to expand some word-say, pasta? At first we
get simply PASTA and SAUCE (that's ALL!). The next stage yields
PASTA and SAUCE (that's ALL!) and SHAD and unusual COFFEE
(eccellente!) (that's a luscious lunch). We could obviously go on expanding
acronyms forever-or at least until we filled the universe up to its very brim with
mouth-watering descriptions of Italian food. But what if we were less ambitious, and
wanted merely to fill half a page or so with such a description? How might we find a
way to halt this seemingly bottomless recursion in midcourse?
         Well, of course, the key word here is "bottomless", and the answer it implies
is: Put in a mechanism to allow the recursion to bottom out. The bottomlessness
comes from the fact that at every stage, every acronym is allowed to expand, that is,
to spawn further acronyms. So what if, instead, we kept tight control of the spawning
process, being generous in the first few "generations" and gradually letting fewer and
fewer acronyms spawn progeny as the generations got later? This would be similar to
a redwood tree in a forest, which begins with a single "branch" (its trunk), and that
branch spawns "progeny", namely, the first generation of smaller branches, and they
in turn spawn ever more progeny-but eventually a natural "bottoming out" occurs as a
consequence of the fact that teeny twigs simply cannot branch further. (Somehow,
trees seem to have gotten their wires crossed, since for them, bottoming out generally
takes place at the top.)
         If this process were completely regular, then all redwood trees would look
exactly alike, and one could well agree with former Governor Reagan's memorable
dictum, "If you've seen one redwood tree, then you've seen them all." Unfortunately,
though, redwood trees (and some other things as well) are trickier than Governor
Reagan realized, and we have to learn to deal with a great variety of different things
that all go by the same name. The variety is caused by the introduction of randomness
into the choices as to whether to branch or not to branch, what angle to branch at,
what size branch to grow, and so on.




Lisp: Recursion and Generality                                                    437

        Similar remarks apply to the "trees" of mutually recursive acronyms in
expanding tomatoes we always made exactly the same control decisions about which
acronyms to expand when, then there would be one and only one type of rhubarb
expansion, so that here too, it would make sense to say
        "If you've seen one rhubarb, you've seen them all." But if we allow some
randomness to enter the decision-making about spawning, then we can get many
varieties of rhubarb, all bearing some telltale resemblance to one another, but at a
much more elusive level of perception.
        How can we do this? The ideal concept to bring to bear here is that of the
random-number generator, which serves as the computational equivalent of a coin flip
or throw of dice. We'll let all decisions about whether or not to expand a given
acronym depend on the outcome of such a virtual coin flip. At early stages of
expansion, we'll set things up so that the coin will be very likely to come up heads (do
expand); at later stages, it will be increasingly likely to come up tails (don't expand).
The Lisp function rand will be employed for this. It takes no arguments, and each
time it is called, it returns a new real number located somewhere between 0 and 1,
unpredictably. (This is an exaggeration-it is actually 100 percent predictable if you
know how it is computed; but since the algorithm is rather obscure, for most purposes
and to most observers the behavior of this function will be so erratic as to count as
totally random. The story, of random number generation is itself quite a fascinating
one, and would be an entire article in itself.)
        If we want an event to happen with a probability of 60 percent, first we ask
rand for a value. If that value turns out to be 0.6 or below, we go ahead, and if not,
we do not. Since over a long period of time, rand sprays its outputs uniformly over
the interval between 0 and 1, we will indeed get the go-ahead 60 percent of the time.

                                       *   *   *

        So much for random decisions. How will we get an acronym to expand when
told to? This is not too hard. Suppose we let each acronym be a Lisp function, as in
the following example: .

(def tomatoes (lambda 0
       '(tomatoes on macaroni (and tomatoes only), exquisitely spiced)))

        The function tomatoes takes no arguments, and simply returns the list of
words that it expands into. Nothing could be simpler.
        Now suppose we have a variable called acronym whose value is some
particular acronym-but we don't know which one. How could we get that acronym to
expand? The way we've set it up, that acronym must act as a function call. In order for
any atom to invoke a function, it must be the car of a list, as in the examples (plus 2
2), (rand), and (rhubarb). Now if we were




Lisp: Recursion and Generality                                                       438

to write (acronym), then the literal atom acronym would be taken by the genie as a
function name. But that would be a misunderstanding. It's certainly not the atom
acronym that we want to make serve as a function name, but its value, be it
macaroni, cheese, or what-have-you.
        To do this, we employ a little trick. If the value of the atom acronym is
rhubarb and if I write (list acronym), then the value the Lisp genie will return to me
will be the list (rhubarb). However, the genie will simply see this as an inert piece of
Lispstuff rather than as a little command that I would like to have executed. It cannot
read my mind. So how do I get it to perform the desired operation? Answer: I
remember the function called eval, which makes the genie look upon a given data
structure as a wish to be executed. In this case, I need merely say (eval (list
acronym)) and I will get the list (ravioli, heated under butter and rhubarb
(basta!)). And had acronym had a different value, then the genie would have handed
me a different list.
        We now have just about enough ideas to build a function capable of expanding
mutually recursive acronyms into long but finite phrases whose sizes and structures
are controlled by many "flips" of the rand coin. Instead of stepping you through the
construction of this function, I shall simply display it below, and let you peruse it. It is
modeled very closely on the earlier function replace.

(def expand (lambda (phrase probability)
       (cond ((atom phrase) phrase)
              ((is-acronym (car phrase))
              (cond ((Iessp (rand) probability)
                       (append
                       (expand (eval (list (car phrase))) (lower probability))
                       (expand (cdr phrase) probability)))
                    (t
                       (cons (car phrase) (expand (cdr phrase) probability)))))
              (t (cons (expand (car phrase) (lower probability))
                       (expand (cdr phrase) probability))))))

        Note that expand has two parameters. One represents the phrase to expand,
the other represents the probability of expanding any acronyms that are top-level
members of the given phrase. (Thus the value of the atom probability will always be
a real number between 0 and 1.) As in the redwood-tree example, the expansion
probability should decrease as the calls get increasingly recursive. That is why lines
that call for expansion of (car phrase) do so with a lowered probability. To be exact,
we can define the function lower as follows:

(def lower (lambda (x) (times x 0.8)))




Lisp: Recursion and Generality                                                         439

Thus each time an acronym expands, its progeny are only 0.8 times as likely to
expand as it was. This means that sufficiently deeply nested acronyms have a
vanishingly small probability of spawning further progeny. You could use any
reducing factor; there is nothing sacred about 0.8, except that it seems to yield pretty
good results for me.
         The only remaining undescribed function inside the definition above is is-
acronym. Its name is pretty self-explanatory. First the function tests to see if its
argument is an atom; if not, it returns nil. If its argument is an atom, it goes on to see
if that atom has a function definition-in particular, a definition with the form of an
acronym. If so, Is-acronym returns the value t; otherwise it returns nil. Precisely how
to accomplish this depends on your specific variety of Lisp, which is why I have not
shown it explicitly. In Franz Lisp, it is a one-liner.
         You may have noticed that there are two cond clauses in close proximity that
begin with t. How come one "otherwise" follows so closely on the heels of another
one? Well, actually they belong to different cond's, one nested inside the other. The
first t (belonging to the inner cond) applies to a case where we know we are dealing
with an acronym but where our random coin, instead of coming down heads, has
come down tails (which amounts to a decision not to expand); the second t (belonging
to the outer cond) applies to a case where we have discovered we are simply not
dealing with an acronym at all.
         The inner logic of expand, when scrutinized carefully, makes perfect sense.
On the other hand, no matter how carefully you scrutinize it, the output produced by
expand using this famiglia of acronyms remains quite silly. Here is an example:

(rich Italian green and red linguini and shad and unusual choice of fine flavors,
especially excellent, strong, powerful, rich, espresso, suppressing sleep
outrageously (eccellentel) and green and red lasagne in cheese (noodles
everywhere!) in cheddar, havarti, emmenthaler (especially sharp emmenthaler)
noodles (oodles of delicious linguini), elegantly served (oodles of delicious
linguini), elegantly served (oodles of delicious linguini and sauce and garlic
(noodles (oodles of delicious linguini), elegantly served everywhere!) and meat
balls, casually heaped onto pasta and sauce (that's sill) and sauce (that's a
luscious lunch) sauce (including noodles (oodles of delicious linguini), elegantly
served), gotten usually In Northern Italy), elegantly served over crisp cheese,
heated immediately and tomatoes on macaroni and cheese (a repast of Naples,
Italy) (and tomatoes only), exquisitely sweetly pickled in cheese endive dressing
(or noodles instead) and vongole in oil, lavishly Introduced, heated under butter
and rich Italian gnocchi and tomatoes (or noodles instead) and vongole in oil,
lavishly Introduced, heated under butter and rigatoni and vongole in oil, lavishly
Introduced, heated under butter and ravioli, heated under butter and rich
Italian garlic noodles over crisp cheese, heated Immediately and tomatoes (or
noodles Instead) and vongole In oil, lavishly Introduced, heated under butter and
ravioli, heated under butter and




Lisp: Recursion and Generality                                                        440

rhubarb (bastal) (basta! (bastal) (bastal) (belly all stuffed (tummy ache!))
(basta!)

      Oh, the glories of recursive spaghetti! As you can see, Lisp is hardly the
computer language to learn if you want to lose weight. Can you figure out which
acronym this gastronomical monstrosity grew out of?

                                       *   *   *

The expand function exploits one of the most powerful features of Lisp -that is, the
ability of a Lisp program to take data structures it has created and treat them as pieces
of code (that is, give them to the Lisp genie as commands). Here it was done in a most
rudimentary way. An atom was wrapped in parentheses and the resulting minuscule
list was then evaluated, or evaled, as Lispers' jargon has it. The work involved in
manufacturing the data structure was next to nothing, in this case, but in other cases
elaborate pieces of structure can be "consed up then handed to the Lisp genie for
evaling. Such pieces of code might be new function definitions, or any number of
other things. The main idea is that in Lisp, one has the ability to "elevate" an inert,
information-containing data structure to the level of "animate agent", where it
becomes a manipulator of inert structures itself. This program-data cycle, or loop, can
continue on and on, with structures reaching out, twisting back, and indirectly
modifying themselves or related structures.
        Certain types of inert, or passive, information-containing data structures are
sometimes referred to as declarative knowledge-"declarative" because they often have
a form abstractly resembling that of a declarative sentence, and "knowledge" because
they encode facts about the world in some way, accessible by looking in an index in
somewhat the way "book-learned" facts are accessible to a person. By contrast,
animate, or active, pieces of code are referred to as procedural knowledge-
"procedural" since they define sequences of actions ("procedures") that actually
manipulate data structures, and "knowledge" since they can be viewed as embodying
the program's set of skills, something like a human's unconscious skills that were once
learned through long, rote drill sessions. (Sometimes these contrasting knowledge
types are referred to as "knowledge that" and "knowledge how".)
        This distinction should remind biologists of that between genes--relatively
inert structures inside the cell-and enzymes, which are anything but inert. Enzymes
are the animate agents of the cell; they transform and manipulate all the inert
structures in indescribably sophisticated ways. Moreover, Lisp's loop of program and
data should remind biologists of the way that genes dictate the form of enzymes, and
enzymes manipulate genes (among other things). Thus Lisp's procedural-declarative
program-data loop provides a primitive but very useful and tangible example of one
of the most fundamental patterns at the base of life: the ability of passive structures




Lisp: Recursion and Generality                                                       441

to control their own destiny, by creating and regulating active structures whose form
they dictate.

                                        *   *   *

        We have been talking all along about the Lisp genie as a mysterious given
agent, without asking where it is to be found or what makes it work. It turns out that
one of Lisp's most exciting properties is the great ease with which one can describe
the Lisp genie's complete nature in Lisp itself. That is, the Lisp interpreter can be
easily written down in Lisp. Of course, if there is no prior Lisp interpreter to run it, it
might seem like an absurd and pointless exercise, a bit like having a description in
flowery English telling foreigners how best to learn English. But it is not so silly as
that makes it sound.
        In the first place, if you know enough English, you can "bootstrap" your way
further into English; there is a point beyond which explanations written in English
about English are indeed quite useful. What's more, that point is not too terribly far
beyond the beginning level. Therefore, all you need to acquire first, and
autonomously, is a "kernel"; then you can begin to lift yourself by your own
bootstraps. For children, it is an exciting thing when, in reading, they begin to learn
new phrases all by themselves, simply by running into them several times in
succession. Their vocabulary begins to grow by leaps and bounds. So it is once there
is a Lisp kernel in a system; the rest of the Lisp interpreter can be-and usually is-
written in Lisp.
        The fact that one can easily write the Lisp interpreter in Lisp is no mere fluke
depending on some peculiarly introverted fact about Lisp. The reason it is easy is that
Lisp lends itself to writing interpreters for all sorts of languages. This means that Lisp
can serve as a basis on which one can build other languages.
        To put it more vividly, suppose you have designed on paper a new language
called "Flumsy". If you really know how Flumsy should work, then it should not be
too hard for you to write an interpreter for it in Lisp. Once implemented, your Flumsy
interpreter then becomes, in essence, an intermediary genie to whom you can give
wishes in Flumsy and who will in turn communicate those wishes to the Lisp genie in
Lisp. Of course, all the mechanisms allowing the Flumsy genie to talk to the Lisp
genie are themselves being carried out by the Lisp genie. So is this a mere facade? Is
talking Flumsy really just a way of talking Lisp in disguise?
        Well, when the U.S. arms negotiators talk to their Soviet counterparts through
an interpreter, are they really just speaking Russian in disguise? Or is the crux of the
matter whether the interpreter's native language was English or Russian, upon which
the other was learned as a second tongue? And suppose you find out that in fact, the
interpreter's native language was Lithuanian, that she learned English only as an
adolescent and then learned Russian by taking high-school classes in English? Will
you then feel that when she speaks Russian, she is actually speaking English in
disguise, or worse, that she is actually speaking Lithuanian, doubly disguised?




Lisp: Recursion and Generality                                                         442

        Analogously, you might discover that the Lisp interpreter is in fact written in
Pascal or some other language. And then someone could strip off the Pascal facade as
well, revealing to you that the truth of the matter is that all instructions are really
being executed in machine language, so that you are fooling yourself completely if
you think that the machine is talking Flumsy, Lisp, Pascal, or any higher-level
language at all!

                                        *      *       *

        When one interpreter runs on top of another one, there is always the question
of what level one chooses not to look below. I personally seldom think about what
underlies the Lisp interpreter, so that when I am dealing with the Lisp system, I feel as
if I am talking to "someone" whose "native language" is Lisp. Similarly, when dealing
with people, I seldom think about what their brains are composed of; I don't reduce
them in my mind to piles of patterned neuron firings. It is natural to my perceptual
system to recognize them at a certain level and not to look below that level.
        If someone were to write a program that could deal in Chinese with simple
questions and answers about restaurant visits, and if that program were in turn written
in another language-say, the hypothetical language "SEARLE" (for "Simulated East-
Asian Restaurant-Lingo Expert"), I could choose to view the system either as
genuinely speaking Chinese (assuming it gave a creditable and not too slow
performance), or as genuinely speaking SEARLE. I can shift my point of view at will.
The one I adopt is governed mostly by pragmatic factors, such as which subject area I
am currently more interested in at the moment (Chinese restaurants, or how grammars
work), how closely matched the program's speed is to that of my own brain, and -not
least-whether I happen to be more fluent in Chinese or in SEARLE. If to me, Chinese
is a mere bunch of "squiggles and squoggles", I may opt for the SEARLE viewpoint;
if on the other hand, SEARLE is a mere bunch of confusing technical gibberish, I will
probably opt for the Chinese viewpoint. And if I find out that the SEARLE interpreter
is in turn implemented in the Flumsy language, whose interpreter is written in Lisp,
then I have two more points of view to choose from. And so on.
        With interpreters stacked on interpreters, however, things become rapidly very
inefficient. It is like running a motor off power created through a series of electric
generators, each one being run off power coming from the preceding one: one loses a
good deal at each stage. With generators there is usually no need for a long cascade,
but with interpreters it is often the only way to go. If there is no machine whose
machine language is Lisp, then you build a Lisp interpreter for whatever machine you
have available, and run Lisp that way. And Flumsy and SEARLE, if you wish to have
them at your disposal, are then built on top of this virtual Lisp machine. This
indirectness can be annoyingly inefficient, causing your new "virtual Flumsy
machine" or "virtual SEARLE machine" to run dozens of times more slowly than you
would like.




Lisp: Recursion and Generality                                                       443

                                        *   *   *

        Important hardware developments have taken place in the last several years,
and now machines are available that are based at the hardware level on Lisp. This
means that they "speak" Lisp in a somewhat deeper senselet us say, "more fluently"-
than virtual Lisp machines do. It also means that when you are on such a machine,
you are "swimming" in a Lisp environment. A Lisp environment goes considerably
beyond what I have described so far, for it is more than just a language for writing
programs. It includes an editing program, with which one can create and modify one's
programs (and text as well), a debugging program, with which one can easily localize
one's errors and correct them, and many other features, all designed to be compatible
with each other and with an overarching "Lisp philosophy".
        Such machines, although still expensive and somewhat experimental, are
rapidly becoming cheaper and more reliable. They are put out by various new
companies such as LMI (Lisp Machine, Inc.), Symbolics, Inc., both of Cambridge,
Massachusetts, and older companies such as Xerox. Lisp is also available on most
personal computers-you need merely look at any issue of any of the many small-
computer magazines to find ads for Lisp.
        Why, in conclusion, is Lisp popular in artificial intelligence? There is no
single answer, nor even a simple answer. Here is an attempt at a summary:

   (1) Lisp is elegant and simple.
   (2) Lisp is centered on the idea of lists and their manipulation-and lists are
      extremely flexible, fluid data structures.
   (3) Lisp code can easily be manufactured in Lisp, and run.
   (4) Interpreters for new languages can easily be built and experimented with in
      Lisp.
   (5) "Swimming" in a Lisp-like environment feels natural for many people.
   (6) The "recursive spirit" permeates Lisp.

       Perhaps it is this last rather intangible statement that gets best at it. For some
reason, many people in artificial intelligence seem to have a deep sense that
recursivity, in some form or other, is connected with the "trick" of intelligence. This is
a hunch, an intuition, a vaguely felt and slightly mystical belief, one that I certainly
share-but whether it will pay off in the long run remains to be seen.

Post Scriptum.
       In March of 1977, 1 met the great AI pioneer Marvin Minsky for the first time.
It was an unforgettable experience. One of the most memorable remarks he made to
me was this one: "Gödel should just have thought up




Lisp: Recursion and Generality                                                        444

Lisp; it would have made the proof of his theorem much easier." I knew exactly what
Minsky meant by that, I could see a grain of truth in it, and moreover I knew it had
been made with tongue semi in cheek. Still, something about this remark drove me
crazy. It made me itch to say a million things at once, and thus left me practically
speechless. Finally today, after my seven-year itch, I will say some of the things I
would have loved to say then.
        What Minsky meant, paraphrased, is this: "Probably the hardest part of
Godel's proof was to figure out how to get a mathematical system to talk about itself.
This took several strokes of genius. But Lisp can talk about itself, at least in the sense
Gödel needed, directly. So why didn't he just invent Lisp? Then the rest would have
been a piece of cake." An obvious retort is that to invent Lisp out of the blue would
have taken a larger number of strokes of genius. Minsky, of course, knew this, and at
bottom, his remark was clearly just a way of making this very point in a facetious
way.
        Still, it was clear that Minsky felt there was some serious content to the
remark, as well. (And I have heard him make the same remark since then, so I know it
was not just a throwaway quip.) There was the implicit question, "Why didn't Gödel
invent the idea of direct self-reference, as in Lisp?" And this, it seemed to me, missed
a crucial point about Gödel’s work, which is that it showed that self-reference can
crop up even where it is totally unexpected and unwelcome. The power of Gödel’s
result was that it obliterated the hopes for completeness of an already known system,
namely Russell and Whitehead's Principia Mathematica; to have destroyed similar
hopes for some newly concocted system, Lisp-like or not, would have been far less
significant (or, to be more accurate, such a result's significance would have been far
harder for people to grasp, even if it were equally significant).
        Moreover, Godel's construction revealed in a crystal-clear way that the line
between "direct" and "indirect" self-reference (indeed, between direct and indirect
reference, and that's even more important!) is completely blurry, because his
construction pinpoints the essential role played by isomorphism (another name for
coding) in the establishment of reference and meaning. Gödel’s work is, to me, the
most beautiful possible demonstration of how meaning emerges from and only from
isomorphism, and of how any notion of "direct" meaning (i.e., codeless meaning) is
incoherent. In brief, it shows that semantics is an emergent quality of complex
syntax, which harks back to my earlier remark in the Post Scriptum to Chapter 1,
namely: "Content is fancy form." So the serious question implicit in Minsky's joke
seemed to me to rest on a confusion about this aspect of the nature of meaning.

                                            *   *   *

        Now let me explain this in more detail. Part I of Gödel’s insight was to realize
that via a code, a number can represent a mathematical symbol 11 e.g., the integer
eleven can represent the left parenthesis, and the integer




Lisp: Recursion and Generality                                                        445

thirteen the right parenthesis 13. The analogue of this in human languages is the
recognition that certain orally produced screeches or manually produced scratches
(such as "word", "say", "language", "sentence", "reference", "grammar", "meaning",
and so on) can stand for elements of language itself (as distinguished from screeches
or scratches such as "cow" and "splash", which stand for extralinguistic parts of the
universe). Here we have pieces of language talking about language. Maybe it doesn't
seem strange to you, but put yourself, if you can, back in the shoes of barefoot cave
people who had barely gotten out of the grunt stage. How amazingly magical it must
have felt to the beings in whose minds such powerful concepts as words about words
first sparked! In some sense, human consciousness began then and there.
        But a language can't get very far in terms of self-reference if it can talk only
about isolated symbols. Part II of Gödel’s insight was to figure out how the system
(and here I mean Principia Mathematica "and", as Gödel’s paper's title says, "related
systems") could talk about lists of symbols, and even lists of lists of symbols, and so
on. In the analogy to human language, making this step is like the jump from an
ability to talk about people's one-word utterances ("Paul Revere said `Land!'.") to the
ability to talk about arbitrarily long utterances, and nested ones at that ("`Douglas
Hofstadter wrote, "Paul Revere said, `Land!'."."').
        Gödel found a way to have some integers stand for single symbols and others
stand for lists of symbols, usually called strings. An example will help. Suppose the
integer 1 stands for the symbol '0', and as I mentioned earlier, 11 and 13 for
parentheses. Thus to encode the string "(0)" would require you to combine the
integers 11, 1, and 13 somehow into a single integer. Gödel chose the integer
7500000000000-not capriciously, of course! This integer can be viewed as the
product of the three integers 2048, 3, and 1220703125, which in turn are,
respectively: 2", 3', and 5's. In other words, the three-symbol string whose symbols
individually are coded for by 11 and 1 and 13 is coded for in toto by the single integer
2"3'5'3. Now 2, 3, and 5 are of course the first three primes, and if you want to encode
a longer string, you use as many primes as you need, in increasing order. This simple
scheme allows you to code strings of arbitrary length into large integers, and
moreover-since large integers can be exponents just as easily as small ones can-it
allows for recursive coding. In other words, strings can contain the integer codes for
other strings, and this can go on indefinitely. An example: the list of strings "0", "(0)",
and "((0))" is coded into the stupendously large integer


       The proverbial "astute reader" might well have noticed a possible ambiguity:
How can you tell if an integer is to be decomposed via prime factorization into other
integers or to be left alone and interpreted as a code




Lisp: Recursion and Generality                                                         446

for an atomic symbol? Gödel’s simple but ingenious solution was to have all atomic
symbols be represented by odd integers. How does that solve the matter? Easy: You
know you should not factorize odd integers, and conversely, you should factorize
even ones, and then do the same with all the exponents you get when you do so.
Eventually, you will bottom out in a bunch of odd integers representing atomic
symbols, and you will know which ones are grouped together to form larger chunks
and how those chunks are nested.
        With this beautifully polished scheme for encoding strings inside integers and
thereby inside mathematical systems, Gödel had discovered a way of getting such a
system to talk-in code-about itself. He had snuck self-reference into systems that were
presumed to be as incapable of self-reference as are pencils of writing on themselves
or garbage cans of containing themselves, and he had done so as wittily as the Greeks
snuck a boatload of unacceptable soldiers into Troy, "encoded" as one single large
acceptable structure.
        Historically, the importance of Godel's work was that it revealed a plethora of
unexpected self-references (via his code, to be sure, but that fact in no way diminishes
their effect) within the supposedly impregnable walls of Russell and Whitehead's
Troy, Principia Mathematica. Now in Lisp, it's possible to construct and manipulate
pieces of Lisp programs. The idea of quoted code is one of those deep ideas that make
Lisp so appealing to Al people. Okay, but-when you have a system constructed
expressly to have self-referential potential, the fact that it has self-referential
structures will amaze no one. What is amazing and wonderful is when self-reference
pops up inside the very fortress constructed expressly to keep it out! The
repercussions of that are enormous.
        One of the clear consequences of Gödel’s revelation of this self-referential
potential inside mathematical systems was that the same potential exists within any
similar formalism, including computer languages. That is simply because computers
can do all the standard arithmetic operations-at least in theory-with integers of
unlimited size, and so coded representations of programs are being manipulated any
time you are manipulating certain large integers. Of course, which program is being
manipulated depends on what code you use. It was only after Gödel’s work had been
absorbed by a couple of generations of mathematicians, logicians, and computer
people that the desirability of inserting the concept of quotation directly into a formal
language became obvious. To be quite emphatic about it, however, this does not
enhance the language's potential in any way, except that it makes certain constructions
easier and more transparent. It was for this reason of transparency that Minsky made
his remark.
        Oh yes, I agree, Gödel’s proof would have been easier, but by the time Gödel
dreamt it up, it would have long since been discovered (and called "Snoddberger's
proof") had Gödel been in a mindset where inventing Lisp




Lisp: Recursion and Generality                                                       447

was natural. I'm all for counterfactuals, but I think one should be careful to slip things
realistically.

                                        *   *   *

         After this diatribe, you will think I am crazy when I turn around and tell you:
Gödel did invent Lisp! I am not trying to take anything away from John McCarthy,
but if you look carefully at what Gödel did in his 1931 article, written roughly 15
years before the birth of computers and 27 years before the official birth of Lisp, you
will see that he anticipated nearly all the central ideas of Lisp. We have already been
through the fact that the central data structure of Lisp, the list, was at the core of
Gödel’s work. The crucial need to be able to distinguish between atoms and lists-
something that modern-day implementors of Lisp systems have to worry about-was
recognized and cleverly resolved by Gödel, in his odd-even distinction. The idea of
quoting is, in essence, that of the Gödel code. And finally, what about recursive
functions, the heart and soul of Lisp programming technique? That idea, too, is an
indispensable part of Gödel’s paper! This is the astounding truth.
         The heart of Gödel’s construction is a chain of 46 definitions of functions,
each new one building on previous ones in a dizzying spire ascending toward one
celestial goal: the definition of one very complex function of a single integer, which,
given any input, either returns the value 1 or goes into an infinite loop. It returns 1
whenever the input integer is the Gödel number -the code number-of a theorem of
Principia Mathematica, and it loops otherwise. This is Gödel’s 46th function, which
he named "Bew", short for "beweisbar", meaning "provable" in German.
         If we could calculate the value of function 46 swiftly for any input, it would
resolve for us any true-false question of mathematics that a full axiomatic system
could resolve. All we would need to do is to write down the statement in question in
the language of Principia Mathematica, code the resulting formula into its Godel
number (the most mechanical of tasks), and then call function 46 on that number. A
result of 1 means true, looping forever means false. Do I hear the astute reader
protesting again? All right, then: If we want to avoid any chance of having to wait
forever to find out the answer, we can encode the negation of the statement in
question into its Gödel number as well, and also call function 46 on this second
number. We'll let the two calculations proceed simultaneously, and see which one
says `1'. Now, as long as Principia Mathematica has either the statement or its
negation as a theorem, one of the two calls on function 46 will return 1, while the
other will presumably spin on unto eternity.
         How does function 46 work? Oh, easy-by calling function 45 many times. And
how does function 45 work? Well, it calls functions 44 and others, and they call
previously defined functions, some of which call themselves (recursion!), and so on
and so forth-all of these complex calls eventually bottoming out in calls to absolutely
trivial functions such as "S", the




Lisp: Recursion and Generality                                                        448

successor function, which returns the value 18 when you feed it the integer 17. In
short, the evaluation of a high-numbered function in Gödel’s paper sets in motion the
calling of one subroutine after another in a hierarchical chain of command, in
precisely the same way as my function expand called numerous lower-level functions
which called others, and so on. Gödel’s remarkable series of 46 function definitions
is, in my book, the world's first serious computer program-and it is in Lisp. (The
Norwegian mathematician Thoralf Skolem was the inventor of the study of recursive
functions theoretically, but Gödel was the first person to use recursive functions
practically, to build up functions of great complexity.)
        It was for all these reasons that Minsky's pseudo joke struck my intellectual
funnybone. One answer I wanted to make was: "I disagree: Gödel shouldn't have and
couldn't have invented Lisp, because his work was a necessary precursor to the
invention of Lisp, and anyway, he was out to destroy PM, not Lisp." Another answer
was: "Your innuendo is wrong, because any type of reference has to be grounded in a
code, and Godel's way of doing it involved no more coding machinery grounding its
referential capacity than any other efficient way would have." The final answer I
badly wanted to blurt out was: "No, no, no-Gödel did invent Lisp!" You see why I
was tongue-tied?

                                        *      *       *

        One reason I mention all this about Gödel is that I wish to make some
historical and philosophical points. There is another reason, however, and that is to
point out that the ideas in Lisp are intimately related to the basic questions of
metamathematics and metalogic, and these, translated into a more machine-oriented
perspective, are none other than the basic questions of computability-perhaps the
deepest questions of computer science. Michael Levin has even written an
introduction to mathematical logic using Lisp, rather than a more traditional system,
as its underlying formal system. For this type of reason, Lisp occupies a very special
place inside computer science, and is not likely to go away for a very long time.
        However ... (you were waiting for this, weren't you?), there is a vast gulf
between the issues that Lisp makes clear and easy, and the issues that confront one
who would try to understand and model the human mind. The way I see it is in terms
of grain size. To me, the thought that Lisp itself might be "more conducive" to good
Al ideas than any other computer language is quite preposterous. In fact, such claims
remind me of certain wonderfully romantic but woefully antic claims I have heard
about the Hopi language. The typical one runs something like this: "Einstein should
just have invented Hopi; then the discovery of his theory of relativity would have
been much easier." The basis for this viewpoint is that Hopi, it is said, lacks terms for
"absolute time" in it. Supposedly, Hopi (or a language with similar properties) would
therefore be the ideal language in which to speak of relativity, since absolute time is
abandoned in relativity.




Lisp: Recursion and Generality                                                       449

        This kind of claim was first put forth by the outstanding American linguist
Edward Sapir, was later polished by his student Benjamin Whorf, and is usually
known these days as the Sapir-Whorf hypothesis. (I have already made reference to
this view in Chapters 7 and 8 on sexist language and imagery.) To state the Sapir-
Whorf thesis explicitly: Language controls thought. A milder version of it would say:
Language exerts a powerful influence upon thought.
        In the case of computer languages, the Sapir-Whorf thesis would have to be
interpreted as asserting that programmers in language X can think only in terms that
language X furnishes them, and no others. Therefore, they are strapped in to certain
ways of seeing the "world", and are prevented from seeing many ideas that
programmers in language L can easily see. At least this is what Sapir-Whorf would
have you believe. I will have none of it!
        I do use Lisp, I do think it is very convenient and natural in many ways, I do
advocate that anyone seriously interested in AI learn Lisp well; all this is true, but I do
not think that deep study of Lisp is the royal road to AI any more than I think that
deep study of bricks is the royal road to understanding architecture. Indeed, I would
suggest that the raw materials to be found in Lisp are to AI what raw materials are to
architecture: convenient building blocks out of which far larger structures are made.
        It would be ridiculous for anyone to hope to acquire a deep understanding of
what Al is all about without first having a clear, precise understanding of what
computers are all about. I know of no shorter cut to that latter goal than the study of
Lisp, and that is one reason Lisp is so good for AI students. Beginners in Lisp
encounter, and are in a good position to understand, fundamental issues in computer
science that even some advanced programmers in other languages may not have
encountered or thought about. Such concepts as lists, recursion, side effects, quoting
and evaluating pieces of code, and many others that I did not have the space to present
in my three columns, are truly central to the understanding of the potential of
computing machinery. Moreover, without languages that allow people to deal with
such concepts directly, it would be next to impossible to make programs of subtlety,
grace, and multi-level complexity. Therefore I advocate Lisp very strongly.
        It would similarly be next to impossible to build structures of subtlety, grace,
and multi-level complexity such as the Golden Gate Bridge and the Empire State
Building out of bricks or stone. Until the use of steel as an architectural substrate was
established, such constructions were unthinkable. Now we are in a position to erect
buildings that use steel in even more sophisticated ways. But steel itself is not the
source of great architects' inspiration; it is simply a liberator. Being a super-expert on
steel may be of some use to an architect, but I would guess that being quite
knowledgeable will suffice. After all, buildings are not just scaled-up girders. And so
it is with Lisp and Al. Lisp is not the "language of thought" or the `language of the
brain"-not by any stretch of the imagination. Lisp is,




Lisp: Recursion and Generality                                                         450

however, a liberator. Being a super-expert on Lisp may be of some use to a person
interested in computer models of mentality, but being quite knowledgeable will
suffice. After all, minds are not just scaled-up Lisp functions.
        Let me switch analogies. Is it possible for a novelist to conceive of plot ideas,
characters, intrigues, emotions, and so on, without being channelled by her own or his
own native language? Are the events that take place in, say, Anna Karenina
specifically determined by the nature of the Russian language and the words that it
furnished to Tolstoy? If that were the case, then of course the novel would be
incomprehensible to people who do not know the Russian language. It would simply
make no sense at all. But that is not even remotely the case. English-language readers
have read that novel with great pleasure and have just as fully fathomed its
psychological twists and turns as have Russian-language readers. The reason is that
Tolstoy's mind was concerned with concepts that float far above the grain size of any
human language. To think otherwise is to reduce Tolstoy to a mere syntactician, is to
see Tolstoy as pushed around by low-level quirks and local flukes of his own
language.
        Now please understand, I am not by any means asserting that Tolstoy
transcended his own culture and times; certainly he belongs to a particular era and a
particular set of circumstances, and those facts flavor what he wrote. But "flavor" is
the right word here. The essence of what he did-the meat of it, to prolong the "flavor"
metaphor-is universal, and has to do with the fact that Tolstoy had profoundly
experienced the human condition, had felt the pangs of many conflicting emotions in
all sorts of ways. That's where the power of his writing comes from, not from the
language he happened to inherit (otherwise, why wouldn't all Russians be great
novelists?); that's why his novels survive translation not only into other languages (so
they reach other cultures), but also into other eras, with different sensibilities. If
Tolstoy manages to reach further into the human psyche than most other writers do, it
is not the Russian language that deserves the credit, but Tolstoy's acute sensitivity and
empathy for people.
        The analogous statement could be made about AI programs and AI
researchers. One could even mechanically substitute "AI program" for "novel", "Lisp"
for "Russian", and-well, I have to admit that I would be hard pressed to come up with
"the Tolstoy of Al". Oh, well. My point is simply that good Al researchers are not in
any sense slaves to any language. Their ideas are as far removed from Lisp (or
whatever language they program in, be it Lisp itself, a "super-Lisp" (such as n-Lisp
for any value of n), Prolog, Smalltalk, and so on) as they are from English or from
their favorite computer's hardware design. As an example, the AI program that has
probably inspired me more than any other is the one called Hearsay-II, a speech-
understanding program developed at Carnegie-Mellon University in the mid-1970's
by a team headed up by D. Raj Reddy and Lee Erman. That program was written not
in Lisp but in a language called SAIL, very




Lisp: Recursion and Generality                                                       451

different in spirit from Lisp. Nonetheless, it could easily have been written in Lisp.
The reason it doesn't matter is simply that the scientific questions of how the mind
works are on a totally different level from the statements of any computer language.
The ideas transcend the language.

                                        *   *   *

        To some, I may seem here to be flirting dangerously with an anti-mechanistic
mysticism, but I hasten to say that that is far from the case. Quite the contrary. Still, I
can see why people might at first suspect me of putting forth such a view. A
programmer's instinct says that you can cumulatively build a system, encapsulating all
the complexity of one layer into a few functions, then building the next layer up by
exploiting the efficient and compact functions defined in the preceding layer. This
hierarchical mode of buildup would seem to allow you to make arbitrarily complex
actions be represented at the top level by very simple function calls. In other words,
the functions at the top of the pyramid are like "cognitive" events, and as you move
down the hierarchy of lower-level functions, you get increasingly many ever-dumber
subroutines, until you bottom out in a myriad calls on trivial, "subcognitive" ones. All
this sounds very biological -even tantalizingly close to being an entire resolution of
the mind-brain problem. In fact, for a clear spelling-out of just that position, see
Daniel Dennett's book Brainstorms, or perhaps worse, see parts of my own Gödel,
Escher, Bach !
        Yes, although I don't like to admit it, I too have been seduced by this recursive
vision of mechanical mentality, resembling nothing so much as an army, with its
millions of unconscious robot privates carrying out the desires of its top-level
cognitive general, as conveyed to them by large numbers of obedient and semi-bright
intermediaries. Probably my own strongest espousal of this view is found in Chapter
X of GEB, where a subheading blared out, loud and clear, "AI Advances Are
Language Advances". I was arguing there, in essence, for the orthodox Al position
that if we could just find the right "superlanguage"-a language presumably several
levels above Lisp, but built on top of it as Flumsy or SEARLE are built on top of it-
then all would be peaches and cream. We would be able to program in the legendary
"language of thought". AI programs would practically write themselves. Why, there
would be so much intelligence in the language itself that we could just sit back and
give the sketchiest of hints, and the computer would go off and do our tacit bidding!
        This, in the opinion of my current self, is a crazy vision, and my reasons for
thinking so are presented in Chapter 26, "Waking Up from the Boolean Dream". I am
relieved that I spent a lot more time in GEB knocking down that orthodox vision of Al
rather than propping it up. I argued then, as I still do now, that the top-level behavior
of the overall system must emerge statistically from a myriad independent pieces,
whose actions are almost as likely to cancel each other out as to be in phase with each
other. This picture,




Lisp: Recursion and Generality                                                         452

most forcefully presented in the dialogue Prelude ... Ant Fugue and surrounding
chapters, was the dominant view in that book, and it continues to be my view. In this
book, it is put forth in Chapters 25 and 26.
        In anticipation of those chapters, you might just ponder the following question.
Why is it the case that, after all these millennia of using language, we haven't
developed a single word for common remarks such as "Could you come over here and
take a look at this?" Why isn't that thought expressed by some minuscule
epigrammatic utterance such as "Cycohatalat"? Why are we not building new layer
upon new layer of sophistication, so that each new generation can say things that no
previous generation could have even conceived of? Of course, in some domains we
are doing just that. The phrase "Lisp interpreter" is one that requires a great deal of
spelling out for novices, but once it is understood, it is a very useful shorthand for
getting across an extremely powerful set of ideas. All through science and other
aspects of life, we are adding words and phrases.. Acronyms such as "radar", "laser",
"NATO", and "OPEC", as well as sets of initials such as "NYC", "ICBM", "MIT",
"DNA", and "PC", are all very common and very wordlike.
        Indeed, language does grow, but nonetheless, despite what might be
considered an exponential explosion in the number of terms we have at our disposal,
nobody writes a novel in one word. Nobody even writes a novel in a hundred words.
As a matter of fact, novels these days are no shorter than they were 200 years ago.
Nor are textbooks getting shorter. No matter how big an idea can be packed into a
single word, the ideas that people want to put into novels and textbooks are on a
totally different scale from that. Obviously, this is not claiming that language cannot
express the ideas of a novel; it is simply saying that it takes a heap o' language to do
so, no matter how the language is built. That is the issue of grain size that I alluded to
before, and I feel that it is a deep and subtle issue that will come up more often as
theoretical Al becomes more sophisticated.

                                        *       *      *

       Those interested in the Sapir-Whorf thesis might be interested to learn of
Novelflo, a new pseudo-natural language invented by Rhoda Spark of Golden,
Colorado. Novelflo is intended as a hypothetical extension of, or perhaps successor to,
English. In particular, it is a language designed expressly for streamlining the writing
of novels (or poetry). You write your novel in Novelflo in a tiny fraction of the
number of words it would take in full English; then you feed it into Spark's
copyrighted "Expandatron" program, which expands your concise Novelflo input into
beautiful, flowing streams of powerful and evocative English prose. Or poetry, for
that matter -you simply set some parameters defining the desired "shape" of the poem
(sonnet, limerick, etc.; free verse or rhyme; and similar decisions), and out comes a
beautifully polished poem in mere seconds. Some of Spark's advertised features for
Novelflo are:




Lisp: Recursion and Generality                                                        453

   • Plot Enhancement Mechanisms (PEM's)
   • Default Inheritance Assumptions
   • Automatic Character Verification (checks for consistency of each character's
     character)
   • Automatic Plot Verification (checks to be sure plot isn't self-contradictory-an
     indispensable tool for any novel writer)
   • SVP's (Stereotype Violation Mechanisms)-allow you to override default
     assumptions with maximal ease)
   • Sarcasm Facilitators (allows sarcasm to be easily constructed)
   • VuSwap (so you can shift point of view effortlessly)
   • AEP's (Atmosphere Evocation Phrases)-conjure up, in a few words, whole
     scenes, feelings, moods, that otherwise would take many pages if not full
     chapters

         This entire Post Scriptum, incidentally, was written in Novelflo (it was my
first attempt at using Spark's language), and before being expanded, it was only 114
words long! I 'must admit, it took me over 80 hours to compose those nuggets of
ideas, but Spark assures me that one gets used to Novelflo quickly, and hours become
minutes.
         Spark is now hard at work on the successor to Novelflo, to be called
Supernovelflo. The accompanying expansion program will be called, naturally
enough, the Superexpandatron. Spark claims this advance will allow a further factor
of compression of up to 100. (She did not inform me how the times for. writing a
given passage in Supernovelflo and Novelflo would compare.) Thus this whole P.S. -
yes, all of it-would be a mere three words long, when written in Supernovelflo. It
would run as follows (so she tells me):

               SP 91 pahC TM-foH

Now I'd call that jus-telegant!




Lisp: Recursion and Generality                                                     454

               Heisenberg's Uncertainty Principle
              and the Many Worlds Interpretation
                    of Quantum Mechanics
                                       July, 1981


YOU’VE dropped a coin between some cushions in a fancy old chair. You're very
anxious to retrieve your coin, so you gingerly try to reach between the cushions and
grab the coin. But the very act of sticking your hand in there widens the crevice and
the coin slips farther in. You can see hat any more of this reaching, and your coin will
be lost forever in the innards of that chair. What to do? This commonplace little
drama illustrates a feeling we all know: that striving for something can have the effect
of reducing that thing's availability.
        A good friend is visiting from far away and before she returns home, you want
to capture her infectious smile on film. But she is terribly camera-shy. The instant you
bring out your camera, she freezes: spontaneity is lost, and there is no way to record
that smile. The act of trying to capture this elusive phenomenon completely destroys
the phenomenon.
        Examples such as these are sometimes erroneously attributed to the
uncertainty principle. That notorious principle of quantum mechanics was first
enunciated by Werner Heisenberg in about 1927. Careless paraphrases since then,
however, have eroded and obscured the true meaning of the principle in the popular
mind. I would like to clarify matters a bit by discussing the genuine uncertainty
principle and its phony imitators.
        Let me first exhibit a typical imitation version clearly, so that you know what I
am attacking. The standard pseudo-uncertainty principle states:




Heisenberg's Uncertainty Principle and the                                           455
Many Worlds Interpretation of Quantum Mechanics

       I
       The observer always interferes with the phenomenon under observation.

It tends to be cited heavily in particular domains, often where the phenomenon
involves a reciprocal observer-someone who can observe back. But even in such
cases, this pseudo-principle is too simplistic. It rests on a misunderstanding of how
experimentation proceeds, and how science explains. The main thing to keep in mind
is \_that science is about classes of events, not particular instances. Science explains
through abstractions that underlie a potentially unlimited number of concrete
phenomena.
        Consider the following example. I recently read of a woman who remarked,
"Rosa always date shranks." She had meant to say, "Rosa always dated shrinks." But
the tense marker somehow got shifted from the verb "date" onto the noun "shrink",
which was then conjugated as if it were functioning as a verb: "dated shrinks" became
"date shranks". It would be fascinating to know, exactly what was going on in the
woman's brain as she made this bizarre transformation. We would like to know
exactly how things went awry. Something went down the wrong track: what, and
why?
        But this was a one-shot phenomenon; it will probably never be repeated. We
can't expect a scientific explanation of those details. Instead, we have to abstract some
general phenomenon that, we think is the essential component of this particular event.
We have to be able to imagine other events in the same general class. We have to be
able to imagine some way to provoke them or to detect them when they happen, so
that we can study the patterns. Perhaps the appropriate level of abstraction is:
"grammatical errors in the speech of woman W". Or perhaps it is: "shifts of tense
markers from verbs to nouns". In any case, we will have to plan a course of
experimentation suitable to the way we choose to abstract this event.
        In the case of the camera-shy friend, presumably her smile is a repeatable
phenomenon; in missing it once, you haven't missed it forever. And with sufficient
patience and ingenuity, you could set up a telephoto lens on a distant camera
controlled remotely by a button you can carry in your hand. You could put the camera
in an unlikely window a few dozen yards from a table where you sneakily take your
friend one day, and then snap her smile without her ever suspecting it.
        In the case of the coin in the cushion, with some effort you could make a
special tool to retrieve it with. In fact, in any such everyday case, even those involving
reciprocal observers, by investing sufficient effort and time and ingenuity-and most
likely money-into a revised version, you will find you can isolate the phenomenon,
you can render it impervious to the fact that you are observing it. You will never get a
perfect replay of some specific event, but as long as it's a general phenomenon and
not a one-shot event that you're interested in, then you can always reduce the effect of
the observer (yourself) to as close to nil as you want. A budget of a trillion dollars
would suffice for most purposes.




Heisenberg's Uncertainty Principle and the                                            456
Many Worlds Interpretation of Quantum Mechanics

       Points such as these bear repeating, because many people think the quantum-
mechanical uncertainty principle actually applies to everyday phenomena. Nothing
could be further from the truth! What, then, is Heisenburg´s principle about?

                                        *   *   *

        To explain it, we have to go back to one of Albert Einstein's three fundamental
papers of the year 1905: the paper in which he postulated that light is made up of the
discrete entities he called photons. It was in this paper that the window onto the
mysterious world of quantum mechanics was first opened. Two centuries of careful
experimentation and observation had demonstrated unequivocally that visible light
acts like a wave with an exceedingly short wavelength (some 10-4 centimeter). Light
waves had been observed interfering with themselves, canceling or reinforcing
themselves. Such behavior is analogous to phenomena seen on lakes or other bodies
of water, such as the momentary canceling of one part of a speedboat's wake by
another part reflected off a jetty, or the shimmering pattern created on a still lake by
the crisscrossing circular ripples emanating from the successive bounces of a skipped
rock.
        In some ways, light waves are simpler than water waves. Whereas water
waves of different wavelengths travel at different speeds, all light waves travel at one
speed: c, or 3 X 1010 centimeters per second. In water, waves of long wavelength
travel faster than waves of short wavelength. Water is thus said to be a dispersive
medium. A single circular ripple, as it expands, breaks up into its various components.
The outer edge, traveling fastest, consists of long-wavelength components, while the
inner edge consists of short-wavelength components. Gradually, because of this
dispersion, the leading and trailing edges of the ripple get so far apart that the ripple
can no longer be perceived. By contrast, the medium that light waves travel through
through is nondispersive: all wavelengths travel at exactly the same speed But what is
that medium? The rather crazy fact of the matter is that light waves need no medium-
or, if you prefer, vacuum is light's medium. But how very peculiar it is for waves to
wave even when there's nothing to wave?
        This anomaly persistently puzzled the young Einstein, and in 1905 his fertile
mind came across two fundamental elements of the resolution. One element was the
counterintuitive theory of special relativity, and the other was the counterintuitive idea
of particle-like quanta out of which light waves would somehow be constituted. But
where did this curious flash of insight come from?

                                        *   *   *

       The classical theory of light as an electromagnetic wave had left a mystery
concerning the way light of various colors, or wavelengths, is emitted from




Heisenberg's Uncertainty Principle and the                                            457
Many Worlds Interpretation of Quantum Mechanics

a "black body". The term is somewhat misleading; it merely means any object that
absorbs light of all frequencies and does not reflect light of any frequencies. As a
black body heats up, it begins to glow: first dull red, then bright red, then orange,
eventually white, and then, surprisingly enough, bluish! (Think of the glowing burner
on an electric stove.) The unsolved problem was to determine how much light of each
wavelength is put out by a black body at a given temperature. In short, how does
intensity depend on wavelength (at a fixed temperature)? In the water-wave analogy,
this would correspond roughly to predicting how deep the leading, central, and
trailing parts of a ripple created by a falling stone would be, as a function of, say, the
kinetic energy of the stone as it hit the water surface.
        Now the actual black-body spectrum at many temperatures had been carefully
measured by experimental physicists, and the characteristic shape of the curve of
intensity versus wavelength (at a fixed temperature) was familiar. At very long and
very short wavelengths, the intensity died away toward zero, and at an intermediate
value determined by the temperature, the intensity hit its maximum. This disagreed
sharply with the prediction of classical physics concerning the intensity of the various
colors. Classical physics predicted that at very short wavelengths, no matter what the
temperature, the intensity would approach infinity. In modern terminology, this
amounts to saying that every object, even, an ice cube, is constantly radiating lethal
gamma rays at arbitrarily great intensities! This is obviously preposterous. Up to
1900, however, no one had any idea of how to patch up the classical theory.
        In that year, Max Planck invented a sort of hybrid formula that looked like a
mathematical splicing-together of two different components, one pertaining to long
wavelengths and the other to short wavelengths. At the longer wavelengths, the
formula agreed with the classical prediction and also with the measured data. At the
shorter wavelengths, Planck's formula diverged from the classical prediction but
stayed in agreement with the data. The long and the short of it was that Planck's
equation seemed right on the money for all wavelengths and temperatures-but it had
not been derived from the first principles. It was a lucky guess, although much more
than luck was involved, since Planck's intuition had guided him like a bloodhound to
this formula.
        Planck himself was particularly baffled by the fact that he'd had to throw a
strange quantity he called "the elementary quantum of action", h, into his formula.
What h represented physically was unclear. It was just a constant that, with a suitable
value, would make the formula exactly reproduce the observed spectrum. It seemed
therefore to be a universal constant of nature.
        But what in the world was it doing in this equation? What did it mean?
Einstein was the first to postulate a physical reason for the appearance of Planck's
constant h in the equation. Einstein began with the concept that the energy content of
light waves is deposited in tiny "lumps"--photons-whose size has to do with h and
their wavelength. For example, if the light is red,




Heisenberg's Uncertainty Principle and the                                            458
Many Worlds Interpretation of Quantum Mechanics

the photons carry always 3.3 X 10-12 erg of energy. Green photons carry 4 X 10-12 erg.
AM radio-wave photons carry somewhere between 3 X 10221 and 9 X 10-21 erg
(depending on what station you're listening to). The amount of energy per photon was
postulated to be invariant, given its color that is, its wavelength).
        In the water-wave analogy, you can try to envision ripples that, when they
reach the shore, suddenly disappear and are replaced by frogs who hop up the bank
where the waves, had they landed, would have lapped. The longer the wavelength of
the ripple, the tinier the frog that jumps out, and conversely: delicate ripples with very
short wavelengths, when they reach the shore, suddenly become thundering monster-
frogs who knock eucalyptus trees down and send boulders crashing into the lake (this
is the infamous phrogo-eucalyptic effect, so yclept by reason of its analogy with the
famous photoelectric effect, in which incoming photons of sufficient energy knock
electrons out of a metal surface).
        Einstein's interpretation of Planck's formula implied that a frog's energy -or
rather, a photon's energy-and its wavelength must be inversely proportional. The
equation linking them is:

                       E = hc / λ

Here, E is the photon's energy, h is Planck's newly discovered constant, c is the speed
of light, and λ is the photon's wavelength. E and λ are the only variables. This mixing
of wave and particle viewpoints was one of the most baffling aspects of quantum
mechanics, and it has continued to plague the intuitions of physicists ever since,
although mathematically it was greatly cleared up by the blossoming of the field in
the 1920's and 1930's

                                           *   *   *

        The next step en route to Heisenberg's uncertainty principle came in 1924,
when Prince Louis-Victor de Broglie was reflecting on the mysterious particle-like
nature of light waves. He asked himself: Why should only light waves be particle-
like? Why not the reverse? That is, mightn't particles also have wavelike properties?
De Broglie's intuition was more or less as follows: If you want to generalize Einstein's
equation so that it holds for particles other than photons, you have to get rid of the one
direct reference in it to light, namely c. Hence de Broglie thought about how he might
most elegantly and relativistically recast the equation in a c-less form.
        This proved to be not too hard, because by then it was known that photons
have both energy E and momentum p, and that they are related by the equation E =pc.
If you combine the two equations, you can cancel out the c , and the result is:

                              p = h / λ.




Heisenberg's Uncertainty Principle and the                                            459
Many Worlds Interpretation of Quantum Mechanics

Mathematically speaking, this equation of de Broglie's is new, but physically
speaking, its content is no different from that of Einstein's original equation -at least
when it is applied to photons. De Broglie's conceptual bravery was to propose-without
any experimental evidence for it-that this equation should be universal. It should
apply to all matter: not just photons, but also electrons, protons, atoms, billiard balls,
people-even frogs! Thus Kermit the Frog would have a quantum-mechanical
wavelength whose value would depend on how fast he's hopping.
        What would this mean physically? What can a hopping frog's wavelength
mean? Well, if you calculate it, you will find that Kermit's wavelength comes out far
shorter than the radius of a proton-yet Kermit himself is considerably bigger than a
proton. If Kermit were very, very small-small enough that his wavelength and his own
size were comparable-then his wavelength would make him diffract around objects
the way water waves and sound waves do. But since Kermit is macroscopic, his
having a microscopic wavelength is all but irrelevant.
        For electrons, though, it is entirely another matter. They are smaller than their
own wavelengths. (In fact, as far as anyone knows, electrons are perfect point
particles, with zero radius.) Shortly after de Broglie's suggestion, experiment and
theory thoroughly confirmed his notion. Electron waves were soon being diffracted in
laboratories around, the world, just like light waves. But now there arises a puzzle.
Are electrons spread out in space in the way waves must be, or are they localized? If
they are truly points, how can they be diffracted? If they are truly waves, where is
their electric charge carried?

                                        *   *   *

Experiments have shown that even a single electron can be diffracted. Richard
Feynman, in his little book The Character of Physical Law, describes it beautifully. In
an idealized experiment, one electron is released in the direction of a barrier with two
slits in it. On the far side of the barrier is a detecting screen. The electron follows
some trajectory and hits the screen somewhere. One such event simply results in one
dot being made on the screen. Suppose we repeat the experiment many times, each
time releasing just one electron. We get a buildup of dots on the screen. Intuition,
building on our experience with such things as bullets fired from a gun, tells us clearly
to expect the dots to be clustered directly behind each of the two slits, with their
distribution tailing off with distance from the center of each cluster. In other words,
we would expect to find two clusters of dots and no other kind of distribution. (See
Figure 20-1a.)
         But if the de Broglie wavelength of the electron is close to the distance
between the slits, the pattern on the screen after thousands of arrivals will look very
different. It will be a complex regular structure characteristic of waves interfering
with each other. In fact, it will reproduce the intensity ,pattern created by a wave that
splits itself into two pieces, which pass through the two slits and interfere with each
other on the far side of the




Heisenberg's Uncertainty Principle and the                                            460
Many Worlds Interpretation of Quantum Mechanics

barrier. (See Figures 20-1b and 20-Ic.) It must be inferred that each electron, as it flew
in its trajectory from source to screen,, somehow "sensed" both slits and interfered
with itself in the manner of a wave and yet deposited itself froglike (that is, in a point)
on the screen without a trace of its schizophrenia.
         The dilemma is, then, that electrons act as if they are both spread out and
localized-as if they were both waves and particles. This kind of wishy-washiness is
inconceivable in the macroscopic realm. Most of us have no trouble distinguishing
between, say, ripples on a pond, and frogs. For those who do, however, it might be
useful to clip out the following handy frog-ripple distinguisher:

      Test 1: Is the candidate solid, tangible, and above all, always somewhere?
      If your answers to these three questions are yes, you are probably dealing
      with a frog.
      Test 2: Is the candidate massless, intangible, and spread out?
      If your answers to these three questions are yes, it is probably a ripple

       If you are hungry for frog's legs and want to know where a frog is, you can just
look around, and as soon as you sense some froglike photons entering your eyes, you
will have found it. Those photons bounced off the frog and into your eyes. But
suppose the frog somehow grew smaller and smaller. After it got down to the size of a
mitochondrion in a living cell, its diameter would be about the wavelength of frog-
green light. Then it would diffract light, and you would not be able to find it so easily.
If it grew even smaller, something terrible would begin to happen. The individual
photons hitting it would, with their momentum and energy, begin to jostle it around.
The particle-like quality of photons would start to enter the picture. Indeed, a frog the
size of an electron would probably be very hard to find. So if you were starved for
frog's legs, you would do better to look around for a bigger one.
Unfortunately, though, no matter how starved you might be for electron's legs, you
cannot find a bigger electron! To find an electron, you cannot do anything but
bombard it with other particles or with photons. Since particles and photons have both
particle-like and wavelike aspects, either bombardment will lead to similar
consequences. If you want to pinpoint a particle, you need waves whose wavelength
is about the size of that particle (or shorter). To understand this intuitively, think of
the way water waves would be affected by a floating piece of wood. If they have a
very long wavelength, they will not even "notice" the wood. Only if their wavelength
gets down to the size of the object will they begin to be affected by it.
Consequently, in order to find our electron, we need photons of very short
wavelength. But wavelength is inversely proportional to momentum.




Heisenberg's Uncertainty Principle and the                                             461
Many Worlds Interpretation of Quantum Mechanics

Heisenberg's Uncertainty Principle and the        462
Many Worlds Interpretation of Quantum Mechanics

That is the deadly import of de Broglie's equation. You pay for your short wavelength
 by having a lot of momentum. And so, as you try to diffract waves ever so gently off
      your particle, hoping not to move it, you will not -be able to do so without
 transmitting momentum to it. Either you are gentle (using long-wavelength photons)
 and do not see the electron well, or you are violent (using short-wavelength photons)
                    and throw the electron completely off its course.
        Heisenberg made a careful study of this perversity, which follows from de
Broglie's equation, and, to the bewilderment of epistemology lovers the world over,
he discovered that to know the position of a particle perfectly is to give up any hope
of knowing its momentum, and that to know the momentum is to give up any hope of
knowing its position. And knowing either one imprecisely still imposes bounds on the
precision with which you could know the other. The principle can even be
summarized in an inequality, which Heisenberg deduced. If you are trying to
determine the location of the particle, there will be an uncertainty, conventionally
denoted Ax. There will also be an uncertainty in the value of the momentum, -denoted
Op. Heisenberg's uncertainty principle is the following inequality:

                       AxAp > h/4π.

There are a couple of things to point out here. First, note the presence of h, Planck's
mysterious constant. This tells you that the effect is due to the wave-particle duality of
matter (and of photons), and has nothing to do with ,the notion of an observer
disturbing the thing under observation. Second,

FIGURE 20-1. Three related two-slit experiments, two classical and one quantum-
mechanical.

        In (a), a wildly swinging machine gun sprays bullets toward a wall with two
holes in it. Occasionally, a bullet will pass through one of the two holes, and will hit
the backstop and make a mark. Eventually, the buildup of marks looks as shown. It
has two peaks, one for each hole.
        In (b), a bobbing buoy creates ripples that spread out toward a jetty with two
breaks in it. When the ripples hit the jetty, new circular ripples emanate from each of
the two breaks, and those ripples, crisscrossing each other, interfere constructively at
some points and destructively at others. On a vertical barrier parallel to the jetty, areas
of highly constructive interference are dark, and areas of highly destructive
interference are white. This characteristic interference pattern is due to two facts: first,
that any ripple passes through both holes, rather than just one, and second, that the
phases at the two holes are correlated.
        In (c), a wildly swinging electron gun sprays electrons toward a wall with two
holes in it. Beyond the wall there is a backstop made of some material that emits a
flash whenever an electron hits it. There is no classical way to describe what happens
to any electron en route, but that what is certain is that, when it comes in for a landing
on the backstop, its local spot of arrival is clearly visible, just as in (a) (thus reminding
us of the corpuscular, or bullet-like, nature of electrons); and yet, if those flashes are
tallied up over a period of time, they are found to be distributed in an interference
pattern just like the one formed in (b) (thus reminding us of the undulatory, or ripple-
like, nature of electrons). Any attempt to ascertain which of the two holes the
electrons pass through ends up in destruction of the interference pattern. [Drawing by
David Moser, after Richard Feynman. ]

Heisenberg's Uncertainty Principle and the                                              463
Many Worlds Interpretation of Quantum Mechanics

notice that even with this epistemological restriction, arbitrarily accurate measurement
of either position or momentum is possible; you just can't get both.
        In short, it is a total misinterpretation of Heisenberg's uncertainty principle to
suppose that it applies to macroscopic observers making macroscopic measurements.
For example, it does not follow from Heisenberg's principle that psychologists
studying the phenomena of human cognition are somehow limited in principle by the
fact that the conscious human beings they are observing are capable of the same kind
of observation. What psychologists are limited by is their knowledge of the human
brain, their ingenuity, and, of course, their funding.
        If you wanted to know more about grammatical anomalies in the speech of
woman W, there are all sorts of ways that you could, in principle, go about it without
making her self-conscious. For just a few thousand dollars, for instance, you could
secretly install a bug in her home and monitor all her conversations. For a few
hundred thousand dollars, you could have tiny radio transmitters manufactured and
secretly sewn into all her lapels. For, say, a few million dollars, you might be able to
convince her she needed minor surgery of some sort, and then while she was
anesthetized you could open up her skull and have harmless electrodes implanted in
her brain to monitor her speech areas-all without her knowing. If you fear that such
blatant physical interference with her brain might disturb her grammatical habits, then
you may have to wait a while longer until we figure out how neural activity can be
examined remotely. These possibilities are clearly extravagant, even ridiculous, but
the point is that, in principle, we can study macroscopic phenomena with an arbitrary
degree of precision.
        To recapitulate: The uncertainty principle states not that the observer always
interferes with the observed, but rather that at a very fine grain size, the wave-particle
duality of the measuring tools becomes relevant. It is a consequence of the fact that
Planck's constant is not zero, rather than an epistemological law about observation
that would have been discovered with or without the discovery of quantum
mechanics.

                                        *   *   *

       The uncertainty principle is not an axiom of quantum physics; it is a deduced
principle, just as Einstein's most famous equation E = mc2 was deduced from the
more fundamental equations of special relativity-a fact that most non-physicists do
not appreciate. Both equations are useful (and famous) because they are so pithy. For
example, the uncertainty principle is often applied by physicists as a rule of thumb. If
you want to estimate the approximate momentum a neutron will have when it is
emitted by a nucleus decaying from an excited state, a seat-of-the-pants estimate is
given by p = h/d, where d is on the order of the dimensions of the confining nucleus.
You can think of the confinement within the nucleus as making the position




Heisenberg's Uncertainty Principle and the                                            464
Many Worlds Interpretation of Quantum Mechanics

uncertainty very small, so that the neutron is bouncing around inside its ``cage" with a
compensatingly large momentum uncertainty. When it escapes, a rough estimate of
the momentum it will have is given by the uncertainty value.
        When you examine the foundations of quantum mechanics, it becomes clear
that the uncertainty principle is more than an epistemological restriction on human
observers; it is a reflection of uncertainties in nature itself. Quantum-mechanical
reality does not correspond to macroscopic reality. It's not just that we cannot know a
particle's position and momentum simultaneously; it doesn't even have definite
position and momentum simultaneously!
        In quantum mechanics, a particle is represented by a so-called wave function
describing the probabilities that the particle is here, there, or somewhere else; that the
particle is heading east, west, north, or south; and so on. For each point in space, there
is what is called a probability amplitude of finding the particle there, and this number
is given by the wave function. Alternatively, one can read the wave function through
different "mathematical glasses" and obtain a probability amplitude for each possible
value of momentum. All the facts about the particle are wrapped up in its wave
function. In more modern terminology, the term "state" is often used instead of "wave
function".
        In classical physics, quantities such as x and p-position and momentum
directly enter the equations governing a particle's behavior. The values of x and p are
definite at any one moment, and they change according to :the forces that are acting
on the particle. With such equations of motion, physicists can plot in advance the
positions and momenta of particles in -simple, stable systems with incredible
accuracy. An example is the motions of the planets, which even the ancients learned
to predict with considerable accuracy. A more contemporary example is provided in
computer space games, where rockets and planets are affected by a star's gravity and
can go into orbit right before your eyes, swinging about in perfect ellipses on a screen.
The underlying equations of such motion are differential equations, and one obvious
property they have-we take it for granted-is that the motions they describe are smooth.
Planets and rocket ships do not jump out of their orbits. There are no sudden
discontinuities in their motion.
        In quantum mechanics; x and p do not enter into the equations of motion as
they do in classical mechanics. Instead, it is the wave function (in nonrelativistic
quantum mechanics) that evolves in time according to a differential equation:
Schrödinger’s equation, named for Heisenberg's contemporary, the quantum-
mechanical pioneer Erwin Schrödinger. As time progresses, the values of the wave
function ripple through space just the way a water wave ripples on a lake's surface.
This would seem to imply that quantum phenomena, like nonquantum ones, proceed
smoothly and with no jumps. In one sense, that is right. A well-known example is the
smooth precession of a spinning charged particle in a magnetic field. It is a kind of




Heisenberg's Uncertainty Principle and the                                            465
Many Worlds Interpretation of Quantum Mechanics

electromagnetic analogue to the precession of a spinning top on a table. The
parameters that characterize the state of the spinning top or spinning particle do
indeed change smoothly, without any jumps.
        HOWEVER-a big however-there are exceptions to this smooth behavior, and
they seem to form just as central a part of quantum theory as does the smooth
evolution of states. The exceptions occur in the act of measurement, or the interaction
of a quantum system with a macroscopic one. As quantum mechanics is usually cast,
it accords a privileged causal status to certain systems known as "observers", without
spelling out what observers are (in particular, without spelling out whether
consciousness is a necessary ingredient of being an observer). To clarify this, I now
present a quick overview of the measurement problem in quantum mechanics, and I
will use the metaphor of the "quantum water faucet" for that purpose.

                                        *   *   *

        Imagine a water faucet with two knobs, one labeled "H" and one labeled "C",
each of which you can twist continuously. Water comes streaming out of the faucet,
but there is a strange property to this system: the water is always either totally hot or
totally cold; there is no in-between. These are called the two temperature eigenstates
of the water. (The prefix eigen- can be translated from the German as "particular."
Here it refers to the fact that the temperature has a particular value.) The only way
you can tell which eigenstate the water is in is by sticking your hand in and feeling it.
Actually, in orthodox quantum mechanics it is trickier than that. It is the act of putting
your hand in the water that throws the water into one or the other eigenstate. Up till
that very instant, the water is said to be in a superposition of states (or more
accurately, a superposition of eigenstates).
        Depending on the settings of the knobs, the likelihood of your getting cold
water will vary. Of course, if you open only the "H" valve, then you'll get hot water
always, and if you open only the "C" valve, then you'll get cold water for sure. But if
you open both valves, you'll create a superposition of states. By trying it out over and
over again with one setting, you can measure the probability of getting cold water
with that setting. After that you can change the setting and try again. There will be
some crossover point where hot and cold are equally likely. It will then be like
flipping a coin. (This quantum water faucet is sadly reminiscent of many a bathroom
shower ...) You can eventually build up enough data to draw a graph of the probability
of cold water as a function of the knobs' settings.
        Quantum phenomena are like this. Physicists can twiddle knobs and put
systems into superpositions of states, analogous to the superpositions of the hot-cold
system. As long as no measurement is made of the system, a physicist cannot know
which eigenstate the system is in. Indeed, it can be shown that in a very fundamental
sense, the system itself does not "know" which eigenstate it is in, and only decides-at
random-at the moment the




Heisenberg's Uncertainty Principle and the                                            466
Many Worlds Interpretation of Quantum Mechanics

observer's hand is stuck in to "test the water", so to speak. (Note that a nonsexist
reader is in a superposition of states at this very moment, not knowing if this
hypothetical observer (or for that matter, the hypothetical nonsexist reader) is male or
female!) Up to the moment of observation, the system acts as if it were not in an
eigenstate. For all practical purposes, for all theoretical purposes-in fact for all
purposes-the system is not in an eigenstate.
        You can imagine doing a lot of experiments on the water coming out of a
quantum water faucet to determine whether the water is actually hot or actually cold
without sticking your hand in it. (We're of course assuming that there are no telltale
clues to the temperature of the water, such as steam rising from it.) For example, run
your washing machine on water from the quantum faucet. Still, you won't know
whether your wool sweater has shrunk or not until the moment you open the machine
(a measurement made by a conscious observer). Make some tea with water from the
faucet. Still, you won't know whether you've got hot tea or not until you taste it (again
a measurement made by a conscious observer). The critical point here is that the
sweater and the tea, not having conscious-observer status themselves, have to play
along with the gag and, just as the water did, enter superpositions of states: shrunk
and non-shrunk, hot tea and cold tea.
        All this may sound as if it has nothing to do with physics per se, but merely
with ancient philosophical conundrums such as: "Does a tree in a forest make a noise
when it falls, if there's nobody there to hear it?" But the quantum-mechanical twist on
such riddles is that there are observational consequences of the reality of the
superpositions, consequences diametrically opposite to those that would ensue if a
seemingly mixed state were in reality always a true eigenstate, merely hiding its
identity from observers until the moment of measurement. In crude terms, a stream of
maybe-hot maybe-cold water would act differently from a stream of water that is
actually hot or actually cold, because the alternatives "interfere" with each other. This
would become manifest only after a large number of sweater-washings or tea-
makings, just as in the two-slit experiment it takes a large number of electron-landings
to reveal the interference pattern of the alternative trajectories. (Quantum-mechanical
interference resembles the classical phenomenon of two notes beating against each
other, except that in quantum mechanics, instead of producing a chord of sounds, the
superposition produces a distribution of probability-a "chord of possibilities".)
Interested \_,readers should consult either Feynman's The Character of Physical Law
or, for an account with more detail, Volume III of The Feynman Lectures on Physics.

                                        *      *       *

       The plight of "Schrödinger’s cat" carries this idea further: it suggests that even
a cat might be in a quantum-mechanical superposition of states until




Heisenberg's Uncertainty Principle and the                                           467
Many Worlds Interpretation of Quantum Mechanics

FIGURE 20-2. Schrodinger's cat in a superposition of states, partly alive and partly
dead. [From The Many-Worlds Interpretation of Quantum Mechanics, edited by
Bryce S. DeWitt and Neill Graham.]

a human observer intervened. The tale of this unfortunate cat goes like this. A box is
prepared for a cat's occupancy. Inside this box, there is a small sample of radium.
Also in the box is a detector of radiation, which will detect any decays of radium
nuclei in the sample. The sample has been chosen so that there is a 50-50 probability
that within any hour-long period, one decay will occur. On the occurrence of such a
decay, a circuit will close, tripping a switch that will break a beaker filled with a
deadly liquid, spilling the liquid onto the floor of the box, and killing the cat. (See
Figure 20-2.)
        The cat is now placed in the box, the lid firmly shut, and an hour ticks away.
At the hour's end, a human observer approaches the box and opens the lid to see what
has happened. According to one extreme view of quantum mechanics (and the reader
should bear in mind that it is not the usual view), only at that moment will the system
be forced to "jump" into one of the two possible eigenstates-cat alive and cat dead-
that are represented together as a superposition in the wave function of the system.
(Notice that it is necessary that the randomness be of a clearly quantum-mechanical
origin: the decay of the radium nucleus. This thought experiment would not pack any
punch if there were a spinning roulette wheel in the box instead of a radium sample.)
        One might object and say, "Wait a minute! Isn't a live cat as much of a
conscious observer as a human being is?" Probably it is, but notice that this cat is
possibly a dead cat, and in that case certainly not a conscious observer. We have in
effect created in Schrodinger's cat a superposition of two eigenstates, one of which
has observer status, the other of which lacks it. Now what do we do? This situation is
reminiscent of a Zen riddle (recounted in Zen Flesh, Zen Bones by Paul Reps):




Heisenberg's Uncertainty Principle and the                                         468
Many Worlds Interpretation of Quantum Mechanics

   Zen is like a man hanging in a tree by his teeth over a precipice. His hands grasp
   no branch, his feet rest on no limb, and under the tree another person asks him:
   "Why did Bodhidharma come to China from India?" If the man in the tree does
   not answer, he fails; and if he does answer, he falls and loses his life. Now what
   shall he do?

                                         *   *   *

        The idea that consciousness is responsible for the "collapse of the wave
function"-a sudden jump into one randomly chosen pure eigenstate leads to further
absurdities. For instance, it would imply that nothing ever happened for the first
umpteen billion years of the universe, until one day, a million or so years ago, some
human being woke up-and at that instant the enormously swollen universal wave
function collapsed down into one world-and this person blinked, peered around, and
saw Mesopotamia or Kenya ...
        The alternative left to us is that observers-things that make a wave function
collapse-need not be conscious, but merely macroscopic. However, isn't a
macroscopic object just a collection of microscopic objects? How would a wave
function "know" it was dealing with a macroscopic object? More concretely, what is it
about a screen that forces an electron to reveal itself?
        To many physicists, the distinction between systems with observer status and
those without has seemed artificial, even repugnant. Moreover, the idea that an
intervention of an observer causes a "collapse of the wave function" introduces
caprice into the ultimate laws of nature. "God does not play dice" (Der Herrgott
wurfelt nicht) was Einstein's lifelong belief.
        A radical attempt to save both continuity and determinism in quantum
mechanics is known as the many-worlds interpretation, first proposed in 1957 by
Hugh Everett III. According to this very bizarre theory, no system ever jumps
discontinuously into an eigenstate. What happens is that the superposition evolves
smoothly with its various branches unfolding in parallel. Whenever it is necessary, the
state sprouts further branches that carry the various new alternatives. For example,
there are two branches in the case of Schrödinger’s cat, and they develop in parallel.
"Well, what happens to the cat? Does it feel itself to be alive or does it feel itself to be
dead?" One must wonder. Everett would answer: "It depends on which branch you
look at. On one branch, the cat feels itself to be alive, and on the other, there is no cat
to feel anything." With intuition beginning to rebel, one then asks: "Well, what about
a few moments before the cat on the fatal branch died? How did the cat feel then?
Surely the cat can't feel two ways at once! Which of the two branches contains the
genuine cat?"
        The problem becomes even more intense as you realize the implications of this
theory as applied to you, here and now. For every quantum-mechanical branch point
in your life (and there have been billions upon billions), you have split into two or
more you's riding along parallel but disconnected




Heisenberg's Uncertainty Principle and the                                             469
Many Worlds Interpretation of Quantum Mechanics

branches of one gigantic "universal wave function". (By this term is meant the
enormous wave function representing all the particles in all the parallel universes.) At
the critical spot in his article where this difficulty arises, Everett calmly inserts the
following footnote:

           At this point we encounter a language difficulty. Whereas before the
   observation we had a single observer state, afterwards there were a number of
   different states for the observer, all occurring in a superposition. Each of these
   separate states is a state for an observer, so that we can speak of the different
   observers described by different states. On the other hand, the same physical
   system is involved, and from this viewpoint it is the same observer, which is in
   different states for different elements of the superposition (i.e., has had different
   experiences in the separate elements of the superposition). In this situation we
   shall use the singular when we wish to emphasize that a single physical system
   is involved, and the plural when we wish to emphasize the different experiences
   for the separate elements of the superposition. (E.g., "The observer performs an
   observation of the quantity A, after which each of the observers of the resulting
   superposition has perceived an eigenvalue.")

        All said with a poker face. The problem of how it feels subjectively is not
treated; it is not even swept under the rug. It is probably considered meaningless.
        And yet . one simply has to wonder: "Why, then, do I feel myself to be in just
one world?" Well, according to Everett's view, you don't-you feel all the alternatives
simultaneously. It's just this you going down this branch who doesn't experience all
the alternatives. This is completely shocking. In his story "The Garden of Forking
Paths", the Argentinian writer Jorge Luis Borges describes a fantastic vision of the
universe in this way:

           ... a picture,- incomplete yet not false, of the universe as Ts'ui Pen
   conceived it to be. Differing from Newton and Schopenhauer, ..: [he] did not
   think of time as absolute and uniform. He believed in an infinite series of times,
   in a dizzily growing, ever spreading network of diverging, converging and
   parallel times. This web of time-the strands of which approach one another,
   bifurcate, intersect, or ignore each other through the centuries-embraces every
   possibility. We do not exist in most of them. In some you exist and not I, while
   in others I do, and you do not, and in yet others both of us exist. In this one, in
   which chance has favored me, you have come to my gate. In another, you,
   crossing the garden, have found me dead. In yet another, I say these very same
   words, but am an error, a phantom.

This quotation is featured at the beginning of the book The Many-Worlds
Interpretation of Quantum Mechanics: A Fundamental Exposition, edited by Bryce S.
Dewitt and Neill Graham. The ultimate question is this: "Why is this me in this
branch, then? What makes me-I mean this me-feel itself-I mean myself-unsplit?"

                                        *   *   *




Heisenberg's Uncertainty Principle and the                                             470
Many Worlds Interpretation of Quantum Mechanics

        The sun is setting one evening over the ocean. You and a group of friends are
standing at various points along the wet sand. As the water laps at your feet, you
silently watch the red globe drop nearer and nearer to the horizon. As you watch,
somewhat mesmerized, you notice how the sun's reflection on the wave crests forms a
straight line composed of thousands of momentary orange-red glints-a straight line
pointing right at you! "How lucky that I am the one who happens to be lined up
exactly with that line!" you think to yourself. "Too bad not all of us can stand here
and experience this perfect unity with the sun." And at the same moment, each of your
friends is having precisely the same thought ... or is it the same?
        Such musings are at the heart of the "soul-searching question". Why is this
soul in this body? (Or on this branch of the universal wave function?) Why, when
there are so many possibilities, did this mind get attached to this body? Why can't my
"I-ness" belong to some other body? It is obviously circular and unsatisfying to say
something like "You are in that body because that was the one made by your parents."
But why were they my parents, and not sometwo else? Who would have been my
parents if I had been born in Hungary? What would I have been like if I had been
someone else? Or if someone else had been me? Or-am I someone else? Am I
everyone else? Is there only one universal consciousness? Is it an illusion to feel
oneself as separate, as an individual? It is rather eerie to find these bizarre themes
reproduced at the core of what is supposedly our stablest and least erratic science.
        And yet in a way it is not so surprising. There is a clear connection between
the imaginary worlds of our minds and the alternate worlds evolving

FIGURE 20-3. A robot in an anxious superposition of mental states. [Drawing by
Rick Granger. ]




Heisenberg's Uncertainty Principle and the                                        471
Many Worlds Interpretation of Quantum Mechanics

in parallel with the one we experience. The proverbial young man picking apart a
daisy and muttering, "She loves me, she loves me not, she loves me, she loves me not
. . ." is clearly maintaining in his mind at least two different worlds based on two
different models for his beloved. (See Figure 20-3.) Or would it be more accurate to
say that there is one mental model of his beloved that is in a mental analogue of a
quantum-mechanical superposition of states?
         And when a novelist simultaneously entertains a number of possible ways of
extending a story, are the characters not, to speak metaphorically, in a mental
superposition of states? If the novel never gets set to paper, perhaps the split
characters can continue to evolve their multiple stories in their author's brain.
Furthermore, it would even seem strange to ask which story is the genuine version.
All the worlds are equally genuine.
         And in like manner, there is a world-a branch of the universal wave function-
in which you didn't make that stupid mistake you now regret so much. Aren't you
jealous? But how can you be jealous of yourself ? Besides which, there's yet another
world in which you made yet stupider mistakes, and in which you are jealous of this
very you, here and now in this world!

                                        * *    *

         Perhaps one way to think of the universal wave function is as the mindor
brain, if you prefer-of the great novelist in the sky, God, in which all possible
branches are being simultaneously entertained. We would be mere subsystems of
God's brain, and these versions of us are no more privileged or authentic than our
galaxy is the only genuine galaxy. God's brain, conceived in this way, evolves
smoothly and deterministically, as Einstein always maintained. The physicist Paul
Davies, writing on just this subject in his recent book Other Worlds, says: "Our
consciousness weaves a route at random along the ever-branching evolutionary
pathway of the cosmos, so it is we, rather than God, who are playing dice."
         Yet this leaves unanswered the most fundamental riddle that each of us must
ask: "Why is my unitary feeling of myself propagating down this random branch
rather than down some other? What law underlies the random choices that pick out
the branch I feel myself tracing out? Why doesn't my feeling of myself go along with
the other me's as they split off, following other routes? What attaches me-ness to the
viewpoint of this body evolving down this branch of the universe at this moment in
time?" These questions are so basic that they almost seem to defy clear formulation in
words. And their answers do not seem to be forthcoming from quantum mechanics. In
fact, this is exactly the collapse of the wave function reappearing at the far end of the
rug it wasn't swept under by Everett ...
         It turns it into a problem of personal identity, no less perplexing than the
problem it replaces.
         One can fall even more deeply into the pit of paradox when one realizes




Heisenberg's Uncertainty Principle and the                                           472
Many Worlds Interpretation of Quantum Mechanics

that there are branches of this one gigantically branching universal wave function on
which there is no Werner Heisenberg, no Max Planck, no Albert Einstein, branches on
which there is no evidence for quantum mechanics whatsoever, branches on which
there is no uncertainty principle or many-worlds interpretation of quantum mechanics.
There are branches on which the Borges story did not get written, branches in which
this column did not get written. There is even a branch in which this entire column got
written just as you see it here, except for one noun which was replaced by its exact
antonym, at the column's very beginning.

Post Scriptum.
                                   Quantum particles: the dreams that stuff is made of.
                                                                         David Moser

         If this was your introduction to the weirdness of quantum mechanics (which I
doubt), then may I say how delighted I am to have been your guide. But in that case, I
also must say that you really deserve a more complete introduction. This article was
aimed mostly at people who already have at least a nodding acquaintance with
quantum phenomena. The Feynman books alluded to in the article are ideal
introductions. There are other books that purport to explain quantum mechanics to
novices, and in some cases they may do a fairly good job of it, but some of them have
the serious drawback of trying to link quantum-mechanical reality with Eastern
mysticism, a connection I find superficial and misleading. I cannot fault people who
wish to make some observations about the worldview of ancient Buddhists and to
point out that a few statements written thousands of years ago can, if very liberally
interpreted, be taken to say things that are not inconsistent with discoveries of modern
physics, but to claim that "Western science is only now catching up with the ancient
wisdom of the East", as most of those authors do (and in roughly those words), is, in
my view, both silly and anti-intellectual.
         I call it "anti-intellectual" because most Western people infatuated with
Eastern mysticism hold a grudge against the encroachment of science on territory they
consider beyond science. This attitude may be a holdover from the bitterly anti-
scientific, anti-intellectual mood that gripped the United States during much of the
Viet Nam War era. Those people have some sort of axe to grind, perhaps
subconsciously; they want to see science "put in its place". Curiously, many of them
are scientists themselves and revel in a kind of self-deprecation, thinking that they are
lifting themselves up to transcendent heights and seeing things from a "higher plane
of enlightenment" than science affords. Usually, at that point, their prose




Heisenberg's Uncertainty Principle and the                                           473
Many Worlds Interpretation of Quantum Mechanics

abruptly changes mood, moving from precise terms to mushy, vague, poetic terms
(such as "mushy", "vague", and "poetic"). Don't you just hate that sort of thing?
        These are the sorts of people who propagate misinformation about the
discoveries of modern physics (such as the pseudo-uncertainty principle). They
encourage people to think that any wild theory explaining any mystery (or alleged
mystery) might well be correct, as long as it uses voguish technical terms from
physics-terms like "tachyon", "Bell's inequality", "EPR paradox", "gravitational
waves", and so on. A typical abuser of physics in this way is Arthur Koestler; in his
book The Roots of Coincidence, he purports to explain "psi phenomena" in terms of
some five-dimensional theory of particle physics that includes a host of hypothetical
particles called "psitrons".
        To me, a very troubling aspect of an "explanation" such as this (which,
actually, Koestler didn't invent himself but borrowed from a physicist named Adrian
Dobbs) is that very similar explanations are used by physicists .themselves-not so
often of "psi phenomena", but of currently unexplained real phenomena in particle
physics. When I was a graduate student in particle physics, quite a number of years
ago, I read paper after paper in which not only new particles were invoked to explain
some observation, but new families of particles were routinely postulated. As a matter
of fact, one of those papers was the straw that broke the camel's back, as far as I was
concerned. In that three-author paper, the authors had the audacity to invent some
totally off-the-wall superfamily of particles that consisted of a large number of
families, each containing quite a few particles on its own. As I recall, there were
something like 140 new particles introduced in one fell swoop-and, mind you, this
was done merely to explain some rather small discrepancies between things measured
and things predicted by previous theories. A far cry from the days when it was a
highly daring step to introduce even one new particle! It was at that point that I
decided I should bow out of that branch of theoretical physics.

                                       *   *   *

         I am not really trying to castigate the whole field of particle physics, because
all I learned for sure from my long, gruelling, and ultimately broken engagement with
that field was that I personally was not cut out to be a particle physicist. However, I
did learn one disillusioning thing about science in general, and that is that large
segments of it-including, very often, the most forbidding and technically prickly
papers-are just as nonsensical and empty as the pseudo-scientific papers that try to
shore up "psi phenomena", "remote viewing", "telekinesis", or the like. (Is it
reasonable for me to continue using quotation marks around those words? I think so. I
don't like using words in such a way that I help to lend them legitimacy when I think
there is nothing behind them.) Bad science




Heisenberg's Uncertainty Principle and the                                           474
Many Worlds Interpretation of Quantum Mechanics

permeates good science the way that gristle runs through meat (a "meataphor"
exploited in a different context in Chapter 21).
         I am afraid that this is an example of an inevitable phenomenon: If you are
throwing darts and want not only to hit the bull's-eye every time but also to cover the
entire bull's-eye evenly, so that you are equally likely to hit any point inside the bull's-
eye but totally unlikely to go outside of it, then you are dreaming a pipe dream! You
have to pay in some way for the privilege of filling up that inner circle-and you pay
either by sometimes overflowing the boundaries of the bull's-eye (being too loose, so
to speak), or by covering it unevenly, having a high concentration in the middle of the
bull's-eye and a low concentration near its edges (being too tight or controlled). In
science, this translates to the trade-off between being too speculative and too cautious.
It is impossible for all the papers in a field to be both right and significant. Either
many will be wrong or many will be trivial. The former corresponds, obviously, to
throwing outside the circle, and the latter, a little less obviously, to covering it fully
but unevenly. This inevitable trade-off is very much like that spoken of in Chapter 13,
where in trying to produce all -the truths expressible in a formal system or all the
members of a semantic category, you wind up with either an incomplete system or an
inconsistent system.
         I guess this makes me sound somewhat cynical about science. But I would
make similar noises about human endeavors of any sort that involve skill. For
instance, not all the letters I receive from people who have read things I've written hit
the bull's-eye; some of them are the cat's meow, but a larger number are either old hat,
off base, full of hot air, or some combination thereof. So if I want to get some good
letters, I have no choice but to be willing to wade through a bunch of bad ones, too.
And, regretfully, I must say that this law applies just as much to my own output: not
all of it can be of the same caliber. If it's all correct, then much of it will be mundane;
and if I regularly dare to go far beyond the mundane, then some of it will wind up
being wrong.
         Some people choose to see trade-offs such as these as more examples of a kind
of "uncertainty principle": you can't have both total correctness and total novelty. You
must take your pick. This "either-or" quality, however, has very little to do with the
quantum-mechanical substrate of our world. It just has to do with statistical
phenomena in general.

                                         *   *   *

        I would like to say something about the alienness of quantum-mechanical
reality. It is no accident, I would maintain, that quantum mechanics is so wildly
counterintuitive. Part of the nature of explanation is that it must eventually hit some
point where further probing only increases opacity rather than decreasing it. Consider
the problem of understanding the nature of solids. You might wonder where solidity
comes from. What if




Heisenberg's Uncertainty Principle and the                                             475
Many Worlds Interpretation of Quantum Mechanics

someone said to you, "The ultimate basis of this brick's solidity is that it is composed
of a stupendous number of eensy-weensy bricklike objects that themselves are rock-
solid"? You might be interested to learn that bricks are composed of micro-bricks, but
the initial question-"What accounts for solidity?"-has been thoroughly begged. What
we ultimately want is for solidity to vanish, to dissolve, to disintegrate into some
totally different kind of phenomenon with which we have no experience. Only then,
when we have reached some completely novel, alien level will we feel that we have
really made progress in explaining the top-level phenomenon.
        That's the way it is with quantum-mechanical reality. It is truly alien to our
minds. Who can fathom the fact that light-that most familiar of daily phenomena-is
composed of incredible numbers of indescribably minuscule "particles" with zero
mass, particles that recede from you at the same speed no matter how fast you run
after them, particles that produce interference patterns with each other, particles that
carry angular momentum and that bend in a gravitational field? And I have barely
scratched the surface of the nature of photons! I like to summarize this general
phenomenon in the phrase "Greenness disintegrates." It's a way of saying that no
explanation of macroscopic X-ness can get away with saying that it is a result of
microscopic X-ness ("just the same, only smaller"); macroscopic greenness, solidity,
elasticity-X-ness, in short-must, at some level, disintegrate into something very, very
different.
        I first saw this thought expressed in the stimulating book Patterns of Discovery
by Norwood Russell Hanson. Hanson attributes it to a number of thinkers, such as
Isaac Newton, who wrote, in his famous work Opticks: "The parts of all homogeneal
hard Bodies which fully touch one another, stick together very strongly. And for
explaining how this may be, some have invented hooked Atoms, which is begging the
Question." Hanson also quotes James Clerk Maxwell (from an article entitled
"Atom"): "We may indeed suppose the atom elastic, but this is to endow it with the
very property for the explanation of which .... the atomic constitution was originally
assumed." Finally, here is a quote Hanson provides from Werner Heisenberg himself:
"If atoms are really to explain the origin of color and smell of visible material bodies,
then they cannot possess properties like color and smell." So, although it is not an
original thought, it is useful to bear in mind that greenness disintegrates.

                                       *   *   *

        One of the most beautiful features of the quantum-mechanical description of
reality is how a bridge is erected between the microscopic and the macroscopic. The
nature of that bridge is characterized by the correspondence principle, which states:

   In the limit of large sizes, quantum-mechanical phenomena must look
   indistinguishable from their classical counterparts.




Heisenberg's Uncertainty Principle and the                                           476
Many Worlds Interpretation of Quantum Mechanics

This can be converted into a more mathematical statement, as follows:

   In the limit of large quantum numbers, quantum-mechanical equations must
   reproduce their classical counterparts.

A physicist does not have to work to make an equation describing quantum
phenomena obey this principle; if the equation is correct, it will obey it automatically.
However, a physicist cannot always be sure that a proposed equation is correct.
Therefore, the correspondence principle provides a very useful check on any proposed
equation-for if it fails to yield the familiar classical equation in the limit of large sizes
(or more accurately, large quantum numbers), it is surely wrong. Of course, merely
passing this test is no guarantee that an equation is right, but it is a confirming piece
of evidence.
        Quantum-mechanical phenomena are characterized by "quantum numbers",
which are always integers. When those integers are small-less than 5 or so-you have
quintessentially quantum phenomena. But when you plug fairly large values such as
20 into the equations, you get behavior that floats midway between the quantum style
and the classical style. And when ou take the limit of infinitely large values, you
should get back the familiar old equations from the pre-quantum era: such things as
Newton's laws of motion, for instance.
        A striking example of this idea is furnished by so-called "Rydberg atoms",
highly excited atoms whose outermost electrons have very large quantum numbers,
and which are consequently tethered so loosely to their central nucleus that their
orbits begin to be somewhat less "cloud-like" (i.e., less quantum-mechanical), and
more like the familiar planetary orbits that electrons used to follow, back in the short-
lived "semiclassical" era of physics, after Ernest Rutherford's discovery of nuclei, but
before Schrödinger and Heisenberg. These bridges between the alien world and the
familiar world help provide the intuitions necessary for macroscopic people to
imagine how jolly giant greenness could emerge from murky, unfathomable
microdepths. .




Heisenberg's Uncertainty Principle and the                                              477
Many Worlds Interpretation of Quantum Mechanics

                                              Section IV

                                    Spirit and Substrate




Review of Alan Turing: The Enigma                    481

                                                                       Section IV

                                                      Spirit and Substrate
         The world has traditionally been divided into the animate and the inanimate.
Inanimate things do not have feelings or wills of their own, and can therefore be
smashed, burned, or harnessed by animate ones without the animate ones having to
feel guilty. This borderline, so long taken for granted by people, is gradually
becoming blurrier with the advent of computers, especially as programs acquire more
and more flexibility-and with that flexibility, a seeming mentality or personality. How
and when could mind and emotions-surely the essence of the animate-emerge from
complex inanimate substrates? What does it take to make spirit out of pure matter
pattern? A number of recent artificial-intelligence programs have been touted as
"thinking". Yet no one who looked closely could fail to see that there remains a huge
gap between human self-aware fluidity and such programs. Even the best of them is
still relatively rigid and unaware of anything, let alone itself. But where is the
borderline between the highest inanimate flexibility and the lowest animate sentience?
When does a system or organism have the right to call itself "I", and to be called
"you" by us? Will we be able to recognize systems deserving of our respect when they
come along, or will we abuse them? Will such systems have as much free will as we
don't? These and other philosophically motivated questions about mind and
mechanism, free will and determinism, randomness and rule-following, are examined
in the following six chapters.




Review of Alan Turing: The Enigma                                                  482


                         Review of Alan Turing:
                              The Enigma
                                   November, 1983

AN true intelligence be embodied in any sort of substrate-organic, electronic, or
otherwise? Is mind more than pattern? How can we distinguish between a genuine
mind and a clever facade? Is free will compatible with a materialist, mechanistic view
of living beings? Is there a contradiction in the notion of rule-bound creativity? Do
our emotions and intellects belong to separate compartments of our selves? Could
machines have emotions? Could machines be enchanted by ideas, by le, by other
machines? Could machines be attracted to each other, fall in love? What would be the
social norms for machines in love? Would there and improper types of machine love
affairs? Could a machine be frustrated and suffer? Could a frustrated machine release
its pent-up feelings by going outdoors and self propelling ten miles? Could a machine
learn to enjoy the sweet pain of marathon running? Could a machine with a seeming
zest for life destroy itself purposefully one day, planning the episode so as to fool its
mother machine into "thinking" (which of machines cannot do) that it had perished by
accident?
        These are the sorts of questions that burned in Alan Turing's brain, and, taken
at another level, they reveal highlights of Turing's troubled life. It would require
someone who shares much with Turing to plumb his story deeply enough to do it
justice, and Andrew Hodges, a young British writer with a doctorate in mathematics,
has wonderfully succeeded in doing so. His 500 page biography of Turing
painstakingly put together from innumerable sources, including conversations with
scores of people who knew Turing at various stages of his life, provides as vivid a
picture as one hope of a most complex and intriguing individual. And it's about time,
for not only was Turing a very significant person in the science of this century, but his
fascinating and difficult life illustrates serious problems that yet has not yet grappled
with successfully.
        Hedges' rich and engrossing portrait is not the first book about Turing,




Review of Alan Turing: The Enigma                                                    483

since his mother, Sara Turing, wrote a sketchy memoir a few years after her son's
death, which presents an image of Turing as a lovable, eccentric boy of a man, filled
with the joy of ideas and driven by an insatiable curiosity about questions concerning
mind and life and mechanism. Hodges goes far more deeply into Turing's mind, body,
and soul than Sara Turing ever dared, for she wore conventional blinders and did not
want to see how poorly her son fit into the standard molds of British society. Alan
Turing was homosexual, a fact that he took no particular pains to hide, especially as
he grew older. And for a boy growing up in the 1920's and for a man in the next few
decades, being homosexual-especially if one was British and belonged to the upper
classes-was an unmentionable, terrible, and mysterious affliction.
        Alan Turing, an atheist, homosexual, eccentric English mathematician, was in
large part responsible not only for the concept of computers, incisive theorems about
their powers, and a clear vision of computer minds, but also for the cracking of
German ciphers during World War 11. It is fair to say that we owe much to Alan
Turing for the fact that we are not under Nazi rule today. And yet this salient figure in
world history has remained, as the book's subtitle says, an enigma.
        Turing was born in London in 1912 of relatively well-to-do parents in the civil
service in India. Not long after his birth, his father returned to India, followed by his
mother, and they spent the next few years there, leaving young Alan in England. Then
they decided to return closer to England, and for a time lived in France, which gave
Alan the opportunity to take school vacations there and learn French. As a boy, he
was inquisitive and humorously inventive but definitely not a child prodigy. At age
thirteen, he was sent off to a boys' private boarding school called Sherborne, in the
west of England. He made quite a hit his first day, for he arrived on bicycle, having
pedaled the 60 miles from Southampton, where the ferry from France had left him on
a day of general strikes and no trains. However, as the weeks passed, his hero status
declined as he revealed himself to be a rather untidy pupil prone to getting ink all over
himself, and one who did not distinguish himself in most of his classes. Alan was a
solitary boy and his first venture into serious friendship came to an unexpected and
tragic ending, when his friend and idol, Christopher Morcom, succumbed to bovine
tuberculosis.
        Alan never forgot this first and perhaps deepest of all his human contacts, for
it was in fact a mixture of friendship and love. Although Alan never confided his love
to Chris, it is apparent that Alan was in love with him. Later in life, Alan would have
romances with other men, as well as more numerous and more sordid one-night
stands, but the purity of his love for Christopher Morcom was never surpassed. A
flame in Alan's heart continued to burn for




Review of Alan Turing: The Enigma                                                    484

Iris Morcom, and he faithfully visited the Morcom family for years afterward, seeking
some sort of spiritual communion with his lost friend. j!erhaps this was one of the key
sources of Alan's abiding interest in the connection between the elusive human soul
and its mortal incarnation.
        At Sherborne, Alan excelled at mathematics to the exclusion of pretty much
everything else. In the end, his school recognized his great talent and awarded him
several science prizes. At age twenty, Alan went on to Cambridge. This was 1933,
and the scientific world was charged with the excitement of absorbing several
revolutionary discoveries of the previous decade. Relativity, one of Turing's early
obsessions, was now old hat, while quantum mechanics and mathematical logic were
in their heyday. Quantum mechanics made a deep impression on Alan's mind. In
quantum systems such as an atom, an electron can jump from one orbit (or "state") to
another without occupying any intermediate position between them. It would be as if
a space satellite jumped from one orbit to another without traveling between them.
        Equally striking to Turing was the mechanization of mathematical reasoning,
which he first read about in a philosophical book by Bertrand Russell. Later he
studied the ambitious "Hilbert program", whose aim was to demonstrate the
possibility of capturing in a single system all the valid principles of mathematical
reasoning. In that system, all possible true consequences would flow out of a small set
of axioms by means of a well-defined set of rules, like automobiles from an assembly
line, or physical systems jumping from one state to another. This image of a machine
that jumped from one state to another according to a finite set of rules became
uppermost in Turing's mind. What fascinated him was the idea that such meaningless
actions could also be viewed as having meanings. For instance, one rule-obeying
machine might be viewed as making moves of chess, ' another as producing truths of
mathematics, and yet another as writing poetry.
        In 1931, the Austrian logician Kurt Gödel devastated Hilbert's and Russell's
hopes of creating a perfect formalization of all mathematical reasoning. Gödel had
demonstrated that there were undecidable propositions in any consistent axiomatic
system of the Hilbert-Russell sort, propositions based on famous paradoxes of logic
that had plagued logicians ever since the Greeks. (The sentence "I am lying" is a good
example, as is the fruitless exercise of trying to catch a glimpse in the mirror of what
you look like with your eyes perfectly closed.) What GOdel had left unsettled,
however, was the question of whether, given an axiomatic system' and an arbitrary
proposition in it, one could determine mechanically whether that proposition was
undecidable in that system. If this were possible, then one could discard undecidable
propositions as easily as one trims pieces of fat from a steak; if it were impossible,
then mathematics would resemble a piece of steak riddled with gristle, so that no
matter how you slice it, what remains will contain some gristle.




Review of Alan Turing: The Enigma                                                   485

                                       *   *   *

        Alan Turing chose to work on this question of whether the gristle could be
cleanly lopped off the rest of mathematics, leaving the "meat of mathematics intact
and mechanizable, and the rest just a collection of quirky Gödelian curiosities,
sideshow freaks contrasting with the vast world of normal mathematical propositions.
To his surprise, he discovered that for very Gödelian reasons, no machine could be
built that could infallibly recognize undecidable propositions.
        He began by trying to specify exactly the most general possible notion of what
a "machine" is. In fact, the concept he arrived at, now called a "Turing machine",
forms a central part of Turing's contribution to the theory of computing. Although
fundamentally all a Turing machine can do is jump from one discrete state to another
by means of very simple transition rules, Turing was able to show that such machines
could do anything that one could reasonably expect of any machine or any human
following well-defined rules. He went further and showed that a very complex type of
Turing machine, called a "universal" Turing machine, was capable of being fed a
single number that encoded the structure of any other Turing machine, much in the
way that DNA codes for the structure of an organism. The universal machine could
then act indistinguishably from that machine. The discovery of such a universal
Turing machine made all "specialist" Turing machines obsolete. For instance, if some
Turing machine could play chess, then the universal Turing machine could also play
chess, simply by being fed the code number of the chess-playing Turing machine.
Ditto for theorem-producing and poetry-writing. If Turing were still alive, he would
probably have relished Woody Allen's recent character Leonard Zelig, the "Human
Chameleon"-a living, breathing universal Turing machine, one that could perfectly
simulate any other, if fed the right code number.
        Turing's death blow to the hopes of logicians such as Russell and Hilbert was
delivered in two st ges. First, he supposed that a machine for recognizing undecidable
propositions exists; then he showed how that assumption leads to self-contradiction.
He began by showing that any such machine-if it existed-would closely resemble a
universal Turing machine, in that it could accept the description number of any
machine and simulate it. Slyly, then, he proposed feeding into this hypothetical
machine its own description number. This action, he showed, would instantly send it
into a dizzy loop and make it perish of computational vertigo. In other words, the idea
of such a machine is self-contradictory. The tantalizing upshot of this twisty-argument
is the discovery that undecidable propositions run through mathematics like
ineradicable threads of gristle that crisscross a steak in such a dense way that they
cannot be cut out without the entire steak's being destroyed. In short, through Gödel’s
and Turing's work, mathematics was revealed to be unmechanizable-or, more
precisely, incompletely mechanizable, no matter how complex the machine involved.




Review of Alan Turing: The Enigma                                                  486

        Though on the surface this defeat for mechanism might seem to imply that
human reasoning can always outwit or transcend mechanical imitations, -©n deeper
analysis it turns out that Turing's argument can be applied to 'humans as well.
Consider the yes-no question, "Will your answer to this particular question be `no'?"
You will find that you too go into a sort of computational vertigo in trying to answer
it with a "yes" or a "no". This Question exemplifies the sort of undecidability
problems that Turing showed machines and mathematical systems are subject to.
Though the example is simplistic, it reminds us of an essential fact of the human
c3ndition-that people, no matter how aware they are of their minds, cannot fully take
their own complexity into account in attempting to understand themselves, and, quite
like Turing machines baffled by their own descriptions, may be plunged into a vertigo
of the psyche when they attempt to calculate their own hypothetical or future acts.
        Just as people can be surprised by their own complexity, so can machines, in
that they can't predict their own behavior. People attribute this feature of themselves
to "free will", and speak of "making choices". Turing's observation that machines will
go into endless loops when trying to predict their own behavior suggests that a
sufficiently complex machine might also cc' ne to suffer from that seemingly
inevitable human delusion: believing that. one has free will and is able to make
choices that transcend physical law.
        Thus Turing's seemingly negative result about machines can be seen as a
positive result, in that it sheds new light on how physical objects might reflect on
themselves and even consider themselves to be conscious, deliberating beings. A
mechanical approach to the mysteries of consciousness was Alan Turing's dream, and
probably by the late 1930's he was a believer in the possibility that a properly
organized machine could be intelligent, conscious, and have free will-at least to the
extent that we or `a y physical object cnN do so.

                                         * * *
       The war came as an interruption to young Turing's budding career as a
mathematical logician-he had by this time held fellowships at both Cambridge and
Princeton (at the Institute for Advanced Study, where he enjoyed such august
company as that of Einstein, Gödel, and John von Neumann)-and he was pressed into
service as a code breaker. It actually turned out not so badly for Turing, on a personal
level. At Bletchley Park, midway between Cambridge and Oxford, he and a small
cohort of powerful mathematical minds turned their powers of analysis to furthering
the already impressive work done by Polish codebreakers. It was known that the
German high command was sending orders in a code to its forces, including its vast
submarine network, by means of a machine called the "Enigma". exact construction of
the Enigma was known, but this was not enough, as Alan Turing so well knew: the
code breakers also needed to know the Machine's internal state, which could be any
one of an astronomical number.




Review of Alan Turing: The Enigma                                                   487

Any configuration of several independently turning rings in the enciphering machine
constituted a state; only when they knew that configuration could the code breakers
quickly decipher a message.
        Turing, Gordon Welchman, and a few others worked in close collaboration.
Together, they analyzed strategies for using intercepted coded messages and high-
speed searching machines to pinpoint the Enigma's state. Feverishly they worked as
British ships were sunk one after another by the German U-boats. It was clear that the
Nazis would bring Britain's war effort to an end unless the Enigma could be
outwitted.
        At first, they were able to decipher messages only a couple of weeks after
receiving them-obviously far too late. As they began to succeed, they reduced the gap
to a few days, then one day, and finally they reached the break-even point of decoding
messages in minutes. However, it then turned out that the Germans were referring to
places by special code names and unusual coordinates, so a second layer of decoding
was needed. Fortunately, this could be done by watching where ships actually were
sunk and correlating that information with the mysterious coordinates in the decoded
messages. Once the second layer of code had been peeled off, it was as if, all of a
sudden, the German fleet in the Atlantic were simply displayed on a screen in front of
them.
        There was an immediate dramatic increase in the number of British ships
getting through the U-boats' offensive network. To the Germans, this ought to have
been a dead giveaway that their code had been broken, but ironically, they were so
certain of the undecipherability of the Enigma machine that their own logic forced
them to conclude instead that the British must have very good spies, and so they
looked for the spies instead of inventing a new coding machine. There was
nonetheless a fragile, touch-and-go quality to the decipherment operation, because the
Germans would occasionally alter the Enigma machine in various ways, precipitating
desperate scrambles for new theories in Bletchley Park. But Turing and his associates
always came up with the theories, and the British government knew regularly and
with certainty what the Nazi command was up to.
        Meanwhile, the figure at the center of this activity, Alan Turing, was running
in long races and riding his dumpy bicycle to and from work, seemingly oblivious to
rain. People noticed that every so often he would stop to adjust his bicycle chain so
that it wouldn't fall off. Characteristically, he knew just when such a stop was needed,
for he had observed that his bicycle had an internal state, just like an Enigma
machine, determined by the relative positions of several independently turning gear
wheels in the mechanism. As long as he monitored the state of this "Turing machine",
he was able to forestall disaster. And meanwhile, on a more global scale, he was
forestalling disaster no less.

                                       *   *   *




Review of Alan Turing: The Enigma                                                   488

         When the war came to an end, Turing's ideas about how machines could
imitate the mind had matured considerably. He had been in contact for several years
now with electronic machinery, and many of the tasks he had been involved in had
fertilized his brain with new ideas. The problem now was that there was no longer any
war to make his abilities and his ideas seem crucial to anyone with money. He tried to
find funding to build his universal Turing machine, but his awkward ways with people
and his tendency to advocate long-term philosophical goals along with nearer-term
practical ones seemed to put people off. Rather than gaining respect, he became
known as something of an oddball. His powerful vision of the best way to go about
creating a universal machine, based on his deep preference that all flexibility come
from software (internal programs) rather than hardware, was gradually
circumnavigated, and he found himself left out in the cold. Eventually, a British
computer was built at Manchester University in the late 1940's, but not along the lines
Turing had advocated.
         Fortunately, while Turing was out of favor in the "proper" intellectual circles,
he was able to concentrate on philosophical issues connected with mechanical
thought, and in 1950, at age 38, he put his reflections into one of the classic articles on
that subject, entitled "Computing Machinery and Intelligence". In it he proposed what
has come to be known as the "Turing Test". The idea is to get around emotionally
charged questions like "Can a machine think?". Taking his cue from operationalism,
he replied in effect, "You want to know if that machine can think? Put it behind a
curtain and see if it can fool people into thinking it is human on the basis of what it
types to them." This has its parallel, interestingly, in the way some orchestra
conductors do auditions: they have each candidate stand behind a curtain, hidden, and
play from there, so they will not be swayed by age, sex, dress, or other external
aspects.
         The Turing Test (or "Imitation Game", as Turing called it) involved
communication between a human interrogator and an unknown languageusing
"being". Knowing that there was ferocious resistance to the image that computing
machinery might soon, or indeed, ever, think, Turing took pains to point out the
remarkable generality of the probing allowed by his test, by presenting a pair of short
sample dialogues in which it was shown how a skillful human interrogator might try
to elicit odd and recondite knowledge, subtle judgments, and even emotional
responses from the unknown "being". But most people remain skeptical about the
Turing Test even after- reading these dialogues, probably because they fear that they
might be easily taken in by the wiles of a superficial machine. They do not appreciate
how deeply and broadly the Turing Test potentially would allow them to probe.
         In his article, Turing raised nine plausible objections to his own Imitation
Game approach to the question of mechanical thought, and answered them cogently
one by one. The most serious one seems to be "Lady Lovelace's objection": that
computers cannot originate anything, but can do only what




Review of Alan Turing: The Enigma                                                      489

we explicitly tell them to do. Turing's answer to this-that one does not know what one
has programmed a machine to do, except in the most superficial and general way-has
a depth that eludes many good minds. I suspect that the Turing Test's profundity as an
examination of an alleged "thinking machine" will only gradually seep into our
culture as we collectively absorb the subtle and many-layered complexities of
computers.
        A sad footnote: In the early 1950's, the BBC recorded several radio interviews
with Turing on the subject of minds and machines, but for some reason did not
preserve any of them, and so we are left without a trace of a voice that, by all
accounts, was quite peculiar and revealing. Even though Turing's own Imitation Game
stressed the power of the printed word to convey all the nuances of personality, it
seems poignant to think that the voice of such a recent figure is forever lost, and all
we have to go on is the written word.

                                       *   *   *

        For his entire life, Alan Turing had been fascinated by the problem of
morphogenesis: how whole organisms synchronize and coordinate their growth. An
example is the fivefold symmetry of a starfish-how in the world does a cell know
what part of the organism it is in, and how do various cells communicate with one
another to plan the tricky overall pattern that they eventually wind up forming? It is as
if the card stunt section in a football stadium had to coordinate complex patterns
entirely by having nearest neighbors talk with each other. Turing's mathematically-
based theories, developed in the early 1950's, were typically ahead of their time and
even today they hold up well.
        His long-time enjoyment of long-distance running remained with him, and he
could look forward, so it might seem, to a happy life of pursuing his intellectual
dreams and his romantic hopes in a more peaceful world. Unfortunately, the Britain of
those days was as troubled politically as the United States, and homosexuality was
seen as a "dangerous" disease, symptomatic of mental instability. And ironically, just
as the anti-homosexual attack was getting more virulent, Alan Turing was becoming
increasingly courageous and vocal about his own sexual nature, often ignoring the
advice of friends to be more cautious.
        Turing's house was burglarized in 1952, and it was quickly clear to him that
one of his occasional lovers was somehow involved. In the course of making
depositions to the police, Turing indirectly revealed his homosexuality. Instantly, the
course of his life was irrevocably changed. No longer just a victim, he was now a
criminal in his own right-and rather than protest his innocence of the "crime" of
homosexuality, he talked freely about his "crime".
        At that time in Britain, there was a movement to look upon homosexuality as a
disease caused by hormone imbalances, and various physicians had proposed various
"cures". In America, castration had been a very popular




Review of Alan Turing: The Enigma                                                    490

"cure" for males (Hodges cites figures to the effect that by 1950, at least 50,000
castrations had been performed), but in Britain this was eschewed for less violent but
no less barbaric treatment. Turing was found guilty and sentenced to "treatment":
regular injections of female sex hormone, supposedly to quell his sex drive. This was
the way that British society thanked the person most responsible for the safety of its
ships during the World War. Of course, they had no way of knowing what role Turing
had taken during the war, since it was top secret and would remain so for many years
more. And in any case, Turing's wartime role should not have been seen as a
mitigating factor in his "crime", since that would have meant that the millions of other
more ordinary British homosexuals would have still been guilty of the same "crime".
Turing saw this, and did not want to try to use any of his connections in government
or the academic world to mitigate his sentence, and he simply endured it, growing
breasts and being rendered impotent.
        After one year, the sentence was over and he was free to return to a more
normal "state of affairs". But torment like that leaves permanent scars, and deep
inside Alan Turing something had changed. For the next couple of years, he appeared
for the most part quite happy to his friends, and he joked and chatted about his future.
But one day in 1954, he prepared a cyanide-coated apple, just as he had once seen the
Wicked Witch do in the, Walt Disney movie of Snow White and the Seven Dwarfs.
Unlike her, he bit into his own apple. "Dip the apple in the brew, Let the sleeping
death seep through." And he was found dead the next day. He had planned it in such a
way that his mother would interpret it as an "accident with chemicals", but others
knew better. Although today all evidence strongly suggests that the machine known as
Alan Mathison Turing halted itself of its own free will, the ultimate reason remains an
enigma to us, an undecidable question.

                                    *       *       *

        Andrew Hodges has painted, in his book, a beautiful portrait of a multifaceted
man whose honesty and decency were too much for his society and his times, and
who brought about his own downfall. Beyond the evident empathy that Hodges feels
for Turing, there is another level of depth and understanding, one that makes all the
difference in a biography of a scientific figure: scientific accuracy. Hodges has done
an admirable job of presenting each idea in detail to the lay reader, but moreover, it is
obvious that he is passionately intrigued by all the ideas that fascinated Turing. This
book is a first-rate presentation of the life of a first-rate scientific mind, and because
this particular mind was attached to a body that had a mind of its own, the book is a
very important document for social reasons as well. Alan Turing would have
shuddered if he had ever known that his life story would be made public, but he is in
good hands: it is hard to imagine a more thoughtful and warm portrait of a life than
this one.




Review of Alan Turing: The Enigma                                                     491


                    A Coffeehouse Conversation
                        on the Turing Test
                                      May, 1981



                            Participants in the dialogue:
    Chris, a physics student; Pat, a biology student; Sandy, a philosophy student.

CHRIS: Sandy, I want to thank you for suggesting that I read Alan Turing's article
   "Computing Machinery and Intelligence". It's a wonderful piece and certainly
   made me think-and think about my thinking.
SANDY: Glad to hear it. Are you still as, much of a skeptic about artificial
   intelligence as you used to be?
CHRIS: You've got me wrong. I'm not against artificial intelligence; I think it's
   wonderful stuff perhaps a little crazy, but why not? I simply am convinced that
   you Al advocates have far underestimated the human mind, and that there are
   things a computer will never, ever be able to do. For instance, can you imagine a
   computer writing a Proust novel? The richness of imagination, the complexity of
   the characters
SANDY: Rome wasn't built in a day!
CHRIS: In the article, Turing comes through as an interesting person. Is he still alive?
SANDY: No, he died back in 1954, at just 41. He'd be only 70 or so now, although he
   is such a legendary figure it seems strange to think that he could still be living
   today.
CHRIS: How did he die?
SANDY: Almost certainly suicide. He was homosexual, and had to deal with some
   pretty barbaric treatment and stupidity from the outside world. In the end, it got
   to be too much, and he killed himself.
CHRIS: That's horrendous, especially in this day and age.
SANDY: I know. What really saddens me is that he never got to see the amazing
   progress in computing machinery and theory that has taken place since 1954. Can
   you imagine how he'd have been wowed?




A Coffeehouse Conversation on the Turing Test                                        492

CHRIS: Yeah...
PAT: Hey, are you two going to clue me in as to what this Turing article is t?
SANDY: It is really about two things. One is the question "Can a machine think?"-or
   rather, "Will a machine ever think?" The way Turing answers the question-he
   thinks the answer is yes, by the way-is by batting down a series of objections to
   the idea, one after another. The other point he tries to make is that, as it stands,
   the question is not meaningful. It's too full of emotional connotations. Many
   people are upset by the suggestion that people are machines, or that machines
   might think. Turing tries to defuse the question by casting it in less emotional
   terms. For instance, what do you think, Pat, of the idea of thinking machines?
PAT: Frankly, I find the term confusing. You know what confuses me? It's those ads
   in the newspapers and on TV that talk about "products that "intelligent ovens" or
   whatever. I just don't know how seriously to take them.
SANDY: I know the kind of ads you mean, and they probably confuse a lot people.
   On the one hand, we're always hearing the refrain "Computers are really dumb;
   you have to spell everything out for them in words of one syllable"-yet on the
   other hand, we're constantly bombarded with advertising hype about "smart
   products".
CHRIS: That's certainly true. Do you know that one company has even taken to
   calling its products "dumb terminals" in order to stand out from the crowd?
SANDY: That's a pretty clever gimmick, but even so it just contributes to the trend
   toward obfuscation. The term "electronic brain" always comes to mind when I'm
   thinking about this. Many people swallow it com y, and others reject it out of
   hand. It takes patience to sort out the issues and decide how much of it makes
   sense.
PAT: Does Turing suggest some way of resolving it, some kind of IQ test for -
   machines?
SANDY: That would be very interesting, but no machine could yet come close taking
   an IQ test. Instead, Turing proposes a test that theoretically be applied to any
   machine to determine whether or not it can think. . Does the test give a clear-cut
   yes-or-no answer? I'd be skeptical if it claimed to.
SANDY: No, it doesn't claim to. In a way that's one of its advantages. It shows how
   the borderline is quite fuzzy and how subtle the whole question is. PAT; And so,
   as usual in philosophy, it's all just a question of words!
SANDY: Maybe, but they're emotionally charged words, and so it's important, it
   seems to me, to explore the issues and try to map out the meanings of the crucial
   words. The issues are fundamental to our concept of ourselves, so we shouldn't
   just sweep them under the rug.
PAT: Okay, so tell me how Turing's test works.




A Coffeehouse Conversation on the Turing Test                                      493

SANDY: The idea is based on what he calls the Imitation Game. Imagine that a man
   and a woman go into separate rooms, and from there they can be interrogated by
   a third party via some sort of teletype set-up. The third party can address
   questions to either room, but has no idea which person is in which room. For the
   interrogator, the idea is to determine which room the woman is in. The woman,
   by her answers, tries to help the interrogator as much as she can. The man,
   though, is doing his best to bamboozle the interrogator, by responding as he
   thinks a woman might. And if he succeeds in fooling the interrogator ...
PAT: The interrogator only gets to see written words, eh? And the sex of the author is
   supposed to shine through? That game sounds like a good challenge. I'd certainly
   like to take part in it someday. Would the interrogator have met either the man or
   the woman before the test began? Would any of them know any of the others?
SANDY: That would probably be a bad idea. All kinds of subliminal cueing might
   occur if the interrogator knew one or both of them. It would certainly be best if
   all three people were totally unknown to one another.
PAT: Could you ask any questions at all, with no holds barred?
SANDY: Absolutely. That's the whole idea!
PAT: Don't you think, then, that pretty quickly it would degenerate into sex-oriented
   questions? I mean, I can imagine the, man, overeager to act convincing, giving
   away the game by answering some very blunt questions that most women would
   find too personal to answer, even through an anonymous computer connection.
SANDY: That's a nice observation. I wonder if it's true ...
CHRIS: Another possibility would be to probe for knowledge of minute aspects of
   traditional sex-role differences, by asking about such things as dress sizes and so
   on. The psychology of the Imitation Game could get pretty subtle. I suppose
   whether the interrogator was a woman or a man would make a difference. Don't
   you think that a woman could spot some telltale differences more quickly than a
   man could?
PAT: If so, maybe the best way to tell a man from a woman is to let each of them play
   interrogator in an Imitation Game, and see which of the two is better at telling a
   man from a woman!
SANDY: Hmm ... that's a droll twist. Oh, well. I don't know if this original version of
   the Imitation Game has ever been seriously tried out, despite the fact that it
   would be relatively easy to do with modern computer terminals. I have to admit,
   though, that I'm not at all sure what it would prove, whichever way it turned out.
PAT: I was wondering about that. What would it prove if the interrogator say a
   woman-couldn't tell correctly which person was the woman? It certainly wouldn't
   prove that the man was a woman!
SANDY: Exactly! What I find funny is that although I strongly believe in the idea of
   the Turing Test, I'm not so sure I understand the point of its basis, the Imitation
   Game.




A Coffeehouse Conversation on the Turing Test                                      494

CHRIS: As for me, I'm not any happier with the Turing Test as a test for thinking
   machines than I am with the Imitation Game as a test for femininity.
PAT: From what you two are saying, I gather the Turing Test is some kind of
   extension of the Imitation Game, only involving a machine and a person instead
   of a man and a woman.
SANDY: That's the idea. The machine tries its hardest to convince the interrogator
   that it is the human being, and the human tries to make it clear that he or she is
   not the computer.
PAT: The machine tries ? Isn't that a loaded way of putting it?
SANDY: Sorry, but that seemed the most natural way to say it.
PAT: Anyway, this test sounds pretty interesting. But how do you know that it will
   get at the essence of thinking? Maybe it's testing for the wrong things. Maybe,
   just to take a random illustration, someone would feel that a machine was able to
   think only if it could dance so well that you couldn't tell it was a machine. Or
   someone else could suggest some other characteristic. What's so sacred about
   being able to fool people by typing at them?
SANDY: I don't see how you can say such a thing. I've heard that objection before,
   but frankly, it baffles me. So what if the machine can't tap-dance or drop a rock
   on your toe? If it can discourse intelligently on any subject you want, then it has
   shown that it can think-to me, at least! As I see it, Turing has drawn, in one clean
   stroke, a clear division between thinking and other aspects of being human.
PAT: Now you're the baffling one. If you couldn't conclude anything from a man's
   ability to win at the Imitation Game, how could you conclude anything from a
   machine's ability to win at the Turing Game?
CHRIS: Good question.
SANDY: It seems to me that you could conclude something from a man's win in the
   Imitation Game. You wouldn't conclude he was a woman, but you could
   certainly say he had-good insights into the feminine mentality (if there is such a
   thing). Now, if a computer could fool someone into thinking it was a person, I
   guess you'd have to say something similar about it-that it had good insights into
   what it's like to be human, into "the human condition" (whatever that is).
PAT: Maybe, but that isn't necessarily equivalent to thinking, is it? It seems to me that
   passing the Turing Test would merely prove that some machine or other could do
   a very good job of simulating thought.
CHRIS: I couldn't agree more with Pat. We all know that fancy computer programs
   exist today for simulating all sorts of complex phenomena. In theoretical physics,
   for instance, we simulate the behavior of particles, atoms, solids, liquids, gases,
   galaxies, and so on. But no one confuses any of those simulations with the real
   thing!
SANDY: In his book Brainstorms, the philosopher Daniel Dennett makes a similar
   point about simulated hurricanes.




A Coffeehouse Conversation on the Turing Test                                        495

CHRIS: That's a nice example, too. Obviously, what goes on insid, computer when it's
   simulating a hurricane is not a hurricane, for machine's memory doesn't get torn
   to bits by 200-mile-an-hour win the floor of the machine room doesn't get
   flooded with rainwater, and on.
SANDY: Oh, come on-that's not a fair argument! In the first place, i programmers
   don't claim the simulation really is a hurricane. It's mer a simulation of certain
   aspects of a hurricane. But in the second pla you're pulling a fast one when you
   imply that there are no downpours 200-mile-an-hour winds in a simulated
   hurricane. To us there aren't ai but if the program were incredibly detailed, it
   could include simulat people on the ground who would experience the wind and
   the rain ji as we do when a hurricane hits. In their minds-or, if you'd rather, their
   simulated minds-the hurricane would be not a simulation, buy genuine
   phenomenon complete with drenching and devastation.
CHRIS: Oh, my-what a science-fiction scenario! Now we're talking abc simulating
   whole populations, not just a single mind!
SANDY: Well, look-I'm simply trying to show you why your argument th a simulated
   McCoy isn't the real McCoy is fallacious. It depends on tl tacit assumption that
   any old observer of the simulated phenomenon equally able to assess what's
   going on. But in fact, it may take an observ with a special vantage point to
   recognize what is going on. In td hurricane case, it takes special "computational
   glasses" to see the ra and the winds.
PAT: "Computational glasses"? I don't know what you're talking about.
SANDY: I mean that to see the winds and the wetness of the hurricane, yc have to be
   able to look at it in the proper way. You
CHRIS: No, no, no! A simulated hurricane isn't wet! No matter how mu( it might
   seem wet to simulated people, it won't ever be genuinely wet! An no computer
   will ever get torn apart in the process of simulating wind
SANDY: Certainly not, but that's irrelevant. You're just confusing levels. Th laws of
   physics don't get torn apart by real hurricanes, either. In the cap of the simulated
   hurricane, if you go peering at the computer's memor expecting to find broken
   wires and so forth, you'll be disappointed. Bi look at the proper level. Look into
   the structures that are coded for i memory. You'll see that many abstract links
   have been broken, man values of variables radically changed, and so on. There's
   your flood, yoI devastation-real, only a little concealed, a little hard to detect.
CHRIS: I'm sorry, I just can't buy that. You're insisting that I look for a ne, kind of
   devastation, one never before associated with hurricanes. Th: way you could call
   anything a hurricane as long as its effects, seen throug your special "glasses",
   could be called "floods and devastation".
SANDY: Right-you've got it exactly! You recognize a hurricane by its effect You
   have no way of going in and finding some ethereal "essence c hurricane", some
   "hurricane soul" right in the middle of the storm's ey4 Nor is there any ID card to
   be found that certifies "hurricanehood". It's




A Coffeehouse Conversation on the Turing Test                                       496

   just the existence of a certain kind of pattern-a spiral storm with an eye and so
   forth-that makes you say it's a hurricane. Of course, there are a lot of things you'll
   insist on before you call something a hurricane.
PAT: Well, wouldn't you say that being an atmospheric phenomenon is one
   prerequisite? How can anything inside a computer be a storm? To me, a
   simulation is a simulation is a simulation!
SANDY: Then I suppose you would say that even the calculations computers do are
   simulated-that they are fake calculations. Only people can do genuine
   calculations, right?
PAT: Well, computers get the right answers, so their calculations are not exactly fake-
   but they're still just patterns. There's no understanding going on in there. Take a
   cash register. Can you honestly say that you feel it is calculating something when
   its gears mesh together? And the step from cash register to computer is very
   short, as I understand things.
SANDY: If you mean that a cash register doesn't feel like a schoolkid doing
   arithmetic problems, I'll agree. But is that what "calculation" means? Is that an
   integral part of it? If so, then contrary to what everybody has thought up till now,
   we'll have to write a very complicated program indeed to perform genuine
   calculations. Of course, this program will sometimes get careless and make
   mistakes, and it will sometimes scrawl its answers illegibly, and it will
   occasionally doodle on its paper ... It won't be any more reliable than the store
   clerk who adds up your total by hand. Now, I happen to believe that eventually
   such a program could be written. Then we'd know something about how clerks
   and schoolkids work.
PAT: I can't believe you'd ever be able to do that!
SANDY: Maybe, maybe not, but that's not my point. You say a cash register can't
   calculate. It reminds me of another favorite passage of mine from Dennett's
   Brainstorms. It goes something like this: "Cash registers can't really calculate;
   they can only spin their gears. But cash registers can't really spin their gears,
   either; they can only follow the laws of physics." Dennett said it originally about
   computers; I modified it to talk about cash registers. And you could use the same
   line of reasoning in talking about people: "People can't really calculate; all they
   can do is manipulate mental symbols. But they aren't really manipulating
   symbols; all they are doing is firing various neurons in various patterns. But they
   can't really make their neurons fire; they simply have to. let the laws of physics
   make them fire for them." Et cetera. Don't you see how this reductio ad absurdum
   would lead you to conclude that calculation doesn't exist, that hurricanes don't
   exist-in fact, that nothing at a level higher than particles and the laws of physics
   exists? What do you gain by saying that a computer only pushes symbols around
   and doesn't truly calculate?
PAT: The example may be extreme, but it makes my point that there is a vast
   difference between a real phenomenon and any simulation of it. This is so for
   hurricanes, and even more so for human thought.
SANDY: Look, I don't want to get too tangled up in this line of argument,




A Coffeehouse Conversation on the Turing Test                                        497

   but let me try one more example. If you were a radio ham listening to another
   ham broadcasting in Morse code and you were responding in Morse code, would
   it sound funny to you to refer to "the person at the other end"?
PAT: No, that would sound okay, although the existence of a person at the other end
   would be an assumption.
SANDY: Yes, but you wouldn't be likely to go and check it out. You're prepared to
   recognize personhood through those rather unusual channels. You don't have to
   see a human body or hear a voice. All you need is a rather abstract manifestation-
   a code, as it were. What I'm getting at is this. To "see" the person behind the dits
   and dahs, you have to be willing to do some decoding, some interpretation. It's
   not direct perception; it's indirect. You have to peel off a layer or two to find the
   reality hidden in there. You put on your "radio-ham's glasses" to "see" the person
   behind the buzzes. Just the same with the simulated hurricane! You don't see it
   darkening the machine room; you have to decode the machine's memory. You
   have to put on special "memory-decoding" glasses. Then what you see is a
   hurricane.
PAT: Oh, ho ho! Talk about fast ones-wait a minute! In the case of the shortwave
   radio, there's a real person out there, somewhere in the Fiji Islands or wherever.
   My decoding act as I sit by my radio simply reveals that that person exists. It's
   like seeing a shadow and concluding there's an object out there, casting it. One
   doesn't confuse the shadow with the object, however! And with the hurricane
   there's no real storm behind the scenes, making the computer follow its patterns.
   No, what you have is just a shadow-hurricane without any genuine hurricane. I
   just refuse to confuse shadows with reality.
SANDY: All right. I don't want to drive this point into the ground. I even admit it is
   pretty silly to say that a simulated hurricane is a hurricane. But I wanted to point
   out that it's not as silly as you might think at first blush. And when you turn to
   simulated thought, then you've got a very different matter on your hands from
   simulated hurricanes.
PAT: I don't see why. You'll have to convince me.
SANDY: Well, to do so, I'll first have to make a couple of extra points about
   hurricanes.
PAT: Oh, no! Well, all right, all right.
SANDY: Nobody can say just exactly what a hurricane is-that is, in totally precise
   terms. There's an abstract pattern that many storms share, and it's for that reason
   we call those storms hurricanes. But it's not possible to make a sharp distinction
   between hurricanes and non-hurricanes. There are tornados, cyclones, typhoons,
   dust devils ... Is the Great Red Spot on Jupiter a hurricane? Are sunspots
   hurricanes? Could there be a hurricane in a wind tunnel? In a test tube? In your
   imagination, you can even extend the concept of "hurricane" to include a
   microscopic storm on the surface of a neutron star.
CHRIS: That's not so far-fetched, you know. The concept of "earthquake" has actually
   been extended to neutron stars. The astrophysicists say that the tiny changes in
   rate that once in a while are observed in the pulsing of a pulsar are caused by
   "glitches"-starquakes-that have just occurred' on the neutron star's surface.
SANDY: Oh, I remember that now. That "glitch" idea has always seemed eerie to me-
   a surrealistic kind of quivering on a surrealistic kind of surface.
CHRIS: Can you imagine-plate tectonics on a giant sphere of pure nuclear matter?
SANDY: That's a wild thought. So, starquakes and earthquakes can both be subsumed
   into a new, more abstract category. And that's how science constantly extends


A Coffeehouse Conversation on the Turing Test                                       498

   familiar concepts, taking them further and further from familiar experience and
   yet keeping some essence constant. The number system is the classic example-
   from positive numbers to negative numbers, then rationals, reals, complex
   numbers, and "on beyond zebra", as Dr. Seuss says.
PAT: I think I can see your point, Sandy. In biology, we have many examples of close
   relationships that are established in rather abstract ways. Often the decision about
   what family some species belongs to comes down to an abstract pattern shared at
   some level. Even the concepts of "male" and "female" turn out to be surprisingly
   abstract and elusive. When you base your system of classification on very
   abstract patterns, I suppose that a broad variety of phenomena can fall into "the
   same class", even if in many superficial ways the class members are utterly
   unlike one another. So perhaps I can glimpse, at least a little, how to you, a
   simulated hurricane could, in a funny sense, be a hurricane.
CHRIS: Perhaps the word that's being extended is not "hurricane", but "be"
PAT: How so?
CHRIS: If Turing can extend the verb "think", can't I extend the verb "be"? All I mean
   is that when simulated things are deliberately confused with genuine things,
   somebody's doing a lot of philosophical wool-pulling. It's a lot more serious than
   just extending a few nouns, such as "hurricane".
SANDY: I like your idea that "be" is being extended, but I sure don't agree with you
   about the wool-pulling. Anyway, if you don't object, let me just say one more
   thing about simulated hurricanes and then I'll get to simulated minds. Suppose
   you consider a really deep simulation of a hurricane-I mean a simulation of every
   atom, which I admit is sort of ridiculous, but still, just consider it for the sake of
   argument.
PAT: Okay.
SANDY: I hope you would agree that it would then share all the abstract structure
   that defines the "essence of hurricanehood". So what's to keep you from calling it
   a hurricane?
PAT: I thought you were backing off from that claim of equality.




A Coffeehouse Conversation on the Turing Test                                        499

SANDY: So did I, but then these examples came up, and I was forced back to my
   claim. But let me back off, as I said I would do, and get back to thought, which is
   the real issue here. Thought, even more than hurricanes, is an abstract structure, a
   way of describing some complex events that happen in a medium called a brain.
   But actually, thought can take place in any one of several billion brains. There
   are all these physically very different brains, and yet they all support "the same
   thing": thinking. What's important, then, is the abstract pattern, not the medium.
   The same kind of swirling can happen inside any of them, so no person can claim
   to think more "genuinely" than any other. Now, if we come up with some new
   kind of medium in which the same style of swirling takes place, could you deny
   that thinking is taking place in it?
PAT: Probably not, but you have just shifted the question. The question now is: How
   can you determine whether the "same style" of swirling is really happening?
SANDY: The beauty of the Turing Test is that it tells you when! Don't you see?
CHRIS: I don't see that at all. How would you know that the same style of activity
   was going on inside a computer as inside my mind, simply because it answered
   questions as I do? All you're looking at is its outside.
SANDY: I'm sorry, I disagree entirely! How do you know that when I speak to you,
   anything similar to what you call thinking is going on inside me? The Turing
   Test is a fantastic probe, something like a particle accelerator in physics. Here,
   Chris-I think you'll like this analogy. Just as in physics, when you want to
   understand what is going on at an atomic or subatomic level, since you can't see
   it directly, you scatter accelerated particles off a target and observe their
   behavior. From this, you infer the internal nature of the target. The Turing Test
   extends this idea to the mind. It treats the mind as a "target" that is not directly
   visible but whose structure can be deduced more abstractly. By "scattering"
   questions off a target mind, you learn about its internal workings, just as in
   physics.
CHRIS: Well ... to be more exact, you can hypothesize about what kinds of internal
   structures might account for the behavior observed-but please remember that they
   may or may not in fact exist.
SANDY: Hold on, now! Are you suggesting that atomic nuclei are merely
   hypothetical entities? After all, their existence (or should I say hypothetical
   existence?) was proved (or should I say suggested?) by the behavior of particles
   scattered off atoms.
CHRIS: I would agree, but you know, physical systems seem to me to be much
   simpler than the mind, and the certainty of the inferences made is
   correspondingly greater. And the conclusions are confirmed over and over again
   by different types of experiments.
SANDY: Yes, but those experiments still are of the same sort-scattering, detecting
   things indirectly. You can never handle an electron or a quark. Physics
   experiments are also correspondingly harder to do and to interpret. Often they
   take years and years, and dozens of collaborators' are involved. In the Turing
   Test, though, just one person could perform many highly delicate experiments in
   the course of no more than an hour. I maintain that people give other people
   credit for being conscious simply because of their continual external monitoring
   of other people-which is itself something like a Turing Test.
PAT: That may be roughly true, but it involves more than just conversing with people
   through a teletype. We see that other people have bodies, we watch their faces
   and expressions-we see they are human beings, and so we think they think.


A Coffeehouse Conversation on the Turing Test                                      500

SANDY: To me, that seems a narrow, anthropocentric view of what thought is. Does
   that mean you would sooner say a mannequin in a store thinks than a wonderfully
   programmed computer, simply because the mannequin looks more human?
PAT: Obviously I would need more than just vague physical resemblance to ,the
   human form to be willing to attribute the power of thought to an entity. But that
   organic quality, the sameness of origin, undeniably lends a degree of credibility
   that is very important.
SANDY: Here we disagree. I find this simply too chauvinistic. I feel that the key
   thing is a similarity of internal structure-not bodily, organic, chemical structure
   but organizational structure-software. Whether an entity can think seems to me a
   question of whether its organization can be described in a certain way, and I'm
   perfectly willing to believe that the Turing Test detects the presence or absence
   of that mode of organization. I would say that your depending on my physical
   body as evidence that I am a thinking being is rather shallow. The way I see it,
   the Turing Test looks far deeper than at mere external form.
PAT: Hey, now-you're not giving me much credit. It's not just the shape of a body
   ,that lends weight to the idea that there's real thinking going on inside. It's also, as
   I said, the idea of common origin. It's the idea that you and I both sprang from
   DNA molecules, an idea to which I attribute much depth. Put it this way: the
   external form of human bodies reveals that they share a deep biological history,
   and it's that depth that lends a lot of credibility to the notion that the owner of
   such a body can think.
SANDY: But that is all indirect evidence. Surely you want some direct evidence.
   That's what the Turing Test is for. And I think it's the only way to test for
   thinkinghood.
CHRIS: But you could be fooled by the Turing Test, just as an interrogator could
   mistake a man for a woman.
SANDY: I admit, I could be fooled if I carried out the test in too quick or too shallow
   a way. But I would go for the deepest things I could think of.
CHRIS: I would want to see if the program could understand jokes-or better yet, make
   them! That would be a real test of intelligence..
SANDY: I agree that humor probably is an acid test for a supposedly intelligent
   program, but equally important to me-perhaps more so




A Coffeehouse Conversation on the Turing Test                                          501

   would be to test its emotional responses. So I would ask it about its reactions to
   certain pieces of music or works of literature-especially my favorite ones.
CHRIS: What if it said, "I don't know that piece", or even, "I have no interest in
   music"? What if it tried its hardest (oops!-sorry, Pat!) ... Let me try that again.
   What if it did everything it could, to steer clear of emotional topics and
   references?
SANDY: That would certainly make me suspicious. Any consistent pattern of
   avoiding certain issues would raise serious doubts in my mind as to whether I
   was dealing with a thinking being.
CHRIS: Why do you say that? Why not just conclude you're dealing with a thinking
   but unemotional being?
SANDY: You've hit upon a sensitive point. I've thought about this for quite a long
   time, and I've concluded that I simply can't believe emotions and thought can be
   divorced. To put it another way, I think emotions are an automatic by-product of
   the ability to think. They are entailed by the very nature of thought.
CHRIS: That's an interesting conclusion, but what if you're wrong? What if I
   produced a machine that could think but not emote? Then its intelligence might
   go unrecognized because it failed to pass your kind of test.
SANDY: I'd like you to point out to me where the boundary line between emotional
   questions and non-emotional ones lies. You might want to ask about the meaning
   of a great novel. This certainly requires an understanding of human emotions!
   Now is that thinking, or merely cool calculation? You might want to ask about a
   subtle choice of words. For that, you need an understanding of their connotations.
   Turing uses examples like this in his article. You might want to ask for advice
   about a complex romantic situation. The machine would need to know a lot about
   human motivations and their roots. If it failed at this kind of task, I would not be
   much inclined to say that it could think. As far as I'm concerned, thinking,
   feeling, and consciousness are just different facets of one phenomenon, and no
   one of them can be present without the others.
CHRIS: Why couldn't you build a machine that could feel nothing (we all know
   machines don't feel anything!), but that could think and make complex decisions
   anyway? I don't see any contradiction there.
SANDY: Well, I do. I think that when you say that, you are visualizing a metallic,
   rectangular machine, probably in an air-conditioned room-a hard, angular, cold
   object with a million colored wires inside it, a machine that sits stock still on a
   tiled floor, humming or buzzing or whatever, and spinning its tapes. Such a
   machine can play a good game of chess, which, I freely admit, involves a lot of
   decision-making. And yet I would never call it conscious.
CHRIS: How come? To mechanists, isn't a chess-playing machine rudimentarily
   conscious?
SANDY: Not to this mechanist! The way I see it, consciousness has got to come from
   a precise pattern of organization, one we haven't yet figured out how to describe
   in any detailed way. But I believe we will gradually come to understand it. In my
   view, consciousness requires a certain way of mirroring the external universe
   internally, and the ability to respond to that external reality on the basis of the
   internally represented model. And then in addition, what's really crucial for a
   conscious machine is that it should incorporate a well-developed and flexible
   self-model. And it's there that all existing programs, including the best chess-
   playing ones, fall down.



A Coffeehouse Conversation on the Turing Test                                      502

CHRIS: Don't chess programs look ahead and say to themselves as they're figuring
   out their next move, "If my opponent moves here, then I'll go there, and then if
   they go this way, I could go that way ..."? Doesn't that usage of the concept "I"
   require a sort of self-model?
SANDY: Not really. Or, if you want, it's an extremely limited one. It's an
   understanding of self in only the narrowest sense. For instance, a chess-playing
   program has no concept of why it is playing chess-, or of the fact that it is a
   program, or is in a computer, or has a human opponent. It has no idea about what
   winning and losing are, or
PAT: How do you know it has no such sense? How can you presume to say what a
   chess program feels or knows?
SANDY: Oh, come on! We all know that certain things don't feel anything or know
   anything. A thrown stone doesn't know anything about parabolas, and a whirling
   fan doesn't know anything about air. It's true I can't prove those statements-but
   here, we are verging on questions of faith.
PAT: This reminds me of a Taoist story I read. It goes something like this. Two sages
   were standing on a bridge over a stream. One said to the other, "I wish I were a
   fish. They are so happy." The other replied, "How do you know whether fish are
   happy or not? You're not a fish!" The first said, "But you're not me, so how do
   you know whether I know how fish feel?"
SANDY: Beautiful! Talking about consciousness really does call for a certain amount
   of restraint. Otherwise, you might as well just jump on the solipsism bandwagon
   ("I am the only conscious being in the universe") or the panpsychism bandwagon
   ("Everything in the universe is conscious! ").
PAT: Well, how do you know? Maybe everything is conscious.
SANDY: Oh, Pat, if you're going to join the club that maintains that stones and even
   particles like electrons have some sort of consciousness, then I guess we part
   company here. That's a kind of mysticism I just can't fathom. As for chess
   programs, I happen to know how they work, and I can tell you for sure that they
   aren't conscious. No way!
PAT: Why not?
SANDY: They incorporate only the barest knowledge about the goals of chess. The
   notion of "playing" is turned into the mechanical act of comparing a lot of
   numbers and choosing the biggest one over and over




A Coffeehouse Conversation on the Turing Test                                    503

   again. A chess program has no sense of disappointment about losing, or pride in
   winning. Its self-model is very crude. It gets away with doing the least it can, just
   enough to play a game of chess and nothing more. Yet interestingly enough, we
   still tend to talk about the "desires" of a chess-playing computer. We say, "It
   wants to keep its king behind a row of pawns" or "It likes to get its rooks out
   early" or "It thinks I don't see that hidden fork".
PAT: Yes, and we do the same thing with insects. We spot a lonely ant somewhere
   and say, "It's trying to get back home" or "It wants to drag that dead bee back to
   the colony". In fact, with any animal we use terms that indicate emotions, but we
   don't know for certain how much the animal feels. I have no trouble talking about
   dogs and cats being happy or sad, having desires and beliefs and so on, but of
   course I don't think their sadness is as deep or complex as human sadness is.
SANDY: But you wouldn't call it "simulated" sadness, would you?
PAT: No, of course not. I think it's real.
SANDY: It's hard to avoid use of such teleological or mentalistic terms. I believe
   they're quite justified, although they shouldn't be carried too far. They simply
   don't have the same richness of meaning when applied to present-day chess
   programs as when applied to people.
CHRIS: I still can't see that intelligence has to involve emotions. Why couldn't you
   imagine an intelligence that simply calculates and has no feelings?
SANDY: A couple of answers here. Number one, any intelligence has to have
   motivations. It's simply not the case, whatever many people may think, that
   machines could think any more "objectively" than people do. Machines, when
   they look at a scene, will have to focus and filter that scene down into some
   preconceived categories, just as a person does. And that means seeing some
   things and missing others. It means giving more weight to some things than to
   others. This happens on every level of processing.
PAT: I'm not sure I'm following you.
SANDY: Take me right now, for instance. You might think I'm just making some
   intellectual points, and I wouldn't need emotions to do that. But what makes me
   care about these points? Just now-why did I stress the work "care" so heavily?
   Because I'm emotionally involved in this conversation! People talk to each other
   out of conviction-not out of hollow, mechanical reflexes. Even the most
   intellectual conversation is driven by underlying passions. There's an emotional
   undercurrent to every conversation-it's the fact that the speakers want to be
   listened to, understood, and respected for what they are saying.
PAT: It sounds to me as if all you're saying is that people need to be interested in
   what they're saying, otherwise a conversation dies.
SANDY: Right! I wouldn't bother to talk to anyone if I weren't motivated by interest.
   And "interest" is just another name for a whole constellation of subconscious
   biases. When I talk, all my biases work together, and what




A Coffeehouse Conversation on the Turing Test                                       504

   you perceive on the surface level is my personality, my style. But that style arises
   from an immense number of tiny priorities, biases, leanings. When you add up a
   million of them interacting together, you get something that amounts to a lot of
   desires. It just all adds up! And that brings me to the other answer to Chris'
   question about feelingless calculation. Sure, that exists-in a cash register, a
   pocket calculator. I'd say it's even true of all today's computer programs. But
   eventually, when you put enough feelingless calculations together in a huge
   coordinated organization, you'll get something that has properties on another
   level. You can see it in fact, you have to see it-not as a bunch of little calculations
   but as a system of tendencies and desires and beliefs and so on. When things get
   complicated enough, you're forced to change your level of description. To some
   extent that's already happening, which is why we use words such as "want",
   "think", "try", and "hope" to describe chess programs and other attempts at
   mechanical thought. Dennett calls that kind of level-switch by the observer
   "adopting the intentional stance". The really interesting things in AI will only
   begin to happen, I'd guess, when the program itself adopts the intentional stance
   toward itself!
CHRIS: That would be a very strange sort of level-crossing feedback loop.
SANDY: It certainly would. When a program looks at itself from the outside, as it
   were, and tries to figure out why it acted the way it did, then I'll start to think that
   there's someone in there, doing the looking.
PAT: You mean an "I"? A self?
SANDY: Yes, something like that. A soul, even-although not in any religious sense.
   Of course, it's highly premature for anyone to adopt the intentional stance (in the
   full force of the term) with respect to today's programs. At least that's my
   opinion.
CHRIS: For me an important related question is: To what extent is it valid to adopt
   the intentional stance toward beings other than humans?
PAT: I would certainly adopt the intentional stance toward mammals.
SANDY: I vote for that.
CHRIS: Now that's interesting. How can that be, Sandy? Surely you wouldn't clairrr
   that a dog or cat can pass the Turing Test? Yet don't you maintain the Turing
   Test is the only way to test for the presence of consciousness? How can you have
   these beliefs simultaneously?
SANDY: Hmm ... All right. I guess that my argument is really just that the Turing
   Test works only above a certain level of consciousness. I'm perfectly willing to
   grant that there can be thinking beings that could fail at the Turing Test-but the
   main point that I've been arguing for is that anything that passes it would be a
   genuinely conscious, thinking being.
PAT: How can you think of a computer as a conscious being? I apologize if what I'm
   going to say sounds like a stereotype, but when I think of conscious beings, I just
   can't connect that thought with machines. To me, consciousness is connected
   with soft, warm bodies, silly though it may sound.
CHRIS: That does sound odd, coming from a biologist. Don't you deal with




A Coffeehouse Conversation on the Turing Test                                          505

   life so much in terms of chemistry and physics that all magic seems to vanish?
PAT: Not really. Sometimes the chemistry and physics simply increase the feeling
   that there's something magical going on down there! Anyway, I can't always
   integrate my scientific knowledge with my gut feelings. CHRIS: I guess I share
   that trait.
PAT: So how do you deal with rigid preconceptions like mine?
SANDY: I'd try to dig down under the surface of your concept of "machine" and get
   at the intuitive connotations that lurk there, out of sight but deeply influencing
   your opinions. I think we all have a holdover image from the Industrial
   Revolution that sees machines as clunky iron contraptions gawkily moving under
   the power of some loudly chugging engine. Possibly that's even how the
   computer inventor Charles Babbage saw people! After all, he called his
   magnificent many-geared computer the "Analytical Engine".
PAT: Well, I certainly don't think people are just fancy steam shovels or electric can
   openers. There's something about people, something that that-they've got a sort of
   flame inside them, something alive, something that flickers unpredictably,
   wavering, uncertain-but something creative !
SANDY: Great! That's just the sort of thing I wanted to hear. It's very human to think
   that way. Your flame image makes me think of candles, of fires, of vast
   thunderstorms with lightning dancing all over the sky in crazy, tumultuous
   patterns. But do you realize that just that kind of thing is visible on a computer's
   console? The flickering lights form amazing chaotic sparkling patterns. It's such
   a far cry from heaps of lifeless, clanking metal! It is flamelike, by God! Why
   don't you let the word "machine" conjure up images of dancing patterns of light
   rather than of giant steam shovels?
CHRIS: That's a beautiful image, Sandy. It does tend to change my sense of
   mechanism from being matter-oriented to being pattern-oriented. It makes me try
   to visualize the thoughts in my mind-these thoughts right now, even!-as a huge
   spray of tiny pulses flickering in my brain.
SANDY: That's quite a poetic self-portrait for a mere spray of flickers to have come
   up with!
CHRIS: Thank you. But still, I'm not totally convinced that a machine is all that I am.
   I admit, my concept of machines probably does suffer from anachronistic
   subconscious flavors, but I'm afraid I can't change such a deeply rooted sense in a
   flash.
SANDY: At least you sound open-minded. And to tell the truth, part of me
   sympathizes with the way you and Pat view machines. Part of me balks at calling
   myself a machine. It is a .bizarre thought that a feeling being like you or me
   might emerge from mere circuitry. Do I surprise you?
CHRIS: You certainly surprise me. So, tell us-do you believe in the idea of an
   intelligent computer, or don't you?




A Coffeehouse Conversation on the Turing Test                                      506

SANDY: It all depends on what you mean. We've all heard the question "Can
   computers think?" There are several possible interpretations of this (aside from
   the many interpretations of the word "think"). They revolve around different
   meanings of the words "can" and "computer".
PAT: Back to word games again ...
SANDY: I'm sorry, but that's unavoidable. First of all, the question might mean,
   "Does some present-day computer think, right now?" To this I would
   immediately answer with a loud no. Then it could be taken to mean, "Could some
   present-day computer, if suitably 'programmed, potentially think?" That would
   be more like it, but I would still answer, "Probably not". The real difficulty
   hinges on the word "computer". The way I see it, "computer" calls up an image
   of just what I described earlier: an air-conditioned room with cold rectangular
   metal boxes in it. But I suspect that with increasing public familiarity with
   computers and continued progress in computer architecture, that vision will
   eventually become outmoded.
PAT: Don't you think computers as we know them will be around for a while?
SANDY: Sure, there will have to be computers in today's image around for a long
   time, but advanced computers-maybe no longer called "computers"-will evolve
   and become quite different. Probably, as with living organisms, there will be
   many branchings in the evolutionary tree. There will be computers for business,
   computers for schoolkids, computers for scientific calculations, computers for
   systems research, computers for simulation, computers for rockets going into
   space, and so on. Finally, there will be computers for the study of intelligence.
   It's really only these last that I'm thinking of-the ones with the maximum
   flexibility, the ones that people are deliberately attempting to make smart. I see
   no reason that these will stay fixed in the traditional image. They probably will
   soon acquire as standard features some rudimentary sensory systems-mostly for
   vision and hearing, at first. They will need to be able to move around, to explore.
   They will have to be physically flexible. In short, they will have to become more
   animal-like, more self-reliant.
CHRIS: It makes me think of the robots R2D2 and C3PO in the movie Star Wars.
SANDY: Not me! In fact, I don't think of anything remotely like them when I
   visualize intelligent machines. They are too silly, too much the product of a film
   designer's imagination. Not that I have a clear vision of my own. But I think it's
   necessary, if people are realistically going to try to imagine an artificial
   intelligence, to go beyond the limited, hard-edged picture of computers that
   comes from exposure to what we have today. The only thing all machines will
   always have in common is their underlying mechanicalness. That may sound cold
   and inflexible, but then just think -what could be more mechanical, in a
   wonderful way, than the workings of the DNA and proteins and organelles in our
   cells?




A Coffeehouse Conversation on the Turing Test                                     507

PAT: To me, what goes on inside cells has a "wet", "slippery" feel to it, and what
   goes on inside machines is dry and rigid. It's connected with the fact that
   computers don't make mistakes, that computers do only what you tell them to do.
   Or at least that's my image of computers.
SANDY: Funny-a minute ago, your image was of a flame, and now it's of something
   wet and slippery. Isn't it marvellous, how contradictory we can be?
PAT: I don't need your sarcasm.
SANDY: No, no, I'm not being sarcastic-I really do think it's marvellous. PAT: It's
   just an example of the human mind's slippery nature-mine, in this case.
SANDY: True. But your image of computers is stuck in a rut. Computers certainly
   can make mistakes-and I don't mean on the hardware level. Think of any present-
   day computer predicting the weather. It can make wrong predictions, even
   though its program runs flawlessly.
PAT: But that's only because you've fed it the wrong data.
SANDY: Not so. It's because weather prediction is too complex. Any such program
   has to make do with a limited amount of data-entirely correct data-and
   extrapolate from there. Sometimes it will make wrong predictions. It's no
   different from a farmer gazing at the clouds and saying, "I reckon we'll get a little
   snow tonight." In our heads, we make models of things and use those models to
   guess how the world will behave. We have to make do with our models, however
   inaccurate they may be, or evolution will prune us out ruthlessly-we'll fall off a
   cliff or something. And for intelligent computers, it'll be the same. It's just that
   human designers will speed up the evolutionary process by aiming explicitly at
   the goal of creating intelligence, which is something nature just stumbled on.
PAT: So you think computers will be making fewer mistakes as they get smarter?
SANDY: Actually, just the other way around! The smarter they get, the more they'll
   be in a position to tackle messy real-life domains, so they'll be more and more
   likely to have inaccurate models. To me, mistake-making is a sign of high
   intelligence!
PAT: Wow-you throw me sometimes!
SANDY: I guess I'm a strange sort of advocate for machine intelligence. To some
   degree I straddle the fence. I think that machines won't really be intelligent in a
   humanlike way until they have something like your biological wetness or
   slipperiness to them. I don't mean literally wet-the slipperiness could be in the
   software. But biological-seeming or not, intelligent machines will in any case be
   machines. We will have designed them, built them-or grown them! We'll
   understand how they work-at least in some sense. Possibly no one person will
   really understand them, but collectively we will know how they work.
PAT: It sounds like you want to have your cake and eat it too. I mean,




A Coffeehouse Conversation on the Turing Test                                       508

   you want to have people able to build intelligent machines and yet at the same
   time have some of the mystery of mind remain.
SANDY: You're absolutely right-and I think that's what will happen. When real
   artificial intelligence comes
PAT: Now there's a nice contradiction in terms!
SANDY: Touche! Well, anyway, when it comes, it will be mechanical and yet at the
   same time organic. It will have that same astonishing flexibility that we see in
   life's mechanisms. And when I say mechanisms, I mean mechanisms. DNA and
   enzymes and so on really are mechanical and rigid and reliable. Wouldn't you
   agree, Pat?
PAT: Sure! But when they work together, a lot of unexpected things happen. There
   are so many complexities and rich modes of behavior that all that mechanicalness
   adds up to something very fluid.
SANDY: For me, it's an almost unimaginable transition from the mechanical level of
   molecules to the living level of cells. But it's that exposure to biology that
   convinces me that people are machines. That thought makes me uncomfortable in
   some ways, but in other ways it is exhilarating.

CHRIS: I have one nagging question ... If people are machines, how come it's so hard
   to convince them of the fact? Surely a machine ought to be able to recognize its
   own machinehood!
SANDY: It's an interesting question. You have to allow for emotional factors here. To
   be told you're a machine is, in a way, to be told that you're nothing more than
   your physical parts, and it brings you face to face with your own vulnerability,
   destructibility, and, ultimately, your mortality. That's something nobody finds
   easy to face. But beyond this emotional objection, to see yourself as a machine,
   you have to "unadopt" the intentional stance you've grown up taking toward
   yourself-you have to jump all the way from the level where the complex lifelike
   activities take place to the bottom-most mechanical level where ribosomes chug
   along RNA strands, for instance. But there are so many intermediate layers that
   .they act as a shield, and the mechanical quality way down there becomes almost
   invisible. I think that when intelligent machines come around, that's how they
   will seem to us-and to themselves! Their mechanicalness will be buried so deep
   that they'll seem to be alive and conscious just as we seem alive and conscious ...
CHRIS: You're baiting me! But I'm not going to bite.
PAT: I once heard a funny idea about what will happen when we eventually have
   intelligent machines. When we try to implant that intelligence into devices we'd
   like to control, their behavior won't be so predictable. SANDY: They'll have a
   quirky little "flame" inside, maybe? PAT: Maybe.
CHRIS: And what's so funny about that?
PAT:. Well, think of military missiles. The more sophisticated their target-tracking
   computers get, according to this idea, the less predictably they will function.
   Eventually, you'll have missiles that will decide they are




A Coffeehouse Conversation on the Turing Test                                     509

   pacifists and will turn around and go home and land quietly without blowing up.
   We could even have "smart bullets" that turn around in midflight because they
   don't want to commit suicide!
SANDY: What a nice vision!
CHRIS: I'm very skeptical about all this. Still, Sandy, I'd like to hear your predictions
   about when intelligent machines will come to be.
SANDY: It won't be for a long time, probably, that we'll see anything remotely
   resembling the level of human intelligence. It rests on too awesomely
   complicated a substrate-the brain-for us to be able to duplicate it in the
   foreseeable future. Anyhow, that's my opinion.
PAT: Do you think a program will ever pass the Turing Test?
SANDY: That's a pretty hard question. I guess there are various degrees of passing
   such a test, when you come down to it. It's not black and white. First of all, it
   depends on who the interrogator is. A simpleton might be totally taken in by
   some programs today. But secondly, it depends on how deeply you are allowed
   to probe.
PAT: You could have a range of Turing Tests-one-minute versions, five minute
   versions, hour-long versions, and so forth. Wouldn't it be interesting if some
   official organization sponsored a periodic competition, like the annual computer-
   chess championships, for programs to try to pass the Turing Test?
CHRIS: The program that lasted the longest against some panel of distinguished
   judges would be the winner. Perhaps there could be a big prize for the first
   program that fools a famous judge for, say, ten minutes.
PAT: A prize for the program, or for its author?
CHRIS: -For the program, of course!
PAT: That's ridiculous! What would a program do with a prize?
CHRIS: Come now, Pat. If a program's human enough to fool the judges, don't you
   think it's human enough to enjoy the prize? That's precisely the threshold where
   it, rather than its creators, deserves the credit, and the rewards. Wouldn't you
   agree?
PAT: Yeah, yeah-especially if the prize is an evening out on the town, dancing with
   the interrogators!
SANDY: I'd certainly like to see something like that established. I think it could be
   hilarious to watch the first programs flop pathetically!
PAT: You're pretty skeptical for an AI advocate, aren't you? Well, do you think any
   computer program today could pass a five-minute Turing Test, given a
   sophisticated interrogator?
SANDY: I seriously doubt it. It's partly because no one is really working at it
   explicitly. I should mention, though, that there is one program whose inventors
   claim it has already passed a rudimentary version of the Turing Test. It is called
   "Parry", and in a series of remotely conducted interviews, it fooled several
   psychiatrists who were told they were talking to either a computer or a paranoid
   patient. This was an improvement over an earlier version, in which psychiatrists
   were simply handed




A Coffeehouse Conversation on the Turing Test                                        510

FIGURE 22-1. In (a), a program is enjoying the great reward for passing the Turing
   Test: an evening out on the town, dancing with the interrogator. Can the reader
   spot the program? In (b), an interrogator is enjoying the great reward for
   successfully unmasking an unthinking robot: an evening out on the town, dancing
   with the robot. Can the robot spot the reader? NOTE: One of these two
   photographs was not taken by David,J. Moser. Can the interrogator tell which
   one?

   transcripts of short interviews and asked to determine which ones were with a
   genuine paranoid and which ones were with a computer simulation.
PAT: You mean they didn't have the chance to ask any questions? That's a severe
   handicap-and it doesn't seem in the spirit of the Turing Test. Imagine someone
   trying to tell which sex I belong to, just by reading a transcript of a few remarks
   by me. It might be very hard! I'm glad the procedure has been improved.




A Coffeehouse Conversation on the Turing Test                                     511

CHRIS: How do you get a computer to act like a paranoid?
SANDY: Now just a moment-I didn't say it does act like a paranoid, only that some
   psychiatrists, under unusual circumstances, thought so. One of the things that
   bothered me about this pseudo-Turing Test is the way Parry works. "He", as the
   people who designed it call it, acts like a paranoid in that "he" gets abruptly
   defensive and veers away from undesirable topics in the conversation. In effect,
   Parry maintains strict control so that no one can truly probe "him". For reasons
   like this, simulating a paranoid is a whole lot easier than simulating a normal
   person.
PAT: I wouldn't doubt that. It reminds me of the joke about the easiest kind of human
   being for a computer program to simulate.
CHRIS: What is that?
PAT: A catatonic patient-they just sit and do nothing at all for days on end. Even I
   could write a computer program to do that!
SANDY: An interesting thing about Parry is that it creates no sentences on its own-it
   merely selects from a huge repertoire of canned sentences the one that in some
   sense responds best to the input sentence.
PAT: Amazing. But that would probably be impossible on a larger scale, wouldn't it?
SANDY: You better believe it (to use a canned remark)! Actually, this is something
   that's really not appreciated enough. The number of sentences you'd need to store
   in order to be able to respond in a normal way to all possible turns that a
   conversation could take is more than astronomical-it's really unimaginable. And
   they would have to be so intricately indexed, for retrieval ... Anybody who thinks
   that somehow, a program could be rigged up just to pull sentences out of storage
   like records in a jukebox, and that this program could pass the Turing Test, hasn't
   thought very hard about it. The funny part is that it is just this kind of
   unrealizable "parrot program" that most critics of artificial intelligence cite, when
   they argue against the concept of the Turing Test. Instead of imagining a truly
   intelligent machine, they want you to envision a gigantic, lumbering robot that
   intones canned sentences in a dull monotone. They set up the imagery in a
   contradictory way. They manage to convince you that you could see through to
   its mechanical level with ease, even as it is simultaneously performing tasks that
   we think of as fluid, intelligent processes. Then the critics say, "You see! A
   machine could pass the Turing Test and yet it would still be just a mechanical
   device, not intelligent at all." I see things almost the opposite way. If I were
   shown a machine that can do things that I can do-I mean pass the Turing Test-
   then, instead of feeling insulted or threatened, I'd chime in with philosopher
   Raymond Smullyan and say, "How wonderful machines are!"
CHRIS: If you could ask a computer just one question in the Turing Test, what would
   it be?
SANDY: Uhmm ...




A Coffeehouse Conversation on the Turing Test                                       512

PAT: How about this: "If you could ask a computer just one question in the Turing
   Test, what would it be?"?


Post Scriptum.
     In 1983, I had the most delightful experience of getting to know a small group of
extremely enthusiastic and original students at the University of Kansas in Lawrence.
These students, about thirty in number, had been drawn together by Zamir Bavel, a
professor in the Computer Science Department, who had organized a seminar on my
book Gödel, Escher, Bach. He contacted me and asked me if there was any chance I
could come to Lawrence and get together with his students. Something about his way
of describing what was going on convinced me that this was a very unusual group and
that it would be worth my while to try it out. I therefore made a visit to Kansas and
got to know both Zamir and his group. All my expectations were met and surpassed.
The students were full of ideas and. warmth and made me feel very much at home.
The first trip was so successful that I decided to do it again a couple of months later.
This time they threw an informal party at an apartment a few of them shared. Zamir
had forewarned me that they were hoping to give me a demonstration of something
that had already been done in a recent class meeting. It seems that the question of
whether computers could ever think had arisen, and most of the group members had
taken a negative stand on the issue. Rod Ogborn, the student who had been leading
the discussion, had asked the class if they would consider any of the following
programs intelligent:

  (1) a program that could pass a course in beginning programming (i.e., that
     could take informal descriptions of tasks and turn them into good working
     programs);
  (2) a program that could act like a psychotherapist (Rod gave sample dialogues
     with the famous "Doctor" program, also known as "Eliza", by Joseph
     Weizenbaum);
  (3) a program called "Boris", written at Yale by Michael Dyer, that could read
     stories in a limited domain and answer questions about the situation which
     required filling in many unstated assumptions, and making inferences of
     many sorts based on them.

       The class had come down on the "no" side of all three of these cases, although
they got progressively harder. So Rod, to show the class how difficult this decision
might be if they were really faced with a conversational program, managed to get a
hookup over the phone lines with a natural




A Coffeehouse Conversation on the Turing Test                                       513

language program called "Nicolai" that had been developed over the last few years by
the Army at nearby Fort Leavenworth. Thanks to some connections that Rod had, the
class was able to gain access to an unclassified version of Nicolai and to interact with
it for two or three hours. At the end of those hours, they then reconsidered the
question of whether a computer might be able to think. Still, only one student was
willing to consider Nicolai intelligent, and even that student reserved the right to
switch sides if more information came in. About half the others were noncommittal,
and the rest were unwilling, under any circumstances, to call Nicolai intelligent. There
was no doubt that Rod's demonstration had been effective, though, and the class
discussion had been one of the most lively.
        Zamir told me all of this on our drive into Lawrence from the Kansas City
airport, and he explained that the group had been so stimulated by this experience that
they were hoping to get reconnected to Nicolai over the phone lines, and to let me try
it out during the party. I thought it sounded quite amusing, and since I have tried out
and watched a lot of natural-language programs in my time, I thought I would have an
easy time coming up with good probes into the weaknesses of Nicolai. "Besides," I
thought (somewhat chauvinistically), "how good can a program be that is being
developed by the Army in an out-of-the-way place like Fort Leavenworth? To reveal
its mechanicalness ought to be a piece of cake for me!" So I was looking forward to
the challenge.
        When we first arrived at the party, everyone was drinking beer and munching
on various chips and dips and talking excitedly about what had happened in the
previous class meeting with Nicolai. There was some consternation, though, because
it seemed that the connection to Nicolai was hard to re-establish. It had to go through
a computer at the University of Kansas, and thence to Fort Leavenworth, all of which
meant that it was a slow and tenuous link. But after about 45 minutes of trying to set
up the link, someone announced with pleasure that Nicolai was responding. One of
the students was sitting at the terminal and typing in various simple questions and
getting back simple answers. I watched over his shoulder, and felt the itch to sit down
and try to unmask Nicolai for what "he" really was. After a couple of minutes, the
student relinquished the terminal and I got my chance, with a small group of people
looking over my shoulder as I typed. Nicolai's answers were quite slow in coming
over the phone lines, and the wait was made all the more suspenseful by the fact that
they were transmitted only in full lines of text, so there were some pretty long delays
between successive lines of response. But that didn't seriously dampen my enthusiasm
for trying to quickly poke holes in the program's facade of understanding.
        Herewith follows a somewhat abridged but accurate transcript of the
"conversation" I had with Nicolai over the next hour or so. For convenience's sake,
what Nicolai typed is shown in capitals, and my inner thoughts at the time (as best I
can reconstruct them now) are in italics.




A Coffeehouse Conversation on the Turing Test                                       514

What would you like to talk about?
(Let's start out easy.)
I CAN TALK ABOUT ANYTHING.
(Yeah, yeah-tell it to the Marines!)
Is there anything you specialize in?
WHAT DOES "SPECIALIZE" MEAN? PLEASE USE A COMPLETE SENTENCE.
(For a supposedly smart program, Nicolai doesn't seem to have a very big
vocabulary.)
"Specialize" means "to be good at".
(I wonder if Nicolai likes those quotation marks, or if they will throw it for a loop.)
I UNDERSTAND.
(Izzatso?)
When were you written?
LAST COMPILATION: "4/27/83 19.247"
(Oh, look at that stupid, rigid answer, complete with pointless quotation marks and
that extra, unexplained number at the end. I'm exposing Nicolai pretty well, so far ...)
What language is Nicolai written in?
(I think I'll see if it knows its own name and recognizes third person references to
itself.)
NICOLAI IS WRITTEN MOSTLY IN ADA.
(Hmm ... Not bad. It handled both "you " and "Nicolai " when I talked about it. Those
Army programmers aren't so bad!)
What are feet?
(Choose an ambiguous word, see what it comes up with ... )
FEET ARE 12.0 INCHES.
(Wow, weird! A person would never put it that way. They, d say, `Afoot is a unit of
length, twelve inches long", or something like that. And anyway, feet are primarily
parts of the body, secondarily units of length. Why didn't it even seem to recognize
that ambiguity? Oh, well. I guess that kind of dumbness is the kind of thing I wanted
to expose. Score one for me.)
What are ears?
(Why not try another ambiguous word? After all, I did well on the last one.)
EARS ARE AUDITORY ORGANS FOUND ON ANIMALS.
(This is really funny! Why were feet not thought of primarily as body parts, then? This
body-parts stuff seems to give me lots of points-I think I'll try it again.)
What are arms?
THAT INFORMATION IS CLASSIFIED.




A Coffeehouse Conversation on the Turing Test                                       515

(Oh, no! That is too stupid for words! Or is this somebody's idea of a joke?)
I like to understand up without seventeen clams.
(Let's toss some nonsense at Nicolai and see what it does. )
WHAT DO CLAMS HAVE TO DO WITH IT?
(Oh, come on, Nicolai! Do you mean you understood all the rest, except for the clams
bit? You are really dumb!)
Why are you so dumb?
I AM A SMART PROGRAM.
(Well, that's kind of touching, actually. A little self-defense is appealing, in a way.)
You are pretty confident of yourself, aren't you?
WHAT DOES "PRETTY CONFIDENT" MEAN?
(Oh, yeah, I should have known- "pretty " as an adverb might be too colloquial. )
Did you understand up without my last sentence?
(I figure, since earlier it didn't complain about my phrase "understand up ', let's see
what it does with it this time.)
YOUR SENTENCE IS NONSENSE.
(Oh, really? Well, then, why wasn't the other one with that same phrase in it also
considered nonsense?)
Do you enjoy nonsense?
IF YOU INSIST ON TALKING NONSENSE, I HAVE OTHER TASKS TO
PERFORM.
(A strange response. Who ever insisted on talking nonsense? I was just asking a
simple question. Oh, well, let me try out another angle.)
What is your opinion of M. C. Escher?
M. C. ESCHER, ARTIST, TWENTIETH CENTURY
(Oh, not again-another vacuous dictionary look-up. Can't it do better than that,
especially when I've asked it a probing question about taste? Let's see how it handles
a complaint ... )
You did not answer my question.
WHICH QUESTION DID I NOT ANSWER?
(Hey, this Nicolai is a pretty slick character. Not so different from Parry. 1 must say, I
am really pretty impressed by the robustness of this program. I mean, we have
covered a lot of ground so far, and at least it makes 'pseudo-sensible answers " most
of the time. Not at all bad!)
Do you own any money?
I DO NOT OWN ANYTHING.




A Coffeehouse Conversation on the Turing Test                                         516

(Fair enough.)
Have you seen the latest issue of "Scientific American"?
(Uh-oh--once again, I don't know if I'm going to help it or confuse it by putting in
those quotes. Let's see.)
I DO NOT HAVE ANY "SCIENTIFIC AMERICAN".
(That sounds awkward. Very computerish. I doubt if it knows what I meant. Oh, well.)
YOU'VE BEEN TALKING TO ME FOR QUITE A WHILE. WHAT IS YOUR
NAME?
(Hmm ... That was an unexpected remark by Nicolai. Should I answer it? Sure, why
not?)
My name is "Doug".
DO YOU HAVE A LAST NAME?
(Inquisitive little devil, isn't it?)
Of course I do.
(Ask a stupid question, get a stupid answer. Let's see how it fields this one.)
WOULD YOU CARE TO TELL ME WHAT YOUR LAST NAME IS?
(Damn, it just totally ignored the tone of my remark. And why is it being so polite,
suddenly?)
Only if you'll tell me yours.
(That ought to challenge it.)
I DO NOT HAVE A LAST NAME.
(I don't seem to be able to foil its defenses very easily. It always walks around me.
Those programmers sure anticipated a lot of ideas. I guess they must have debugged
Nicolai for years, watching typical dialogues with people who just come in off the
street. For a dumb program, Nicolai is pretty smart!)
Well then, I won't tell you mine.
THAT'S NOT BEING VERY FRIENDLY.
(You know, sometimes, I hate to admit it, this stupid program really does seem to
understand something!)
I'm sorry. How would you like me to give you a last name?
IF YOU WOULD LIKE TO DO THAT, THAT WOULD BE FINE WITH ME.
(Now that is obviously a canned sentence, or at least the latter half must certainly be.
Score one more point for me. And yet, I'm not doing nearly as well as I'd hoped I
would be, by now. What's wrong?)
Your last name is "Rimsky-Korsakov".
THANK YOU.




A Coffeehouse Conversation on the Turing Test                                       517

(Well, that's a strange kind of thing to thank somebody for. Just goes to show how the
rigidity of programs makes them say the strangest things.)
Don't mention it.
(I'll play along with this game for a second.)
WHAT IS YOUR LAST NAME?
(My god, it remembered, after all this while, that as soon as it had told me its last
name, I was going to tell it mine. Well, I d better make good on that promise.)
My last name is "Hofstadter".
ARE YOU "(32) DOUGLAS R. HOFSTADTER"?
(Now why, in the midst of so much elegance, does it suddenly act so crude, putting my
name in quotes and including that silly extraneous index number? I can 't for the life
of me figure out why this program is so mixed in quality.)
One and the same.
(Try out a strange idiomatic phrase.)
(83) "PARSING ERROR IN PREVIOUS INPUT"
(Aha fooled it again. All of a sudden, a whole lot of mechanicalness is showing
through.)
How many letters are there in your last name?
(I think I'll see whether it can figure out ambiguous words by context.)
14 PLUS A HYPHEN.
(Good. Now I'll let "letters" switch meaning. What will it do?)
How many letters are there in an ordinary envelope?
18.
(Oh, no! It didn't recognize the shift! And it merely counted the number of letters in
"an ordinary envelope"! But if that's the case, then how come in the earlier case it
didn't count the letters in the phrase "your last name'; but instead retrieved the
referent? There's something so inconsistent and uneven about this program. It just
goes to show how far we have to go before mechanical intelligence bears any
resemblance to human fluidity.)
ARE YOU WORKING ON A NEW BOOK?
(I wonder what will happen if I simply sidestep that question. Anyway, I d like to
probe its ability to deal with humor.)
Tell me a joke.
I WANT A COOOKIE.
(What kind of a joke is that? Hey, it spelled "cookie" wrong! Hmm ... is that its joke?
Or what is going on here?)

By the way, I am sort of working on a new book. Care to hear about it?

(Can it handle weird phrases like "sort of "? And questions with no subject?)




A Coffeehouse Conversation on the Turing Test                                      518

NO, BUT IF YOU WOULD GIVE ME A COOKIE, I WOULD GLADLY PAY
YOU TUESDAY.
(That sounds like a familiar line. Have I heard it somewhere before? Is that supposed
to be the punch line of its 'joke "?)
Tell me another one.
HOW MANY COMPUTER PROGAMMERS DOES IT TAKE TO SCREW IN A
LIGHT BULB?
(Okay, I see. It's got a library of jokes and riddles, and it just fishes them out at
random. I'll play along.)
4.
INCORRECT. THAT IS A HARDWARE PROBLEM.
(Okay, so it told its joke. But I bet it has no idea what any of the terms in it mean, such
as "programmer" . Hey, it left an 'r' out of "programmer"! I guess whoever input the
joke misspelled the word. So let's see what happens when I ask about it ...
What is a computer progammer?
WOULD YOU LIKE ME TO EXPLAIN THE JOKE?
(Hey, damn it, it turned the tables on me! Does it think I don't know what a '
progammer" (or a programmer) is? I don't like this slipperiness. Every time I try to
jet it to act really dumb, it somehow slips out of it. Strange ... All right, let me probe
some more about its understanding of the terms in that joke. How about a potentially
off-color word?)
What does "to screw" mean?
IN WHAT CONTEXT?
(Aha, at least it knows it's an ambiguous term.)
In bed.
"TO SCREW IN BED" MEANS "TO HAVE SEX".
(Okay, fine, it has some definitions stored according to context. I guess this word must
come up all the time in conversations. But can it go any further than that definition?)
What is having sex all about?
I CAN RECOMMEND AN EXCELLENT BOOK ON THAT SUBJECT.
(What a great stock answer! One point for the Army programmers!)
Tell me, tell me!
THE BOOK IS CALLED "GODEL, ESCHER, BACH".
(Huh? Say, wail a minute, this is too weird. I'm not so sure this is really a program.
There's too damn much cleverness in these weird responses. I can't imagine bizarre
mechanical devices coming up with this much coherence, even if it is completely off
the wall, at times.)
I like to call it "GEB", myself.
THAT IS VERY NICE, DOUG. WHAT ELSE DO YOU CALL "GEB"?
(All right, the jigs up. What's going on here?)




A Coffeehouse Conversation on the Turing Test                                          519

        At this point, everybody started laughing. They had all been trying hard to
keep from doing so the whole time, but at this point, they couldn't contain themselves.
It was time to reveal to me what had been going on. They took me downstairs and
showed me that three of the students-Roy Leban, Howard Darsche, and Garrick
Stangle-had been collectively playing Nicolai. There was in reality no such program,
and there hadn't ever been one. (In retrospect, I am reminded of the famous French
mathematician Nicolas Bourbaki-a hypothetical person, actually an amalgam of over
a dozen eminent mathematicians writing under that collective pen name.) There had
indeed been a similar demonstration for the class a few days earlier, and the class, like
me, had been taken in for a long time. In my case, Roy, Howard, and Garrick had
worked very hard to give the impression of mechanicalness by spewing back "parsing
error" and other indications of rigidity, and also by sending what looked very much
like canned phrases from time to time. That way, they could keep sophisticates like
me believing that there was a program behind it all. Only by that point I was
beginning to wonder just how sophisticated I really was.
        The marvelous thing about this game is that it was, in many ways, a Turing
Test in reverse: a group of human beings masquerading as a program, trying to act
mechanical enough that I would believe it really was one. Hugh Kenner has written a
book called The Counterfeiters about the perennial human fascination with such
compounded role flips. A typical example is Delibes' ballet Coppelia, in which human
dancers imitate life-sized dolls stiffly imitating people. What is amusing is how
Nicolai's occasional crudeness was just enough to keep me convinced it was
mechanical. Its "willingness" to talk about itself, combined with its obvious
limitations along those lines (its clumsy revelation of when it was last compiled, for
instance), helped establish the illusion very strongly.

                                       *   *   *

        In retrospect, I am quite amazed at how much genuine intelligence I was
willing to accept as somehow having been implanted in the program. I had been
sucked into the notion that there really must be a serious natural-language effort going
on at Fort Leavenworth, and that there had been a very large data base developed,
including all sorts of random information: a dictionary, a catalogue containing names
of miscellaneous people, some jokes, lots of canned phrases to use in difficult
situations, some self-knowledge, a crude ability to use key words in a phrase when it
can't parse it exactly, some heuristics for deciding when nonsense is being foisted on
it, some deductive capabilities, and on and on. In hindsight, it is clear that I was
willing to accept a huge amount of fluidity as achievable in this day and age simply
by putting together a large bag of isolated tricks kludges and hacks, as they say.




A Coffeehouse Conversation on the Turing Test                                        520

       Roy Leban, one of the three inside Nicolai's mind, wrote the following about
the experience of being at the other end of the exchange:

   Nicolai was a split personality. The three of us (as well as many kibitzers)
   argued about practically every response. Each of us had a strong preconceived
   notion about what (or who) Nicolai should be. For example, I felt that certain
   things (such as "Douglas R. Hofstadter") should be in quotation marks, and that
   feet should not be 12 inches, but 12.0. Howard had a tendency for rather flip
   answers. It was he who suggested the "classified" response to the "arms"
   question. And somehow, when he suggested it, we all knew it was right.

   Several times during our conversation, I felt quite amazed at how fluently Nicolai
was able to deal with things I was bringing up, but each time I could postulate some
not too sophisticated mechanical underpinning that would allow that particular thing
to happen. As a strong skeptic of true fluidity in machines at this time, I kept on trying
to come up with rationalizations for the fact that this program was doing so well. My
conclusion was that it was a very vast and quite sophisticated bag of tricks, no one of
which was terribly complex. But after a while, it just became too much to believe.
Furthermore, the mixture of crudity and subtlety became harder and harder to
swallow, as well.
   My strategy had been, in essence, to use spot checks all over the map: to try to
probe it in all sorts of ways rather than to get sucked into some topic of its own
choice, where it could steer the conversation. Daniel Dennett, in a paper on the depth
of the Turing Test called "Can Machines Think?", likens this technique to a strategy
taught to American soldiers in World War II for telling German spies from genuine
Yankees. The idea was that even if a young man spoke absolutely fluent American-
sounding English, you could trip him up by asking him things that any boy growing
up in those days would be expected to know, such as "What is the name of Mickey
Mouse's girlfriend?" or "Who won the World Series in 1937?" This expands the
domain of knowledge necessary from just the language itself to the entire culture-and
the amazing thing is that just a few well-placed questions can unmask a fraud in a
very brief time-or so it would seem.
   The problem is, what do you do if the person is extremely sharp, and when asked
about Minnie Mouse, responds in some creative way, such as, "Hah! She ain't no
girlfriend-she's a mouse !"? The point is that even with these trick probes that should
ferret out frauds very swiftly, there can be clever defensive countermaneuvers, and
you can't be sure of getting to the bottom of things in a very brief time.
   It seems that a few days earlier, the class had collectively gone through something
similar to what I had just gone through, with one major difference. Howard Darsche,
who had impersonated (if I may use that peculiar choice of words!) Nicolai in the first
run-through, simply had acted himself, without trying to feign mechanicalness in any
way. When asked




A Coffeehouse Conversation on the Turing Test                                         521

what color the sky was, he replied, "In daylight or at night?" and when told "At
night", he replied, "Dark purple with stars." He got increasingly poetic and creative in
his responses to the class, but no one grew suspicious that this Nicolai was a fraud. At
some point, Rod Ogborn simply had to stop the demonstration and type on the screen,
"Okay, Howard, you can come in now." Zamir (who was not in cahoots with Rod and
his team) was the only one who had some reluctance in accepting this performance as
that of a genuine program, and he had kept silent until the end, when he voiced a
muted skepticism.
   Zamir summarizes this dramatic demonstration by saying that his class was willing
to view anything on a video terminal as mechanically produced, no matter how
sophisticated, insightful, or poetic an utterance it might be. They might find it
interesting and even surprising, but they would find some way to discount those
qualities. Why was this the case? How could they do this for so long? And why did I
fall for the same kind of thing?
   In interacting with me, Nicolai had seemed to waver between crude
mechanicalness and subtle flexibility, an oscillation I had found most puzzling and
somewhat disturbing. But I was still taken in for a very long time. It seems that, even
armed with spot checks and quite a bit of linguistic sophistication and skepticism,
unsuspecting humans can have the wool pulled over their eyes for a good while. This
was the humble pie I ate in this remarkable reverse Turing Test, and I will always
savor its taste and remember Nicolai with great fondness.

                                        *   *   *

   Alan Turing, in his article, indicated that his "Imitation Game" test should take
place through some sort of remote teletype linkup, but one thing he did not indicate
explicitly was at what grain size the messages would be transmitted. By that, I mean
that he did not say whether the messages should be transmitted as intact wholes, or
line by line, word by word, or keystroke by keystroke. Although I don't think it
matters for the Turing Test in any fundamental sense, I do think that which type of
"window" you view another language-using being through has a definite bearing on
how quickly you can make inferences about that being. Clearly, the most revealing of
these possibilities is that of watching the other "person" operate at the keystroke level.
   On most multi-user computer systems, there are various ways for different users to
communicate with each other, and these ways reflect different levels of urgency. The
slowest one is generally the "mail" facility, through which you can send another user
an arbitrarily long piece of text, just like a letter in an envelope. When it arrives, it
will be placed in the user's "mailbox", to be read at their leisure. A faster style of
communicating is called, on Unix systems, "write". When this is invoked, a direct
communications link is set up between you and the person you are trying to reach
(provided they are




A Coffeehouse Conversation on the Turing Test                                         522

logged on). If they accept your link, then any full line typed by either of you will be
instantly transmitted and printed on the other party's screen-where a lineful is signaled
by your hitting the carriage-return key. This is essentially what the Nicolai team used
in communicating with me over the Kansas computer. Their irregular typing rhythm
and any errors they might have made were completely concealed from me this way,
since all I saw was a sequence of completely polished lines (with the two spelling
errors "coookie" and "progammer", which I was willing to excuse because Nicolai
generated them in a "joke" context).
    The most revealing mode is what, on Unix, is called "talk". In this mode, every
single keystroke is revealed. You make an error, you are exposed. For some people,
this is too much like living in a glass house, and they prefer the shielding afforded by
"write". For my part, I like living dangerously. Let the mistakes lfy! In computer-
mediated conversations with my friends, I always opt for "talk". I have been amused
to watch their "talk" styles and my own slowly evolve to relatively stable states.
    When we in the Indiana University Computer Science Department first began
using the "talk" facility, we were all somewhat paranoid about making errors, and we
would compulsively fix any error that we made. By this, I mean that we would
backspace and retype the character. The effect on the screen of hitting the- backspace
key repeatedly is that you see the most recently typed characters getting eaten up, one
by one, right to left, and if necessary, the previous line and ones above it will get
eaten backwards as well: Once you have erased the offending mistakes, you simply
resume typing forwards. This is how errors are corrected. We all began in this finicky
way, feeling ashamed to let anything flawed remain "in print", so to speak, visible to
others' eyes. But gradually we overcame that sense of shame, realizing that a typo
sitting on a screen is not quite so deathless as one sitting on a page in a book.
    Still, I found that some people just let things go more easily than others. For
instance, by the length of the delay after a typo is made, you can tell just how much its
creator is hesitating in wondering whether to correct it. Hesitations of a fraction of a
second are very noticeable, and are part of a person's style. Even if a typo is left
uncorrected, you can easily spot someone's vacillations about whether or not to fix it.
    The counterparts of these things exist on many levels of such exchanges. There are
the levels of word choice (for instance, some people who don't mind having their
typos on display will often backtrack and get rid of words they now repudiate),
sentence-structure choice, idea choice, and higher. Hesitations and repairs or restarts
are very common. I find nothing so annoying as someone who has gotten an idea
expressed just fine in one way, and who then erases it all on the screen before your
eyes and proceeds to compose it anew, as if one way of suggesting getting together
for dinner at Pagliai's at 6 were markedly superior to another!
    There are ways of exploiting erasure in "talk" mode for the purposes of




A Coffeehouse Conversation on the Turing Test                                        523

humor. Don Byrd and I, when "talk"ing, would often make elaborate jokes exploiting
the medium in various ways. One of his I recall vividly was when he hurled a nasty
insult onto the screen and then swiftly erased it, replacing it by a sweetly worded
compliment, which remained for posterity to seeat least for another minute or so. One
of our great discoveries was that some "arrow" keys allowed us to move all over the
screen, thus to go many lines up in the conversation and edit earlier remarks by either
of us. This allowed some fine jokes to be made.
   One hallmark of one's "talk" style is one's willingness to use abbreviations. This is
correlated with one's willingness to abide typos, but is not by any means the same. I
personally was the loosest of all the "talkers" I knew, both in terms of leaving typos
on the screen and in terms of peppering my sentences with all sorts of silly
abbreviations. For instance, I will now retype this very sentence as I would have in
"talk mode", below.

   F ins, I will now retype is very sent as I wod hv in "talko mode", below.

Not bad! Only two typos. The point is, the communication rate is raised considerably-
nearly to that of a telephone-if you type well and are willing to be informal in all these
ways, but many people are surprisingly uptight about their unpolished written prose
being on exhibit for others to see, even if it is going to vanish in mere seconds.

                                        *   *   *

   All of this I bring up not out of mere windbaggery, but because it bears strongly on
the Turing Test. Imagine the microscopic insights into personality that are afforded by
watching someone-human or otherwise typing away in "talk" mode! You can watch
them dynamically making and unmaking various word choices, you can see
interferences between one word and another causing typos, you can watch hesitations
about whether or not to correct a typo, you can see when they are pausing to work out
a thought before typing it, and on and on. If you are just a people-watcher, you can
merely observe informally. If you are a psychologist or fanatic, you can measure
reaction times in thousandths of a second, and make large collections and catalogue
them. Such collections have really been made, by the way, and make for some of the
most fascinating reading on the human mind that I know of. See, for instance, Donald
Norman's article "Categorization of Action Slips" or Victoria Fromkin's book Errors
of Linguistic Performance: Slips of the Tongue, Ear, Pen, and Hands.
   In any case, when you can watch someone's real-time behavior, a real live
personality begins to appear on a screen very quickly. It is far different in feel from
reading polished, post-edited linefuls such as I received from Nicolai. It seems to me
that Alan Turing would have been most intrigued and pleased by this time-sensitive
way of using his test, affording so many




A Coffeehouse Conversation on the Turing Test                                         524

lovely windows onto the subconscious mind (or pseudo-mind) of the being (or
pseudo-being) under examination.
    As if it were not already clear enough, let me conclude by saying that I am an
unabashed pusher of the validity of the Turing Test as a way of operationally defining
what it would be for a machine to genuinely think. There are, of course, middle
grounds between real thinking and being totally empty inside. Smaller mammals and
in general, smaller animals, seem to have "less thought" going on inside their
craniums than we have inside ours. Yet clearly animals have always done, and
machines are now doing, things that seem to be best described using Dennett's
"intentional stance". Donald Griffin, a conscious mammal, has written thoughtfully on
these topics (see, for instance, his book The Question of Animal Awareness). John
McCarthy has pointed out that even electric-blanket manufacturers use such phrases
as "it thinks it is too hot" to explain how their products work. We live in an era when
mental terms are being both validly extended and invalidly abused, and we are going
to need to think hard about these matters, especially in face of the onslaught of
advertising hype and journalese. Various modifications of the Turing Test idea will
undoubtedly be suggested as computer mastery of human language increases, simply
to serve as benchmarks for what programs can and cannot do. This is a fine idea, but
it does not diminish the worth of the original Turing Test, whose primary purpose was
to convert a philosophical question into an operational question, an aim that I believe
it filled admirably.




A Coffeehouse Conversation on the Turing Test                                      525

                      On the Seeming Paradox of
                       Mechanizing Creativity
                                    September, 1982


It is a commonly heard statement that there is such a thing as the "creative spark",
that an "unanalyzable leap of the imagination" takes place when a great mind comes
up with a new idea or work of art. Great creators are sometimes said to be a "quantum
leap" away from ordinary mortals. People like Mozart are held to be somehow
divinely inspired, to have magical insights for which they could no more be expected
to be able to account than spiders for the wondrous webs they weave. It is ail felt to be
somehow too deep down, too hidden, too occult a gift, to be mechanical in any sense.
Creativity, in fact, is perhaps one of the last refuges of the soul. "You may mechanize
your logic, " says the English professor to the computer scientist, "but you'll never lay
a finger on poetry. " (You may substitute music or any other domain of artistic
creation for poetry.)
        Is this kind of statement irrational? Is it a reflection of a deep-seated fear that
even this most sacred aspect of humanity is doomed to be taken over soon by metallic
machines, or by silicon chips? Why make such a big deal out of an activity of the
human mind which, like every other activity in life, has shades and degrees? After all,
the creative blurs with the mundane so much that it would be hopeless, would it not,
to try to cull what is truly creative from what is not? Or-is there some clean dividing
line that distinguishes the run-of-the-mill workaday deviser of ditties from the Great
Composer of Eternal Symphonic Masterpieces? And if so, is it possible that here lies
the elusive difference between the living and the dead, the human and the machine,
the mental and the mechanical?
        With such a "magical" view of creativity, there is, of course, a problem. It
would seem to imply that the poor composer of ditties is actually dead and mechanical
inside; that only certified geniuses like Mozart are qualitatively different from
machines-and that even old Mozart was nonmechanical only when he was composing
(certainly not when he was merely sipping ale at a tavern!). Probably most people
who believe in the magical view of




On the Seeming Paradox of Mechanizing Creativity                                       526

creativity would dispute this way of portraying their position. They' could maintain
that Mozart was nonmechanical all the time; moreover that you and I, no less than
Mozart, are also nonmechanical all the time. No matter at some, even many, human
abilities have already been mechanized or will be mechanized someday.
        About the touchy question of the mechanization of the mental, many educated
people feel that, although a machine may now or someday be able -to do a creditable
job of acting like a person, any machine's performance will always remain lackluster
and dull, and that after a while, this dullness will always shine through. You'll simply
be able to tell that it is unoriginal, that its ideas and thoughts are all being drawn from
some storehouse of formulas and cliches, that ultimately there is nothing alive and
dynamic-no flan vital-behind its facade. If it comes up with a bon mot now and then,
well, taut mieux-but even the best will just be an automaton par excellence. There
may be nothing specific to point to other than the "vibes" you pick up of its dullness
and unoriginality, but after a while they will inevitably start to come in loud and clear.
(Incidentally, I would be delighted if some of the more vocal antimechanists felt that
way, instead of insisting, as they more often do, that operational tests are of no use in
deciding who or what possesses "genuine mental states".)
        This sense that you will eventually be able to "just tell", from its inevitable
lack of sparkle, that you're dealing with a machine and not a person, seems to depend
upon a tacit assumption about human thought, one with which
I fully agree: namely, that "creative spark" is not the exclusive property of just a few
rare individuals down the centuries, but quite to the contrary, it
is an intrinsic ingredient of the everyday mental activity of everyone, even the most
run-of-the-mill people. In short, it seems that people who feel that machines-even
intelligent ones-will always remain duller than minds are tacitly relying on the
following thesis: Creativity is part of the very fabric of all human thought, rather than
some esoteric, rare, exceptional, and fluky by-product of the ability to think, which
every so often surfaces in places spread far and wide.
        With this thesis I agree. Where I differ with the antimechanists is over the
fitter of whether creativity lies beyond intelligence. I see creativity and Might, for
machines no less than for people, as intimately bound up with intelligence, so that I
cannot imagine a noncreative yet intelligent machine -something that, in order to
make a point about what is essentially human, `+ seem to be willing and able to do.
To me, "noncreative intelligence" a flat-out contradiction in terms.

*   *   *

       In this column, I would like to describe some ideas I have about how CM-
Ativity is founded on mechanisms, mechanisms that, to be sure, lie deeply hidden in
the depths of the structure of our brains, but mechanisms that




On the Seeming Paradox of Mechanizing Creativity                                       527

nonetheless exist and can perhaps be approximated using the hardware and software
of the machines we have today, crude though they are in certain ways. The gist of my
notion is that having creativity is an automatic consequence of having the proper
representation of concepts in a mind. It is not something you add on afterward. It is
built into the way concepts are. To spell this out more concretely: If you have
succeeded in making an accurate model of concepts, you have thereby also succeeded
in making a model of the creative process, and even of consciousness.
        Another way of talking about concepts is to talk about memory, which is the
"place" where concepts are stored. It is the organization of memory that defines what
concepts are. Incidentally, when I first wrote the preceding sentence, it ended
differently. It said, "It is the organization of memory that defines what concepts will
be accessible under what conditions." But on rereading it, I felt it was too weak that
way. It took for granted the notion that all readers have a clear concept of what a
concept is. But that is hardly takable-for-granted! Granted, we all have some concept
of what a concept is, but a clear one?
        So I dropped the phrase beginning with "will be accessible" and replaced it
with a stark "are". This way, the sentence does more than simply state that memory is
a storehouse of some things called concepts. It emphasizes that what establishes the
"concepthood" of something is the way it is integrated into memory. Or to put it the
other way 'round, nothing is a concept except by virtue of the way it is connected up
with other things that are also concepts. In other words, the property of being a
concept is a property of connectivity, a quality that comes from being embedded in a
certain kind of complicated network, and from nowhere else. Put this way, concepts
sound like structural or even topological properties of vast tangly networks of sticky
mental spaghetti.
        That's more or less the image I feel it is important to convey: namely, that
concepts derive all their power from their connectivity to one another. And now,
having expressed that idea, I can return to the sentence as it was originally put: It is
the organization of memory that defines what concepts will be accessible under what
conditions-and surely, the happy choice of the right concept at the right time is the
essence of the creative. Therefore it is imperative to study deeply the nature of that
network-to ask the question "What is a concept?".
        Some questions that come to mind are: What is the relationship between a
general, or Platonic, concept, such as that of "tree", and the concept you form of some
specific tree? That is, what is the distinction between semantic or perceptual
categories and the representations of individual instances of them? How is a given
situation filed away in memory so that one has access to it under an enormous variety
of future situations-access that is often via analogy or other abstract pathways, rather
than by simplistic superficial traits? Or, to flip that coin, how does a given situation
cause the highly selective retrieval from memory of a small number of previous
situations




On the Seeming Paradox of Mechanizing Creativity                                    528

that seem relevant? Only through a deep understanding of the organization of
memory-which is to say, only by answering the question "What is a concept?"-will we
be able to make models of the creative process. This will be a long and arduous
process, not one that will yield answers overnight, or even in a few decades.
Nonetheless, we have the right beginnings, in the sciences of cognitive psychology
and artificial intelligence. Philosophers of mind and neuroscientists will undoubtedly
contribute as well. The union of all these disciplines is called "cognitive science".

                                      *   *   *

        A question that arises at the outset is: "What kinds of objects have concepts
stored inside them, and what kinds do not?" One of my favorite passages that opens
this question wide is in Dean Wooldridge's book Mechanical Man: The Physical Basis
of Intelligent Life, and it runs this way:

  When the time comes for egg laying, the wasp Sphex builds a burrow for the
  purpose and seeks out a cricket which she stings in such a way as to paralyze
  but not kill it. She drags the cricket into the burrow, lays her eggs alongside,
  closes the burrow, then flies away, never to return. In due course, the eggs hatch
  and the wasp grubs feed off the paralyzed cricket, which has not decayed,
  having been kept in the wasp equivalent of a deepfreeze. To the human mind,
  such an elaborately organized and seemingly purposeful routine conveys a
  convincing flavor of logic and thoughtfulness-until more details are examined.
  For example, the wasp's routine is to bring the paralyzed cricket to the burrow,
  leave it on the threshold, go inside to see that all is well, emerge, and then drag
  the cricket in. If the cricket is moved a few inches away while the wasp is inside
  making her preliminary inspection, the wasp, on emerging from the burrow, will
  bring the cricket back to the threshold, but not inside, and will then repeat the
  preparatory procedure of entering the burrow to see that everything is all right.
  If again the cricket is removed a few inches while the wasp is inside, once again
  she will move the cricket up to the threshold and reenter the burrow for a final
  check. The wasp never thinks of pulling the cricket straight in. On one occasion
  this procedure was repeated forty times, with the same result.

        One can make the obvious remark that perhaps not the wasp but the
experimenter was the one in the rut-but humor aside, this is a rather shocking
revelation of the mechanical underpinning, in a living creature, of what looks like
quite reflective behavior. There seems to be something supremely unconscious about
the wasp's behavior here, something totally opposite to what we feel we are all about,
particularly when we talk about our own consciousness. I propose to call the quality
here portrayed sphexishness, and its opposite antisphexishness (a vexish word to
pronounce!), and then I propose that consciousness is simply the possession of
antisphexishness to the highest possible degree. The point is that sphexishness and
antisphexishness are two extremes along a




On the Seeming Paradox of Mechanizing Creativity                                    529

continuum. Let me give a few examples distributed along that continuum, starting at
the most sphexish and finishing with the most antisphexish:

  1. A stuck record. This can be especially ironic if it's a recording of something
     that has a vibrant, lifelike dynamism to it (such as the music of contemporary
     composer Steve Reich), and then the illusion is shattered by the mechanical
     repetition of the jumping needle.
  2. The Sphex wasp herself, and other examples from the insect world. For
     instance, suppose you have a mosquito in your bedroom. You try to swat it,
     and miss. It takes off and flies around the room, losing you. But after a while,
     it settles down and you spot it somewhere on the wall. Again you try to swat
     it and miss. As this cycle progresses, is the mosquito aware of the repetition?
     Does it begin to sense that there is an organized conspiracy against it, or does
     each new swat attempt come as fresh and unexpected as the previous one?
     Does the mosquito formulate some such notion as "the animate agent trying
     to wipe me out"? Sadly for the mosquito (but fortunately for you), it seems
     highly doubtful.
  3. A herd of cattle in a corral, waiting to get branded. There is general
     commotion and hubbub, caused by the noise each cow makes at the moment
     of branding, and propagated outward by the cows closest to it. But does each
     cow in the corral recognize the overall pattern? Is its increased state of
     agitation due to the fact that the cow sees what is coming, or is it rather just a
     kind of vague apprehension, perhaps merely a raised adrenaline level without
     any specific meaning or referential quality?
  4. A dog who is fooled every time by a faking motion in which you pretend to
     throw a ball, but instead don't release it. Actually, I don't know any dog who
     would fall for such an elementary trick. However, I do know a dog (who
     shall remain nameless-although he does happen to be an Airedale) who did
     not catch on when I threw his toy to an upstairs landing instead of down the
     hall (where he expected it). I led him up the stairs and showed him where it
     was. I expected he would know to go upstairs the next time. But no such
     luck. He just ran down the hallway again. Even after I had thrown his toy
     upstairs fifteen times more, he still ran down the hallway, then came back
     looking confused. Poor doggie! True, some of those seventeen painful times
     he did start going up the stairs, but each time he got only partway up, then
     turned around, and hightailed it down the hallway. To me, it was a
     disappointingly sphexish kind of behavior for a dog.
  5. Glassy-eyed gamblers in Las Vegas, glued to their slot machines. To this can
     be added glassy-eyed teen-agers and college students glued to video games
     and pinball machines. Is there not some kind of deadening rut here? And yet
     so many people do this over and over again with seeming pleasure.




On the Seeming Paradox of Mechanizing Creativity                                      530

  6. A happy-go-lucky person who sings or whistles all the time-and if you listen
     closely, you notice that it's always the same little refrain, day in, day out,
     year in, year out: never any variety.
  7. People who make what seems to be the same joke, only in slightly different
     guises, over and over and over again. Or inveterate punsters, who simply
     cannot stop making one pun after another.
  8. Junior-high-school students who fill each other's yearbooks with those same
     pat phrases and corny poems as your junior-high class did.
  9. A mathematician who exploits one single technique to advantage in paper
     after paper, making advances in many different branches in mathematics, yet
     always with a distinct, idiosyncratic touch, and always, in some deep sense,
     just doing "the same old trick" again and again.
  10. People whose rut-stuck behavior leads them down harmful pathways in their
     lives, for instance in their romances or their jobs. We all know people who
     "blow it" in the same 'way each time when faced with a situation that
     matters.
  11. Social trends that become completely stylized and predictable, such as the
     endless trashy sitcoms that television networks keep churning out, the
     movies one after another based on some gimmick exploited in slightly
     different ways. For instance, one could perceive the movies Breaking Away,
     The Black Stallion, and Chariots of Fire as simply three ways of plugging
     specific values for variables into one successful formula-an upcoming
     championship race, a lovable underdog, a rival, and, of course, ultimate
     victory. And these are sophisticated, compared to some books and movies
     that much more blatantly exploit famous predecessors.
  12. Styles in art that become dated and routinized to the point of no longer being
     creative. This happens to-every style, but at the moment of its happening,
     there are always some people who are breaking out of the rut and creating
     totally new styles. However, there are others who become technically
     proficient at an old style, and who continue to create in an old-fashioned
     vein.

How different are these last few examples from the stuck record, or from the Sphex
wasp? What is the real difference we feel as we progress down this list?
        I would summarize it by saying that it is a general sensitivity to patterns, an
ability to spot patterns of unanticipated types in unanticipated places at unanticipated
times in unanticipated media. For instance, you just spotted an unanticipated pattern-
five repetitions of a word. And I'm sure you picked up on all the French phrases
crowded together earlier on in this chapter. Neither in your schooling nor in your
genes was there any explicit preparation for such acts of perception. All you had
going for you is an ability to see sameness. All human beings have that readiness, that
alertness, and that is what makes them so antisphexish. Whenever they get into some
kind of




On the Seeming Paradox of Mechanizing Creativity                                    531

"loop", they quickly sense it. Something happens inside their heads-a kind of "loop
detector" fires. Or you can think of it as a "rut detector", a "sameness detector"-but no
matter how you phrase it, the possession of this ability to break out of loops of all
sorts seems the antithesis of the mechanical. Or, to put it the other way around, the
essence of the mechanical seems to be in its lack of novelty and its repetitiveness, in
its trappedness in some kind of precisely delimited space. This is why the wasp, the
dog, even some humans seem so mechanical.

                                         *   *   *

        How many computers do you know that would react with outrage (or guffaws)
to the simultaneous occurrence on a single mailing list of "Bernie Weinreb", "Bernie
W. Weinreb", "Mr. Bernie Weinreb, R.M.", "Barnie Weinrab", and so forth?
Computers do not have automatic sensitivity to patterns in the data that they deal
with. And of course, how could they be expected to? As one old saw goes, they do
only what they are programmed to do. Computers are not inherently bored by adding
long columns of numbers, even when all the numbers are the same. But people are.
What is the difference?
        Clearly there is something lacking in the machine that allows it to have this
unbounded tolerance for repetitive actions. This thing that is lacking can be described
in a few words: It is the ability to watch oneself as one deals with the world, to
perceive in one's own activities a pattern, and to be able to do so at many levels of
abstraction. Thus, consider the case of a hypothetical self-watching computer. To be
sensitive in this way, it should get bored whenever it is forced to add a long column of
identical numbers together. Wouldn't you? It should get bored whenever it is forced to
do just adding over and over again, even when the numbers are different. Wouldn't
you? It should even get bored when asked to do many arithmetic operations in any
sort of repetitive pattern! Wouldn't you? Any loop of any sort should become tedious!
Wouldn't it?
        But where does it stop? Surely if a computer could perceive that all it ever
does is pull up one instruction after another from memory (a piece of hardware, not to
be confused with human memory), execute those instructions, and change various
registers, it would yawn very boredly and probably soon go to sleep. And by the same
token, you or I, if we ever gained access to the firings of our neurons, would find
watching the activity to be one of the most stultifying things imaginable.
        But this is not the kind of self-watching I mean. Watching one's own internal
microscopic patterns is bound to be boring, because any complex system is bound to
be made up out of thousands, millions, or even more copies of small elements (such
as gears, transistors, cells, and so on). What is critical is to be able to watch activities
on a completely different level the collective level, in which huge patterns of activity
of these many




On the Seeming Paradox of Mechanizing Creativity                                       532

components assume regular behaviors perceptible on their own. A hurricane is a huge
pattern of activity of tiny atoms, but one that has such regularity and pattern that we
can predict hurricanes without ever thinking of their constituent atoms. A thought is a
huge pattern of activity of tiny cells, of which much the same can be said.
        Antisphexishness has to do with self-perception at this kind of level. Rather
than watching its neurons or transistors or registers, an antisphexish toeing watches its
own high-level patterns, looking for similarities somewhat the way meteorologists
might look for one hurricane following another in a regular way.
        Thus we should not expect or even want a self-watching computer to be able
to see down to the level of its circuitry; it would not watch itself doing machine-
language operations such as ADD, STORE, and JUMP in loop-like .patterns. The
effects of such operations are to change larger things called “data structures" in
memory. Self-watching involves monitoring those changes as they happen, filtering
out the dull ones, and recording certain aspects of the interesting ones in other data
structures. (The fact that such monitoring, filtering, and recording would, on a more
microscopic level, involve the very same kinds of elementary machine-language
operations would be invisible to the computer, since it should be shielded from that
detailed a view of itself.) Thus patterns in the changes taking place in one set of data
structures would get recorded in another set of data structures. Should we then not set
up a third level of data structures, to watch the second level, should patterns occur in
it? And a fourth, to watch the third? This seems prime territory for an infinite regress:
an endless hierarchy of structures, each one monitoring changes in the level below it.
        Now that is quite true, and it is because you are a self-watching human being
that you caught onto this pattern, and probably before I had spelled out. It is in the
nature of human pattern perception to be able to detect such infinite regresses, and to
stop them short before they ever get anywhere. But what about the hypothetical self-
watching computer, with its infinitely many layers of watchers?
        Well, surely one of the most salient features-no, definitely the most salient
feature-of what I have just described is the pattern of the data 'structures themselves:
the hierarchy stretching upwards repetitively towards infinity. Shouldn't this pattern
be as blatant to a self-watcher as it is to us? Indeed yes, it should. If we were to label
the bottom level `0' and the first watching level 'I', then logically we should label the
further levels `2', `3', and so on. Each level in this potentially infinite set can be
identified with a natural number. Once the pattern is perceived by a watcher, that
watcher can form the general concept of "all the levels seen at once", associated with
the concept of "all the natural numbers conceived of at once". The Conventional name
for the set of all natural numbers is `co' (omega), which we can take as the name of a
new watching level that looks out for patterns in this potentially infinite tower of
watchers.




On the Seeming Paradox of Mechanizing Creativity                                      533

        You need not worry, by the way, that in proposing such a self-watching
computer I am presupposing an infinite machine. Precisely the opposite. The whole
purpose of stopping infinite regress in its tracks is so that we will not need to actually
build an infinite tower of data structures and watching processes, a feat that would
clearly be impossible, aside from being monumentally sphexish. At any stage, only a
finite amount of recording would have been done, so that only a finite number-in fact,
a small number -of levels of structure would exist. The only requirement is that there
should exist the potential to extend it further.
        It would be the co-watcher that would perceive (as you and I and any human
being would) the infinite-regress pattern of attempts to build the w-tower itself. The
w-watcher would catch any such infinite regress before it could start. If a change in
level 0 caused a change in level 1 that caused yet another change in level 2, and if
these changes seemed to be patterned in such a way that an inevitable infinite ripple
upwards would ensue, the co-watcher, ever alert for such patterns in the other
watchers, would come to the rescue, shouting "Wait! Enough! Halt!" Thus in fact, no
infinite regress would actually occur; it would be nipped in the bud by the same sorts
of mechanisms that allow you to cut off a bore at a party. "Excuse me, I think I'll go
get some more punch."

*   *   *

         The problem is, there's nothing to prevent the co-level itself from going into
loops-so if we're going to obviate that, we have to have a higher watcher-
conventionally called "w+ I". Uh-oh! Before I even had a chance to begin spelling it
out, you sniffed a new infinite regress! (You ruin all my fun!) Well, I'm going to spell
it out, anyway. Level w+ l needs to be watched by level w+2, and that level by level
w+3. Thus we have a second potentially infinite tower of watchers, all of whom will
be watched over by the Grand Watcher: level 2w. But if there can be two towers, then
why not three? And so, of course, it goes. Wheels within wheels, patterns of patterns
of patterns. We get watchers 2w, 3w, and now our tower of towers needs a new Great-
Grand Watcher: W2. And then
         Excuse me; I think I'll go get some more punch. There is a problem once you
start getting into infinite regresses composed of other infinite regresses -the whole
thing just never stops, and it becomes a bore. Or not exactly a bore, but a very
complex and confusing thing, whose reality and relevance become ever more
questionable. And yet, when you bring it back to the domain of sphexishness, it
becomes the very real and very relevant question of how to build a machine that can
sense unanticipated patterns in its own behavior.
         This is related to a classic problem in the theory of computability, called the
halting problem: It is the question of whether there exists any computer program that
can inspect other programs before they run, and reliably




On the Seeming Paradox of Mechanizing Creativity                                      534

predict whether or not they will go into infinite loops ("going into an infinite loop"
means, of course, never coming to a halt-and conversely, "halting" means avoiding
any infinite loop). The answer turns out to be "Definitely not", and for elegant, deep
reasons. (Recall Chapter 21.) Of course, the thing hinges on getting this halting
inspector to try to predict its own behavior when looking at itself trying to predict its
own behavior when looking at itself trying to predict its own behavior when ... Excuse
me; I think I'll go get some more punch.
         This halting-problem idea is closely related to our question about self-
watching programs, but it is not really the same thing. First of all, the halting problem
is concerned with an inspection to be carried out on programs before they are running,
like looking at blueprints of buildings before they are built to see if they are
earthquake-proof. Here we are talking about a program that is observing some
program while it is running-and what's more, it's not just "some program" that it is
watching, but itself. Of course, not all of its attention is being devoted to seeing if it's
gotten into a rut (for that would itself constitute ruttish behavior!), but while it's doing
other things, it's keeping its eye peeled, so to speak, for signs of ruttishness inside
itself.
         In computability theory, when a program or system of any sort turns back on
itself in this manner, the turning-back-on-itself is known as diagonalization. To some
people, diagonalization seems a bizarre exercise in artificiality, a construction of a
sort that would never arise in any realistic context. To others, its flirtation with
paradox is tantalizing and provocative, suggesting links to many deep aspects of the
universe. Now here we see a dynamic diagonalization-a self-watching program-that
seems t6 be closely connected with what makes a human being so utterly different
from a stuck record or a Sphex wasp. Surely that is not such a bizarrely artificial thing
to ponder!
         Probably the most significant difference between the halting problem and the
idea of a self-watching program is that in trying to build an artificial intelligence, we
are not really so concerned with the mathematical perfection of our self-watching
system as with its likelihood of survival in a complex world; after all, that's what
intelligence is about. So if there is a mathematical theorem telling us that no program
whatsoever will be a perfect self-watcher, able to catch itself in any conceivable kind
of infinite regress, well, that is simply a statement that perfect intelligence is
unreachable something that ought to please us rather than dismay us, since it would be
ether horrible and disappointing if someone came up with some finite program after a
while, and could legitimately announce, "Well, folks, here it is at last: the end-all of
intelligence, a perfectly intelligent program."
         But don't worry about that. The metamathematical work of Kurt Gbdel, Alan
Turing, Stephen Kleene, and others, on such things as the halting problem and the
theory of infinite ordinals (such as the towers of numbers and w's), tells us that this
scenario will not come to pass, for neither is there




On the Seeming Paradox of Mechanizing Creativity                                       535

a perfect halting inspector, nor is there any ultimate scheme for naming ordinals.
What this latter result means is that there is no finite mechanism that can possibly
detect all patterns, patterns of patterns, patterns of patterns of patterns of patterns
(aha!-fooled you that time, didn't l?), and so on.

                                       *   *   *

        In his famous paper "Minds, Machines, and Gödel", the English philosopher J.
R. Lucas attempted to capitalize on these sorts of "negative" results of
metamathematics by claiming that they provided the key element in a proof that no
machine could ever be conscious in the way that humans are. Let Lucas speak for
himself
           At one's first and simplest attempts to philosophize, one becomes
   entangled in questions of whether when one knows something one knows that
   one knows it, and what, when one is thinking of oneself, is being thought about,
   and what is doing the thinking. After one has been puzzled and bruised by this
   problem for a long time, one learns not to press these questions: the concept of a
   conscious being is, implicitly, realized to be different from that of an
   unconscious object. In saying that a conscious being knows something, we are
   saying not only that he knows it, but that he knows that he knows it, and that he
   knows that he knows that he knows it, and so on, as long as we care to pose the
   question: there is, we recognize, an infinity here, but it is not an infinite regress
   in the bad sense, for it is the questions that peter out, as being pointless, rather
   than the answers. The questions are felt to be pointless because the concept
   contains within itself the idea of being able to go on answering such questions
   indefinitely. Although conscious beings have the power of going on, we do not
   wish to exhibit this simply as a succession of tasks they are able to perform, nor
   do we see the mind as an infinite sequence of selves and super-selves and super-
   super-selves. Rather, we insist that a conscious being is a unity, and though we
   talk about parts of the mind, we do so only as a metaphor, and will not allow it
   to be taken literally.
           The paradoxes of consciousness arise because a conscious being can be
   aware of itself, as well as of other things, and yet cannot really be construed as
   being divisible into parts. It means that a conscious being can deal with
   Gödelian questions in away in which a machine cannot, because a conscious
   being can both consider itself and its performance and yet not be other than that
   which did the performance. A machine can be made in a manner of speaking to
   `consider' its performance, but it cannot take this `into account' without thereby
   becoming a different machine, namely the old machine with a .new part' added.
   But it is inherent in our idea of a conscious mind that it can reflect upon itself
   and criticize its own performances, and no extra part is required to do this: it is
   already complete, and has no Achilles' heel.

        Somehow-and I think understandably-Lucas was under the impression that
human beings are endowed with powers that are equivalent to a self-watcher of
infinite depth, someone who will detect and terminate any




On the Seeming Paradox of Mechanizing Creativity                                    536

and all patterned behavior: the ultimate in antisphexishness. I call this hypothetical
ability "Breaking Out Of Loops Everywhere"-" BOOLE" for short, in honor of
George Boole, who wrote one of the most influential books of the nineteenth century,
The Laws of Thought, surely a forerunner of today's artificial intelligence work.
        Lucas seems to think that to be human is to be endowed with this "BOOLE"
ability-this total and perfect antisphexishness-intrinsically. On reflection, however,
one realizes this surely is not the case. Despite not being Sphex wasps or Airedales,
we humans are all still vulnerable to getting caught in ruts, as I attempted to point out
in the dozen-item list above. None of us is immune. Each of us-even the Mozarts
among us-exhibits a "cognitive style" that in essence defines the ruts we are
permanently caught in.
        Far from being a tragic flaw, this is what makes us interesting to each other. If
we limit ourselves to thinking about music, for instance, each composer exhibits a
"cognitive style" in that domain-a musical style. Do we take it as a sign of weakness
that Mozart did not have the power to break out of his "Mozart rut" and anticipate the
patterns of Chopin? And is it because he lacked spark that Chopin could not see his
way to inventing the subtle harmonic ploys of Maurice Ravel? And from the fact that
in "Bolero" Ravel does not carry the idea of pseudo-sphexish music to the
intoxicating extreme that Steve Reich has, should we conclude that Ravel was less
than magical?
        On the contrary. We celebrate individual styles, rather than seeing them
negatively, as proofs of inner limits. What in fact is curious is that those people who
are able to put on or take off styles in the manner of a chameleon seem to have no
style of their own and are simply saloon performers, amusing imitators. We accord
greatness to those people whose "limitations", if that is how you want to look at it, are
the most apparent, the most blatant. If you are familiar with his style, you can
recognize music by Maurice Ravel any time. He is powerful because he is so
recognizable, because he is trapped in that inimitable "Ravel rut". Even if Mozart had
jumped that far out of his Mozart system, he still would have been trapped inside the
Ravel system. You simply can't jump infinitely far!
        The point is that Mozart and Ravel, and you and I, are all highly antisphexish,
but not perfectly so, and it is at that fuzzy boundary where we can no longer quite
maintain the self-watching to a high degree of reliability that our own individual
styles, characters, begin to emerge to the world.
        Although Lucas has been roundly criticized, and rightly so, I believe, by many
philosophers, logicians, and computer scientists for failing to see many important
subtleties of the GSdel argument on which he bases his paper, most of his critics have
failed to see the crucial aspect of mind that Lucas was one of the first to point out.
Lucas correctly observes that the degree of nonmechanicalness that one perceives in a
being is directly related to its ability to self-watch in ever more exquisite ways.
Unfortunately, too




On the Seeming Paradox of Mechanizing Creativity                                     537

many artificial-intelligence people are ready to pooh-pooh the Lucas article on the
grounds that its central thesis-the impossibility of mechanizing mind-is wrong. What
they miss is that it is pointing at very deep issues that have much to do with the very
core of intelligence and creativity.

                                       *   *   *

        Earlier I stressed the importance of the organization of memory and the
pressing need to come at the question "What is a concept?" Critical to the way our
memory is organized is our automatic mode of storing and retrieving items, our
knowledge of when we know and do not know, of how we know or why we wouldn't
know. Such aspects of what is sometimes called "metaknowledge" are fluidly
integrated into the way our concepts are meshed together. They are not some sort of
"extra layer" added on top by a second-generation programmer who decided that
metaknowledge is a good thing, over and above knowledge! No, metaknowledge and
knowledge are simmering together in a single stew, totally fused and flavoring each
other richly. This makes self-watching an automatic consequence of how memory is
structured. How is this wondrous stew of antisphexishness realized in the human
brain?
        And how can we create a program that, like a human brain, is all "of a piece",
a program that is not simply a stack of ever-higher "other-watchers", but is truly a
seamless "self-watcher", where all levels are collapsed into one? If we wish to have a
program that breaks out of the extremely sphexish mold that all programs seem to be
in today, we have to figure out how a flexible perception program might exploit its
own flexibility to look at itself. Of course, no such program will be written as I just
stated. That is, it will not come into being in the following way:

       Step 1. We write a flexible perception program.
       Step 2. We turn that program back on itself as a self-watcher.

Rather, to achieve the results desired in Step 1, we must have incorporated the goals
of Step 2 into the design from the start! In other words, these two goals are
intertwined, more in the following sense:

       Goal 1. Flexible perception.
       Goal 2. Self-watching.

There is no chronological priority here, for the two goals are too intertwined to have
one precede the other. This is a tricky foldback, quite a bit more elaborate than the
one involved in the halting problem, yet in spirit related to it.
       It is interesting that Lucas' argument was based on Gödel’s Theorem, whose
proof depends on making one of these seemingly impossible (or at




On the Seeming Paradox of Mechanizing Creativity                                   538

least highly counterintuitive) foldbacks-this one where a mathematical system of
reasoning folds back on itself and subsumes itself as an object of study. What is
fascinating in that proof is how, in such a system, there is a kind of level-collapse that
ensues from the ability of a system to see itself. Rather than there being towers of
watchers, then towers of those towers, and so on ad infinitum in the worst possible
sort of multiply infinite regress, all those degrees and levels of self-perception are
achieved at once by the fact that the system can mirror itself. Not that it mirrors itself
in every aspect, mind you-for that would entail contradiction-but it does so at all
levels of complexity.
        The seemingly distinct levels of watcher and watched are totally fused, in the
Godel construction, exactly as Lucas would have it occurring in the minds of all
conscious beings. The only thing that Lucas failed to understand is that the ability to
fold around and see oneself in the wonderfully circular GOdelian way does not-in
fact, cannot-bring with it total antisphexishness. That, fortunately or unfortunately,
depending on your point of view, is a chimera.

*   *   *

        Back in 1952, the philosopher and composer John Myhill wrote a lyrical
article entitled "Some Philosophical Implications of Mathematical Logic: Three
Classes of Ideas". The three classes are borrowed from mathematical logic, and
Myhill's names for them are the effective, the constructive, and the prospective. In
logic, they are known more technically as the recursive, the renotrec (short for
"recursively enumerable but not recursive"), and the productive. Their essence is
described below.
        A category is effective provided that there is a way, given a candidate for
membership, of deciding without any doubt whether that object is or is not a member.
Is Ronald Reagan a KGB agent? Is the Pope Catholic? Although these two questions
are easy to answer, which would seem to imply that being a KGB agent and being
Catholic are examples of the effective, this is slightly misleading. Was Lee Harvey
Oswald a KGB agent? Is an excommunicated bishop Catholic? Examples like these
show that these categories are not genuinely effective categories-but then nothing in
the real world is as clean as it is in logic. I could have asked, "Is 29 prime?" but I
wanted to show how these notions extend beyond the mathematical realm. In natural
languages, grammaticality (syntactic well-formedness) is a rather fuzzy property, but
in an idealized language or formal system, it would be a perfect example of an
effective property.
        We pass on to the constructive. A property that is constructive is more elusive
than one that is effective. The idea here is that some means exists whereby members
of the category can be churned out one by one, so that you will eventually see any
particular member if you wait long enough, but no means exists for doing the
complementary operation-namely, churning out nonmembers, one by one.
Unfortunately, although this kind of set in




On the Seeming Paradox of Mechanizing Creativity                                      539

mathematics is an I extremely important one, easily definable examples of it are rather
hard to come by. The set of all theorems in any formal axiomatic system is always
recursively enumerable, but very often its complement is also, which turns the set into
an effective one rather than a constructive one. You have to be dealing with a formal
system whose non theorems are not themselves producible by some complementary
formal system. Only then do you have a renotrec, or constructive, set. The set of
theorems of any formalized version of number theory turns out (by Gödel’s theorem)
to have this property.
        So much for the "constructive". We finally come to the prospective, also
known as the productive. Myhill's characterization of it is this: "A prospective
character is one which we cannot either recognize or create by a series of reasoned but
in general unpredictable acts." Thus it is neither effective nor constructive. It eludes
production by any finite set of rules. However-and this is important-it can be
approximated to a higher and higher degree of accuracy by a series of bigger and
better sets of generative rules. Such rules tell you (or a machine) how to churn out
members of this prospective category. In mathematical logic, works by Tarski and
Gödel establish that truth has this open-ended, prospective character. This means that
you can produce all sorts of examples of truths-unlimitedly many-but no set of rules is
ever sufficient to characterize them all. The prospective character eludes capture in
any finite net. (See Chapter 13 for a discussion of Platonic notions such as
"chairness", 'A'-ness, etc.)
As his prime example outside of mathematical logic of this quality, Myhill suggests
beauty. As he puts it:

  Not only can we not guarantee to recognize it [beauty] when we encounter it,
  but also there exists no formula or attitude, such as that in which the romantics
  believed, which can be counted upon, even in a hypothetical infinitely
  protracted lifetime, to create all the beauty that there is.

Thus beauty admits of a succession of ever-better approximations, but is never fully
attainable. Beauty and irrationality are often linked. Is it coincidental that the first
example of such a notion of something approximable but never attainable in a finite
process is called an "irrational" number?
        Myhill is bold enough to speculate as follows: "The analogue of GOdel's
theorem for aesthetics would therefore be: There is no school of art which permits the
production of all beauty and excludes the production of all ugliness." To each coin
there are two sides; and the obverse side of beauty is ugliness. By a rather ironic
coincidence, the complementary set to a productive (or prospective) set is called, in
the jargon of mathematical logic, creative. It must be admitted that it would take a
stupendously brilliant, if perverse, sort of creativity to produce all possible ugly
objects.
        If we see the aim of art as the production of all possible objects of beauty




On the Seeming Paradox of Mechanizing Creativity                                    540

(which is doubtless an oversimplification, but let us adopt that view nonetheless), then
each individual artist contributes objects in a particular style. That style is a product of
the artist's heredity and formation, and becomes a hallmark. To the extent of having
an individual style, any artist is sphexish-trapped within invisible, intangible, but
inescapable boundaries of mental space. But that is nothing to lament. Artists in
groups form movements or schools or periods, and what limits one artist need not
limit another. Thus, by the fact that its boundaries are wider, a school is less sphexish-
more conscious-than any of its members.
          But even the collective movement of a school of art has its limits, shows its
finitude, after a period of time. It begins to wind down, to lose fertility, to stagnate.
And a new school begins to form. What no individual can make out clearly is perhaps
seen collectively, on the level of a society. Thus art progresses towards an ever wider
vision of beauty-a "prospective" vision of beauty-by a series of repeated
"diagonalizations": processes of recognizing and breaking out of ruts. As I like to put
it, this is the process of jootsing (jumping out of the system) to ever wider worlds.
          This endless jootsing is a process whose totality (so says GSdel) cannot be
formalized, either in a computer or in any finite brain or set of brains. Thus one need
not fear that the mechanization of creativity, if ever it comes about, will mark the end
of art. Quite the contrary: It is a day to look forward to, for on that day our eyes will
open-as will those of computers, to be sure -onto whole new worlds of beauty. It will
be a happy day when, hand in hand with our new computer friends, we take an
unanalyzable leap out of the system and go get some.more punch.

                                         *   *   *

        Do you know the Saint-Saens Violin Concerto No. 3? Its middle movement
happens to be based on a ravishingly beautiful melody-long, sinuous, flowing, lyrical.
I suggest you get a hold of it and listen to it! Where do such melodies come from? Did
they always exist? Are some people just lucky to have picked them up, these pretty
pebbles lying on the musical beach?
        Well, I hardly want to get into the discovery-invention-existence quagmire
here. I have my own opinions, to be sure, but what I am more concorned with is
where such inspiration comes from. One can point with a fair degree of objectivity to
certain composers as being the most melodically gifted. These names come to my
mind, for instance: Chopin, Rachmaninoff, Saint-Saens, Tchaikovsky, Brahms, Bach,
Mendelssohn, Puccini-and, switching gears somewhat, Cole Porter, Richard Rodgers,
Jerome Kern, and George Gershwin. Obviously there are others. Some people
undoubtedly would strike some off this list and would suggest others-perhaps
Schubert,




On the Seeming Paradox of Mechanizing Creativity                                       541

Dvorak, Prokofiev, Scott Joplin, Fats Waller, Frederick Loewe, the Beatles, Carole
King ... It's hard to draw the line.
         The main point is that certain rare people seem to be able to tap into some
magic vein in which flow incredibly catchy patterns, deeply intoxicating to the human
spirit. Leonard Bernstein once wrote a lively dialogue encatchily titled "Why Don't
You Run Upstairs and Write a Nice Gershwin Tune?". In it, he talks about why that
vein is so hard to tap. Bernstein should know, of course, since he too is one of the
great melodic inventors of our time.
         The problem is that melody invention, like every other art, looks so easy after
the fact. In fact, in many ways it looks easier than creating other kinds of beauty,
because melodies are such small, easily described structures. Making a beautiful turn
on skis at least involves a continuum of possibilities, whereas a melody usually
involves a very restricted, discrete alphabet (the notes within a two-octave range or
so), and isn't even very long!
         It is tempting, therefore, to imagine that good melodies are producible from
some sort of recipe or mathematical formula, or, what comes to nearly the same thing,
to think that the amount of beauty in a melody could be measured by some sort of
machine, just as the amount of radioactivity in a sample of ore can be measured by a
scintillation counter. You would stick your proposed string of notes into a machine
and out would come a number called its "CQ" ("catchiness quotient").
         If you doubt that the very idea of such a number is coherent, just remember
that attached to every piece of existent music there really is a measure of its
catchiness-namely, how often it actually is listened to, at the present time. Pieces can
be rank-ordered according to this very cold, linear measure. This is not to suggest that
the top piece is the best, but only to point out that the idea of a single, one-
dimensional "catchiness index" applying to every possible string of notes is by no
means absurd. Admittedly, under the present circumstances, it seems to take an entire
society of millions of people to calculate the value for any string of notes, but could
all that not be simulated? Perhaps the catchiness-quotient machine could be built to
accept a set of parameters characterizing the target culture and its general musical
mood at the time, and then it would predict how the given tune would fare in the
given society under the specified musical circumstances. Is that not an engaging
notion?
         Are the musical receptivities of a culture truly characterizable in purely
mathematical terms relating only to the syntactical structures of melodies? Ultimately,
of course, the answer has got to be "yes", if by "syntactical structures" you mean
structures whose recognition might require bringing in arbitrary amounts of external
information. Sufficiently deep syntactic probing is tantamount to semantic probing, a
motto from Chapter I's P.S. The question is, then, just how complex a "syntax
machine" that creates, or at least measures, melodic beauty would be. (Let's assume
that it contains adjustable parameters for culture and mood.) Need it be as complex as
a human society or a human brain? Can wonderful, lyrical, sinuous, and




On the Seeming Paradox of Mechanizing Creativity                                    542

rapturous melodies come pouring out of a black box that can do nothing but that?
Readers of Godel, Escher, Bach (especially pages 676-680) might recall that I am
extremely skeptical on that score. Yet how solid is the ground I am standing on?
Could music not yield to brute computational power as swiftly as chess skill has
(something which, in the same passages in GEB, I also was very skeptical about)?

                                       *   *    *

        It is funny how certain fads catch on, seemingly for no reason, while other
things die, again for no clear reason. We all laugh at the Edsel today-yet what exactly
is there to laugh at, except the fact that it did so poorly? What exactly was wrong with
the Edsel? What is wrong with those thousands upon thousands of melodies that are
composed every year and go nowhere? What made Michael Jackson and Pachelbel's
simple Canon all the rage? Why did the typeface Helvetica catch on like wildfire
when it was first invented, when a dozen extremely similar ones died on the vine?
Why did the typographical gimmick of symmetrically capitalizing both the first and
the last letter of a word or title, as in


       G ATEWA Y                                       PRINCE
           INN                                 SPAGHETTI

become a sudden vogue about four years ago?
        Why is it now faddish to write run-on words such as "Intelligenetics" or
"PEOPLExpress"? What makes words like "Da-glo", "Turbomatic", and "Rayon"
seem slightly dated? Why is "Qantas" still modern-sounding? What is poor about
brand names like "Luggo" and "Flimp"? Why are 'x's now so popular in brand names?
And yet why would "Goxie" be a weak name compared with, say, "Exigo" or
"Xigeo"? Why are the ordinaryseeming names that nasal-voiced comedians Bob and
Ray come up withfor example, "Wally Ballou", "Hudley Pierce", "Bodin Pardew",
and "John W. Norbis"-apt to evoke snickers? How come Norma Jean Baker changed
her name to "Marilyn Monroe"? Why would it not do for a movie star to be named
"Arnold Wilberforce"? Why is the name "Tiffany" popular today, and why was "Lisa"
so popular a few years earlier? Is something wrong with "Agnes", "Edna", or
"Thelma"? With "Clyde", "Lance", or "Bartholomew"? Mere length certainly cannot
be the answer (think of "Elizabeth"). Nor can the sound, in any simple sense. (Why is
"Lance" bad if "Vance" is okay?)
        All this may seem a far, far cry from sphexishness and self-watching
computers and brains. But what I am getting at is the unbelievable number of forces
and factors that interact in our unconscious processing of even very tiny structures
composed of discrete parts, such as words and names only a few letters long, let alone
melodies several dozen notes long. Most of us




On the Seeming Paradox of Mechanizing Creativity                                    543

could not put our finger on the answers to any of these questions. In fact, nobody
could really answer these questions definitively. If we are going to try to get machines
to do the subtlest of cognitive tasks, we had jolly well better be able to explain how
mere words are appealing or repelling!

                                        *   *   *

         There are currently some efforts in artificial intelligence to imbue programs
with a certain type of introspective capacity. Such a capacity is usually termed
"reflection", a self-explanatory name that harks back to mathematical logic. A formal
system is said to be capable of reflection if it can reason about itself. Godel was the
first person to discuss such things in detail. Nowadays reflective systems are the bread
and butter of many a logician. However, computer modeling of logic is just now
reaching the point where reflection is being seriously explored.
         The idea is very enticing, but I think it has less to do with genuine progress in
AI than it does with progress in elegant formal systems. It all has to do with one's
ultimate view of what thought is. If you believe that thought is intimately tied up with
some strict notion of truth and reasoning, and that exquisitely honed deductive
capacities are the centrepiece of mentality, then you will naturally be drawn toward
reflective reasoning systems. If, on the other hand, you believe, as I do, that reasoning
is a far, far cry from the core of thought, then you will not be too inclined to jump
toward such systems.
         One way of looking at things is this. Imagine you have a set of rules that are
supposed to capture the way people think in some domain-say that of melody
composition. Now you try them out, and you find that most of the time they fail for
complex reasons, but reasons that you have some intuitions about. How should you
proceed now? There are two main rival avenues, the way I see it.
         One avenue says, "Add meta-rules! Then add meta-meta-rules! Then ... ad
infinitum!" This might be called the "meta-meta" school of Al. The strategy is to
improve the performance of a given set of rules by having higher-order meta-rules
that help determine when and how to apply the ordinary rules. And this process
knows no bounds, even to the point that one can formalize the progression from one
level to its meta-level, so that in principle, an infinite number of meta-levels now are
"there" to be consulted if needed.
         The alternate avenue is to sidestep the topless tower of bureaucracies and
meta-bureaucracies above by making rule-like behavior emerge out of a multi-level
bubbling broth of activity below. This means that you give up the idea of trying to
explicitly tell the system as a whole how to run itself. Instead, you content yourself
with defining explicit micro-behaviors that will interact in vast numbers, and then you
just let them go, carefully watching




On the Seeming Paradox of Mechanizing Creativity                                      544

what ensues and noting what you like and what you don't like. After the run, you
theorize about what might have made the system's top-level behaviour more closely
resemble your ultimate goals, and you go back and tinker around with the micro-
elements whose micro-behavior you have explicit control over, using your best guess
as to what sorts of changes will improve overall performance. Then you run the
system again.
        I remember a long time ago seeing a television show-perhaps you have seen it,
too-in which someone set up a bathtub full of spring-loaded mousetraps holding ping-
pong balls. Then they threw a single ping-pong ball in, and WHAM! The whole thing
exploded madly, in parallel chain reactions. It was all over in a few seconds, but you
can imagine running a film of it in slow motion. There are numerous large-scale
features of the explosion that one could aim at creating, such as how long the pop
takes, how high the average ping-pong ball flies, what the envelope of the flying balls
looks like, and so on. If there were more types of micro-element and their interactions
were more variegated, then you can imagine how multi-dimensional the system's
macrobehavior would be, and how hard it would be to predict even its most basic
features.
        Yet when certain vast ensembles grow sufficiently big, the statistical principle
called "the law of large numbers" sets in, in essence guaranteeing that there will be so
much cancellation in the chaos that ultimately, a kind of order will emerge. It is for
reasons like this that the National Safety Council can predict fairly accurately how
many deaths there will be on a Labor Day weekend, even though they have no idea
where any particular one will occur. Somehow, amazingly, the drivers cooperate and
produce just about the predicted number each time, usually even on the state-by-state
level, although less accurately.
        The difference between such statistically emergent macrobehavior and rigidly
constrained macrobehavior is best made by contrasting the mousetrap system with a
huge domino-chain network, involving branching and rejoining paths, paths that
climb hills and go back down, anything you ran imagine as long as it's entirely self-
determined (i.e., no unanticipated external events start chains falling). In this kind of
system, you know how everything is going to work beforehand. It's true that you may
not be able to predict which of two "rival" pathways will reach a certain point first,
but this kind of unpredictability is not nearly as hard to correct as that of the
mousetrap system. If on one run the result is not what you want, you can just set it up
again the same way, change some specific region, and you know what will happen.
You can program this kind of system, but you cannot program a statistical system in
the same sense. You can only tailor its micro-elements, and then release them and see
what happens.
        Which approach to mind is superior? Is the mind more like a fancy system of
domino chains or a bathtub full of spring-loaded mousetraps? I'm betting on the latter.
More will be found on this topic in Chapters 25 and 26 and their postscripts.




On the Seeming Paradox of Mechanizing Creativity                                     545

                                       *   *   *

         I received a letter from Thomas P. Laubert, in which he expressed
considerable perplexity over a paragraph he had come across containing the following
sentence: "Experience had taught the du Pont engineers to provide ... flexibility in the
design, wherever possible, to meet unforeseen problems that were sure to arise."
Laubert mused: "But if the nature of the problems was unforeseen, then what
parameters were used to determine these built-in flexibilities?" Another reader, whose
letter I have unfortunately misplaced, brought up a similar point about engineering.
What I remember vividly is his term "UNK-UNK"s-meaning the unknown unknowns
that plague all complex systems. He was asking, rather skeptically, as I recall, how
one can ever hope to build a system that anticipates all possible problems.
         These simple-seeming questions hit the nail on the head. An intelligence is, by
definition, a system supposed to be able to deal with the unpredictable. But how can
any set of rules "frozen" into a machine's design do that? Doesn't the very fact of
being frozen make any foreordained system/program/machine/organism vulnerable in
some way that actually follows from the rules themselves? This, of-course, is the
Gödelian point that J. R. Lucas was trying to make in the article I quoted from in the
column. And the only satisfactory answer that I can see is to admit that, yes, all
intelligences are indeed vulnerable-including biological ones, and that means people
no less than Sphex wasps. Natural selection has looked favorably upon organisms
with highly abstract kinds of vulnerability, highly abstract kinds of sphexishness. And
so for the time being, humans are doing all right. But as for there being a fixed recipe
that would allow an organism to cope with all the curves that the universe at large
might throw at it, that is a vain and crazy hope!




On the Seeming Paradox of Mechanizing Creativity                                    546

                                          24.
                      Analogies and Roles
                in Human and Machine Thinking
                                   September, 1981


IN our research in artificial intelligence, my graduate students Gray Clossman and
Marsha Meredith and I have been looking at typical human thought processes in
everyday life as well as in more limited domains, and everywhere we look, we seem
to find that within the internal representations of concepts there are substructures that
have a kind of independence of the structures of which they are part. Such a
substructure is modular-exportable from its native context to alien contexts. It is an
autonomous structure in its own right, and we call these modules roles. A role, then, is
a natural "module of description" of something, a sort of bite-sized chunk that seems
to be comfortable moving out of its first home and finding homes in other places,
some of them unlikely at first glance.
        One intriguing example is the "First Lady" role. Probably most Americans use
this term more flexibly than they realize. They would most likely say, if asked, that
the term means "the wife of the president", and not think any more about it. But if
they were asked about the First Lady of Canada, what would almost surely pop into
their mind is the name or image of Margaret Trudeau. They might reject the thought
as soon as it occurred to them, but for us the important thing is that the thought of her
would arise at all. First of all, people know her as the former wife of Pierre Elliott
Trudeau. Second, Trudeau is not the president of Canada but its prime minister. How,
then is "former wife of the prime minister" the same as "wife of the president"?
        Before you answer, "Well, `wife' and `former wife' are related concepts, as are
`prime minister' and `president"', consider who might be said to be the current First
Lady of Britain. Whose name comes to your mind? Margaret Thatcher? Queen
Elizabeth? They are women, but do they really play the role of First Lady? How about
Denis Thatcher or Prince Philip? At first these suggestions seem silly, but in a strange
way they start to seem compelling, particularly the thought of Denis Thatcher. In fact,
I once




Analogies and Roles in Human and Machine Thinking                                    547

clipped a newspaper article that portrayed Denis Thatcher as Britain's First Lady.
         What kind of sense does this make? How can a male be a lady? Well, language
is far slipperier than dictionary definitions would have you believe. Its slipperiness
comes from the underlying slipperiness of concepts, in particular these elusive things
we are calling roles.
         Of course, you could argue that what "First Lady" really means is "spouse of
the head of state", and so the First Lady role goes over without any trouble into
"husband of the prime minister". But this won't do either. In Haiti, until recently the
title of First Lady belonged to Simone Duvalier, the wife of the late former president,
Francois ("Papa Doc") Duvalier. She is also the mother of the current president, Jean-
Claude ("Baby Doc") Duvalier. Not long ago there was a bitter power struggle
between Simone Duvalier and her daughter-in-law Michelle Bennett Duvalier, the
wife of Baby Doc, for the title of First Lady. In the end, the younger woman
apparently gained the upper hand, taking the title "First Lady of the Republic" away
from her mother-in-law, who in compensation was given the lifetime title "First Lady
of the Revolution".
         Do you want to amend your suggestion so that it will say "spouse or parent,
present or former, of a head of state, present or former"? You know perfectly well that
we'll be able to come up with other exceptions. For example, imagine a meeting of the
Pooh-Bah Club at which the Grand Pooh-Bah's favorite aunt was introduced as the
First Lady of the club. Of course, the Grand Pooh-Bah is hardly a head of state, and so
you could amend your definition to say "spouse or favorite relative, present or former,
of the head, present or former, of any old organization". But suppose ... Actually, I
think I'll let you go on inventing exceptional cases. For any rule you propose, there is
bound to be some conceivable way to get around it.
         Worse yet, something terrible is happening to the concept as it gets more
flexible. Something crucial is gradually getting buried, namely the notion that "wife
of the president" is the most natural meaning, at least for Americans in this day and
age. If you were told only the generalized definition, a gigantic paragraph in legalese,
full of subordinate clauses, parenthetical remarks, and strings of or's -the end product
of all these bizarre cases-you would be perfectly justified in concluding that Sam
Pfeffenhauser, the former father-in-law of the corner drugstore's temporary manager,
is just as good an example of the First Lady concept as Nancy Reagan is. When this
happens, something is wrong. The definition not only should be general, but also
should incorporate some indication of what the spirit of the idea is.

                                       *   *   *

        Computers have a hard time getting the spirit of things; they prefer to know
things to the letter. And so people spend an enormous amount of time




Analogies and Roles in Human and Machine Thinking                                   548

talking to computers, writing long and detailed descriptions of ideas they could get
across in one good example to anyone with half a brain. So a challenging question is
how to get a computer to understand what is meant by "First Lady". For this we need
to examine the idea of "roles" in detail.
        In order to illustrate how the notion of "role" can be modeled in domains more
formal than that of political protocol, I now will switch to one of my favorite
domains: the natural numbers. I will present some puzzles that Gray, Marsha, and I
have been thinking about. Each of them has a set of possible answers with varying
degrees of plausibility or defensibility. We are working on a computer program that is
able to see the rationale behind each possible answer, and thus is able to come up with
the same set of "feelings" as a typical person would have, about what is a good
answer and what is a bad one.
        The domain of natural numbers might sound at first like a hard-edged,
objective mathematical world, but actually it is a domain in which problems requiring
extremely subtle subjective judgments can be formulated. We have given our program
very little detailed arithmetical knowledge about the integers. The program does not,
for example, recognize 9 as a square; in fact, it doesn't even know about
multiplication! It does not know that 6 is even and 7 is odd. So what does it know? It
knows how to count up or down -that is, it has a knowledge of successorship and
predecessorship. Thus it recognizes that the sequence of numerals "12345" represents
an upward counting process. It is also able to apply the notion of counting to
structures it is looking at, as in "44444", which it could recognize as a group of five
copies of the numeral `4'. It knows that 9 is bigger than 4, although it has no idea how
much bigger. (Subtraction and other arithmetical operations are unknown to it.) You
can think of our computer program as having the arithmetical sophistication of a five-
year-old and an avid curiosity about number patterns. (By the way, it is not tied to or
affected by decimal notation. The number 10 is not considered any more special than
the number 9.)
        Here is the first problem (invented, as were many of the following ones, by
Gray). Consider the following structure, which we'll call A:

              A: 1234554321

Now consider the structure called B:

              B:    12344321

The question is: What is to B as 4 is to A? Or, to use the language of roles: What plays
the role in B that 4 plays in A?
        Note that by asking it this way, we leave it to the puzzle solver to decide what
role 4 actually does play in A. It would be analogous to asking "Who is the Nancy
Reagan of Britain?", leaving it to the listener to figure out what




Analogies and Roles in Human and Machine Thinking                                   549

conceptual role Nancy Reagan fills, and then to try to export that role to Britain. I
have found that many people who balk at calling Denis Thatcher the "First Lady of
Britain" are quite content with calling him the "Nancy Reagan of Britain". A curious
point that this illustrates, and to which we will return, is this: If the role is left implicit,
nonverbalized, it has more fluidity in the way it transfers than if it is "frozen" in an
English phrase.
        As a matter of fact, most analogies crop up in this type of nonverbal way.
Seldom does someone say to you explicitly: "What is the counterpark of Central Park
in San Francisco?" Usually it happens through a more implicit channel. When you are
visiting San Francisco for the first time, you are driven through Golden Gate Park,
and somehow it reminds you of Central Park. After the fact, you can point out some
shared features: both are long thin rectangles; both contain lakes, curving roads, and
excellent museums; and so on. Most analogies arise similarly-as a result of
unconscious filterings and arrangings of perceptions, rather than as consciously
sought solutions to cooked-up puzzles. To put it another way, to be reminded of
something is to have unconsciously formulated an analogy.
        Incidentally, when I first thought of writing about roles and analogies, I had in
mind both the First Lady example and the numerical examples. As my thoughts
evolved, I realized I was unconsciously developing a parallel in my mind between the
First Lady example and the numerical examples. I'll call it a "meta-analogy", since it
is an analogy between analogies. In this meta-analogy, I see structure A as
corresponding to the United States, structure B to Britain, 4 to Nancy Reagan and the
unknown number to the unknown person. We'll come back to the meta-analogy later
on.

                                           *   *   *

       Let us now look at some possible answers to the first number-analogy
problem. The most sensible answer is 3-and fortunately, it is also the most frequently
given one. The usual justification is that 4 precedes the central pair (55) in A, and the
corresponding central pair in B is 44, which is preceded by 3. Well, then, what would
you say for C? What is to C as 4 is to A?

                C:   12345666654321

        The central pair in C is 66, which is flanked by 6's. Is 6, therefore, to C what 4
is to A? Well, most people probably would prefer 5, although it is perfectly logical to
insist on 6. The preference for 5 comes, nonetheless, from a very sensible (and also
logical) instinct to generalize the notion of "central pair" (itself, to be sure, a role) to
"central plateau" (or whatever you want to call it). There are competing urges: first, to
stay with the exact original concept, and second, to flex and bend when it "feels
right", when it would seem rigid and stodgy to insist on established conventions over




Analogies and Roles in Human and Machine Thinking                                          550

simple and "natural" extensions. But it is just these sorts of terms-"flex", "bend",
"feels right", "rigid", "natural", and so on-that are so extraordinarily hard to put into
programs, logical though programming might be.
        Now let us investigate some other ways to make the role of 4 slip. Consider
this structure:

               D:    11223344544332211

         Here is a curious kind of reversal; now there is no central pair-yet everything
else is in pairs. Some people might still pick 4, since it is next to the center. But what
about 44, a pair rather than a single number? After all, as long as "pair" and
"singleton" have switched places, we might as well go all the way and give an answer
that reflects this perceptual turnabout. In fact, it would seem rigid and unimaginative
to insist on sticking with single numbers when it is so obvious that the easiest way to
perceive D is in terms of pairs:

               1-1 2-2 3-3 4-4 5 4-4 3-3 2-2 1-1

        Not just 4 but every part of A has a role, and there are corresponding roles in
D. As you can see, within each role the concepts of pair and singleton have been
switched.
        Now is as good a time as any to return to my meta-analogy and to point out
some correspondences between these problems and the First Lady problem. If you
think of the president as "the highest, most central figure in the land" and his wife as
"the one standing next to him", you will see that this characterization carries over
almost literally to the numerical problems. In structure A, the highest, most central
figure-the "president"-is 5 (or possibly the pair of 5's) and his "wife", standing next to
him, is 4. In B, the president is 4 (or the pair of 4's) and his wife is 3. In C, the
president is 6 (or the group of 6's), and his wife is 5. In D, the president is (for once)
unambiguous (5), but to compensate, there is a dilemma concerning the identity of his
wife. If you think of pairs as males and singletons as females, then D presents us with
a case where the sexes are reversed, exactly as in the First Lady of Britain problem.
The most reasonable answer seems to be the "spouse" (in this case, the husband) of 5,
namely the pair 44.
        Consider now the following couple of curious cases:

               E: 12345678

               F: 87654321

       What can we make of these? A very rigid person might cling to the idea
captured in the phrase "number to the left of the central pair", despite the fact that
nothing at all distinguishes the central pair in either of these




Analogies and Roles in Human and Machine Thinking                                     551

examples. Such a person would give the inane answers of 3 for E, and 6 for F. Such a
person would do better to take up football instead of analogies, as Lewis Carroll's
Tortoise once remarked to Achilles.
        But what would be a wiser view of, say, E? How should one map E onto A?
Any mapping is doomed to be imperfect, so how can we do it best (that is, with the
least pain or frustration)? We might think of E, since it rises uniformly, as mapping
onto the left half of A. This would involve a tacit judgment that it is all right to
abandon the attempt to map E onto all of A, in return for the ease of mapping E onto a
"natural" portion of A. That is a pretty subtle step to take, I would say. It would
suggest 7 as the answer.
        Well, what about F, then? Do we prefer 2 or 7? It depends on whether we
choose to map F onto the left half or the right half of A. Mapping F onto A's left half
involves mapping a descending sequence onto an ascending one. But either choice
requires a willingness to let go of qualities that had seemed important, a willingness to
bend gracefully under pressure. Fluid analogies are not a game for rigid minds!
        These kinds of situations are difficult because in essence they call for splitting
the role of 4 in A into two rival facets. In the mapping of A onto F, one of the rival
facets sees 4's role in A as "one less than the president", whereas the other facet sees
4's role as "the next-to-rightmost element of a staircase". Thus one facet is primarily
concerned with magnitude and the other primarily with position. The facet you find
more important will determine your answer to F.
        Pretty much this kind of split happened when you tried to decide whether the
First Lady of Britain was Queen Elizabeth or Margaret Thatcher-or one of their
husbands. Is being a figurehead or being a head of state more likely to make
someone's spouse a First Lady? In the United States, these features coincide in one
person (the president), but in Britain they do not. Consider the following target
structures:

               G: 5432112345
               H: 123465564321

In G, what is most central is simultaneously lowest, and what is highest is
simultaneously most peripheral! (G can be pictured as a valley and A as a mountain
peak.) We have a "ceremonial figure" (the 5's flanking the structure) and we have a
"head of state" (the two central l's). Which one's spouse would better fill the role of
First Lady? Or, to put it most simply: "What in G plays the role of 4 in A?" I
personally would opt for 2 because it stands next to the central group. To me,
centrality seems more important here than magnitude, just as political power seems
more substantive than ceremonial show. Correspondingly, I would opt for Denis
Thatcher rather than Prince Philip as "Britain's Nancy Reagan".
        Now what happens when we tackle H? There are three "reasonable"




Analogies and Roles in Human and Machine Thinking                                     552

possibilities (in the sense of appealing to law's proverbial "reasonable human"): 6
(flanking the central pair of 5's), 5 (being the next-to-largest number) and 4 (flanking
the central "crater" 6556). Once again there is no Gloriously Right Answer, but there
are certainly ideas that seem good and ideas that seem shaky. For instance, if someone
suggested, "The answer is 4, because 4 is the fourth term of H, just as it is the fourth
term of A", I would be nonplussed. That would be a childish generalization based on
the most superficial of scans of both structures. It would be as childish as inferring
that, since our national holiday falls on the fourth of July, other countries' national
holidays would also have to fall on the fourth of various months. To see 4 as no more
than the fourth element of A is to ignore all of A's structure. It is to see A as nothing
richer than this:

               ooo4oooooo

A good answer must take A's structure into account in a full, rich, and yet simple way.
This means that, to the extent it is possible, all of A must be perceived in terms of
interacting, mutually intertwining conceptual structures-roles that are mutually
dependent, in the way that "family", "husband", "wife", "mother", "father",
"daughter", "son", "brother", "sister", "relative", "in-laws", and so on are all
interdependent concepts.
        The word "role" makes us think of the theater. In a play, the various roles all
mingle together in scenes. A scene is a larger-scale structure than an individual role; it
is a place where several roles coexist and interact. In our analogy problems, one might
try to conceive of the two structures involved as if they were two enactments of a
single scene, portrayed by different directors working with different actors. Thus the
core roles would exist and would be filled in both presentations, but at the same time
each presentation would have minor aspects, or roles, unique to it. For example, the
adaptation of the Greek legend of Orpheus and Eurydice into a contemporary context
of carnival time in Rio de Janeiro is the basis for the movie Black Orpheus. Many
original features cannot be directly exported, but with poetic modification they can be,
and the director, Marcel Camus, met the challenge with great flair. In the movie there
are, of course, many minor parts-extras -that add Brazilian flavor, yet they do not
impair the analogy at all; in fact, they enrich it. This is the kind of thing that appeals
deeply to human .sensibilities, both intellectually and emotionally.

                                        *   *   *

       Now that you have seen some variations, I would like to return to our first
puzzle and point out some of its hidden subtlety. First of all, the "central pair" notion,
which functions as the keystone of structure A, is actually just a kind of by-product,
an accidental artifact of the structure of A. To see what I mean, consider this question.
How would you efficiently describe the




Analogies and Roles in Human and Machine Thinking                                     553

structure of A (without quoting it digit by digit)? You would probably say it rises
from 1 to 5 and then falls from 5 to 1, making two halves that are mirror images.
Nowhere in- this description was there any mention of some kind of central pair or
central plateau. It was not needed; one will just appear automatically there when
anyone follows your description. In fact, anyone who constructs a copy of A is very
likely to see a central plateau, even if the concept was never suggested to them. To the
mind's eye it appears something like this:

               1 2 3 4 5-5 4 3 2 1

       Somehow a new percept has been born in the center. It is, as I remarked above,
the keystone of A. (Note that the concept of "the keystone of A" depends on, or
implies, a mapping of A onto an arch-yet another analogy.)
       Why don't we perceive the pair of 3's, say, as a unit as well? Probably simply
because they do not touch. And consider this structure:


        The central pair-"5 1"-doesn't pop out as being salient or important, does it? In
A, though, the combination of adjacency and equality, particularly when
supplemented by centrality, somehow makes the two central 5's merge into a unit in
the perceiver's mind, albeit usually not at a conscious level. Still, if this perceptual
shift did not happen, then the answer of 5 for C, based largely on equating the
plateaus in A and C, would be considerably less compelling.
        In the first puzzle, both A and B had obvious central plateaus. This suggested
a good starting point for an overall mapping of A onto B: central plateau onto central
plateau, start onto start, finish onto finish, and so on. But if we tried to complete this
mapping, we would obviously run into trouble:



        We must have 1 in A mapping onto 1 in B, no? And the centers have to match
up too, don't they? But where between 1 and 5 does the analogy break down? It seems
that some kind of mapping of 4 onto 3, as is shown above, is satisfying to many
people. But press them one step more, and they will shrug, grin, and give up.
        Similarly, although you can ask for "the Nancy Reagan of Britain", it makes
less sense to ask, "Who is the Maureen Reagan of Britain?" (Remember that Maureen
Reagan is Nancy Reagan's stepdaughter.) Suppose the Thatchers had a biological
daughter. Would she be the counterpart of Maureen Reagan? Or suppose Margaret
Thatcher had a stepdaughter. Would she be the counterpart? Then again, suppose that




Analogies and Roles in Human and Machine Thinking                                     554

Margaret Thatcher had no daughter but that Denis Thatcher had twin stepdaughters.
Would these twins, taken together, constitute the counterpart of Maureen Reagan?
(How can two people fill a role defined by one person? Well, think of example D,
where the pair of 4's played the role of a single 4 in A. Or think of many European
countries, which have both a president and a prime minister.)
         Issues like this crop up all the time in the pursuit of good analogies, and facing
up to mismatches leads occasionally to productive insights. One could go on and press
for even more detailed correspondences between entities in Britain and in the United
States. What is the British counterpart of Watergate? Who plays the part of Richard
Nixon? Of John Mitchell? Of Senator Sam Ervin? Of Senator Daniel Inouye? Of G.
Gordon Liddy? Of Judge John Sirica? Of John Dean? Of Officer Ulasewicz? Of
Alexander Butterfield? The less salient an object is inside a larger structure, the harder
it is to characterize in an exportable way.
         But what makes something salient? As a rule, it is its proximity, in some
sense, to a "distinguished" element of the larger structure. Consider the following long
structure:


The central 4 is probably the most distinguished individual numeral. Then, depending
on how you perceive the sequence, different features will leap out at you. For nstance,
do you see it as "letters" or as "words" (larger-scale chunks of the sequence)? When I
see it at the "word" level, the central group "3334333" seems just a shade less salient
than the 4, and after that, perhaps, the two flanking groups of l's. The two groups of
3's by themselves come next. Only then do the groups of 2's get recognized. On the
other hand, when I perceive the sequence at the "letter" level, what is salient is quite
different. After the central 4, probably the next most salient numbers to me are the
first and last l's, since they are very easy to describe-then maybe the first and last 2's.
After that, the two 3's flanking the central 4but at this point it starts to get a little
harder to specify various items without resorting to such uninspired descriptions as
"the fourth term".
        A distinguished item is something we can get at via an elegant, crisp,
exportable-sounding description. A nearly distinguished item is something we can get
at by first pointing to a distinguished item, and then, in an exportable way, describing
a short "jog" that leads to it. Just as in giving someone directions, some places are
more salient, others are less so. Some buildings in New York City are inherently
difficult to direct someone to, others are inherently easy. In the same way, some roles
in a complex conceptual structure are highly distinguished and easily exportable,
others are very hard to describe. Although they may have certain idiosyncratic
qualities in their local context, nothing makes them stand out globally.
        As you move progressively farther away from its central roles, any analogy




Analogies and Roles in Human and Machine Thinking                                      555

becomes increasingly strained. For example, "Who is theJackie Washington of
Britain?" Should we begin by getting out the London telephone book and looking
under "Washington, J."? Or should we look under "London, J."? Or is it a
meaningless question, meaningless even to Jackie's best friend? After all, Jackie's role
may just be too small and idiosyncratic within the structure of the United States. It is
not exportable. The fact that Jackie is the manager of Gearloose's Hot Dog Stand in
Duckburg does not help much, because one still has to figure out the identities of the
British Duckburg and the British Gearloose-not to mention the British equivalent of
hot dog stands!
        The moral is a simple one: Don't press an analogy too far, because it will
always break down. In that case, what good are analogies? Why bother with them?
What is the purpose of trying to establish a mapping between two things that do not
map onto each other in reality? The answer is surely very complex, but the heart of it
must be that it is good for our survival (or our genes' survival), because we do it all
the time. Analogy and reminding, whether they are accurate or not, guide all our
thought patterns. Being attuned to vague resemblances is the hallmark of intelligence,
for better or for worse.

                                       *   *   *

        The fact that we use words and ready-made phrases shows that we funnel the
world down into a fairly constant set of categories. Often we end up with one word,
such as "kitchen". In general, two kitchens will not map onto each other exactly, but
we still are satisfied with the abstraction "kitchen". Generally speaking, a kitchen will
have a sink, a stove, a refrigerator, cupboards, counters, drawers, and so on. In the
United States, it is very common for people to assume that the garbage will be in a
cupboard below the sink. The idea of "the cupboard below the sink" is a perfect
example of an exportable role. In fact, isn't your sink the "president" of your kitchen?
And .. .
        Our language provides for mappings of many degrees of accuracy. Some
people, when they see Bossie, see no further than "cow" and accordingly use that
word; others notice that Bossie is female, and will say "heifer". Still others perceive
the breed as easily as they perceive Bossie's "cowness" and talk about "that Angus
heifer". A famous Dublin zookeeper, Mr. Flood, was once asked the secret of his
great success in breeding lion cubs. "Understanding lions." said he. "And in what does
understanding lions consist?" he was asked. His reply: "Every lion is different." This
curious answer denies the category while taking advantage of it. But that is the nature
of categories. Their validity can at best be partial. No matter at what level of detail
you cut off your scrutiny, your perception amounts to filtering out some aspects and
funneling the remainder into a single conceptual target, a mental symbol often labeled
with just one word (such as "word") or stock phrase (such as "stock phrase"). Each
such mental symbol implicitly stands for the elusive sameness shared by all the things
it denotes.




Analogies and Roles in Human and Machine Thinking                                    556

        Beyond the implicit analogies hidden in individual words or stock phrases,
explicit analogies occur all the time on a larger scale in our sentences. We are quite
uninhibited in comparing unfamiliar things with things we assume are more familiar.
We see grids of all kinds as being similar to checkerboards. We see carefully charted
actions in life as being similar to chess moves. We see the eye as a camera, the atom
as a tiny solar system. Science is constantly being likened to a vast jigsaw puzzle (an
analogy I have never cared for). In their eagerness to stretch and bend concepts,
people turn proper nouns into common nouns, as in the statement "Brigitte Bardot is
the French Marilyn Monroe." In such linguistic flexing, both la Bardot and the
Monroe suffer somewhat in the interests of vivid imagery.
        Then, going one step beyond the explicit linguistic level, there are the
analogies and mappings that we use constantly to guide our thoughts on a larger scale.
The perception of romantic dilemmas is one of the most striking places where
mapping or analogical thinking dominates in an obvious way. When someone tells us
of some romantic woe, we can usually map it immediately onto some experience of
our own. In fact, we can probably draw some parallel between any romantic situation
and any other one, and such a mapping will perhaps yield some insight if it is carried
out well. Yet romances are incredibly detailed and idiosyncratic. The point is that we
throw many details away; we skim off some abstractions and are careful not to try to
carry the resemblance too far. And certainly we ignore the trivial aspects. A romance
between Chris and Sandy can certainly map onto one between Pat and Chris or one
between Sandy and Pat, despite the fact that names, hair colors, and other superficial
features do not match!
        The reason, then, for worrying about human analogical thought is that it is
there. To ignore it would be like ignoring Everest in trying to understand mountain
climbing.
        . Let us get back to some concrete problems in our more formal, numerical
domain. Notice that there is an inherent kind of contradiction in setting up analogy
problems-which, after all, are informal by definition-in a rather formal domain. But
the nice thing is that it shows that the domain is actually just as slippery as any
"informal" domain. Here are four further examples:

       I:    123345676543321

       J:    177654321

       K:    697394166

       L:    123456789789654321

       Example I involves what I enjoy referring to as a "governor", namely the pair
33. Here again, one role in A has been split into parts: 55 in A was not




Analogies and Roles in Human and Machine Thinking                                  557

only the sole pair but also the peak, whereas in I, 33 is the sole pair and 7 plays the
role of the peak. We are forced to choose between 2 (the wife of the governor) and 6
(the wife of the president). Actually, the governor has two "wives"-2 and 4, a "left
wife" and a "right wife"-and so we have to choose between them, unless we go with 6
as being the wife of president 7.
        Example J, beginning as it does with "1776", is a patriotic puzzle. (What is its
British counterpart?) Its interest is primarily in that it draws attention to A's
symmetry, which we had taken for granted. When we chose 4 as the president's wife,
were we taking his left wife or his right wife? In A, of course, they coincided, so it
didn't matter. But in J they differ: 1 would be the left wife, 6 would be the right wife.
Because of a tendency to be influenced by our left-to-right scanning, we probably
would choose the left wife under normal circumstances, but here, there is such great
asymmetry that we pause. The regular descent from 7 to 1 corresponds far better to
A's "staircase" structures than does the abrupt leap upward (from 1 to 7 in one step!).
For that reason, 6 probably wins over 1, in this case.
        Example K is a bit obscure, but it has been led up to by example J. In
particular, example j drew attention to the fact that in A there are two 4's, not just one.
Example K plays on the relation of those two 4's to each other. In A, there were two
elements between the two 4's. We can take that property as defining the role of 4 in A.
To be sure, that is not the only relation between the two 4's, but it is the most obvious.
If you "turn off" everything in A but the 4's, you will get an image something like
this: ooo4oo4ooo. That image makes the size of the interval between the 4's stand out.
Given this way of looking at the role of A's two 4's, what in K corresponds? There is
only one number that occurs exactly twice, and its two appearances are separated by
two numbers. That number is 9. If you turn off everything but the 9's in K, you get
this picture: o9oo9oooo. That may or may not be sufficient reason for you to choose 9
as the K-counterpart of A's 4.
        Finally, consider example L. Here, the central-pair notion gets extended one
further degree of abstraction. We go up, step by step, until we hit the second 7. Jolt! It
takes us a moment to get our bearings, and when we recover, we realize that the
central pair consists not of single integers but of "clumps" or "chunks": namely, two
copies of the unit 789. We can aid the eye this way:

       1 2 3 4 5 6 7-8-9   7-8-9 6 5 4 3 2 1

Now the answer seems glaringly obvious: It is 6! On the other hand, maybe we were
supposed to get the hint offered us generously by the central pair. And what was that
hint? It is that we could perceive the whole structure in triples, not just its center. In
this case, L reparses into

       123 456 789 789 654 321




Analogies and Roles in Human and Machine Thinking                                      558

Now the answer should be obvious-except that we are still left with a minor dilemma.
Do we take the president's right-hand wife (654) or his left-hand wife (456)? I am
biased by left-to-right chauvinism and would choose 456. Many people, however,
refuse to see the sequence in triples and stick with 6.
        Here is an innocent-seeming puzzle that points to still more complex issues:

       M:    123457754321

The way I see it, the best answer is 6. You might object, "Why not 5? 6 isn't even
there!" True, but 6 is conspicuous by its absence. The 4 in A precedes the 5 not only
typographically but also arithmetically: 4 is the numerical predecessor of 5. And what
is 5 in A? It could be seen either as the maximum in A or as the number forming the
central pair of A. Both carry over to M, yielding 7 as M's 5. Now, if you choose to see
4's role in A abstractly and arithmetically rather than concretely and typographically,
you can carry your vision directly over to M. Then candidate 6 must be considered a
strong competitor to 5. In my mind, it wins.
        This example opens up an whole new world of levels of abstraction in the
perception of structures. To illustrate briefly, let me propose the following structures:

       A': 1234445678987654444321
       B’: 1112343211

And here is the puzzle: What in B' plays the role that 7 plays in A'? Well, 7 occurs
twice in A', but certainly it seems to play no salient role. As a numeral in A', 7 has no
outstanding characteristic, and so at first its role seems hard to export. However, 7
enters into the structure of A' in another way, and a salient one at that. One of the
most salient features of A' is its large number of 4's. Count them. How many? Seven.
Aha! Thus 7, in its capacity as an invisible counting number rather than as a visible
numeral, plays a very distinguished role in structure A'. Still, is it possible to export
this role to B'?
        We have to decide how to characterize (in an exportable way) just what it is
that 7 is counting. To insist that it must be the number of 4's seems a little parochial,
to say the least. Who says that 4 is the 4 of B'? Perhaps a deeper and more fruitful
way to look at matters is to see 4 as A"s most frequent term. This leads us to look for
the most frequent term in B'; this is 1. So 1 is the 4 of B'. Therefore the counterpart of
7 would be the number of l's in B'-namely, 5-another "invisible" answer, in that 5 does
not appear as a visible numeral in B'.
        It is narrow-minded to insist that there is a big distinction between being
present as a visible numeral and being present in a more abstract sense, for example,
as a counting number. To put it another way, 5 is invisible in B'




Analogies and Roles in Human and Machine Thinking                                     559

only if you think of vision as having no cognitive component, as if we could perceive
only numerals. In fact, with our eyes we are constantly "seeing" abstract qualities.
When we look at a television program, we see more than flickering dots: we see
people. Of course, somewhere deep down in the processing there are components of
our visual system where the dots themselves are "seen" as dots, but ironically, we
would hesitate to describe as "vision" what retinal and other cells do. Vision implies
going beyond the dots; in other words, beyond the primitive visual level. We can
"see" that a certain chess position is ominous, that a certain painting is by Picasso,
that someone is in a bad mood, that this car won't fit in that garage, and so on. If we
accept this notion that vision is imbued with a cognitive component, then we can
agree to "look beyond the numerals". In that case, 5 is directly visible in B'!
        By the way, I carefully drew up A' so that 7 would appear as a numeral in it
(as well as counting the number of 4's). This threw in a complicating factor,
something one had to ignore. I could have had A' have, say, 12 4's, in which case 12
would have "appeared" in A' only at the abstract level of a counting number and not
as a numeral. But real life is seldom so considerate of would-be analogy-perceivers.
For example, in thinking about the question "Who is the Nancy Reagan of Britain?",
you might have felt that this was much harder than "Who is the First Lady of
Britain?" because you may attach certain uniquely personal attributes to Nancy
Reagan, over and above seeing her as the First Lady of the United States. What if I
had asked for "the Eleanor Roosevelt of Britain"? Or, turning the tables, "Who is the
Moshe Dayan of the United States?" I am almost tempted to answer "Douglas
MacArthur", like Dayan a famous and successful general and a controversial political
figure, but then I remember-MacArthur had two eyes! Dayan's eye patch is perhaps
his most memorable feature.
        It is interesting to go back to earlier examples, mapping them onto A' and
asking, "What here plays the role of 7?" You will perceive those old structures
through new eyes (or glasses). I leave a few challenging examples for you to map
onto A and A':

       P: 5432154321

       Q: 543211234554321

       R: 12349876543

       S:   112233445566771217654321

       T:   1234123121213214321

       U:   211221222291232




Analogies and Roles in Human and Machine Thinking                                  560

You might enjoy making up some examples of your own, which potentially might
lead a solver to further unexpected modifications of the perceived role of 4 in A. For
instance, can you devise an example in which it becomes sophisticated, rather than
childish, to perceive 4 as the fourth element of A?

                                          *   *   *

        One of the purposes of these puzzles is to dispel the notion that the full, rich,
intuitive sense of a role, such as that of 4 in A or that of First Lady, can be easily
captured in words. In fact, it might be more accurate to assert the contrary: that
precisely in its nonverbalizability lies its fluidity, its flexibility. This is a crucial idea.
Consider how you would try to capture in some phrase the precise way you see "what
4 is doing" within A. No matter what phrase you give, someone will be able to
concoct another example in which your phrase does not enable anyone to predict what
you will perceive as being analogous to 4. A frozen verbal phrase is like a snapshot
that gives a perfect likeness at one moment but fails to show how things can slip and
move. There is something much more fluid in the way a mind represents the role
internally. Various features are potentially important in defining the role, but not until
an example comes up and makes one feature explicit does that feature's relevance
emerge.
        We make comparisons all the time. It does not seem particularly noteworthy
when someone walks into your kitchen for the first time and says, "I like the way your
kitchen is laid out better than the way mine is. My kitchen has windows over there
and the stove is right here, so it is less convenient and the light isn't so good in the
morning." Clearly the words "here" and "there" conceal implicit mappings of the two
kitchens, otherwise the statement would be utter nonsense. Words like "this" and
"that" and phrases like "that sort of thing" are even better at picking up intangible,
flexible, implicit meanings that can be transported across the borders of situations
differing widely from each other. And that's the name of the game, in thought.
        Right now it seems that what artificial intelligence needs is a way to go
beyond "delta function" programs: programs that are virtuosos in a very narrow
domain but that have no flexibility or adaptability or tolerance for errors. I call these
programs "AE programs": programs that have Artificial Expertise. The trouble with
them is that they are always brittle and narrow. t seems that a careful study of
judgmental processes in even so simple a domain as these curious number analogies
would afford fascinating insights ,unto how computer programs might be made to
approach the flexibility and generality of our own minds.
        To show what I mean, I would like to conclude with a verbatim transcript of a
conversation I had with a friend a while back. It ran this way:




Analogies and Roles in Human and Machine Thinking                                         561

FRIEND: Last Friday afternoon I was over at the Pooh-Bah Club listening to a piece
  on the radio that I was sure was Shostakovich. When it ended and they announced
  it, sure enough, it was! I was thrilled, because that kind of thing has happened to
  me only a couple of times in my life!
ME: That kind of thing? You mean, being at the Pooh-Bah Club and hearing a piece
  on the radio that you thought was Shostakovich on a Friday afternoon?
FRIEND: You're so dense! When those Scientific American people hear about that,
  they probably won't want you to write any more articles for them.
ME: Yeah, I should have known that it didn't have to be on a Friday afternoon.
FRIEND: You should have known that it didn't have to be Shostakovich!

        Quite coincidentally, a recently perfected natural-language computer program
called CORTEX happened to be eavesdropping on us, and it just could not resist
chiming in at this point, saying, "Oh say, that reminds me -something really similar
happened to me the other day. I was at a club whose name is hyphenated, and the
water cooler broke. Ain't that something!" Well, that kind of thing is what I would
like to see artificial intelligence programs doing more of.


Post Scriptum.
Verdi is the Puccini of music.                The knee is the Achilles' heel of the leg.
               -Igor Stravinsky                      -Pasadena (Calif.) Valley Values

        The AI work out of which this column grew was my "Seek-Whence" project.
My original goal was to develop a program that would take as input a sequence of
integers such as 1, 4, 9, 16, and that would detect the underlying law-thus, it would
"seek whence" the sequence came, and would extend it. Over a period of time, it
became obvious that certain aspects of the goal were more central to mentality than
others. In particular, it became clear that the ability to quickly discover the law behind
highly mathematical sequences (even lowly mathematical ones, like the squares) is a
specialized skill that bears little relation to mind in general, but that the ability to
quickly spot patterns (as in "1 2 2 3 4 4 5 6 6") is absolutely indispensable.
        Thus the Seek-Whence project retargeted itself on structures composed of
smallish integers; and the major effort became one of figuring out how to perceive
and "parse" such structures as these:


The latter two examples nicely illustrate one of the major problems to confront: that
of boundary location. Where do you draw the lines separating




Analogies and Roles in Human and Machine Thinking                                     562

one substructure from another one? And what are good ways of restructuring or
regrouping if a first try fails? That whole area of concern is reflected in the project's
very title, "Seek-Whence", which violates the conventional syllabic structure, namely,
"see-kwents", regrouping it as "seek-wents". This kind of difficulty permeates the
efforts to mechanize continuous speech recognition, and is very familiar to anyone
who has been in the position of listening to a foreign language they have studied but
don't know well. Often sounds will flow by so fast that you have absolutely no idea
what is being said, simply because you cannot tell where the word breaks are; what is
most frustrating is that this can happen even if everything being said would be
perfectly clear if you saw it in writing (where word breaks are handed to you on a sil
verplatter). The "parsing" of visual input is likewise permeated by boundary-location
problems. Music is another such domain, and in fact the discrete multi-level patterns
of melodies were among the biggest inspirations for the Seek-Whence domain.
        At one point, in trying to get across the idea of Seek-Whence to someone who
had a distaste for integers, I simply substituted letters for integers (a for 1, b for 2,
etc.), and made up some parsing and analogy problems. For instance: "What is to
abcddcba as d is to abcdeedcba ?" Some people might say that this problem is similar
or analogous to the first numerical analogy problem given in the column; I would say
it is the same problem. Yes, in different clothing, if you like, but the same all the
same. Numerals, capital letters, smalls-what's the difference? At least that was my
feeling. Yet I found that I could usually awaken more interest in people if these
analogy problems were presented in terms of letters instead of numbers. Groan!
        From potentially infinite sequences and the rules behind them, my focus of
attention gradually shifted to rather short sequences and the roles inside them (as I
emphasized in the column). This concentration on roles and analogies then became so
dominant that the Seek-Whence work revealed itself to be primarily a project on
perception of analogies. Once this was out in the open, I decided to reify that concern
by creating a new project, which I dubbed "Copycat", the idea being that being a
copycat, when you're a child, is a universal and primordial experience in doing simple
analogies. If I touch my nose and say to you, "Do this!", what will you do? Most
people will touch their own nose. But why not touch mine? If I touch your nose, what
will you do? Touch your own, or mine? And so on. A set of variations on this theme
is shown in Figures 24-1 and 24-2. You can think of them as symbolizing the entire
Copycat project.
        One can be more flexible or less in how one interprets "Do this!" What oes
one take literally, what does one see as playing a role in a foreordained and familiar
structure? What kinds of familiar structures is one willing to see as identical to each
other? When is it necessary to start inventing new ways of perceiving a given
situation in order to fit it into pre-existent frameworks, which then allow already-
familiar roles to emerge? What remains firm, and What slips? What sticks, and what
gives? These kinds of questions sound rather abstract, but when real analogies are
manufactured, they are the chief




Analogies and Roles in Human and Machine Thinking                                    563

FIGURE 24-1. The "Do this! "problem. In (a), Tom touches his nose and says to
Annie: "Do this!" She wonders what she should do. In (b) and (c) you see how two
Annie-clones respond. What would you do?




Analogies and Roles in Human and Machine Thinking                           564

FIGURE 24-2. More "do-this " questions. In (a), a three-headed Annie is at a loss for
what the best way to "do this " is. In (b), a long-necked Annie and her giraffe friend
wonder what to do. In (c), Fanny the Fish-handless and noseless, but having gills-
muses how to copy Tom. Finally, in (d), how in the world can poor little Elephannie
do what Tom is doing?




Analogies and Roles in Human and Machine Thinking                                 565

concerns of the analogy-maker, whether at a conscious level or not. Therefore these
issues are the heart and soul of the Copycat project, and the purpose of this P.S. is
mainly to show exactly how that is so.

                                        *   *   *

        One serious problem with studying real-world analogies like the "First Lady"
examples is that they bring along too much baggage-too many complex tie-ins to all
sorts of concepts. Another serious problem is that when we study real-world analogies
that real people actually have made-, we are nearly blinded to their essence because
we have nothing to contrast them with. In order to see what makes human-made
analogies good, we have to see the alternatives: analogies that humans would never
make, even in jest. Can you imagine, for instance, an organism trying to understand
the experience of being pregnant by likening it to being an elevator (or a football
stadium, for that matter) containing one person? Something is very wrong with such
an analogy-but what? Try to formulate principles of analogy-building that would
suppress that type of analogy (whatever that means!) but that would recognize the
insight in this one: "Try speaking English for a while without using the letters `e' or `t'
at all. Then you'll have an inkling of how Japan felt, having two of its major cities
wiped out." Both these analogies involve mapping a thinking organism onto
something very alien to it, and on a totally different scale. Yet one succeeds well and
the other flops well. How come?
        In such cases, filled to the brim with myriads of overlapping and softly
blurring concepts, there is virtually no way to unravel what is really going on in a
human understander's mind. To try to model all that at this early stage of research on
natural and artificial minds would be as sillily ambitious as trying to master the most
rich and idiomatic poetry of a foreign language without ever having bothered to tackle
any prose in it at all. That would be arrogant and intellectually upside-down.
        I believe that it is not merely preferable, but indispensable, to look at the
analogy-making process in a pared-down domain, yet a domain where all the essential
qualities of analogy-making remain. Newton couldn't have discovered his laws of
motion if he had concentrated on trying to understand the laws governing waterfalls
or hurricanes. Instead, he boiled the problem of motion down to the most pristine case
he could imagine planets coasting through a vacuum. This is the typical method of
science: Isolate the crucial phenomena and study them in a pure context; then work
your way upward towards phenomena in which two or more fundamental themes
coexist.
This is what I sought to do in Copycat: To lay bare what I saw' as the central
problems of analogy without any extra clutter of real-world knowledge. Those central
problems, as I see them, are:




Analogies and Roles in Human and Machine Thinking                                      566

• deciding how literally to take references (i.e., deciding which parts of each situation
  already have literal counterparts in the other situation, and which parts need
  "literary" counterparts to be discovered or invented);
• deciding what structures are worth perceiving (i.e., deciding which types of
  abstraction are worth bringing to bear as overarching frameworks to guide
  perception, so as to facilitate mapping pieces of one situation onto pieces in the
  other);
• perceiving roles inside structures (i.e., selecting which aspects of the currently
  preferred organizing frameworks are most relevant and which are less relevant);
• deciding how literally to take roles (i.e., recognizing which roles in each situation
  already have literal counterroles in the other situation, and' which roles need
  "literary" counterroles to be discovered or invented);
• weighing rival ways of viewing a situation against each other and choosing the
  most elegant one (or, if you prefer, the simplest one).

The parallel passages about literary and literal parts and roles may seem a bit obscure,
so let me motivate them briefly.
        An entity in one situation can belong simultaneously to other situations. The
sun is an example. From your point of view and mine, it is just the sun, a unique
object. The sun is what I call a "part" of my world. Its counterpart in your world is the
very same thing, not just something like it. Thus the sun is a part of the world. For
another example, take Groucho Marx. When he died, I didn't have to put myself in my
friends' shoes in order to understand what his death was like for them. His role in their
world and his role in my world were so indistinguishable that to attempt such an
empathetic mapping would have been foolishly extravagant.
        But for a contrast, consider your beloved identical twin sister Glunka. Your
connection to her is certainly far different from mine to her. In fact,' I've never even
met her, whereas you've known her all your life! Glunka is a very big part of your life,
but from my rather distant and external point of view, her main identity is as your
twin-after all, I know nothing else about her. Thus to me, she is the filler of a role in
your life. If I were to learn that Glunka had died, I could hardly be expected to weep
rivers, because she is not a part of my life. But that does not mean that I could not
empathize, because I could project. The obvious mapping would see her counterpart
*s my identical twin sister. However, given that I am a male, no such person exists.
Does that mean I am incapable of empathizing? I would be pretty inhuman if my
ability to empathize were that weak. It is easy to slip from ,identical twin sister" to
"twin sister". Trouble is, I don't have one of those either. What can I do, then, to
empathize? How about slipping along a 4jfferent dimension-to my identical twin
brother? That would be fine if he existed, but he doesn't, poor fellow. (And not
because he died!) So I must loosen up still further, and try mapping your identical
twin sister onto my




Analogies and Roles in Human and Machine Thinking                                    567

twin brother, or just my brother. They don't exist either. Damn! In desperation, I try
slipping to just plain old sister. Aha! This time it works. A sister I have (two, in fact),
and in some loose sense, each of them plays a role in my life analogous to that of your
identical twin sister in your life.
         Of course, the analogy is weaker than it might be if I were twins (like you), but
what can a body do? We make do with the best mapping we can find. The notion "my
sister" is what I call the counterrole in my life to the notion "my identical twin sister"
in your life. Of course, to someone else, the counterrole might be "my Siamese twin
brother", or "my best friend, who I have known all my life", or even, for some people,
"my car". The filler of a counterrole is, of course, the counterpart. In your life there
are many parts and role fillers, as there are in mine. By discovering your life's
counterparts to parts of mine, and your life's counterroles to roles in mine, you can
project and identify.
         How distinct are these concepts of role filler and part? Well, as concepts, they
are quite distinct, but life constantly confronts us with blurry situations where people
(or things) are simultaneously parts and role fillers. Think of your old and dear friend
Millapollie, who is also familiar to me, but only mildly so. If you told me that
Millapollie had died, how would I react? I would have dual approaches to the
situation, one seeing Millapollie as a (very small) part of my life, and the other trying
to find a counterpart in my life to Millapollie's part in your life. Thus I would try to
find the filler of the counterrole in my life to the role that Millapollie plays in your
life. Very probably, I would find myself flitting back and forth between these two
visions of oneand the same person. To effect such a part-role compromise is
sometimes easy, but more often very tricky. Usually we are not even conscious of the
conflicting pressures, but we muddle through anyway.

                                        *   *   *

        Cross-language comparisons may also help to make this idea more vivid. How
eager I was, when learning French, to learn how to talk about baseball. I very much
wanted to know how you say "pitcher", "catcher", "fly ball", "out", and so on. To be
sure, such terms do exist in French, and it's fine to learn them, but it seems to me in
retrospect to have been a misguided obsession for someone whose chief motivation
was sheer fluency. In learning a foreign language, why place a high priority on
learning how to talk about your own culture's idiosyncratic features? Instead, strive to
learn the "corresponding"- aspects of that culture-that is, the things that play
counterroles, rather than the translations of many concepts unique to your culture. In
my case, perhaps the appropriate move would have been to learn all about soccer and
its terminology in French.
        Of course, many terms transcend languages. It is important to know how to
say "moon" in both languages, and it seems reasonable to assume that "the moon of
France" and "the moon of the United States" are really the




Analogies and Roles in Human and Machine Thinking                                      568

same, so that the moon is a shared part rather than a private role-filler. Now in a way
this would seem true of any publicly visible entity, such as Algeria. And yet-are the
Algeria of France and the Algeria of the United States really the same thing? Might
Viet Nam not be "the Algeria of the United States", at least from a French
perspective? Is Algeria an objective part of the world or something that plays a role in
various national perspectives? What about Argentina? Australia? Antarctica? And
such questions apply not only to proper nouns. What about wines, cheeses,
languages?
       Native English speakers are quite easily amused by very crude parodies of
German, such as this sign posted near a computer terminal:

                               Alles Lookenspeepers!

  Das Komputermaschine ist nicht fiir gefingerpoken and mittengrabben. Ist easy
  schnappen der Springewerk, blowenfusen, and pappencorken mit
  Spittzensparken. Ist nicht fur gewerken by das Dummkopfen. Das rubber-'
  necken Sightseeren keepen Hands in das Pockets-relaxen and watchen das
  Blinkenlights.


Our amusement is based on the peculiar way in which our language is rooted in the
Germanic family. We tend to find many aspects of German gawky, comical, and old-
fashioned. Obviously, German speakers will not easily see how their language has this
quality to us. They will certainly not get a sense for our feelings of their language's
gawkiness if they tack Germanic endings onto German words, use lots of "sch"
sounds, and make long compound words-but neither will they do so by tacking on
English endings, suppressing "sch" sounds, and breaking up compounds,, because the
historical and social connection between the two languages is not symmetric, and the
effect, even if humorous, would not be analogous. But what, then, would be the
analogue for native German speakers? What is "the German of German"?
        Note that I seem to be implying that there is one best answer. Actually, I doubt
there is. The connections between English and German are many and variegated. In
some ways, English certainly is to German what German is to English. (This harks
back to some problems of translating self referential sentences, dealt with in Chapter 1
and its Post Scriptum.) In other ways, the assumption of symmetry is completely
wrong. What would a German (or French, or whatever) parody of English (or Dutch,
or whatever) look like?
        How does English sound to a native Mandarin speaker studying it? I bt I could
ever know. Yet I am sure there is some fairly uniform reaction English across the
millions of Chinese people who have heard it. What would it be like to hear my own
native language through ears that could not fathom it, or could penetrate it only
superficially? Such an experience is I vied to me-and yet, is it not exactly the same as
my experience when I 'listen to Mandarin? Yes and no.




Analogies and Roles in Human and Machine Thinking                                   569

What entities in a given situation play the role of fixed stars, of mutually shared
global points of reference? These are, in my terminology, parts. What entities are seen
entirely in terms of their role relative to the perceiver, entirely as local occupiers of
standard "slots"? These are role-fillers. What entities float midway between total
globality and total locality, somewhere between being pure parts and pure role fillers?
The answer is, of course, that nearly all entities are free-floating in this way, which is
why the problems of analogy and translation are so deep and so deeply implicated in
the mystery of mind and consciousness. The linguist George Steiner has provocatively
explored these issues in his book After Babel.
        Since a satisfactory discussion of the nature of analogy would take an entire
book, I will not attempt to lay out the philosophy that the Copycat project is based on.
The column gives you some good ideas (provided you can make that giant conceptual
leap from numbers to letters). But it seems worthwhile presenting at least a few
canonical examples from the Copycat project, since I feel they capture in a nutshell all
that we are trying to do.

                                         *   *   *

         It should go without saying to readers of the column that Copycat deals with
an alphabet of stripped-down letters. In particular, all a letter "knows" about itself is
its predecessor and its successor (if it has such; a and z of course are special in that
regard, each one lacking one). Letters do not know what they look like or sound like,
or whether they are vowels or consonants. Since the Platonic alphabet has a starting
point and a finishing point, unlike the integers, there is a kind of symmetry to it. There
are two distinguished elements, namely, the endpoints a and z. These elements have
identities on their own; they are somewhat like royalty. All other letters derive their
identities, directly or indirectly, from these distinguished letters. Obviously b and y
are like royal viziers, and c and x like vice-viziers.
         I visualize the graph representing letters' "importances" as the arc of a
suspension bridge, suspended at both ends from a and z and descending very steeply
to a minimum at the center of the alphabet (see Figure 24-3). Thus in theory, the very
least distinguished letters are m and n. However, practically speaking, all the letters in
that general vicinity are pretty much equally nondescript. After all, if being
nondescript were a salient property, then we would be caught in a paradox: m and n
would be highly salient by virtue of being maximally undistinguished! But m and n do
not know they are of minimal salience, and hence the paradox is obviated. In fact, any
letters further in from the tips than c and x are pretty bland, and even those two aren't
very exciting.
         In the vast midwestern prairies of the Platonic alphabet, one step this way or
that makes little difference. Poor q hardly knows what role it plays in society, since its
only connections are to p and r, letters of equally little distinction. It's just as in human
communities: most people are recognized in their own neighborhood, but as soon as
they leave, they become




Analogies and Roles in Human and Machine Thinking                                       570

FIGURE 24-3. A graphical representation of the effective saliencies of the letters in
the Platonic alphabet of the Copycat world. The "San Francisco " and "New York " of
the alphabet-a and z-are of course glamorous and salient. Nearby letters get some
reflected glory, but as you leave the two "coasts , you fade into the drab middle
regions, where no letter has much to draw attention to it.

anonymous faces. The only thing you know about random strangers is precisely that:
that they are strangers. Copycat letters are like people in that way.
        For example, q recognizes that it doesn't know k or x, but they are so
unfamiliar that it makes no distinction between their degrees of unfamiliarity. Only
when you come quite close to q does it begin to act as if it recognizes you. But even
then, whatever notice q might take of, say, s would not be direct; it would have to be
mediated by r, which has a direct acquaintance with both letters. Generally speaking,
connections decrease quickly as the number of intermediaries goes up, so that
Platonic q loses virtually all "acquaintance" with letters much further from it than s or
o. Still, there is a sort of exponentially decaying "halo" surrounding Platonic q, a
residue of its interactions with its immediate neighbors (and their -interactions with
their neighbors, etc.), giving it a tapering-off set of "fringe acquaintances" (see Figure
24-4). The same phenomenon applies to all the Platonic letters, of course.
        This "renormalization effect" (so called after the analogous effect in particle
physics) is quite well captured by the following candid remarks made to the author by
Platonic q, when queried about various letters:

"Mercy! I certainly don't recall hearing the name m before."
"Now then ... I believe I've seen n somewhere around."
"Oh yes. I know o, though not terribly well."
"Positively. I'm old friends with p."
"Quit kidding! That's me!"
"Really a fine and true friend, is r."
"Sure, I know s, though not frightfully well."
"That's possible ... Probably I've seen t somewhere around."
"Uhh ... No, I definitely don't recall hearing the name u before.",




Analogies and Roles in Human and Machine Thinking                                     571

FIGURE 24-4. The "renormalization effect " by which any letter (q, here) acquires a
large set of virtual acquaintances despite having only two true acquaintances-its
predecessor and successor. In (a), "bare" q dreams of its two neighbors, bare p and
bare r. But those letters can in turn dream of their neighbors, and so on. This
recursive dreaming-pattern is shown in (b). The upshot of it all is shown in (c): a q's-
eye view of the alphabet. This shows how q has an effective connection to every other
letter of the alphabet, although the strength decays rapidly with distance, since it has
to be mediated by all the intervening letters, and each extra link weakens the chain by
some constant factor. This subjective vista, reminiscent of a sensory homunculus in an
animal's brain, translates into a probabilistic statement for the running Copycat
program: the further two letters are from each other, the less likely is any relationship
they bear to each other to be noticed. Thus close letters (within two or so, as a rule of
thumb) are pretty likely to be thought of in terms of their roles relative to each other,
but further-apart letters are likely to be taken simply at face value, without any
attempt to draw connections between them. This has vast repercussions on how
analogies are made.




Analogies and Roles in Human and Machine Thinking                                    572

        Copycat's alphabetic world includes abstractions that lurk behind the scenes,
and it is they that allow viewers to see coherence in groups of letters. The most basic
of those concepts are two group-types based on the two simple relations that exist: (1)
sameness and (2) successorship (or predecessorship, its mirror image). A group of
neighboring elements linked by sameness is called a copy -group, or C-group. Here
are a few C-groups:

               aaa uuuuu cc wowowowowo

Notice that the concept includes the case where a structure (in this case, wo) is
repeated. Degenerate cases include C-groups of length 1, such as c (or even wo), and
worse yet, C-groups of length 0! Although it sounds perverse, sometimes a group
consisting of zero copies of e is quite different from one consisting of zero copies off.
But this is a fine point, and only for advanced copycats.
        The other group-type is that of S-group (and its mirror twin, P group ). An S-
group is simply a group of neighboring successors, as shown below: .

               abc uvwxy cd pqrs

And if you flip these over so that they run backwards, then they are examples of P-
groups:

       cba yxwvu dc srqp

Needless to say, degenerate S-groups and P-groups of length 1 and 0 exist as well, but
we need not worry over them. Our old friend "1 2 3 4 5 5 4 3 2 1" consists of course
of a numerical S-group and P-group, back to back.
Beyond these most-basic abstract constructs, there are more shadowy entities that
move in the wings and exert intangible yet profound forces on perception of
structures. These go by such names as symmetry, uniformity, good substructures,
boundary strength, and so on. They are the kinds of forces that push you toward
perceiving abcpqr as two groups of three, and aabbcc as three groups of two; they
push you toward breaking aakkkkggee into five equal-length C-groups, and toward
breaking abcdefpqr into three equallength S-groups; they lend support to seeing all
three of the structures abcdcba, abcdabc, and abcdwxyz as symmetric, although in
three very different senses; and they make you feel quite uncomfortable with a
structure such as aaabbbqcc. I will not try to spell out these elusive forces here, firstly
cause they are many, secondly because they are very abstract, and finally cause
people tend to grasp intuitively exactly what sorts of pressures are created by them
anyway.
This concludes the presentation of the Copycat domain, nearly isomorphic to the
Seek-Whence domain (except that in Seek-Whence, there is no analogue to z), and so
without further ado, we may proceed to the analogies themselves. The basic set of
analogies is actually a bunch of




Analogies and Roles in Human and Machine Thinking                                      573

variations on a theme (just as in the column). That theme is the following "event":

               abc changes into abd.

Note that it is up to you to decide what happened to abc ; if you think that the c in it
became a d, that is just fine, but it is your perception and not inherent in the change.
Someone else might interpret it as a replacement of the entire "object" abc, lock,
stock, and barrel, by abd.
        Now, there are a number of possible counterpart situations that I could present
and ask you to do "the same thing" in. It's hard to decide which one to try first, but I'll
just plunge in:

               What does pqrs change to?

Pretty easy, eh? I suppose you said pqrt. So does nearly everybody. That is because it
is-in some rather elusive and wonderful sense-the right answer. Yet there are
numerous other candidates that you might have considered. In fact, because it is
almost impossible to appreciate the intricacy and fascination of the Copycat world
unless you "live" there, I strongly urge you to pause here and think: What else could
pqrs go to?

                                        *   *   *

         One possibility is pgrd ; need I explain why? Then there is pqrs, produced
from the input pqrs just as abd was produced from abc: by substituting a d for every c.
Did we, or did we not, do "the same thing" here? Should you touch your nose, or mine
? For that matter, why isn't abd or pqds the answer to this problem? Such questions of
rigidity versus fluidity recur throughout analogy, and seem to resist formalization.
         Speaking of rigidity versus fluidity, when I gave a lecture on analogies in the
Physics Department at the California Institute of Technology several years ago, one
Richard Feynman sat in the front row and bantered with me all the way through the
lecture. I considered him a "benevolent heckler", in the sense that he would reliably
answer each question "What is to X as 4 is to A?" with the same answer, "4! ", and
insist that it was a good answer, probably the best. It seemed to me that Feynman not
only was acting the part of the "village idiot", but even was relishing it. It was hard to
tell how much he was playing devil's advocate and how much he was sincere. In any
case, I will never forget the occasion, since his arguing with me stimulated me no end,
and at least from my point of view, it wound up being one of the best lectures I have
ever given.
         On a subsequent occasion, when giving another lecture on analogy, I
remarked, quite innocently, "Last time I gave this lecture, Richard Feynman sat right
there", and I pointed at a seat just to the left of center in the front row. No sooner had
I said this than I realized the marvelous analogical




Analogies and Roles in Human and Machine Thinking                                      574

transfer I had done so totally subconsciously. After all, at Caltech it had been a
gigantic auditorium (this was a small classroom); the seats were in tiers (here they
were just in ordinary rows); each row was very wide (here they were quite narrow);
and I had been in California (now I was in Ohio). Yet pointing at one seat and saying
"Feynman was sitting there" seemed to make eminent sense, in that context. (Isn't it
equally sensible as claiming that the light bulb was invented in Dearborn, Michigan,
merely because the New Jersey house that Edison did his work in has been
transported to a historical' park there'?) Furthermore, it now occurs to me that "just to
the left of center" is itself the key concept of many of the analogies in the column.
         The more you look at the question of how to do "the same thing" to pqrs, the
more possibilities you see. For instance, many people seem to like pqst, in which the
first two letters are left alone and subsequent letters are replaced by their successors.
Occasionally, people have suggesed pqtu, a rather ingenious notion based on seeing rs
as a single unit whose successor is the unit tu. Somebody pointed out that qrst is a
possibility, based on the idea of changing all letters but a and b to their successors.
And one time someone sug jested dddd, whose justification resides in the even more
village-idiotic notion of changing all letters but a and b to d !
         Some answers appear almost sick. Consider abce. The defense of this answer
is that you take as many letters at the beginning of the alphabet as are in the target,
then convert the final one to its successor. You can even come up with justifications
for such queer answers as abt and, believe it or not, pqre.

                                       *   *   *

        When I call some answers sick, and others healthy by implication, there is
something behind the metaphor. After all, there is a very serious question that always
arises about analogy, but particularly strongly in such an abstract domain as this, and
that is how one can ever speak with confidence about the rightness or wrongness of
something that is so clearly subjective. The way I have come to view this is in terms
of the survival value that an analogy-making capacity confers on its possessors. After
all, our brains got to be the way they are only by helping our forebears to survive
better than heir rivals in this unforgiving world. And analogy-making is at, or close to,
the pinnacle of our mental abilities.
        It seems to me that people do not generally recognize how deeply mplicated
the analogical capacity is in decisions that affect the course of heir lives. On a global
level, it is evident, once pointed out. Is the embroilment of the United States in
Lebanon "another Viet Nam"? How about in El Salvador? How does the American
invasion of Grenada map onto the Soviet invasion of Afghanistan, or the British
invasion of the Falklands? Is the Soviet Union more like an irrational paranoid person,
or someone rational who has been badly bullied recently? Does the current arms race
have valid precedents in history to which it can be compared?




Analogies and Roles in Human and Machine Thinking                                    575

        On a more local scale, our system of law very obviously sanctifies analogy as
the ultimate justification for making a reasonable and even a wise decision. The term
"precedent" is just a legalistic way of saying "well-founded analogue". Two cases that
at a surface level have nothing to do with each other (a bank robbery, say, and a
kidnapping) may be mapped onto each other in exquisite detail at a more abstract
level, with the napped kid being the loot, for instance. Lawyers attempt to sway the
jury by bringing in new ways of looking at the situation that discredit their opponents'
analogies, as well as by making and buttressing their own rival analogies. (Peter
Suber has written a nice article connecting Copycat and Seek-Whence analogies with
legal reasoning. It is called "Analogy Exercises for Teaching Legal Reasoning" and
can be gotten from him at the Philosophy Department of Earlham College in
Richmond, Indiana 47374.)
        In our private lives, most of our important judgments are made by conscious
or unconscious analogy. Should I fight this bureaucracy or accept some annoying
inconvenience? Should I buy this computer or wait for a better one to come along at
the same price? Should we have children now or wait a few years? Should I retire or
continue working beyond retirement age? Questions concerning what to buy, what to
think of someone, whom to marry, whether to move to a new city, how to talk to
someone who has suffered a calamity, and on and on-all of them are influenced in a
myriad ways by prior experiences of the same general sort. And remember that even
in cases where there is not any obvious analogy guiding the judgment, all the
categorization of the situation is being made by a mind exposed to many thousands of
words, and the purpose of words is to label situation types and thus implicitly to make
use of stored analogical mappings.
        As was discussed in the column, the boundary line between making creative
analogies and recognizing pre-existent categories is very blurry. It is signaled when
we feel a desire to pluralize a proper noun ("your Einsteins and your Mozarts") or to
prefix a proper noun by a definite article ("the Podunk of Albania"). Most common
words hide an enormous degree of analogical abstraction. For instance, the
abstractions "female" and "male" are not nearly as simple as most people think,
especially when you consider how they extend to plants. What makes Middle Eastern
pita, Indian purl, French baguettes, and American Wonder all be examples of the
concept "bread"? When you migrate from nouns toward verbs and prepositions, the
difficulties escalate. What do all "x-is-on-y" situations have in common?
        All of this points out how analogies determine the course of our lives in the
present. But I would go much further than that. In pre-civilized days, when people (or
proto-people) lived in caves and hunted bison, analogy played no less important a
role. Samenesses that we have absorbed into our perception as being obvious were
great insights back then. For instance, the idea that one could chart out a plan for
trapping a wild beast by drawing a map on the ground must have been a fabulous
advance. All that is involved, in some sense, is a change in scale-one of the most
obvious of




Analogies and Roles in Human and Machine Thinking                                   576

analogical transforms, yet when it was first invented, it must have been revolutionary.
On the other hand, proto-humans who tried burying meat underground in an attempt
to imitate squirrels' underground hoarding of acorns might thereby seriously damage
their chances for survival. Some analogies help, others hinder.
        Our current mechanisms for analogy-making must certainly have emerged as a
consequence of natural selection. Good mechanisms were selected for, bad ones were
selected against, way back when, in the old times when you and I were but monkeys
and rodents scampering about in tree branches ('member?). The point, then, is that far
more than being just a matter of taste, variations in analogy-making skill can spell the
difference between life and death. That's why "right answer" means something even
for analogies; it's why analogies are only to some degree a matter of taste.

                                       *   *   *

        This finally gets us back to the rivalry among answers in the pqrs case. The
domain does admittedly appear abstract and of course it is totally decoupled from the
cruel world. People who prefer dddd are not suddenly going to get swallowed by a
tiger or topple off a cliff. But people who genuinely believe in dddd's superiority over
pgrt will still have a rough time in life, because they lack the means to size up a
situation and catch its essence in their mind's mesh, letting the trivial pass through.
Something about their analogy making mechanisms is defective.
        To be sure, there is room for argument about answers in this mini-world, just
as there is in a courtroom. But just as a lawyer who suggested that killing a human
being is analogous to breaking a window because both are nasty or because both can
be done with a brick would lose the case in a snap, so .anyone who prefers dddd or abt
to pqrt can be safely assumed to be totally off base. There are absurd answers, there
are good answers, and there are in-between ones, just as there are degrees of edibility
of food. Some foods lead to no survival, some to bare survival, and others to
comfortable survival; the same is true of analogies.
        One can liken the various levels of quality to the concentric circles
surrounding a bull's-eye on a target (see Figure 24-5). In the middle are the totally
edible foods (or insightful analogies); further out are semi-edible substances; such as
grass, hay, or ants (weak analogies) and worse, leather or wood (dubious analogies);
and then way out are the completely inedible hings, such as nails, shards of glass, or
Anglican cathedrals (these correspond to analogies that lead to disaster, such as
forming a higher-level category that lumps tigers together with zebras simply because
both have stripes). In the very center, to be sure, one can argue about taste and it is
indubitably good for the human race that there are people who see analogies
differently in that region, but you cannot get too far-fetched.




Analogies and Roles in Human and Machine Thinking                                   577

FIGURE 24-5. Targets representing the survival values of various actions taken by
(a) a mind seeking answers to the analogy "If abc changes to abd, what does pqrs
change to?", and (b) a mind choosing things for its body to ingest. To be sure, there is
room for debate in the bull's-eye region (hard to say whether having a fried egg
sandwich or a plate of spaghetti is better for you), but as you edge further out, it gets
less and less debatable. Fried dust just doesn't match up to pea soup! Similarly for
analogies. Among the various answers to any analogy problem, some will be
decidedly weaker than others, even if you find that no one answer emerges as the
clear victor.

There is a radius beyond which analogies will be very likely to bring bad
consequences to their proposers, at least if they are acted upon.
        It is for this kind of reason that I unbudgeably believe that there are better and
worse answers to analogies, whether in life or in the Copycat domain. Elegance is
more than just a frill in life; it is one of the driving criteria behind survival. Elegance
is just another way of talking about getting at the essence of situations. If you don't
trust the word "elegance" in this context, then you may substitute "compactness",
"efficiency", or "generality"-in short, survivability. Insight into the mechanisms
behind this sense of elegance is the goal of the Copycat project. And personally, I
would not shy away from equating elegance with wit. I would venture that one cannot
be a successful Copycat without a sense of humor.

                                        *   *   *

        Having seen numerous wild and woolly (and sometimes witty) answers to the
question "What happens to pqrs ?", you might now wish to see some alternate targets.
They are extremely important, because they bring out a variety of new ways of
perceiving the original "abc goes to abd" change. Here are a few fascinating
challenges that I urge you, once again, to actually devote some time to considering.
All of them are based upon our old stock event, abc goes to abd.




Analogies and Roles in Human and Machine Thinking                                      578

       1. What does cab go' to?
       2. What does cba go to?
       3. What does pct go to?
       4. What does pxgxrx go to?
       5. What does aabbcc go to?
       6. What does aaabbbcck go to?
       7. What does srgp go to?
       8. What does spsqsrss go to?
       9. What does abcdeabcdabc go to?
       10. What does bcdacdabd go to?
       11. What does ace go to?
       12. What does xyz go to?

        Each one of these questions shakes some fundamental assumptions about how
you should perceive the original change. Does it necessarily affect just one letter?
Need the object affected be at the righthand extremity? Does the changed piece
always get replaced by its successor? In short, what should be taken literally, and
what slipperily ? Although I would love to do so, I am not going to discuss all twelve
examples here, for that would take a good long time. Each one merits at least a page
on its own. (A set of "answers" is given at the end of this P.S.) I will discuss just one
of them, number 12, in some depth. It has real beauty and raises all the central issues,
so I hope you will give it some thought before reading on.
        May I go on now? All right, Many people are inclined to say that xyz should
go to xya. But who said the alphabet was circular? To make that leap, you almost
need to have had prior experience with circularity in some form, which we all have.
For instance: The hours of a clock form a closed cycle, as do the days of the week, the
months of the year, the cards in a suit, the digits 0-9, and so on. But not all linear
orders are cyclic. The bottom rung on a ladder is not above the top rung! The Empire
State Building's top floor is not the same as its basement! It is a premise of the
Copycat world that z has no successor. Sure, a machine could posit that a is the
successor to z, but to do so would be an act of far greater creativity than it would be
for you, because you have all these prior experiences with wuWaround structures. The
Copycat program does not. Therefore, let us consider xya as admirable but simply too
daring, and look for something more humble yet no less apt. What else remains?
Again, I urge you to think about this before reading on. This is the crucial point where
there simply is no substitute for your own experimentation.

                                       *   *   *

       Okay. You've thought it over. You've got an answer, perhaps even a few,
ranked more or less according to the pleasure they give you. Great! Some people
suggest xy, pure and simple. Since z has no successor, they just let




Analogies and Roles in Human and Machine Thinking                                    579

the third term drop away, as if it had fallen off the edge of the world. Some think,
"Why not just leave z alone, producing xyz ?" Some say, "Since the rightmost letter
has no successor, why not slip over to the next-to-rightmost one and take its
successor, thus producing xzz ?"
         Those are all right, but far more insightful answers are possible. To find them,
one can let the unexpected "snag" (namely, the problem of trying to take the successor
of a successor-less object) trigger a search for something crucial possibly overlooked
earlier. What people tend to see at this stage is the potential correspondence of a and
z-the two extreme letters of the alphabet, our two twin "monarchs". If z is the a of
xyz, then what is the c ? Quite obviously, it is x. Now the question arises: "What to do
to x ?" Should we take its successor, thereby producing yyz ? To my mind, there is
something almost repulsively rigid about this suggestion. After all, the very fabric out
of which abc was constructed has now been reversed in our new way of looking at
xyz. Leftward motion has seized the role of rightward motion, and with it,
predecessorship has taken on the role that successorship played in abc. Therefore
elegance, in the form of a drive toward abstract symmetry, very strongly pushes for
the answer wyz. Now that's a beautiful answer, in my estimation.
         There is one other answer that I have encountered fairly often, and that I
admire and decry at one and the same time. That is wxz. To be sure, it has the same
inner structure as does abd: a jump of size one followed by a jump of size two. That
much is good, but there is something peculiar about wxz nonetheless. To illustrate my
ambivalence toward this answer, I will relate a micro-allegory.
         Arphabelle Snerxis built a lovely house, ultra-modern in every respect. Then
one day she capriciously removed her snazzy, sleek, new doorknob and replaced it by
a most conspicuous creaky, rusty, old doorknob. Now Zulips Twankler, a great
admirer of Arphabelle Snerxis, happened to have built a lovely house, old-fashioned
in every respect; and when he saw Arphabelle's action, he determined to do "the same
thing" to his house. And how did he do that? You might guess that Zulips Twankler
removed his creaky, rusty, old doorknob and replaced it by a most conspicuous
snazzy, sleek, new one! But no, Zulips left his creaky, rusty, old doorknob intact and
instead he tore down the rest of his lovely house, old-fashioned in every respect. Then
Zulips built another quite different but also lovely house, ultra-modern in every
respect-except for the creaky, rusty, old doorknob. And that's how Zulips Twankler
did "the same thing" to his house as Arphabelle Snerxis did to her house (except that
when he'd finished, it wasn't his house any more-it was a different house altogether).
         In this "analogory", Zulips let the identity of his house slip in a manner
determined by its doorknob's properties, while for most people it would seem more
natural to have the slippabilities reversed. There is a parallel in the fight between
answers wxz and wyz. The former sees preservation of the literal intervals (1, then 2)
as necessary at all costs, and it allows the bulk




Analogies and Roles in Human and Machine Thinking                                    580

FIGURE 24-6. A visual comparison of two good answers to the question "If abc
changes to abd, what does xyz change to?" In (a), wyz is depicted. The total symmetry
of this answer is its virtue, In (b), wxz is depicted. Its virtue is that its spacing imitates
the spacing of abd literally, even slavishly. The question is, which of these answers
creates a better overall analogy? Is it wise or is it a cop-out to intone "De gustibus
non est disputandum" and leave it at that?

of the original entity xyz to be shifted in order to achieve those aims. The latter sees
only one small role (analogous to the doorknob) as singled out for modification, and
keeps the bulk intact. The contrast is vividly portrayed in Figure 24-6. Another way to
see the contrast is this: wyz is based on a better imitation of the change from abc to
abd, while wxz is a better imitation of the end product, abd. Which do you find the
more satisfying or elegant action? For me it is wyz, hands down. Still, I find a strange
charm in the Zulipian answer. This is one of those cases where two different answers
(three, if you count xya, that deeply creative answer that comes so easily to us old
circularity hands) can lie inside the innermost "delicious and nutritious" circle, within
which de gustibus non disputandum est.
         There are a host of other "possible" answers to this problem, but many of them
lie in the risky outer circles and are dangerous to their proposers. (Or, to speak more
precisely, the mental mechanisms that would allow them to be considered seriously
are dangerous to their owners, since those same mechanisms could produce very
untrustworthy suggestions for courses of action in cases that are not decoupled from
consequences.) Some of these answers are so far-fetched that they are actually quite
humorous. In fact, some of them are so outlandish that I would claim no conceivable
survivor of evolution's harsh pruning would ever come up with them, unless
deliberately for humorous purposes. Let us look.
         First of all, not all that funny but very much in the Feynman "village idiot"




Analogies and Roles in Human and Machine Thinking                                        581

spirit, is xyd. Just' plain old dullsville. What can you say to such an answer?
Certainly, it has more merit than pqrd did earlier, for here, at least, there is an excuse
for such rigidity: z is a trouble spot, whereas s was not. Another answer, definitely
more pitiful, is dyz. This one just goes "Thud!" very loudly in my mind's ear. Anyone
who would seriously suggest this answer `has seen the light but then dropped the ball,
mixed-metaphorically speaking. That is, they go quite a long ways in mapping a onto
z, c onto x, but then they seemingly totally forget this level of sophistication and
revert to an infantile literality: "Replace whatever plays the c role by d. " It is so inept
that it is funny. In fact, I would say that in the answer dyz to this analogy problem
there is the germ of a rich theory of humor, based on the idea of level-mixing slips of
this sort. Now if only I could say just exactly what I mean by "slips of this sort", I'd be
in business ...
         Equally scramble-brained is dba. Its hypothetical serious proposer has clearly
seen that, in some abstract sense, xyz is a "mirror image" of abc. An attempt is
therefore made to take a mirror image of abd, but somewhere in the shuffle, the type'
of mirror image involved got forgotten, and for it was -substituted the crudest possible
notion of "reversal". The result is another heavy-handed thud. By the way, nobody has
ever seriously suggested to me either the clunky dyz or dba, or even yyz, interestingly
enough. Oh well, I guess such people's ancestors must have all gotten gobbled up by
dreaded human-eating zebras.

                                         *   *   *

         These analogies, as must be abundantly clear by now, are themselves
analogous to real-world analogies. Or better yet, they are allegories or fables for
analogy-land children. They capture the essence of the dilemmas that analogy-makers
face over and over again. Pressures for literality vie mightily with pressures for
"literarity"-that is, for high-flown abstractions that disdain rigidity of reference. In an
analogy where real insight is needed, often some initial forays are made that reveal
some literal-minded ways of seeing the situation, but they ring false or overly crude,
and so one does not stop there. Instead, one continues searching, guided by a number
of small cues, and at some unpredictable moment, something simply snaps, and a host
of things fall into place in a new conceptual schematization of what is going on. What
was once important becomes suddenly trivial, and a new essence emerges, an
organizing concept or set of concepts that seem far superior.
         But beware! You should not take any of this to imply that being literal-minded
is to be avoided at all costs. If jumping to rarefied levels of abstraction were always
preferable, then we would wind up never making any distinctions between situations.
Every situation's optimal description would be: "Something happens." It would be
better to say "woman" than "Mary", "person" than "woman", "animate object" than
"person",




Analogies and Roles in Human and Machine Thinking                                      582

"physical entity" than "animate object", and ultimately, just "thing" would emerge
triumphant. (Actually, even that choice would be nonexistent, for those different
words would not exist, as they would merely make petty distinctions that vanish at
enlightened higher levels.) Rigidity comes in many forms, and a rigid drive toward
abstraction is no less stupid than a rigid refusal to abstract. The clearest and cleanest
statement of the problem that analogy poses is that there are always fights between
forces pushing for literality (in its many forms) and forces pushing for abstraction (in
its many forms). How, in specific circumstances, those forces compete and interact
and in the end come up with some sort of optimal compromise is the problem. And if
you go away from this chapter with one thought, please let it be this: There is no fixed
mathematical recipe for reconciling all the different forces pushing and pulling you in
analogies.

                                       *   *   *

        In the Copycat world, we have remarkably fine control over how pressures
interact, and over the strengths of various rival answers. For instance, we can "tune" a
given analogy by varying knobs on it, until gradually we home in on the most refined
possible version of it. We can also take two possible answers and make each one seem
preferable by very subtly "tweaking" the analogy itself. It is a highly pleasing esthetic
exercise to seek the perfect balance point where an analogy is teetering on the brink
and can tip either way, so that about half the people to whom we give it see it one way
and half see it the other way.
        A lovely example of this idea is provided by a very simple analogy with a
knob that twiddles the part-vs. -role proportion perceived by a typical person. This
example may help clarify that distinction a little, as well. Consider the following three
variations on a familiar theme:

        1. If abc goes to abd, what does pqrs go to?
        2. If abc goes to abe, what does pqrs go to?
        3. If abc goes to abf, what does pqrs go to?
        Line 1 is familiar by now. There, nearly everyone instantly proposes pgrt ;
very few people think of, let alone prefer, pqrd. The reason is that the c-d conversion
seems so obviously to be a "leap" from one letter to its successor that we ourselves
leap to that conclusion. But consider line 3. Here we are confronted by the much
larger c -f gap, too large for us (or the Copycat program) to have any intuition for. It
seems so arbitrary that instead of seeking any "whence" behind it, wee just accept it at
face value, saying to ourselves, "Okay, replace the rightmost letter by f, eh?" And
indeed, most people are happy with the answer pgrf.
        To some, this answer may seem overly Feynman-like, but I must reiterate that
in the Copycat world, there is no simple connection between distant




Analogies and Roles in Human and Machine Thinking                                    583

letters. You can't just "subtract" c from f the way you can subtract 3 from 6.
Subtraction is an unknown concept, just as in Seek-Whence. The connection between
c and f, to the extent there is one, is a conceptual one rather than a mathematical one: f
is the successor of the successor of the successor of c, and that is just too topheavy a
notion to have much charisma.
        All right, if line l's c-d leap has charisma and line 3's c -f leap lacks it, what
about the intermediate case of line 2? Here we are poised between two analogies that
push us in opposite directions. If we were to follow line l's example, we would see the
c-e leap intensionally, a fancy way of saying that we would see e as a role-filler (as
the successor of the successor of c-a bit gawky but still plausible). But if we were to
follow line 3's example, we would see the same leap extensionally, meaning that we
would see e as a mere part of the event, filling no role other than being itself (which it
could hardly help 'doing). This is not gawky, but it is so literal that it provides W
insight into why the given change occurred. So which do you prefer-the intensional
view of e as gawky role-filler, or the extensional view of e as arbitrary part? The
former leads you to answer pqru, the latter to answer pqre.
        Although line 2 may not be your personal balance point, there is surely a point
at which you will switch over from one view to the other. It would seem highly
compulsive if, given the question

       abc goes to abv; what does pqrs go to?

somebody insisted that the v must somehow "come from" the c, and tried to force a
vision of some connection when there really is none. Furthermore, even violating the
spirit of Copycat and seeing v as the order- 19 successor of c is not of much help, for
what is the order-19 successor of s ? It gets us right back to successor-of-z problems,
very messy territory.
        Seeking the balance point of analogies is an esthetic exercise closely related to
the esthetically pleasing activity of doing ambigrams, where shapes must be
concocted that are poised exactly at the midpoint between two interpretations (see
Figures 13-6 and 13-7). But seeking the balance point is far more than just esthetic
play; it probes the very core of how people perceive abstractions, and it does so
without their even knowing it. It is a crucial aspect of Copycat research.

                                            *   *   *

        A few more choice problems are given below for would-be copycats. I do not
have the space-time to discuss them all here; I propose them simply because each one
has a chance of affording you a small but thrilling moment of blinding insight, if you
look at it in just the right way. Our view of the best answers is given in the Post Post
Scriptum.




Analogies and Roles in Human and Machine Thinking                                     584

       1. If aqc goes to abc, what does pqc go to?
       2a. If efgh goes to fghi, what does mvr go to?
       2b. If efgh goes to fghi, what does uuuuu go to?
       3. If beq goes to bqe, what does abcdefpqr go to?
       4. If xyzabc goes to xyzgabc what does abcxyz go to?
       5. If aaggkkkk goes to zaazqqzkkzkkz, what does abcdefstu go to?
       6. If eeeffghhiii goes to eeeefffgghhhiiii, what does eefhii go to?
       7a. If eqe goes to qeq, what does abcdcba go to?
       7b. If eqe goes to qeq, what does aaabccc go to?
       7c. If eqe goes to qeq, what does eqg go to?

        It must be emphasized that the selection of Copycat problems presented here is
but a tiny fraction of all those that I, together with David Rogers, a postdoctoral
fellow working on the Copycat project, have come up with. This selection is biased
toward analogies that have spice and tang, as opposed to bland ones such as "If bbb
goes to bbbb, what does eee go to?" There are of course innumerable ones of this
boring sort, ones that have obvious answers and that present no serious challenges to
adult humans.
        Now, we do not in the least disdain such analogies in the project;* Indeed, it is
a tremendous challenge to make a program that could handle these seemingly easy
cases reliably. They are amazingly deceptive in their subtlety. But it is not of much
interest to people to go down a long list of (not actually but seemingly) trivial
analogies, so that explains the censorship in my choices for you.
        Still, it must be admitted, analogies that seem to require a deep perceptual shift
after an initially unsatisfactory first stab are the ones that beguile us, for they seem to
promise insight into that mystery of mysteries: insight. I must admit to the belief, or at
least the strong intuition, that all the depth of scientific discovery, even the
profoundest discovery, is wrapped up in the mechanisms for solving these simple
problems in which conflicting pressures push around one's percepts and concepts,
letting things bounce against each other until, all at once, something falls into place
and then, presto! A sense of certainty crystallizes, so powerful that you know you
have found the right way to look at things. I firmly believe, in short, that "mini-
breakthroughs" and "maxi-breakthroughs" have precisely the same texture. That's the
faith underlying Copycat.
        It may seem arrogant or blasphemous to compare the trivial alphabetic insights
of a copycat with the genius of an Einstein discovering special relativity, yet I do not
think the comparison is all that silly. What characterized Einstein's unique view of
space and time was that he had decided that certain things were more unslippable than
others: in particular, that the speed of light was unslippable but the notion of absolute
simultaneity of events separated in space was slippable. To be perhaps more accurate,
Einstein didn't decide that simultaneity was slippable, but was forced into that
conclusion, since his stronger intuitive belief in the invariance of




Analogies and Roles in Human and Machine Thinking                                      585

the speed of light simply compelled him to accept its consequences, strange and
counterintuitive though they might be. (Note that counterintuitive consequences can
flow from intuitive grounding.) Einstein did not begin with the idea of simultaneity
being nonabsolute, but when he had to confront that possibility, he let it slip. This
fluidity of mind, guided by a certainty about the deepest, most unslippable concepts,
gave rise to the creative insights of special relativity.
        There is an old song whose lyrics go this way:

               When an irresistible force such as you,
               Meets an old immovable object like me,
               You can bet, as sure as you live,
               Something's gotta give, something's gotta give, something's gotta give!

Yes, something's gotta give, but what? A reliable nose for what might slip and what
ought not marks the difference between a great mind and a small one. If the Copycat
research can unearth the basis for judgments exhibiting creative, artistic slippability
even in our tiny domain, we will be ecstatic, for in our opinion, that would put us well
on the road to understanding where full-scale artistic creativity comes from. Now that
may sound arrogant, but firstly, we are not expecting it to happen just, around the
corner, and secondly, it is just an expression of our faith that we have not lost the
essence of the larger problem in boiling it down this far. If Newton saw whirling
planets in falling apples, why can we not see great leaps in small slips?

*   *   *

        One can look to literature as well as science to find cases where finding the
right things to slip yields highly creative solutions to hard problems. One example I
gave in the Post Scriptum to Chapter 1 was the problem of translating into French (or
the foreign language of your choice) the title of the book All the President's Men. A
word-for-word translation would be as dull as dishwater, as flat as old ginger ale
whose carbonated kick has long since evaporated. In order to keep the title alive in
French, you must seek out a line well known to readers of French that carries the
same subliminal imagery. Need it be a line in the canonical translation of "Humpty
Dumpty"? Need it even be a line from Mother Goose? Of course not. The essence of
the situation does not reside in those particulars. So slip, baby, slip! But how?
        The crux of the matter is to find a line alluding to a famous irreversible
downfall. If it is a line from Pascal's Pensees, so be it. If it is from a popular song of
recent years, so be it. You may have to go further afield to find an appropriate line.
There may be no line of the sort in the popular French-speaking consciousness, in
which case more radical solutions must be sought. There is no clean, clear recipe
guaranteed to work. By the way, I do not have any idea if that book has ever been
translated into other




Analogies and Roles in Human and Machine Thinking                                     586

languages, and if so, what solutions its translators found. But this type of problem is
absolutely standard, since these days particularly, book titles of that style, making an
oblique allusion to some well-known phrase, are a dime a dozen.
         I must admit to some twinges of shame for having leapt aboard the title-as-pun
bandwagon, when Daniel Dennett and I chose the title The Mind's I for our anthology.
Good but non-native speakers of English usually are confused by this title, and often
read the final "I" as the roman numeral for "one", which makes absolutely no sense,
yet it is the best they can do, being unfamiliar with the idiom "the mind's eye".
         Just to give a hint of how creative solutions can be found for such titles, I'll
give one possible French translation for our title (even though it is not yet certain
whether the book will ever be translated into French). Jacqueline Henry, one of the
two translators into French of Gödel, Escher, Bach, came up with Vues de l’esprit-
literally, "Views of Spirit", which clearly gets across one main purpose of the book: to
focus on the nature of mind from many angles. But at the same time, it has a more
idiomatic meaning, namely: grandiose dreams such as are dreamt by visionaries (and
lunatics)-in short, visions or possibly even hallucinations. This too has its own kind of
appropriateness, since a basic theme of the book is that much of the mystery
surrounding mind, spirit, and soul is caused by a kind of hallucination: the
hallucination that there is some thing called "I". Therefore, the French double
meaning is elegant and, though it does not replicate in French the exact effect of the
English double entendre, it is effective and thought-provoking. What more could you
ask?
         Incidentally, if I were writing this P. S. in French, I would of course talk about
books in French, not ones in English, whose titles are hard to translate. Thus a proper
translation of this very passage would involve a good deal of literarity. In fact, I have
one particular book title in French in mind: Le corps a ses raisons-a book about health
and physical fitness, which actually came out in English under the feeble title The
Body Has Its Reasons. Can you do better? Hint: You need to be familiar with a
famous saying by Pascal, namely, Le ca ur a ses raisons que la raison ne connait
point.

*   *   *

        In a certain way, translation is the quintessential form of analogy. You are
given a fixed overarching framework-the home language and culture-and within it, a
novel structure has been erected-a book title or sentence, for instance. Your task as
translator is to replicate, as best you can, the overall "feel" of that small structure, but
in a different overarching fixed framework -the target language and culture. This
description obviously recalls the allegory of Arphabelle and Zulips, where the
frameworks are their respective houses, and it applies equally clearly to the "Do this!"
examples.
        My own mental image that best gets at the nature of translation involves




Analogies and Roles in Human and Machine Thinking                                       587

FIGURE 24-7. A metaphor for translation. A stream (symbolizing reality) has two
sets of stepping-stones (symbolizing the basic ingredients of a language, such as
words and stock phrases) in it. The black stones (Burmese, say) are arranged in one
way, and the white stones (say, Welsh) in some other way. A pathway linking up a few
black stones (a thought expressed in Burmese) is to be imitated by a "similar"
pathway joining up white stones (translated into Welsh). One possibility is the
speckled pathway, located at nearly the same part of the stream as the original
pathway but not terribly similar in shape to it (a fairly literal translation), while a
rival candidate (a more literary translation, needless to say) is the pathway located a
distance upstream and resembling the original in some more abstract ways, including
patterns in some of the "overstones " of the main stones (the similar archipelagos in
Burmese and Welsh stones running roughly parallel to the far bank).

picturing each language as a fixed set of stepping-stones in a stream (see Figure 24-7).
Suppose you are translating from Burmese to Welsh. A Burmese utterance is a
pathway from one place to another via the black stones. They seem to be located in
convenient enough places, and you can get pretty much wherever you want to go. But
when it comes to translating what you have said into Welsh, you find that the Welsh
stepping-stones the white ones-are often not quite in the same places as the Burmese
ones; and even in the cases where they are just about in the same places, they are
shaped differently, and so you can't treat them as identical to the Burmese stones you
are familiar with. You must tread with great care, and sometimes you will find that
there are gaps in one language's set of stones that don't exist in the other, so that some
routes are easier to mimic than others. The most literal translations involve sticking as
close as you can to the original route, at the stone-by-stone level. Of course, no
mimicking route is exactly the same as the original.
        It may be, however, that the essence of a particular route lies not in where it
starts or ends in the stream, but in its shape. It may be that in the particular region of
the stream where the original path was traced, the Burmese stones very easily allow
many shapes to be traced out but the Welsh stones




Analogies and Roles in Human and Machine Thinking                                     588

happen to be sparse. At various places upstream or downstream, however, the
converse is true. If you believe that the essence of the idea resides more in its shape
than in its absolute location relative to the stream bed, then you won't mind moving
upstream or downstream a bit, in order to gain that flexibility. Less metaphorically
speaking, this means that sometimes the overt topic of a passage can slip as long as
something more central-style, perhaps, or metaphorical allusion-is preserved.
         A critical idea is the following one: The longer a passage is, the less the
graininess of the underlying medium is going to be noticed. In purely geometric
terms, what I am saying is this. The larger the curves of the pathway are in
comparison to typical inter-stone distances, the less it matters which of the two grids
of stepping-stones you are using. This can be illustrated very elegantly by thinking of
trying to approximate a circle by filling in various squares on a normal (8 X 8)
chessboard. Clearly you would make the circle as big as you can within the confines
of the board, so as to round off the effects of the squareness. And if you could make
the board bigger, you would. On a 100X 100 chessboard, you could draw a very fine
approximation to a circle, and on a 1,000,000X 1,000,000 board, no one would know
the difference. Furthermore, nobody would even be able to tell, in such a case,
whether the underlying board was a square lattice, a hexagonal lattice, or what. But if
you go back down to circles whose size is on a par with that of the lattice grain, then
of course the lattice becomes very visible.
         For this reason, I feel safe in suggesting that translating a novel's title may
sometimes be the most challenging aspect of translating the whole novel. After all, the
overall message of most novels is on such a vastly larger scale than the grain size of
either language involved that small jogs here and there (where the idiosyncratic
placing of the stepping-stones forces you to take an awkward zigzag) can be
compensated for in other places, and in the larger picture such jogs will balance or
cancel each other out. Recall that I said something similar about computer languages
and Al programs-the grain size of the ideas in a big program is far larger than that of
any conceivable computer language.
         But a title is another story. A title is tiny. Its grain size is barely above that of
the stepping-stones themselves. It consists of a pathway just a handful of stones long,
and the challenge is great when it contains subtlety of any sort -which is the case for
most titles, as it is for proverbs, epigrams, and so on. As they say in Italian,
Traduttore, traditore-which, literally as well as literarily translated, means "Translator,
traitor." In this curious case, the English version is a perfect counterexample to its
own claim, but generally speaking, the Italian epigram is right on target, and pithily
expresses the idea that no translation-no analogy-is perfect. Perhaps a better English
translation of it would therefore be: "Transductor, treasoner."

                                          *   *   *




Analogies and Roles in Human and Machine Thinking                                        589

FIGURE 24-8. Three lattices on which chess-like games could be played. In (a), a
square lattice; in (b), a hexagonal lattice; in (c), a triangular lattice.




Analogies and Roles in Human and Machine Thinking                           590

FIGURE 24-9. Various ways to think about the knight's move. In (a), it is built out of
rook-move and bishop-move primitives. In (b), it is portrayed as the closest spot not
immediately accessible to rook or bishop. In (c), we make some first stabs at the
knight's move on a hexagonal lattice. Are some of these possibilities more defensible
than others?

        The idea of approximating a given shape (such as a circle) on a coarse-grained
grid (such as a chessboard) provides a wonderful way of framing many analogy
issues. But the target shape need not be subtly curvilinear for the analogy problem to
be deeply challenging. If you are trying to export even a very simple shape from one
grid-its "natural habitat"-into another grid, and if it does not export literally to the
target grid, then something's gotta give, and that is the hallmark of a hard analogy
problem.
        Since we are talking about chessboards, let us use a chess example. The
underlying grid of chess is a square lattice. Suppose we pick as our target grid the
hexagonal or triangular lattice (see Figure 24-8), and ask, "What is the knight's move
on this lattice?" We are immediately forced to confront the question, "What is the
essence of the knight's move in the only case we really know?" There are a number of
ways of thinking about it (see Figure 24-9). Which of the following, if any, is the
most insightful characterization of the knight's move?

(1) a rook step of length two followed by a single perpendicular rook step; (2) a single
rook step followed by a perpendicular rook step of length two;




Analogies and Roles in Human and Machine Thinking                                   591

FIGURE 24-10. Recognizing that the coloring of the board plays a significant role in
defining the moves of chess pieces. In (a), the standard coloring pattern of a square
lattice. In (b), the most natural coloring of a hexagonal lattice. Notice that it involves
three colors. In (c), the most natural coloring of a triangular lattice. Finally, in (d),
the guesses of Figure 24-9 are now shown on a colored-in hexagonal lattice. How
does this affect their plausibilities?




Analogies and Roles in Human and Machine Thinking                                     592

(3) a rook step of length one extended by a bishop step of length one;
(4) a bishop step of length one extended by a rook step of length one;
(5) a normal pawn move followed by a pawn's move in "capture mode";
(6) the shortest move that no other piece can make.

Or are all of these simply aspects of the essence of the knight's move? Which aspects
are more central, then? Which would you be willing to relinquish first? Which never?
Are you sure? How slippable are the following aspects, all of which do apply in the
square grid?

•      When a knight moves, it must land on a square of a different color.
•      A knight must be able to jump over or around pieces.
•      A knight can make a tour of the entire chessboard, landing on every square.
•      The starting and stopping squares of a knight's move should lie on opposite
       sides of a straight line that contains one edge of each.
•      The knight's move should not resemble any other piece's move.
•      All knight moves should be congruent except for rotation and reflection.
•      A knight must have about the same power as a bishop and less power than a
       rook.

Once you have tried extending the concept of the knight's move to another lattice,
then you begin to sense all the subliminal features that add up to define its highly
composite identity, features that you most likely never would have thought about
without this pressure. For example, coloring the lattices was a revelation for me, and
turned out to be the royal road to finding elegant knight's-move solutions (see Figure
24-10).
        While working on these two puzzles, I dreamt up what sounded at first like an
absurd challenge: to compress chess into one dimension. In other words, take a chain
of squares of length N (to be determined) and find moves for rook, bishop, knight,
queen, king, and pawn (see Figure 24-11).




FIGURE 24-11. In (a), a board for unidimensional chess, known as chass. Its optimal
number of squares is to be determined. In (b), a wider board for quasi-unidimensional
chess. This variant is known as chass or cheess.

Analogies and Roles in Human and Machine Thinking                                 593

Also consider how to place the pieces on the board at the game's start, and determine
N. This droll exercise gave rise to a number of very stimulating discussions among Al
colleagues and chess-playing friends of mine. The sense of sheer analogical elegance
had to compete with realism about what choices might make the game more
interesting. Along the way, one unexpected suggestion arose: What about-a board two
squares wide, rather than just one? On that type of grid, the knight's move seemed
trivially obvious-but my "obvious" solution turned out to have a fatal flaw, which we
patched in a most intriguing way.
One of the more amusing spinoffs of those discussions was a quest for the name of the
game (as they say). One-dimensional chess was dubbed chass, since 'a' is the first
vowel. Two-dimensional chess retained its name, 'e' being the second vowel. (But
what would seven-dimensional chess be called?) The 2 X N game, delicately poised
between one- and twodimensionality, was yclept chass. And what about the names for
chess on the nonstandard two-dimensional lattices? Here are some solutions that fell
into a delightful pattern:

       Chesh: chess played on a hexagonal lattice;
       Chest: chess played on a triangular lattice; and of course,
       Chess: chess played on a square lattice.

And please-when you are about to deliver the death blow to your opponent's king in a
game of chass, don't forget to triumphantly exclaim, "Chackmate! "

                                       *   *   *

        I found these chess-extension puzzles to be beautiful not only as puzzles, but
much more rewardingly, as examples of issues in analogy. When, for example, I
settled on my answer for chesh, it felt like not just an answer, but the answer: the
knight's move on a hexagonal lattice. This reminded me strongly of the feeling of
absolute certainty that one gets in mathematics when one sees a familiar phenomenon
recurring in a new way in an unfamiliar domain. One says, "Aha! So this is how
Wiggler's Lemma generalizes! The twistoploppic theomorphism is the same for even
clackdoodles, but becomes a hypertwistoploppic pseudotheomorphism for odd
clackdoodles. That's so beautiful!"
        Examples galore of this feeling must have arisen in the minds of the people
who extended the Magic Cube concept to other polyhedra, other dimensions, other
ways of slicing. And once you have made or acquired a new "cube" (such as the
Skewb or IncrediBall), you will want to know how to export a known algorithm,
broken up into its fundamental operators, from a familiar cube. What is the essence of
each operator? One senses a deep invariant lying somehow "down underneath" it all,
something that one can't quite verbalize but that one recognizes so clearly and
unmistakably in each new example, even though that example might violate some
feature one had thought necessary up to that very moment. In fact, sometimes that
violation




Analogies and Roles in Human and Machine Thinking                                 594

is what makes you sure you're seeing the same thing, because it reveals slippabilities
you hadn't sensed up till that time.
         No better example exists than the way mathematicians extended the concept of
exponentiation-putting x to the y power. At first, y had to be a positive integer. Then
it was realized that x°=1 fits exactly into the pattern, so zero was allowed as an
exponent. Immediately, it was seen how the same pattern would suggest-nay, require-
that x-1 be the reciprocal of x. By then the generalization ball was off and rolling.
Fractional powers came along very quickly: 1/2 as an exponent meant you should take
the square root, 1/3 meant the cube root, and so on. Then on to real numbers. But why
stop there? Various abstract representations of what it meant to exponentiate had now
been formulated, allowing one to transcend earlier, primitive notions of what it meant.
Pretty soon, not only could complex numbers be exponents, but so could n X n
matrices, functional operators, and God knows what else! This conceptual supernova
was still very much centered on one core, and blurry though the implicosphere around
it might be, the vastness of this implicosphere's size only made the conceptual core
stronger, firmer, realer.
         Another example: There is clearly only one sensible 4 X 4 X 4 Magic Cube. It
is the answer; it simply has the right spirit. The same holds for the four-dimensional
cube, discovered by Kamack and Keane. Similarly, Scott Kim once generalized the
concept of the "impossible triangle" (a twodimensional drawing read by three-
dimensional viewers as a three-dimensional object that cannot exist) up one
dimension, so that it became the "impossible skew quadrilateral" (a three-dimensional
sculpture read by four-dimensional viewers as a four-dimensional object that cannot
exist). Later, he found out that Roger Penrose, the inventor of the impossible triangle,
had likewise "added 1" to his three-dimensional illusion and come up with exactly the
same construction as Scott had-only Penrose did it fifteen years earlier. Clearly, then,
this was the corresponding paradox, manifesting the same deep essence as the original
and simpler one. Here again we see the aptness of that wonderful saying, Plus la
change, plus c'est la meme chose.
         That feeling of encountering an absolute and almost divine truth andreality
behind highly abstract analogical connections is particularly prevalent in mathematics,
but it can also arise in other areas Qf life. When a "reminding" feeling becomes so
strong that you want to use the same word, that is when you start getting religious
about your discovery. Golomb's "quarks" on the cube, for instance, seem to have
some "essence of quarkness" about them. Is this one phenomenon manifesting itself in
two different ways, or is it simply a pretty coincidence? Such questions can
occasionally not be answered, but very often our minds come to conclusions on such
matters without our ever noticing it. Reification of new categories in words is a
telltale signal, and one of the most important of mental events.

                                      *    *    *




Analogies and Roles in Human and Machine Thinking                                   595

        Some people might look upon the exercise of translating the knight's move
into an alien grid as an amusing but trifling game, and maintain that such things are
far from real-world concerns. Actually, in recent years, problems not too different
from this have become the stock-in-trade of people working on the computerization of
typefaces, where the idea is to pack as much of the spirit of a typeface (such as
Helvetica) into the smallest possible number of "pixels" (on-off dots, usually arranged
in a square lattice, though that is not necessary). Can one make an 'a' that is
recognizably a Helvetica `a', using just 35 pixels arranged in a 5X7 array? This is
certainly beyond feasibility. But how few can you get away with? When does at least
a hint of "Helveticality" start to appear? (See Figure 24-12.) And just what is this
"Helveticality" spirit that is so elusive? How much harder to capture than "essence of
knight's move" is it?
        Attempting to compress a visual form into smaller and smaller arrays of pixels
forces one to confront ever more deeply the question of its essence. What can one
afford to release, and what must be held onto? An analogous aural analogy problem is
very obvious to state, yet seldom explored: Can one translate a complex piece of
music from major into minor, or vice versa? Musicians will immediately recognize
that the major and minor scales here play the roles of underlying lattices, so that we
are undeniably dealing with a lattice-conversion problem. Mechanical methods will
carry you a certain distance, to be sure, but for any complex piece there will always
remain a lot of sticky and idiosyncratic knots. For instance, what if a piece in a major
key turns minor for a short stretch? Should its minor-key "translation" turn major at
the corresponding point? This example is just the tip of the iceberg in the major-minor
translation game. To get into the right spirit, you might try humming to yourself such
old favorites as the popular song "Awful Days Are Here Again" (traditionally sung by
mournful Democrats right after they lose an election) and Frederic Pichon's celebrated
Baptismal March (from his piano sonata in B-flat major) ...
        Another musical analogy problem arises when one tries to arrange a piece of
music for a new instrument or group. Can George Gershwin's very pianistic preludes
for piano be adapted for guitar, for example? Could one convert the wonderful
Mendelssohn violin concerto into a piano concerto? Each instrument forms a kind of
grid, and inter-grid transfer of essence is the problem.
        From vision and hearing, we now move to a more conceptual domain: pieces
of writing. The task of compressing a piece of text one has written into fewer and
fewer words forces one to struggle to define the essence of what one is trying to get
across. Up to a point, a piece of text may actually be improved by having some fat
trimmed here and there, just as a university or government agency can undoubtedly
benefit now and then from a severe budget crunch-but this can be carried too far, and
meaning will certainly begin to suffer. A fascinating exercise is to try to pack a page
of one's writing into half a page, then into a quarter page, and so on, until one has
gone




Analogies and Roles in Human and Machine Thinking                                   596

FIGURE 24-12. Helveticality emerging from the gloom. Proceeding from bottom to top, we
have a series of increasingly fine-grained dot matrices within which to maneuver. Clearly,
both the `a'-ness and the Helveticality get easier and easier to recognize as you ascend-
especially if you look at the page from a few feet away. Proceeding from left to right, we have
a series of increasingly letter-savvy programs doing the choosing of the pixels to light up. (As
a matter of fact, the rightmost column is a very light touch-up job of the third column, done by
a human.)
          The leftmost column is done by a totally letter-naive program. It takes the curvilinear
outline of the target shape and turns on all pixels whose centers fall within that outline.
          The second and third columns are the work of an algorithm that has information
about zones likely to be characteristic and critical for recognizability. It mathematically
transforms the original outline so that the critical zones are disproportionately enlarged (the
way your nose is enlarged when you look at yourself in a spoon). It then applies the naive
algorithm to this new outline (pixels light up if and only if they fall inside). This amounts to an
interesting trade-off sensitivity in the critical zones is enhanced al the sacrifice of sensitivity
in less critical zones. Consequently, some pixels are turned on that do not fall inside the
letter's true outline, while some that do fall inside that outline remain off It's a gamble that
usually pays off but not always, as you can see by comparing the first and second letters in,
say, the third row.
          The difference between the second and third columns is that in the second column, the
critical zones are crude averages fed to the program and don't even depend on the letter
involved. In the third column, however, the program inspects the curvilinear shape and
determines the zones itself according to its knowledge of standard letter features such as
crossbars, bowls, posts, and so on. Then it uses these carefully worked-out zones just the way
the second algorithm uses its cruder zones: by distorting the true outline to emphasize those
zones, and then applying the naive algorithm to the new outline.
          But no matter how smart a program you are, the problem gets harder and harder as
you descend towards typographical hell• matrices too coarse to capture essential distinctions.
En route to hell, more and more sacrifices are made. Helveticalitygoes overboardfirst, then
'a'-ness; andfrom then on, entropy reigns supreme. But just before that point is the ultimate
challenge-and only people can handle it, so far. [Computer graphics by Phill Apley and Rick
Bryan. I




Analogies and Roles in Human and Machine Thinking                                             597

down to a phrafse of only a few words. This can be seen as both a translation problem
and an analogy problem. It is not usually considered either one, but just think: One is
trying to "say this" in an ever sparser and tighter language, an ever more severely
constricted grid. In a similar vein, learning to write in the language called "Nonsexist"
is a great exercise in translation and analogy, as is trying to become fluent in `e'-less
English, referred to earlier. Both provide you with a somewhat modified set of
stepping-stones, and force you to invent and then get accustomed to many new types
of constructs in order to say things that are easily said in the more prevalent mode of
speaking. It is very hard to become totally fluent in either language.

                                        *   *   *

        A significant problem these days, related to that of capturing "Helveticality" in
a low-resolution grid, is that of producing original and esthetically pleasing low-
resolution typefaces-in other words, instead of imitating a known curvilinear typeface,
inventing a new typeface whose natural habitat is, say, a 5 X 7 or 10X 12 grid, all of
whose letters are in "the same style" within that tiny world. Many human designers
have discovered solutions of great ingenuity, but machine designers? There are none.
        Letter Spirit, an Al project of mine currently on the back burner (it is
impatiently waiting for Copycat to come to a boil), has as its aim to produce a
program that can do just that: Given one or two low-resolution letters as inspiration,
complete the alphabet in "the same style"--or rather, the same spirit. Instead of using
pixels (points) as the primitive components of letters, however, I chose to use short
straight-line segments on a fixed grid containing just vertical, horizontal, and 45-
degree diagonal segments. I call those primitive segments "quanta". Figure 24-13
shows the tiny grid permitted, and the stunning variety of 'a's that one can realize
within it. Actually, I estimate there are well over a thousand ways to realize grid-
bound designs possessing some degree of `a'-ness; some will definitely hit the bull's-
eye while others will clearly be way out on the fringes of the `a'-sphere. Then of
course there will be many shapes that hover simultaneously near the fringes of two or
more Platonic letters' spheres of influence. Such shapes are anathema to the human
visual system, which greatly desires unambiguous category membership; they should
be likewise antithetical to the Letter Spirit program.
        The Letter Spirit grid, although seemingly a constraint, actually inspires flights
of fancy that total freedom would not (a fundamental and general lesson about the
deep connection between constraints and creativity). Figure 24-14 gives a sampling of
a few "gridfonts" inspired by various stylistic quirks in one letter or another. Once
again, this is only the tip of the iceberg. There are thousands of intriguing gridfonts to
be designed and savored. As of this writing, I have designed about 150 of them. You
could




Analogies and Roles in Human and Machine Thinking                                     598

FIGURE 24-13. 87 'a's composed of horizontal, vertical, and 45-degree "quanta" in
the Letter Spirit world. How many more shapes recognizable as 'a's do you suppose
lurk in the given grid? (Compare this figure with Figures 12-2, 12-3, and 12-4.)

say I'm addicted! Seven complete gridfonts by me are exhibited in this book, below
the introductory paragraphs to the seven sections.
        The Letter Spirit project was distilled from a far more ambitious dream: that of
producing a program able to create genuinely artistic, curvilinear, full-fledged
typefaces when inspired by one or more sample letterforms. I




Analogies and Roles in Human and Machine Thinking                                   599

FIGURE 24-14. The horizontal and vertical problems as they arise in the Letter Spirit
world. (Compare this figure with Figure 13-8.) The central problem of the Letter
Spirit project is to characterize what it is that items with "the same spirit " (i.e., in the
same row) have in common with each other, so that in general, given a sample letter
or two, the program can "get the hang of it " and then go ahead and design all the
remaining letters in the same spirit, thus creating an esthetically pleasing ' gridfont ".
Readers are encouraged to try their hand at completing the six gridfonts whose
beginnings are shown here, and inventing their own.




Analogies and Roles in Human and Machine Thinking                                       600

wrote "far more ambitious", yet in a way that's not right. After all, during each
boiling-down step (and there were quite a few between the initial conception of the
project and the final arrival at the grid), I assured myself that it truly preserved the
essence of the full typeface problem and merely eliminated some superficial aspect of
it. So in some sense I do believe that Letter Spirit actually encapsulates the central
problem not only of typeface design, but indeed of art and creativity in general. And I
am prepared to stand behind that claim, as long as you give me some grains of salt to
defend myself with.
        I recently had a fascinating visit at Bitstream, a Cambridge, Massachusetts
firm specializing in the digitization of human-designed typefaces. A typical task they
must do over and over again is to take a given font and adapt it from one high-
resolution lattice to another (for example, from 200x200 pixels per letter to merely
IOOX 100). For this, they have specialized graphics hardware that works just fine.
This grid-to-grid conversion is an easily mechanizable analogy, or translation, task.
However, when they want to take a given font and adapt it from a high-resolution
lattice to a medium resolution one (say, down to 15 X 15), their graphics machine
produces an unacceptably crude solution, filled with ragged edges and spurious pixels
of all sorts. To improve on this, Bitstream purchased an expensive Lisp machine and
developed a complex program for this purpose. Some of the results of that work are
shown in Figure 24-12. The point is, severe compression requires far more brute
hardware and sophisticated software than does gentle compression. Finally, when they
need to compress a font from a high-resolution lattice down into a truly coarse-
grained one (say, IOX 10), they turn the task over to human designers, because people
alone can handle the many interacting perceptual forces that emerge at this level of
resolution.
        At first, this may sound counterintuitive, but really it makes perfect sense.
With high-resolution grids, the graininess of the underlying medium all but
disappears, and it is child's play to convert from one grid to another. It wouldn't even
matter if the target grid were hexagonal, as long as it were sufficiently fine-grained.
But compression down into very coarse grids forces one to deal with the conceptual
and perceptual essence of visual forms-and essence, if anything, is the central problem
of analogy. In fact, a sense of essence, in essence, is, in a sense, the essence of sense,
in effect.

                                        *   *   *

        Any analogy can be viewed as an attempt to reproduce in one metaphorical
grid a form that exists in another metaphorical grid. The more coarse-grained the two
"grids" are, the more ingenious the analogy-maker has to be to perform the mapping.
Roles and substructures must be extracted and weighed and mapped against each
other. Shifts of all sorts,




Analogies and Roles in Human and Machine Thinking                                     601

up and down in abstraction as well as sideways in conceptual similarity, must be able
to take place. The analogy-maker attempts to judge proposed solutions for their
elegance, but in the end, only their performance in the world determines their success.
          The Copycat domain might appear less charming a domain than the nose-
touching and chess domains, less grabbing than the Letter Spirit domain. But that is a
superficial viewpoint. To make progress in science, one has to make sure that the
phenomena under study are truly isolated. I am banking on having carried out the job
of isolation very well, and now comes the stage of making the model. That project is
ongoing, and its method of attack-its vision of how to build a system that would run
on a real machine -is an esoteric and complex one. To relate that would be another
very long story. It is the domain itself that has been the subject of discussion here.
          I feel confident that this tiny alphabetic world allows all the key features of
analogy to make their appearances. In fact I would go further and claim: Not only
does the Copycat domain allow all the central features of analogy to emerge, but they
emerge in a more crystal-clear way than in any other domain I've yet come across,
precisely because of its stripped-down-ness. Paradoxically, Copycat's conceptual
richness and beauty emanate directly from its apparent impoverishedness, just as the
richness of the "ideal gas" metaphor emanates from its absolute simplicity. Time will
tell if this limb I am out on is solid.


Post Post Scriptum.
       In retrospect, it seems that this P. S. probably ought to have been a chapter on
its own. I did not dream that it would grow to this size; I merely wanted to let my
readers know what sorts of issues I am working on currently -and I discovered that
sketching that out takes a good deal of time. The original column was disappointingly
coolly received. I hope that this more complete explanation of the driving forces
behind my research projects will awaken more enthusiasm.
       Below I give our "answers" to the analogy problems given in the P. S. Each
problem merits a much longer discussion, but life is short.

Page 579:
       1. dab (chosen over rival cac)
       2. dba (hands down, over cbb)
       3. hard to decide between pdt (ugh!) and pcu (yuk!)
       4. pxqxsx, of course-not pxqxry or pxqxsy




Analogies and Roles in Human and Machine Thinking                                    602

       5. too obvious to need any comment
       6. hmm ... maybe aaabbbddd, maybe aaabbbddk, maybe aaabbbccl, maybe
          even aaabbbddl
       7. trqp, but maybe srqo (definitely not srqq)
       8. tptgtrit is pretty, but so is spsqsrst
       9. abcdeabcdab, of course
       10. bcdacdabc (it's a figure-ground problem-bcd is a, in "code")
       11. acg is way better than acf
       12. you know....
       Page 585:

       1. pqr, a far more insightful answer than pbc
       2. 2a. nws is the only reasonable answer
       2b. uuuuu (not vvvvv, despite the answer to 2a)
       3. aBc dEf pQr goes to aBc pQr dEf
       4. qabcxyzq is incisive, but has a strong rival in abcqxyz
       5. zabczdefzstuz-certainly better than zabcdefzstuz
       6. eeeffghhiii, and yes, that g in the middle is the whole point
       7a. dcbabcd-not hard
       7b. abbbc, based on seeing 3-1-3 go to 1-3-1
       7c. pfr-a daringly abstract vision of "inside-out-ness"

When a program can do analogies like this, I'll be impressed!!!

Post Post Post Scriptum.
        After I'd completed the P.S. and P.P.S., I ran into Richard Feynman at a
conference. I reminded him of my lecture at Caltech three years earlier; his somewhat
vague recollection of it was that it was "silly". I took that as a charitable way of
saying that he hadn't seen any point to it. Which made me think that maybe his
"village-idiot" stance was due to genuine puzzlement, and not just an act.
        I then told him, with a certain amount of trepidation, that in my new book I
had humorously referred to his blunt way of answering all my analogy problems as
"village-idiotic" a few times. Would this offend him? "Oh, no!" he said. "A while
back, Omni magazine interviewed me, and on their cover they advertised it as an
interview with the `world's smartest man'. I think it's good to counterbalance that-so
now you're calling me a village idiot. That's fine. I think my mother would agree with
you more than with Omni. "




Analogies and Roles in Human and Machine Thinking                                 603


                    Who Shoves Whom Around
                      Inside the Careenium?
                     or, What Is the Meaning
                         of the Word "I"?
                                    March, 1981



         The Achilles symbol and the Tortoise symbol encounter each other
                           inside the author's cranium.
\achilles Fancy meeting you here! I'd thought that our dialogue in Paris was the
  last one we'd ever have.
\tortoise You never can tell with this author. Just when you think he's done with
  you, he drags you out again to perform for his readers.
\achilles I don't see why we should have to perform at his whim. TORTOISE:
  Just try resisting. Then you'll see why. You don't have any choice in the matter!
\achilles I don't?
\tortoise Look-to refuse to perform is tantamount to suicide. Let's face it,
  Achilles-you and I (at least in these Hofstadterian versions of ourselves) come to
  life only when Hofstadter writes dialogues about us. We had it good in Gödel,
  Escher, Bach, but now that that's over and done with, I have a feeling the pickings
  are going to be pretty slim. Hofstadter knows he can't live off of us forever! So
  we'd better take what we can get!
\achilles Yes ... I remember those good old days. Sometimes we had such
  wonderful lines. Like that one you had, something about how the "Achillean flash"
  swoops about my brain "in shapes stranger than the dash of a gnat-hungry
  swallow". Isn't that how it went?
\tortoise Something like that. Hofstadter liked that one well enough




Who Shoves Whom Around Inside the Careenium?                                     604
or, What Is the Meaning of the Word "I"?

  that and he had me say it in at least two dialogues! Pretty strange, eh?
\achilles The way you talk about all this is so bizarre, to my mind. I mean,
  granted that we're figments of someone else's imagination; but still, you know how
  characters in a novel are supposed to "come alive" and have "wills of their own" . .
  . Surely it's not just a cliché?
\tortoise I wouldn't know. I'm not a novelist. Nor is Hofstadter. .
\achilles I mean, am I really just a tool of Hofstadter (however benevolent he is),
  or am I genuinely exerting my own free will here (as I feel I am doing)? What it
  comes down to is: Who pushes whom around inside this cranium?
\tortoise Now there's a planted line, if I ever heard one. That's a direct quote from
  GEB, page 710, where Hofstadter is quoting from Roger Sperry of split-brain
  fame. It's where Sperry's giving his mind-brain-freewill philosophy, which Mr. H
  evidently espouses. But let's get on with the subject matter of this dialogue. I think
  we've done enough introduction. You must have something on your mind,
  Achilles, which Mr. H wants to bring up through you.
\achilles I wish you'd quit putting it in that upside-down way, Mr. T. TORTOISE:
  All right. But am I right? Isn't there something you're just itching to tell me?
\achilles Come to mention it, yes. It's related to a book I saw in the bookstore the
  other day, called Molecular Gods: How Molecules Determine Our Behavior. It was
  the subtitle that intrigued me.
\tortoise In what way?
\achilles My first thought on reading it was, "Oh, that's interesting-I didn't know
  that the molecules inside me could affect me that much."
\tortoise A classic reaction.
\achilles I know it sounds dumb, but what's wrong with it?
\tortoise How can you say that? Molecules is all you are, my friend! Read Francis
  Crick's Of Molecules and Men someday.
\achilles Oh, yes-I know I'm made of molecules. Nobody could deny that. It just
  seems to me that my molecules are at my beck and call-not individually, of course,
  but in large "chunks", such as my fingers, when I play my cello or sign a check. So
  that when I decide to do something, my molecules are forced to come along. So-
  haven't you really got it reversed? Isn't it really the case that I shove those
  molecules around, and not vice versa?
TORTOISE (rather exasperated): What do you mean, "I"? What is this "you"?
\achilles How I feel-let me put it that way. My free will determines what I do.
\tortoise All right. Let me suggest a definition. Let me suggest that the term "free
  will", when you use it, is a shorthand for a complex set of predispositions of your
  brain to act in certain ways. Just a moment ago, you used the word "fingers" as an
  abbreviation for a whole bunch of molecules. In a similar way, the phrase "free
  will" could be thought of




Who Shoves Whom Around Inside the Careenium?                                        605
or, What Is the Meaning of the Word "I"?

  as an abbreviation for a whole bunch of natural tendencies and constraints. So ...
  your free will-your set of preferred pathways for neural activity to flow along-
  constrains the motions of molecules inside your brain, and those motions in turn
  are reflected in the patterns that your fingers will trace out.
\achilles Are you saying that when I say "free will", I'm really using a shorthand
  for a kind of "hedge maze", like the ones on the grounds of Victorian palaces, a
  maze that allows some pathways and disallows others?
\tortoise Yes, that's the idea-only of course this "hedge maze" is inside your skull,
  and is a bit more abstract. For instance, it's a little oversimplified to imagine that
  pathways are rigidly allowed or disallowed. It would be more accurate to think of
  the set of predispositions in terms of a set of pins in a pinball machine. You know
  what I mean by "pins"?
\achilles Those stationary round things with rubber "bumpers" that the shiny
  marbles bounce off of?
\tortoise Correct. Were you to take an average over a million marbles, you could
  find out how each pin statistically affects the way the marbles descend to the
  bottom. Pathways aren't just allowed or disallowed; rather, some are more likely,
  some are less likely, depending on how the pins are arrayed. But if you still like the
  image of the maze of hedges, that's not a bad one to hold in your head. The hedges
  make more rigid constraints -things are more black-and-white than with the pins.
  There are fewer degrees of freedom for the motions in a maze. But I can make the
  maze image richer. Suppose that in your maze, one of the effects of the people
  moving through the maze were that the hedges gradually shifted position. It's
  somehow as if the maze were formed of movable partitions that constrain the maze
  runners, yet the maze runners' paths gradually move the partitions, thus changing
  the maze.
\achilles You mean the maze runners could just decide-by free will-that they want
  to pick up a partition and plop it down somewhere else?
\tortoise Not like that. It's got to be a deterministic outcome of the act of maze
  running itself. Let me go back to the pinball analogy. It's more as if the pins,
  instead of being fastened on the board, were slidable objects like hockey pucks,
  objects that as they get banged around from above and below and all sides, slightly
  slip and change positions. The pins need not be circular; they could be longish so
  that two or more located near each other could act like a channel or a funnel for
  marbles. In any case, they are jounced around by the rapidly moving pinballs.
\achilles As in Brownian motion?
\tortoise Exactly. There are really two scales in time and space operating here,
  each affecting the other. The heavy hockey-puck-like pins appear almost stationary
  to the light marbles. To a casual observer who's following the motions of the
  marbles, the massive pins would appear to be determining the light marbles'
  motions, to be telling the marbles where to go-or in Sperry's words, to be "shoving
  them around".




Who Shoves Whom Around Inside the Careenium?                                        606
or, What Is the Meaning of the Word "I"?
\achilles I like that image. It agrees with my earlier view that I shove my
  molecules around.
\tortoise True-provided you identify -yourself" with the configuration of the pins.
\achilles That's a little strange, I admit.
\tortoise Now imagine a second observer, who's watching a film of the whole
  thing speeded up by a factor of a thousand or more. To her, there is a smooth,
  interesting patterned motion of a bunch of large, variously-shaped pucks. She says
  to herself, "Wonder why they're moving that way-I can't see anything visible
  causing any of it."
\achilles She doesn't see the marbles?
\tortoise No-they are shooting around so fast in this time scale that their tracks all
  blur together into one uniform background color with no apparent motion.
\achilles Ah, yes ... Now the facts about Brownian motion begin to come back to
  me. I remember how people were mystified by the jostling motions of colloidal
  particles in solutions when they looked at them under a microscope. They couldn't
  figure out what was causing such motions. The molecules that were battering them
  constantly were too small to be visible, and besides, they were moving too quickly.
\tortoise Exactly. An observer on this time scale might start to develop a sense for
  the slow drifting patterns of the pucks, even without having any clear notion of
  what's causing the pucks to move about.
\achilles It's a natural human tendency. Why not?
\tortoise The observer could anthropomorphize: "Oh, those two little ones don't
  like to be close together, and those two long thin ones are trying to be parallel"-and
  so on. So she develops a teleology, or a way of describing the heavy pucks'
  motions all on their own. She's quite unaware that they are being bombarded
  constantly by teeny objects, as in Brownian motion. (Let's pretend that the marbles
  are more like BB's -really small.) She doesn't know that something smaller is
  making the pucks swim around in those patterns.
\achilles So you can turn a knob on your movie projector and flip back and forth
  between the fast and slow views? Or even smoothly go between them? That's neat!
  At first, at the slowest setting, the immobile pucks seem to determine the paths of
  the many little bouncing marbles. As you speed up the film, the marbles become
  harder and harder to track, and pretty soon they become just a big blur. Meanwhile,
  you begin to notice that the pucks actually aren't immobile, after all. They're being
  shoved about by the marbles. So-who's shoving whom around really ? Well, it's
  mutual, I now see.
\tortoise Good. Now let me add some more richness to this whole metaphor. Let's
  say that marbles are constantly being shot in from all sides of the table, and also
  leaving on all sides. You can envision something like a pool table, with a lot of
  little marble-launching stations mounted on the walls, and a lot of pockets that act
  as exits for stray




Who Shoves Whom Around Inside the Careenium?                                        607
or, What Is the Meaning of the Word "I"?

I
  marbles that land in them. The inflow and outflow are equal, so there's no net gain
  or loss of marbles. And the bombardment is pretty uniform, but not exactly. The
  marbles are launched according to conditions outside the table. For example, if
  there's a red light near a marble-launching station, that station slows down its firing
  rate; if a green light is near it, it speeds it up. So you have a set of transducers from
  external light to internal marble-shooting. Now if the puck observer watches both
  the lights and the pucks, she'll be able to draw some causal connections between
  light patterns outside and the puck-patterns inside. Using mentalistic language will
  become quite natural. For instance, it would sound quite reasonable to say, "It saw
  the green light-it's moving away from it-I guess it doesn't like green." And so on.
\achilles Now you've got me thinking. I too want to add some strange features. I'll
  propose a physical linkage between one particular puck and an external "arm" that
  can move toward or away from the lights. So, when that puck moves a certain way
  on the table, the arm may push a light away or pull it closer. Of course this is
  primitive-there are no fingers or anything, but at least there's now a two-way link
  between the pucks and the lights. Gosh! I'm almost completely forgetting about
  those marbles careening around down there! I'm just relying on the marble-
  shooters to keep on doing their job without much maintenance or attention needed
  ... All I see now is the seemingly animate interplay-a sort of danceamong the
  pucks, the lights, and the arms ...
\tortoise We're really jumping from one metaphor to another, aren't we? And each
  time, we escalate in complexity ... Oh, well, that's fine with me. No matter how
  complex the scene gets, you can always slow down the projector, unblur the
  marbles and no longer see the pucks moving at all.
\achilles Of course. But there's now something that bothers me. In the brain, there
  aren't these large- and small-sized units-everything's uniform, right? I mean, it's all
  just a dense packing of neurons. So where do the two scales come from? If we go
  back to the maze and partitions, there too we had two levels of objects (maze
  people and maze walls), each kind pushing the other around. But in the brain, this
  isn't so-or is it? What else is there besides neural activity?
\tortoise Let's add, then, a level of detail to our picture that we didn't have before.
  Let's say there are no pucks at all. There are only marbles and a number of larger
  stiff yet malleable mobile metal strips, which I'll also describe as "stiff yet
  malleable membranes" (and you'll soon see why). They can be bent into U's or S's
  or circles ...
\achilles So these things are swimming in the soup of marbles, now, but there are
  no more pucks, eh?
\tortoise Right. Can you guess what might happen now? ACHILLES: I can
  imagine that these strips
\tortoise Would you mind calling them "stiff yet malleable membranes", just to
  please me?




Who Shoves Whom Around Inside the Careenium?                                           608
or, What Is the Meaning of the Word "I"?
\achilles Are you going to pull some acronymic trick off in a moment? Let's see-
  "SYMM" doesn't spell anything, does it? Is that really what you want me to call
  them, Mr. T?
\tortoise In fact you anticipated me, Achilles. Go ahead and do call them
  "SYMM" 's.
\achilles All right. So these SYMM's will now be knocked around along with the
  marbles that are bashing into one another. Will the SYMM's occasionally get
  wrapped around some group of marbles and form a circular membrane, separating
  out a group of marbles from the rest?
\tortoise Just call the circular structure so formed a "SYMM-ball", if you please.
\achilles Oh ... I should have seen it coming. All right. Now I see that in this way,
  structures like pucks are emerging again, only this time as composite structures
  made-up out of many, many marbles. So now, my old question of who pushes
  whom around in the cranium-er, should I say `.`fit-tee careenium "?-becomes one
  of"symmballs versus marbles. Do the :marbles push the symmballs around, or vice
  versa? And I can twiddle the speed control on the projector and watch the film fast
  or slow.
\tortoise I should mention that once a symmball is formed, it might have quite a
  bit of stability, because the marbles inside it get fairly densely packed together, and
  jostle each other around only a little bit when the symmball gets hit by a fast
  marble from the outside. The impact gets spread around and shared among the
  marbles inside, and the symmball won't tend to break up-at least not when you
  watch the film at either of the two speeds we've already mentioned. Perhaps the
  fission of a symmball would occur on a longer time scale than the motions of
  symmballs. And the same for the formation of a symmball.
\achilles Would it be fair to liken a symmball's emergence to the solidifying of
  water into a cube of ice?
\tortoise An excellent analogy. Symmballs are constantly forming and unforming,
  like blocks of ice melting down into chaotically bouncing water molecules-and
  then new ones can form, only to melt again. This kind of "phase transition" view of
  the activity is very apt. And it introduces yet a third time scale for the projector,
  one where it is running much faster and even the motions of the symmballs would
  start to blur. Symmballs have a dynamics, a way of forming, interacting, and
  splitting open and disintegrating, all their own. Symmballs can be seen as
  reflecting, internally to the careenium, the patterns of lights outside of it. They can
  store "images" of light patterns long after the light patterns are gone-thus the
  configurations of symmballs can be interpreted as memory, knowledge, and ideas.
\achilles It seems to me that although you got rid of the pucks, you added another
  structure-the SYMM's. So how is this new system any improvement, as a model of
  a brain; over the old one? Don't you still have two levels of basic physical
  constituents and activity?




Who Shoves Whom Around Inside the Careenium?                                         609
or, What Is the Meaning of the Word "I"?
\tortoise The SYMM's are there only to provide a way for marbles to join up and
  form clusters. There are other conceivable ways I could have done this. I could
  have said, "Imagine that each marble is magnetic, or Velcro-coated, so that they all
  attract each other and stick together (unless jostled too hard)." That suggestion
  would have had a similar effect-namely, of making much larger units grow out of
  smaller onesand so you would have only one kind of basic physical constituent.
  Would that be more satisfying to you, Achilles?
\achilles Yes, but then you'd have lost your pun on "symbols", which would be
  too bad.
\tortoise Not at all! I'd cleverly rename the marbles themselves this time, as
  "small yellow magnetic marbles''--" SYMM's"-and a magnetically bound cluster of
  them would form a "SYMM-ball". No loss.
\achilles That's a relief! I would hate to see a good metaphor go down the drain
  for lack of a pun to illustrate it.
\tortoise Hofstadter would never let that happen! You can take it from me.
  Anyway, you can conceive of the larger units however you want, as long as you
  have it clear in your mind that starting with just one level, you wind up with two
  levels and two time scales-three time scales, in fact, when you take into account
  the slow formation, fission, fusion, and fizzling of the symmballs.
\achilles Now can we go back and talk about whether I control my molecules, or
  my molecules control me? That's where this all started, after all.
\tortoise Certainly! Why don't you try to answer the question yourself?
\achilles The problem is that in all those pulsations inside a careenium, I just don't
  see a "me". I see a lot of activity-I see a lot of internalized representations of things
  "out there"-I mean of light patterns, in this case. And with fancier transducers, we
  could have a careenium in which symmball patterns reflected such things as
  sounds, touches, smells, temperatures, and so on.
\tortoise Let your imagination run wild, Achilles!
\achilles All right. If I stretch my imagination, I can even see a gigantic three-
  dimensional careenium, hundreds of feet on a side, filled with billions upon
  billions of marbles floating in zero gravity, shooting back and forth, and all over
  forming short-lived and long-lived symmballs, and with those symmballs in turn
  governing the marbles' paths. I can see all that, and yet I don't see free will or "I". I
  guess I can't see how I myself could be a system like this inside my cranium. I feel
  alive! I have thoughts, feelings, desires, sensations!
\tortoise Hold on, hold on! One at a time. These are all related, but let's try to talk
  about just one-say, thoughts. Let me propose that the word "thought" is a shorthand
  for the activity of the symmballs that you see when you run the movie fast: the way
  they interact and trigger patterns of motions among themselves (mediated, of
  course, by the constant background swarming of marbles, too fast to make out).




Who Shoves Whom Around Inside the Careenium?                                           610
or, What Is the Meaning of the Word "I"?
\achilles But I feel myself thinking. There's no one inside a careenium to feel
  those "thoughts". It's all just a bunch of silly yellow magnetic marbles bashing into
  each other! It's all impersonal and unalive. How can you call that "thought"?
\tortoise Well, isn't it equally true of the molecules running around in your brain?
  Where's the soul of Achilles that "shoves them around"?
\achilles Oh, Mr. T, that's not a good enough answer. I've just heard it said too
  many times that we're made out of atoms, so there's no room for souls or other
  things-but I know I'm there, it's an undeniable fact, so I need more insight than a
  mere reminder that my body obeys the laws of physics. Where does this feeling of
  "I" come from, a feeling that I have and you have but stones don't have?
\tortoise You're calling my bluff, eh? All right. Let's see what I can do to turn you
  around. Let's add one more feature to the careenium-an artificial mouth and throat,
  just as we added an arm. Let various parameters of them be driven by various
  symmballs. Now suppose we turn on a green light on the right-hand side of the
  careenium. New marble activity near that side begins immediately, and there
  follows a complex regrouping of symmballs. As it all settles down into a new
  steady configuration, the mouth-throat combination makes an audible sound:
  "There's a green light out there." Maybe it even says, "I saw a green light out
  there."
\achilles You're trying to play on my weaknesses. You're trying to get me to
  identify with a careenium by making it more human-seeming, by making it
  simulate talking. But to me, this is merely "artificially signaling" (to borrow one of
  my favorite phrases from Professor Jefferson's Lister Oration). Do you expect me
  also to believe that somewhere out there, there is a conscious person reciting the
  time of day twenty-four hours a day, simply because I can dial a certain number
  and hear a human voice say, over the telephone, "At the tone, Pacific Daylight
  Time will be five forty-two"? A voice uttering sentence-like sounds doesn't
  necessarily signify the presence of a conscious being behind it.
\tortoise Granted. But this careenium voice isn't merely uttering a mechanically
  repetitious sequence of sentences. It is giving a dynamic description of what is
  perceived in the vicinity.
\achilles I have a question about that. Is the thing being perceived located outside
  the careenium, or inside it? Why does the mouth say, "I saw a green light out there
  " rather than say something such as, "Inside me, a new symmball just formed and
  exchanged places with an old one"? Isn't that a more accurate description of what it
  perceived?
\tortoise In a way, yes, that is what it perceived, but in another way, no, it did not
  perceive its own activity. Think about what perception really involves. When you
  perceive something "out there", you cannot help but mirror that event inside you
  somehow. Without that internal mirroring event, there would be no perception. The
  trick is to know what kind of external event triggered it, and to describe what you
  felt out loud in




Who Shoves Whom Around Inside the Careenium?                                        611
or, What Is the Meaning of the Word "I"?

  public language that refers to something external. You subtract one layer of
  transduction. You omit, in your description of what happened, one step along the
  way. You omit mention of the step that converted the green light into internal
  symmball responses. You are not even aware of that step, unless you are something
  of a philosopher or psychologist.
\achilles Why would I or anyone else omit a real level? What's the origin of this
  socially conventional lie? I don't omit levels in my speech!
\tortoise Actually, you do. It's a universal phenomenon. If you live near a railroad
  track and hear a certain kind of loud noise coming from that direction-rumbling,
  bells dinging, and so on-do you say, "I Near a train", or do you say, "I hear the
  sound of a train"?
\achilles I guess that ordinarily, I would tend to say, "I hear the train."
\tortoise Do you see a train, or do you see light hitting your eyes? When you
  touch a chair, do you feel the chair, or do you feel your feeling of the chair?
\achilles I opt for the simpler alternative. I never would think those extra
  philosophical thoughts that go along with it. What point would it serve to say, "I
  hear the sound of the train"?
\tortoise Exactly my point. The most convenient language, the least obfuscatory
  and pedantic, omits the heavy "extra" reference to the medium carrying the signals,
  omits mention of the transducers, and so on. It simply gets straight to the external
  source. This seems, somehow, the most honest way to look at things-and the least
  confusing. You hear and see a train, not an image of a train, not the light reflected
  off a train, not retinal cells firing-and most definitely not your perception of a train!
  We are constructed in such a way as to be unaware of our brain's internal activity
  underlying perception, and therefore we "map it outward".
\achilles Yes, I see that pretty well. I think I see why a careenium with a voice
  might talk about a green light rather than talk about its symmballs. But wait a
  minute. How would it know anything about green lights? It might prefer to refer to
  things in the outside world-but nonetheless, all it knows about is its own internal
  state!
\tortoise True, but its way of verbalizing its internal state employs words that you
  and I think refer to objects and facts outside the careenium. In fact, it too thinks so.
  But you could very well argue that it is just making sounds that mirror its internal
  state in some very complex way. It could be deluding itself. There might be
  nothing out there to refer to!
\achilles True, but that's not exactly my question. What I want to know is, how
  come it uses the right words to describe what's out there? Where did it learn to say
  "green light"? The same question goes for people. How come we all say the same
  sounds for the same things?
\tortoise Oh, that's not so hard. I had thought you were asking whether reality
  exists or not. I quickly tire of such pointless quibbling over solipsism. But let me
  answer the question you did ask. When you were a tot, you saw things-say, rattles-
  and heard certain sounds-namely,




Who Shoves Whom Around Inside the Careenium?                                           612
or, What Is the Meaning of the Word "I"?

  various pronunciations of the word "rattle"-at about the same time. Those sights
  and sounds were transduced from your retinas and eardrums into internal symbol
  states inside your cranium. Now, as a member of the human race, you were
  constructed in such a way as to enjoy mimicry, so you made funny noises
  something like "wattle", which were then automatically picked up by your
  eardrums, and fed back into the interior of your cranium. You heard your own
  voice, to your great delight and thrill! You were then able to compare the sounds
  you'd just made with your memory of the sounds you'd heard. By playing this
  exciting new game, you were learning the English words for objects. Of course you
  started with the nouns for visible objects, but quickly you built on that most
  concrete level and over the next few years you developed a large vocabulary
  including such things as "ball", "pick up", "next to", "splash", "window", "seven",
  "remember", "sort of", "zebra", "maze", "stretch", "of course", "by accident",
  "tongue-twister", "blunder", "confetti", "equilibrium", "analogy", "vis-a-vis ",
  "chortle", "Picasso", "double negation", "few and far between", "neutrino",
  "Weltanschauung ", "n-dimensional vector space", "tRNA-amino-acyl synthetase",
  "solipsism", "careenium"
\achilles Wait a minute! What about "banana split"?
\tortoise Now how did I overlook that? A shameful oversight. But I hope you get
  the point.
\achilles I think I see what you mean. Gradually, I internalized a huge set of
  external, public, aural conventions-namely the English words attached to particular
  states of my own brain, states that were correlated with things "out there".
\tortoise Not just things-actions and styles and relationships and so on.
\achilles To be sure. But instead of conceiving that the words described my brain
  state, it was easier to conceive of them as describing things out there directly. In
  this way, by omitting a level in my interpretation of my own brain's state, I cast
  internal images outward.
\tortoise A careenium would do likewise-casting its internal symmball patterns
  outward, attributing them to some properties of the external world. AI1d if a large
  number of careenia happened to be located near some specific stimulus, they could
  all communicate back and forth by means of a set of publicly recognizable noises
  that are externalizations of their internal states! So it's actually very useful to
  subtract out the references to the transduction, perception, and representation
  levels.
\achilles It all makes sense now. But unfortunately, something else. is bothering
  me! If the system projects all its states outward, talking about "green lights" and
  "red lights" and "traffic jams" and so forth, then how is there any room left for it to
  perceive its own internal state? Will it be able to say, "I'm annoyed" or "I forgot"
  or "I don't know" or "It's on the tip of my tongue" or "I'm in a blue funk"? Or will
  it project all those inner states outward as well, attributing weird qualities to things
  outside




Who Shoves Whom Around Inside the Careenium?                                          613
or, What Is the Meaning of the Word "I"?

  of it? Could there be inward-directed transducers that focus on symmballs and
  come up with a representation of symmball activity? That would be a sort of sixth
  sense-an inward-directed sense. TORTOISE: You could call it the "inner eye".
\achilles That's a perfect name for it.
\tortoise The inner eye wouldn't need to do much transducing, would it? It's the
  easiest thing in the world to monitor because it's right there inside you.
\achilles Now, Mr. T, you always warn me about confusing use and mention; I
  think you yourself are committing that error here. To have a word such as
  "Tortoise" in a text is enough to make somebody conjure up the image of a
  Tortoise, but it is not at all the same as making that person start to think about the
  word "Tortoise", is it? They may not notice it at all.
\tortoise Point well taken. There is a difference between having your symmballs
  in certain states, and being aware of that fact. It's something like the difference
  between using grammar correctly and knowing the rules of grammar.
\achilles Now I sense I could get really confused here-things could get very
  tangled. How can symmballs "watch" other symmballs? The ones that react to
  green lights, I can imagine and understand. There are transducers-the marble
  shooters on the borders. But would there be some symmball that always reacts to,
  say, the fusion of two symmballs? How would it detect such a fusion? What would
  make it react that way? Would it be a sort of satellite or U-2 plane, with an
  overview of the whole terrain of the brain? And what purpose would it serve?
\tortoise Imagine that you were watching an actual careenium, and at a very slow
  speed-so slow that you could reach down with your hand and pick up and remove
  an entire symmball before getting struck by any symms careening toward your
  hand. All of a sudden there would be a vacuum, where before there had been a
  dense mass of marbles. If you switched speeds now and watched the results in the
  symmballs' time scale, you'd see a massive regrouping of symmballs all over the
  careenium, a kind of shudder passing through the whole system as all the various
  symmballs come to occupy slightly different positions.
\achilles You could call such a shudder a "mindquake".
\tortoise An excellent suggestion. Various types of "mindquakes" would have
  characteristic qualities to them. They would have "signatures", so to speak. Now if
  you, Achilles, an observer from the outside, could learn to recognize such a
  signature, then why couldn't the system itself, from within, be even more able to do
  so? Such mindquakes would be, after all, just as tangible to the system as is an
  increase of marble-firings on any side. Both are simply internal events, even
  though the one is triggered by something external, while the other is set off by
  something internal.
\achilles So would there be various "seismometer symmballs", each one




Who Shoves Whom Around Inside the Careenium?                                        614
or, What Is the Meaning of the Word "I"?

  sitting there waiting to feel a specific kind of mindquake, and when that happens it
  would react?
\tortoise Sure. And for each type of mindquake, there is a special symmball there
  just sitting there like a pencil on end-and when its type of mindquake comes along,
  it topples. Of course that "toppling" in itself is just more symmball activity
\achilles Another mindquake?
\tortoise Precisely-and it can set off further reactions inside the careenium. The
  whole thing is very circular-one shudder triggers another one, and that one sets off
  more, and so on.
\achilles It sounds like it would never stop. There would just be constant
  symmball activity rippling back and forth across the careenium.
\tortoise Well, of course! That is what happens with conscious systems, isn't it?
  We're constantly thinking thoughts-some fresh, some staleconstantly mentally alive
  and aware-partly of the external world, partly of our own state-for example, how
  confused or tired we are, what something reminds us of, how bored we are with
  this long monotonous dialogue....
\achilles Hey, wait a minute! The reader may be bored, but I'm not!
\tortoise Only kidding, Achilles. Just trying to liven things up a bit.
\achilles All right. Well anyway, I admit that everything you've been saying is
  true, makes sense, but how is it useful for us to monitor our own state?
\tortoise Well, think first of a simple animal. What it needs most of all is food. Its
  brain-if it has one-is connected to its stomach by nerves, and it transduces an
  emptiness in the stomach into a certain configuration of symbols in the brain.
  Actually, this animal might be so simple that the symmball level doesn't exist.
  There might just be marbles zipping around in its cranium, but no larger-scale
  agglomerations. In any case, the effect of this may then be a shuddering in its
  brain, which produces repercussions on the animal's peripheries. It may move. All
  this is very much at the reflex level. Mostly it involves monitoring the organism's
  hunger state and controlling its limbs. Every organism has to monitor itself in
  terms of hunger. But primitive organisms don't use much information about the
  external environment they're in-they just flap about and "hope"-if that isn't too
  strong a word!-to encounter some food. Pretty unconscious. On the other hand,
  take a more complex animal. It will have an elaborate representation of its
  environment inside itself, so it also has a lot of options when it detects internal
  hunger. The symbolic activity representing the empty stomach has to be dealt with
  in the context of all the other symbols, which might represent danger, priorities
  other than eating, choices of when and what to eat, and so on. The total interaction
  of symbols at that point we might call "consideration" or "deliberation" or
  "reflection"-as distinguished from "reflexes". Now after all this, let me ask you:
  Does this help you to see why such a careenium might have a self?




Who Shoves Whom Around Inside the Careenium?                                      615
or, What Is the Meaning of the Word "I"?
\achilles Well ... I might grant that there's reflection going on in there, I might
  even grant that it's thinking-but there's nobody in there doing the thinking!
\tortoise Would you grant that there's free will inside there? ACHILLES: Hardly!
\tortoise Then I can see that you will need some more persuading. All right. Let
  me suggest that there is free will, and that this notion of a careenium may help you
  understand more clearly what free will truly consists in. We began this discussion
  by talking about whether you can "shove your molecules around" or not. This is a
  central question-in truth, it is the central question, I think. So I'd like to ask you,
  Achilles, can you freely decide to do anything?
\achilles Of course I can! That's precisely what free will is about! I can decide to
  do whatever I want!
\tortoise Really? Can you decide, say, to answer me in Sanskrit? ACHILLES:
  Obviously not. But that has nothing to do with it. I don't speak Sanskrit. How could
  I answer you in it? Your question doesn't make sense.
\tortoise Not so. You can only do what your brain will allow you to do, and that
  is very crucial. Let me ask you another question. Can you decide to kill me right
  now?
\achilles Mr. T! What a suggestion! How could you suggest such a thing, even in
  jest?
\tortoise Could you nevertheless decide to do it?
\achilles Sure! Why not? I can certainly imagine myself deciding to do it.
\tortoise That is beside the point, Achilles. Don't confuse hypothetical or
  fictitious worlds with reality. I'm asking you if you can decide to kill me.
\achilles I guess that in this world, in the real world, I could not carry out such a
  decision, even had I "decided"-or claimed I'd decided-to do it. So I guess I couldn't
  decide to do it, actually.
\tortoise That's right. That innocent-seeming trailer phrase that one tends to tack
  on-exactly as you did-is very telling, after all.
\achilles What innocent phrase? What do you mean?
\tortoise Don't you remember? You insisted vehemently to me, "I can decide to
  do whatever I want. " Now that phrase "whatever I want" may sound like a grand,
  universal, all-inclusive, sweeping phrase-but in fact, it represents quite the
  opposite: a severe constraint. It's not true that you are able to decide to do
  anything; you are limited to being able to decide to do only things you want.
  Worse yet, you are in fact limited to doing, at any time, the one thing that you want
  most to do! Here, "want" is a complex function of the state of the entire system.
\achilles Are you saying that choice is an illusion?
\tortoise Only to the extent that "I" is an illusion. Let me explain. It's quite
  common for people to develop interests that begin to consume them-doing puzzles,
  doing music, thinking about philosophy ...




Who Shoves Whom Around Inside the Careenium?                                         616
or, What Is the Meaning of the Word "I"?

  Sometimes such habits get so strong that they begin to interfere with the rest of
  their lives. A wife may pick up a bad habit-say, twiddling a cube or smoking cigars
  or constantly punning-and then try to get rid of it. Her exasperated husband may
  say to her, "What's this trying? Why can't you just decide to stop cubing? It is
  driving a wedge between you and me. Why don't you just decide to quit?" Yet the
  afflicted wife may, for all her good intentions, be unable to do so. Certainly having
  a modicum of desire is not enough. I would put it this way. The husband is
  appealing to what I would call his wife's "soul"-a coherent set of principles and
  tendencies and interests and personality traits and so forth that represent to him the
  person that he married. They have always before seemed to provide reasons or
  explanations for his wife's character, and he loves her for that aggregate of ways of
  being. So he appeals to this "soul" to put a clamp on its new obsession. But once
  the wife starts twiddling her cube, a part of her takes over. She gets obsessed-or
  should I say "possessed"?-by one of her own subsystems!
\achilles "Possessed" is the word for it. I myself find it very hard to stop
  practicing a piece on my cello once I have gotten into the swing of it. Before I
  start, I think, "Now, I'll just play this piece one time. " (Or, "I'll just eat one potato
  chip", or "I'll just solve the cube one time".) But then, once I've let myself start, I'm
  no longer quite the same person-some things inside me have subtly shifted. And
  the new me thinks, "That guy, said he'd do it only once. That's what he thought.
  But I know better!" There is a kind of inner inertia that makes me want to continue,
  even when there are other things I would also like to do. It's as if some part of you
  just "slips away" from a higher level of control-some subsystem gets "out of
  control" and won't obey the soul on top-like a bucking bronco unwilling to obey its
  rider.
\tortoise A powerful image. In such cases the wife herself may be confused and
  torn. Her inner turmoil is like that of a country in inner strife. There are factions
  battling each other-only in this case, the factions are neural firings, not people, of
  course. On some level, this woman may feel she wants to be able to decide to give
  up her habit-yet she may not have enough neurons on her side! And as in a country
  where the people won't support the government, so here: the "soul" has to have the
  support of its neurons! It can't just arbitrarily "shove them around", in reality.
\achilles I'm all confused. Who is in control, here?
\tortoise We'd like to be able to say that the symmballs can decide to do arbitrary
  things, but they are constrained. They are in a system that "wants" its parts to move
  in some ways but doesn't "want" them to move in others. We could come back to
  the hedge-maze metaphor, to make this more vivid.
\achilles Yes, but that applied to the lower-level objects-marbles, symms, or
  neural firings.
\tortoise Exactly. The "heavyweight" entities-hedges, pins, symmballs--




Who Shoves Whom Around Inside the Careenium?                                           617
or, What Is the Meaning of the Word "I"?

  constrain the "lightweight" entities-maze runners, pinballs, symms; but in revenge,
  the little ones, acting together, control how the high-level ones are arrayed.
\achilles So nobody's free here!
\tortoise Well, from the outside, that's the way it seems. But on the inside, the
  system may feel, just as you did, that it can "decide to do whatever it wants to do".
  But mind you, two symmballs in a careenium aren't free to decide, arbitrarily, on
  their own, to move in parallel-they have to have the cooperation of the marbles.
  The marbles have to do the work for them. Similarly, when the unhappy wife tries
  to "decide" to give up her cubing or punning habit, she can't do it without the
  agreement, the support, of her neurons.
\achilles YOU make this wife's "soul" sound like a general trying to marshall
  unruly neurons, to force them into line when they have their own paths to follow.
  A military general has some degree of power over his soldiers, so he can coerce
  them to some extent-but only so far. Beyond that, they'll mutiny. So the general has
  to go along with the tide. He can't really dictate policy-he can only resonate with it.
\tortoise It's true. However, sometimes an unexpected shift at a higher level can
  precipitate an abrupt "phase transition" of lower levels. A million tiny things
  suddenly find themselves swirling around in unexpected ways, and realigning in
  totally novel higher-level patterns. Once in a while just once in a while-the
  "general" can gain control of those unruly neurons-but only when they themselves
  don't know what they want, haven't reached any kind of consensus, and are instead
  in a malleable, leadable, chaotic state.
\achilles It sounds like you're describing a "snap decision"-an exercise of pure
  will power, such as when I say to myself, "I'm going to quit cubing right now!", or
  "I'm going to stop feeling sorry for myself and go out and get something useful
  done." But if I understand your way of looking at this kind of thing, even a phrase
  like "snap decision" is really just a kind of shorthand for summarizing a lot of low-
  level activity. Is that so? It seems to me it would have to be so, in your picture.
\tortoise You're right, saying something like "snap decision" is really a coarse-
  grained manner of speaking about a huge cloud of neural activity, like a huge
  blurry cloud of symms in a careenium projected at a high speed on the screen. And
  sometimes the activity of neurons inside a cranium, or of symms inside a
  careenium, lends itself admirably to such a high-level, coarse-grained, symbolic
  description-or in the case of a careenium, a "symm-ball-ic" description.
\achilles Not always?
\tortoise Are all ponds always frozen?
\achilles Oh, I see what you mean. If the relevant portions of the careenium are
  chunked into- symmballs, then a symmball-level description can be made. One set
  of symmballs is seen to affect other sets




Who Shoves Whom Around Inside the Careenium?                                         618
or, What Is the Meaning of the Word "I"?

  of symmbal!s regularly and predictably. Whereas if there are no symmballs just a
  lot of stray symms careening around with nothing to constrain them except the
  careenium's boundary-then it's kind of chaotic, and no higher-level description
  applies. But when the whole careenium is "symm-ball-ic"-when the phase
  transitions have taken place-then the person-I mean the careenium!-feels very
  much in control of his or her thoughts.
\tortoise Its thoughts?
\achilles Yeah, yeah-that's what I meant. Its thoughts. But when not enough phase
  transitions have taken place, then there's an indescribable hubbub: random symms
  careening all over the place without orderly constraints. But I wonder what it's like
  when the brain is in sort of a halfway state-when there are lots of symmballs, but at
  the same time still a lot of stray symms that belong to no one. It reminds me of a
  half-frozen lake in early winter, or early spring, when the molecules have only
  halfcoalesced into large blocks of ice.
\tortoise That's a wonderful state to be in. I find I'm most creative when I feel my
  brain consisting of such halfway-coalesced symbols-neurons acting somewhat
  independently, somewhat collectively. It's a happy medium where neural bubblings
  cooperate with symbolic channelings and yield the most creative, fulfilling, semi-
  chaotic sense of aliveness.
\achilles YOU think some of that uncoalesced freedom is essential for creativity?
\tortoise I was convinced of it by Hofstadter, who certainly feels that way. In
  GEB, writing about his plight as a writer, he portrayed himself as suffering from
  "helplessness" of the top level, for although he-or his symbol level-may in some
  sense have decided what to write, still he is entirely and utterly dependent on vast
  cooperating teams of neurons to come up with imagery and ideas and choices of
  words and sentence structures. Those lower-level items feel to the top level as if
  they "bubble up" from nowhere. But in reality they are somehow formed from the
  churning, seething masses of interacting neural sparks just as patterns of symmball
  motions emerge out of the chaotic Brownian motion of the many tiny symms. And
  a few of those ideas make it out through the narrow channel of verbalization, like
  grains of sand passing through the narrow neck of an hourglass. Yet most likely
  Hofstadter will insist that he himself is responsible for this dialogue, will desire the
  credit to accrue to him.
\achilles Hmm ... to the overall system that constrained the marbles to jounce in
  those ways . . . It is hard to assign "credit" or "blame", once you start analyzing
  thought mechanistically. I see that "decision" and "choice" are very subtle concepts
  that somehow have to do with the ways in which constraints on two different levels
  affect each other reciprocally, and at two different time scales, inside a cranium, or
  a careenium.
\tortoise You're getting the idea.




Who Shoves Whom Around Inside the Careenium?                                          619
or, What Is the Meaning of the Word "I"?
\achilles Every time you say "bubbling up", I can't help but think about the
  bubbles in ginger ale-I love the stuff. And I'm thirsty. I'm going to have some. Care
  for a glass yourself?
\tortoise Ah, ginger ale-capital suggestion.
ACHILLES sipping from a tall glass of cool ginger ale): Did it ever occur to you that
  when your leg is asleep, it feels like it's full of ginger ale? TORTOISE: Clever
  observation.
\achilles Not really original, I have to admit. I read it once in a "Dennis the
  Menace" cartoon.
\tortoise Are you sure you read it-or was it Hofstadter? ACHILLES: Spoilsport!
\tortoise It strikes me that having your leg fall asleep is one very weird
  experience. It's as if nature were giving you a chance, every once in a while, to be
  privy to all the tiny goings-on inside your leg, feeling the mingling tingling of
  trillions of cells all buzzing at once ...
\achilles Do you suppose that's what being alive really is like, and most of the
  time we're just numb to it?
\tortoise Precisely. Can you imagine if all of your body were always as fizzy and
  tingly as that? I've always wondered why people say their leg is asleep. We
  Tortoises say, "My leg is awake. " And French speakers say, "I've got ants in my
  leg." Those seem so much more accurate to me.
\achilles Phooey!
\tortoise What's the matter now?
\achilles I just realized that Hofstadter planted all of this. I mean, I'd thought I
  was genuinely thirsty. Now I see I was just being manipulated. He wanted to get
  certain remarks in here, and having me be thirsty was just a convenient avenue for
  him to do so. I should have known better. TORTOISE: Oh, so your ginger ale
  doesn't taste any good? ACHILLES: That's not what I mean. It tastes fine!
\tortoise Well then, what are you grumbling about? You're happy, he's happy.
  Would you have been happier if he were unhappy? That would seem a little
  perverse, even to me.
ACHILLES:. I guess you have a point. You know, now that I think about it,
  sometimes the decisions I make seem to be slow percolating processes, things that
  are utterly out of my control. In fact, a rather gory image that illustrates this idea
  flashed before my mind's eye while we were talking about the difficulty of
  breaking out of mental ruts.
\tortoise What was that?
\achilles I imagined a grim scene where a man's young wife is in a car crash and
  is badly mangled. He will certainly react. Perhaps he will react with love and
  devotion, perhaps with pity. Perhaps, to his own own dismay, he will even react
  with revulsion. But it occurred to me that in such emotionally wrenching cases,
  you can hardly decide what you will feel. Something just happens inside you.
  Subtle forces shift deep inside you, hidden, subterranean. It's quite scary, in a way,
  because in real crises like




Who Shoves Whom Around Inside the Careenium?                                        620
or, What Is the Meaning of the Word "I"?

  that, instead of being able to decide how you'll act, you find out what sort of stuff
  you're made of. It's more passive than active-or more accurately put, the action is
  on levels of yourself that are far lower-far more microscopic-than you have direct
  control over.
\tortoise Correct. You and your neurons are not on speaking terms, any more than
  a country could be on speaking terms with its citizens. There is, in both cases, a
  kind of collective action of a myriad tiny elements on low levels that swings the
  balance-exactly as in a country that "decides" to go to war or not. It will flip or not,
  depending on the polarization of its citizens. And they seem to align in larger and
  larger groups, aided by communication channels and rumors and so on. All of a
  sudden, a country that seemed undecided will just "swing" in a way that surprises
  everyone.
\achilles Or, to shift imagery again, it's like an avalanche caused by the collective
  outcome of the way that billions upon billions of snow crystals are poised. One
  tiny event can get amplified into stupendous proportions -a chain reaction. But the
  crystals have to be poised in the right way, otherwise nothing will happen.
\tortoise In cases of judgment, whether it be of one musical composer over
  another, one potential title or subtitle for a book over another, or whatever, the top
  level pretty much has to wait for decisions to percolate up from the bottom level.
  The masses down below are where the decision really gets made, in a time of
  brooding and rumination. Then the top level may struggle to articulate the seething
  activity down below, but those verbalized reasons it comes up with are always a
  posteriori. Words alone are never rich enough to explain the subtlety of a difficult
  choice. Reasons may sound plausible but they are never the essence of a decision.
  The verbalized reason is just the tip of an iceberg. Or, to change images, conflicts
  of ideas are like wars, in which every reason has its army. When reasons collide,
  the real battleground is not at the verbal level (although some people would love to
  believe so); it's really a battle between opposing armies of neural firings, bringing
  in their heavy artillery of connotations, imagery, analogies, memories, residual
  atavistic fears, and ancient biological realities.
\achilles My goodness, it sounds terrifying! You make the battlefield of the mind
  sound like a vast mined battlefield! Or a treacherous ice field on a steep mountain
  face. I never realized that a mechanistic explanation of thinking could sound so
  organic and living. It's sort of awful and yet it's sort of awe-inspiring as well. But I
  am very confused now about the "soul", the free will.
\tortoise I think that all these strangely evocative images have brought us back to
  your original perplexity, over the question of who pushes whom around in the
  cranium. Would you now be inclined to say, Achilles, that your molecules push
  you around, or that you push them around?
\achilles Actually, I'm not sure how this "I" fits into a cranium-or a




Who Shoves Whom Around Inside the Careenium?                                          621
or, What Is the Meaning of the Word "I"?

  careenium. You've got my head so spinning now that I don't know what's up or
  down.
\tortoise Wonderful! At least now your mind may be malleable. Do you see how
  "free will" in a careenium is actually constrained-physically constrained, I mean-by
  the "wants" of the system?
\achilles Yes, I see that these seemingly intangible "wants" are actually physical
  attributes of the overall system-tendencies to shun certain modes of behavior or to
  repeat certain patterns. So in a way I can see that a careenium has "free will" in this
  constrained sense of freedom. Maybe "free will" should be renamed `free won't".
\tortoise Oh, my, Achilles! Did you just make that clever one up?
\achilles I don't know-it just came to me. I never thought about it. It just "bubbled
  up from nowhere". I don't know who deserves the credit. Maybe Hofstadter made
  it up. Or maybe it just bubbled up inside his brain-although I don't quite see the
  difference.
\tortoise It sounds like the sort of thing Hofstadter's friend Scott Kim would say.
\achilles Hmm ... I still wonder, though-could a careenium's symmballs actually
  decide to do anything on their own?
\tortoise They certainly can't disobey the way the symms push them around-but
  on the other hand, the symms are always poised in just such a way that the one
  internal event that the symmballs most want to happen will happen. Isn't that a
  miraculous coincidence?
\achilles Now that I understand how all this comes about, I can see that it's not at
  all a coincidence. By the definition of "want", the symmballs will get shoved
  around the way they want to be (whether they like it or not)! I guess that the real
  conviction of having free will would arise when, repeatedly and reliably, a
  collection of symmballs wants something and then watches its desire getting
  carried out. It must seem like magic!
\tortoise It's what happens when you decide, say, to sign your name. Your fingers
  begin obeying you, and miraculously, you watch your name just appear before you,
  effortlessly! Is that magic?
\achilles Aha! That brings back that ultimately confusing term, "I". We say "I
  decide to sign my name." But what does that mean? I can see everything in a
  careenium-wants, desires, beliefs-but I just can't seem to take that last step. I
  simply fail to see an "I" in there.
\tortoise I've tried to explain that the word "I" is just a shorthand used by a system
  such as a careenium-a system that perceives itself in terms of symmballs and their
  predispositions to act in certain ways and not others-particularly a careenium that
  has not perceived that it is composed of small yellow magnetic marbles.
\achilles Perceive, shmerceive, Mr. T! There's no one inside a careenium who
  could perceive such a thing. Perception requires awareness, which no careenium
  has. There's no one inside a careenium to feel and experience and enjoy its "free
  will", even if it's there, in your sense. Or maybe the




Who Shoves Whom Around Inside the Careenium?                                         622
or, What Is the Meaning of the Word "I"?

  best way to say it is that there's perception and free will there, but there's nobody
  there to have it.
\tortoise You mean you seriously would grant that a physical system could have
  free will but you wouldn't then feel forced to say there was someone exercising
  that free will? Or that there was perception but no perceiver? Perceiverless
  perception? Agentless, subjectless free will? Soulless, inanimate free will? That's a
  real doozy!
\achilles I know it sounds paradoxical. I could almost agree with youexcept I'm
  still hung up on one thing. Just which perceiver, which agent, which subject, which
  soul would it be? Which person gets to be that careenium? Or maybe I should turn
  the question around: Which careenium gets to have a given soul? Do you see what
  I mean?
\tortoise I think so. You seem to be envisioning a corral of souls up in the sky,
  into which God (or some other Grand Overarching Deity) dips, whenever a new
  cranium or careenium comes into existence, and from which It pulls out a soul,
  imbuing that careenium or cranium with that identity forevermore-almost as if It
  were putting a cherry on top of a sundae.
\achilles Are you mocking me?
\tortoise I don't mean to be. If it sounds that way, it's only because I'm trying to
  take what I think your implicit notion of "soul" is and to characterize it explicitly,
  by putting it into as graphic terms as possible. But if you subtract out the imagery
  of a corral and God and cherries on sundaes, am I not putting into words the gist of
  your view?
\achilles In a way, I suppose so-only you've made it sound so silly that I hesitate
  to adopt that view now.
\tortoise It's so tempting to think that different I's are just "out there", dormant,
  waiting to be attached to structures, like saddles put on horses or cherries on
  sundaes. Then, once they are in place, suddenly there is a consciousness that
  "wakes up". As if the consciousness, and the identity, the "who-ness", were
  provided by the cherry, and without it there would be only a hollow "pseudo-I"-a
  thing possessing free will but with nobody to be ! Isn't that a little sad? Wouldn't
  you feel sorry for such a poor, deprived entity? Oh, no, of course you wouldn't-
  there would be no one to be sorry for, right?
\achilles Well, it's hard to see where a sense of "who I am" could come to a bunch
  of marbles in a careenium, or even to a collection of firing neurons. It seems to me
  that the identity has to be imposed on top of such a structure. A careenium is a
  complex pinball machine-a heap of metallic machinery-even if, unlike pinball
  machines, some of its states represent the world and its workings. But until you
  add some sort of living "flame" to that heap, it's empty-soulless. You need "flame"
  (although I admit I don't know quite what I mean by that term) to turn a physical
  object into a being, just as you need flame to turn a pile of wood into a fire. No
  matter how much lighter fluid you pour on it, without a




Who Shoves Whom Around Inside the Careenium?                                        623
or, What Is the Meaning of the Word "I"?

  flame, it's still inert.
\tortoise Wait a minute! A pile of wood starts burning when you set flame to it-
  but does it acquire a soul at that moment? No-as you said, it simply becomes active
  instead of inert. Any old flame would do. The identity of a fire doesn't come from
  the torch that lit it, but from the combustible materials! It's the transition from
  inactivity to activity that makes the flame seem so critical. But a careenium doesn't
  need to become more active than it is. Yet for some reason, Achilles, you seem to
  balk at my suggestion that in that activity there is as much reason to see an
  individualized soul as in neural firing activity. But what's so special about neurons?
  You know what you remind me of?
\achilles I don't know that I want to know, but tell me anyway.
\tortoise You remind me of somebody who runs into a pile of metal that's merrily
  burning away, and who declares that although it looks mighty like a fire, it surely
  isn't a fire (especially not a genuine fire!), because it's made of metal, and everyone
  knows that fires-especially genuine ones-are always made of burning wood or
  paper.
\achilles That sounds pretty silly and narrow-minded-more so than I am, I should
  hope. I'm not insisting that no careenium could have a genuine soul so much as I
  am wondering, "If a careenium had a soul, which soul would it be? Who would be
  this careenium, who would be that one?" On what basis could a decision be made?
\tortoise Wow, have you got things upside down! (Or backwards-I'm not sure
  which.) The same question goes for people as much as for careenia. Who gets to be
  which body? Do you also have the belief that any body could be anybody ? All it
  takes is the right flame inside? Could there be a "flame transplant", where someone
  else's flame-say mine-got implanted in your body, leaving your brain and body
  intact? Then who would be you? Or, who would you be? Or where would you be?
\achilles And where would you be, Mr. T? Something seems wrong in this
  picture, I admit. If a careenium is actually somebody, where does the decision as to
  who it is originate?
\tortoise I think you've got things backwards. (Or upside down-I'm not sure
  which.) First of all, it's not a decision-it's an outcome. Secondly, which "who" a
  careenium is is an outcome of its structure, particularly the way it represents its
  own structure in itself. The more it is able to see itself as an independent and
  coherent agent, the more of a "who" there is for it to be. Eventually, by building up
  enough of a sense of its unique self, it has built up a complete "who" for it to be: a
  soul, if you will. The continuity and strength of the feeling of "being someone"
  come from identification with past and future versions of the same system, from
  the way the system sees itself as a unitary thing moving and changing through
  time.
\achilles That's a strange idea-a thing whose identity remains stable even though
  that thing changes in time. Is it like a country that changes and




Who Shoves Whom Around Inside the Careenium?                                         624
or, What Is the Meaning of the Word "I"?

  yet remains somehow the same country? I think of Poland, for instance. If any
  country has had its soul-flame tampered with, Poland is it-yet it seems to have
  maintained a continuous "Polish spirit" for hundreds of years.
\tortoise A beautiful example. The sense of "one thing, extending through time"
  is very much at the root of our feeling of "being someone". And in a way it is
  nature's hoax: the illusion of soulsameness. Or, if you prefer not to call it an
  illusion, you can say that the ability of an organism to abstract, to think it sees
  some constant thing, over time, that it considers its self even as it changes, makes
  that organism's soul not an illusion.
\achilles You mean anything that can fool itself-I mean, see itself-as unchanging
  over time has a soul?
\tortoise That's not such a silly notion-provided that the verb "see" has its usual
  abstract meaning, not some dilution of the term. If the organism is as perceptually
  powerful as living ones like you and me, then I would definitely say it has a soul, if
  it sees itself as essentially "the same organism" over time.
\achilles But to see itself as an organism is not a trivial thing! It has to see itself as
  one coherent thing acting for reasons, not just randomly.
\tortoise Now you're talking! I couldn't agree more. Such a way of looking at
  something-namely, ascribing mental attributes to it-has been called. by Daniel
  Dennett "adopting the intentional stance" toward that thing. In the case of you
  looking at a careenium, it would come down to your seeing it at the symmball
  level, and interpreting the symmball configurations and the patterns they go
  through over time as representing the system's beliefs, desires, needs, and so on,
  overlooking the underlying masses of marbles, either deliberately or out of
  ignorance.
\achilles But you're not talking about me looking at a system; you're talking about
  a situation where some system does that to itself, right?
\tortoise Exactly. It looks at its own behavior and, instead of seeing all the little
  marbles deep down there making it act as it does, it sees only its symmballs, acting
  in sensible, rational ways
\achilles The system sees itself just as observers of the fast film see it! It could
  say of itself, "It wants this, believes that", and so on-only now it is ascribing all
  these beliefs and penchants and preferences and desires and so forth to itself, so
  instead it says, "I want this, believe that", and so on. This seems peculiar to me.. It
  makes up a bunch of hypothetical notions about itself simply out of convenience,
  then ascribes them to itself in all seriousness. For God's sake, though-if beliefs and
  desires and purposes and so on really existed inside itself, wouldn't the blasted
  careenium itself have direct access to them?
\tortoise What makes you think those beliefs aren't real? Aren't ice cubes and
  traffic jams and symmballs real? And what makes you think that this




Who Shoves Whom Around Inside the Careenium?                                          625
or, What Is the Meaning of the Word "I"?

  self-perception' isn't direct access to its beliefs? After all, does your perception of
  your own feelings via your "inner eye" differ so wildly from this?
\achilles I suppose not.
\tortoise When an outsider ascribes beliefs and purposes to some organism or
  mechanical system, he or she is "adopting the intentional stance" toward that
  entity. But when the organism is so complicated that it is forced to do that with
  respect to itsef you could say that the organism is "adopting the auto-intentional
  stance". This would imply that the organism's own best way of understanding itself
  is by attributing to itself desires, beliefs, and so on.
\achilles That's a very strange sort of level-crossing feedback loop, Mr. T. The
  system's self-image (as a collection of symmballs) is getting recycled back into the
  system but of course this depends on the very concrete symms themselves to carry
  it out. It's like a television looking at its own screen, recycling a representation of
  itself over and over, building up a pattern of nested self-images on the screen.
\tortoise And that stable pattern becomes a real thing in and of itself. If you were
  a careenium, merely by adopting the auto-intentional stance toward yourself, you
  would create a self-perpetuating delusion. As soon as you create this illusion that
  there is just one.thing there-a unitary self with beliefs and desires rather than a
  mere bunch of goalless, soulless marbles-then that illusion reenters the system as
  one of its own beliefs. The more that illusion of unity is cycled through the system,
  the more established and hardened and locked-in the whole illusion becomes. It's
  like a crystal whose crystallization, once started, somehow has a catalyzing effect
  on its further crystallization. Some sort of vicious closed loop that self-reinforces,
  so that even if it starts out as a delusion, by the time it has locked in, it has so
  deeply permeated the system's structure that no one could possibly explain how or
  why the system works as it does without referring to its "silly, self-deluding" belief
  in itself as a self.
\achilles But by that time it isn't so silly any more, is it?
\tortoise No, by then it has to be taken quite seriously, because it will have a lot
  of explanatory power. Once the self has become so locked-in, or "reified", in the
  system's own set of concepts, this fact determines much of the system's own future
  behavior--or at least if you are restricted to watching the fast projector, to looking
  at the symmball level, that is the easiest way to understand matters. And the
  curious thing is that this same level-crossing feedback loop (of adopting the auto-
  intentional stance) takes place in every careenium of sufficient complexity. So that
  whichever careenium you take, the stable self-image pattern that it finally
  establishes in this loopy way is isomorphic to the stable self-image pattern in every
  other careenium!
\achilles Bizarre! The medium is different, but the abstract phenomenon it
  supports is the same. It's a universal. That's sort of hard to grasp.




Who Shoves Whom Around Inside the Careenium?                                         626
or, What Is the Meaning of the Word "I"?
\tortoise Maybe so, but it's right. They all have isomorphic, identical senses of
  "I". There is just one sense of the word-just one referent just one abstract pattern-
  yet each one seems to feel it knows it uniquely ! There's a kind of fight for sole
  possession of something that everyone shares.
\achilles Soul possession, Mr. T?
\tortoise Very astute, Achilles.
\achilles Do you really believe there is just one "I", Mr. T?
\tortoise Not quite-an exaggeration for rhetorical purposes. The real point is,
  there's only one mechanism underlying "I-ness": namely, the circling-back of a
  complex representation of the system together with its representations of all the
  rest of the world. Which "I" you are is determined by the way you carry out that
  cycling, and the way you represent the world.
\achilles So you mean that all that determines who "I" am is the set of
  experience's some organism has gone through?
\tortoise Not at all. I said "the way things are cycled", not "which things are so
  cycled and represented". You've got to distinguish between the set of objects
  represented, and the overall style with which they are represented. It's that style
  that determines how the loop will loop. That's what creates the uniqueness of each
  "I".
\achilles Well, Mr. T, I think I am beginning to see your point. It's just\_ so hard,
  emotionally, to acknowledge that a "soul" emerges from so physical a system as a
  careenium.
\tortoise The trick is in seeing the curious bidirectional causality operating
  between the levels of the system, and in integrating that vision with a sense of how
  symbols have representational power, including the power to recognize certain
  qualities of their own activity, even though only approximately. This is the crux of
  the mental, and the source of the enigma of "I".


Post Scriptum.
    This piece was inspired by my brief contact with a brilliant meteor, my friend
Randy Read, a psychiatrist and writer who lived in San Diego. Randy died about a
year ago, as I write this, yet I still feel his spirit resonating in mine. I sometimes don't
understand why. I barely knew him, in a way-and yet in another way, for a short time
I think I was his best friend.
    I got to know Randy through his letters, beginning in 1979. The first letters from
him were intriguing and full of freshness, but there was also a definite brashness
about them that made me hold back for quite a while. Over time, however, his
exuberant way with words and his relentless




Who Shoves Whom Around Inside the Careenium?                                            627
or, What Is the Meaning of the Word "I"?

questing for complex truths grew on me, and I came to understand that much of his
brashness was just for fun. What came through loud and clear was a tremendous
passion to know and to revel in nature, and to find beauty. He wrote of the loneliness
of "edge life"-a balance somewhere between the impulsive and the reflective, the
intellectual and the emotional.
   I met Randy Read face to face on only a couple of occasions, but during them I
came to appreciate even more his keen sensitivity for beauty, especially music, and
his intense zest for life. I found him to be a surprisingly warm and strangely
vulnerable human being. We once did a little rock climbing on the Pacific coast, a
very memorable occasion for me, since I could enjoy Randy's outdoor side and his
verbal side at the same time. After that we continued to correspond and I kept his
many idiosyncratic letters and cards.
   In 1981, we tried to collaborate via mail and telephone on an expanded version of
the Careenium dialogue, but it didn't work out as well as we had hoped. Though the
outcome was less than perfect, this joint effort afforded us both much pleasure. After
that, our correspondence tailed off somewhat, but still I thought of Randy quite often.
   When I first learned that Randy Read had died, I was absolutely stunned. I had no
idea how much he had touched me. Here are some lines I wrote to myself just after
the shock.

     RR was a quester, a seeker, a reacher, a flailer, a yearner, a powerful,
  wonderful, spiritual soarer who could not deal, somehow, with flat mundane
  reality. No, that's not accurate. He loved reality, every.tiny cubic inch of it,
  every nook and cranny filled with paradox, purity, poetry, power.

   As a memorial to Randy Read, I would like to show how his ideas inspired the
Careenium piece and triggered a myriad other thoughts spread throughout this book.
Randy was a remarkable extemporaneous muser. His thoughts came out fast and
furious in wonderfully chosen words. Fortunately, he often used a dictating machine
and then sent me letters that were, in essence, transcriptions of his musings. Herewith,
then, I present a few of my favorite selections from "The Randy Reader".

                                       *   *   *

                                                                       March 31, 1980

      I'm sitting here today in the warm sunshine sipping on some papaya juice and
  contemplating the mysteries of life. What better time to inflict you with another
  letter!
      Well, here goes: I've just returned from another one of my forays into the
  mountains. I was up in the Sierra Nevadas in California and had a nice time
  despite the occasional blizzard. This has been a high-risk year for avalanches in
  the Sierras, so that more than the usual circumspection was necessary.
      Avalanches, even small ones, are extraordinarily deadly, yet notoriously hard
  to predict.
      Avalanche prediction is such a black art that meteorology seems a hard
  science by comparison. Snow pack, as you may know, has the broadest range of
  physical properties of any known physical substance. It can range from nevi an
  almost crystalline ice-like mass, to the flour-fine dust of a powder avalanche.


Who Shoves Whom Around Inside the Careenium?                                        628
or, What Is the Meaning of the Word "I"?

  Snow pack is plastic, elastic, rigid, brittle, solid, or liquid, depending on the
  interacting physical factors.
      An avalanche of snow pack is often caused by such a complex net of factors
  that the phenomenon seems almost to possess life. One doesn't have to be a
  fanatical animist to sense, in some complex systems, businesses that resemble
  the commerce of life. One ready analogy is that an avalanche with its almost
  "all-or-nothing" response resembles the triggering of a nerve cell. A small
  irritant input, usually the mountaineer himself, unleashes an enormous orgasmic
  response.
      Aggregation of loosely coupled elements can produce the "slop" needed for
  innovation. Counterfactuality in the mind is permitted by an elastic looseness.
  Likewise, systems such as avalanches defy easy prediction because they don't
  "have to" be one way.
      Primitive peoples often credit such complex systems with a mentality or
  spirit. Indeed, one does not have to be particularly superstitious to at times feel
  the presence of a sort of mind in a corniced snow pack. There is a character, a
  sort of irritable grouchiness at times, a playful perversity at others, but almost
  always a sense of being that comes from billions of flakes of snow that feel each
  other's presence.
      Crowd and flock behaviors resemble avalanches. Changes of public opinion.
  can be swift; a mob's mood can snap into violent action as abruptly as a
  collapsing wall of snow. Nations may even have some of a snow pack's
  perversity, for that matter. \_
      Much of the effort to understand systems like avalanches, the weather, gas
  flow in internal combustion systems, and other swirly things has rested on an
  attempt to analyze microscopic elements. Perhaps some day, we'll have the
  mathematics and measuring devices to pull off this sort of analysis, but such
  reductionism often misses how the artists do it. The most savvy mountaineers,
  the most skilled engine tuners rely on an intuitive sense of the personalities of
  the medium they work with. Measurements have their place, to be sure, but a
  sense of the phenomenon's personality often guides attempts to find solutions.
      I think in time computer sciences will lead the way in developing
  classifications of big system "personalities". Computer hardware readily lends
  itself to precise replication and in time, I'm sure we'll develop the sensitivity - to
  recognize big system minds as readily as we do a friend's face. Even today, this
  is beginning to be true, as quirks in wiring and software can yield distinctive
  personalities in our existing computer systems.

                                       *    *       *




Who Shoves Whom Around Inside the Careenium?                                           629
or, What Is the Meaning of the Word "I"?

                                                                    January 20, 1981

     In the natural environment, one can never achieve a zero error rate. So
  nature, in her delicate brilliance, finds a way to make errors work. It's sort of
  like this: outrageous errors are deadly, but slight errors keep things loose
  enough for the new to appear. If the steps of the error can be kept small enough,
  a near-perfect approximation can gradually be built.
     It's hard for me to put into words, but qualities like "quiveriness" and
  "vulnerability" come to mind when I think of creativity. Maybe another image is
  something like the doping used in the manufacture of semiconductors.
  Introducing small bits of something that doesn't belong there can dramatically
  change the structure of a large-scale array. It's not really randomness that is the
  vital ingredient in creativity, but rather the slight tint of the subtle truths that
  creativity finds that changes the random from white to pink. Michelangelo,
  standing before the marble that was to become David, had in his mind the image
  of a young man, but also allowed himself to be invaded by the block of marble,
  the whole universe itself. Creativity requires a sense of smell, a palate to taste
  the scents that make brilliance.
     "Chance favors the prepared mind."-Louis Pasteur. I like that quote. All life
  feeds upon the random. Creativity is simply the haute cuisine.
     Little ripples that exist or rather persist long enough to be observed are the
  lowest forms that possess what could be called personality. Rocks and snow
  tend to have very little personality because the ripples are very small-scale and
  dominated by randomness. Smoke has more. Indeed I think that's why Bach and
  Magritte liked pipe smoke and why I, too, confess to enjoy the habit.
     A puff of smoke is sort of like a cousin of ours. Little eddies, the loosely
  coupled systems, shear and spin, and we can observe the gentle drift of the
  whole ensemble. Bubbling streams have this quality, as do breaking ocean
  waves. The boiling and roiling, the little pieces, each with their own life, each
  wavelet connected, interacting, and yet participating in the whole.

                                     *   *   *

                                                                   February 23, 1981

                                    BUTTERFLIES

                        The best thoughts are the most delicate,
                              fastest, trickiest to capture.
                         Lepidoptera so different on the wing,
                               than when caught, killed,
                                and proudly displayed.

                                     *   *   *

   On April 10, 1983, Randy Read took his own life. I don't know why. Perhaps these
musings, dancing and sparking in the neurons of a few thousand readers out there,
will keep alive, in scattered form, a tiny piece of his soul



Who Shoves Whom Around Inside the Careenium?                                         630
or, What Is the Meaning of the Word "I"?

